% Generated by build_chapter.py from the GY1986 files. Do not edit by hand; edit the files and rebuild.
\chapter{The academic year 1986--87}\label{ch:1986}
\noindent\textbf{A source-defined year, with two chronological layers.}
This chapter covers all 46 catalogue records 667--712 and the 128 supplied
pages P001257--P001384. It follows manuscript order while retaining catalogue
record numbers and dates; these are identifiers, not a silently corrected chronology. The manuscript
dates in the opening spring sequence often say 1986 where the catalogue
says 1987, and record 667 has an earlier handwritten year. Such conflicts
are recorded at the affected talks. Shared page ranges are not copied
wholesale into both talks: headings, page layout, and the mathematical
continuity determine the attributed passage. In particular, the Gurevich
heading on P001303 qualifies record 684's catalogue attribution, and the
Beilinson heading on P001319 recovers surviving material for record 694,
whose catalogue entry has no linked pages. A boundary fragment at P001257
is linked to its earlier-year occurrence rather than reassigned to
Kontsevich. The catalogue itself was checked
\citep{mccme2026gelfandcatalogue}.

\medskip
\noindent\textbf{The mathematical lines.}
String measures and determinants, arrangements of real and complex
subspaces, quantum groups and their representations, matrix integrals
and map enumeration, integrable systems, modular lattices, automorphic
forms, and cohomological constructions run through the packet. The
apparently separate talks often share a concrete obstruction: deciding
what remains valid after a deformation, identifying the correct quotient
or completion, or distinguishing a formal generating series from an
analytic integral. The later fate is correspondingly heterogeneous.
Several talks can be matched to near-contemporary publications; others
restate earlier results; a few questions remain meaningful only after
their missing hypotheses and intended output are specified.

\medskip
\noindent\textbf{Inventory and assessment.}
There are 349 distinct indexed items, 21
problem-index entries, and 129 corrections or provenance notices.
An item is a mathematical unit, not a page or a sentence.
\begin{center}
\begin{tabular}{cl r}
\hline
Class & Historical assessment & Items\\
\hline
A & Published soon after by the speaker or Gelfand & 39 \\
B & Developed by others & 24 \\
C & Partly developed & 10 \\
D & No identified development after a documented search & 0 \\
E & Known before the talk & 242 \\
F & Not identifiable from the surviving notes & 34 \\
\hline
\end{tabular}
\end{center}
Class records historical fate; it does not certify literal correctness.
The separate fields retain missing hypotheses, false literal formulas,
partial publication matches, and the actual level of source inspection.
F includes unrecoverable fragments and unsupported precise readings;
it does not mean that their surrounding subject is unknown.
No item is assigned D: the searches did not justify an assertion that a precise idea had received no development.

\medskip
\noindent\textbf{Bounded developments.}
Two class C questions support explicit work. For the quantum KdV spectral
problem, a complete level-two calculation exhibits a genuine Jordan block
in an irreducible, nonunitary Virasoro module. Thus an unrestricted
eigenbasis formulation needs an extra hypothesis; the notebook itself
does not assert semisimplicity. For the two-source unitary integral,
probability Haar normalization gives a regular entire \(U(2)\) formula,
an all-rank convergent moment expansion with a proved tail bound, and
an exact \(U(3)\) counterexample to a specified naive permutation-only
extension. These do not settle an all-grade spectrum or a generic
large-\(N\) asymptotic problem. Two further annotations rederive a known
transverse-plane criterion and certify the notebook's existing answer
of 96 cube regions. Independent derivations or algorithms agree in
every reported test; code and complete outputs accompany the chapter.

\medskip
\noindent\textbf{Questions and continuity.}
The problem index separates construction of quantum KdV charges from
their spectral classification; finite multi-source integrals from the
broader reduced gauge integral; and established special cases from
unrestricted arithmetic conjectures. Earlier problem identifiers are
reused only for exact recurrences, notably the determinant question
GP-1982-40. Other earlier occurrences are linked as antecedents where
their quantifiers differ. Broad thematic resemblance is not an identity.
The toroidal reformulation of the Riemann hypothesis requires the
correct spectral condition, including temperedness; neither an
unrestricted unitarity slogan nor a related theorem settles RH.

\medskip
\noindent\textbf{Status and evidence.}
The cutoff is 12 November 2026. Newly circulated 2026 papers are explicitly
identified as preprints, with the version and theorem actually inspected;
their existence is not presented as independent proof certification.
Every bibliography entry carries a verification note, and unresolved
leads remain distinguishable from results read in the primary source.
The report provides counts by speaker and area, records the difficult
attributions and searches, and specifies how a website should expose
source, reconstruction, literature, and present status separately.

\GStalk{667}{1987-02-24}{M. L. Kontsevich}{P001257-P001260}{Formal uniformization of curves and the Virasoro action}
The mathematical contribution headed Kontsevich occupies P001258: a curve with a marked point and formal coordinate is acted on by Laurent vector fields, and the Virasoro extension lifts the action to determinant lines. P001257 is a continuation of Beilinson, not part of this contribution. The catalogue date is retained, although the heading visibly reads \(24/\mathrm{II}\text{-}85\); neighboring headings consistently read 1986. Consequently the record cannot establish priority relative to the paper received in December 1986.

\begin{GSitem}{1986-667.1}{A}{Infinitesimal uniformization by Laurent vector fields}
\GSfield{Kind; source}{construction; \GSpage{P001258}}
\GSfield{The notes}{The notes define \(\widetilde M_g=(C,p,x)\), \(\mathfrak g=\mathbb C((x))\partial_x\), and \(\gamma\) as Laurent expansions of vector fields on \(C\setminus\{p\}\); the stated action has stabilizer \(\gamma\).}
\GSfield{Statement}{For smooth connected projective complex curves with a marked point and a formal parameter, Laurent vector fields act infinitesimally on the corresponding moduli stack. At \((C,p,x)\), the tangent space is \(\mathfrak g/H^0(C\setminus\{p\},T_C)\), where the denominator is embedded by its Laurent expansion. This is a formal infinitesimal uniformization, not a claim that the coarse moduli space is globally a quotient by a finite-dimensional Lie group.}
\GSfield{Reconstruction}{The stack/formal interpretation and smooth projective hypotheses make precise the shorthand “homogeneous space.”}
\GSfield{Fate}{\citet[\S1]{kontsevich1987virasoro} gives this construction and its stabilizer. The paper was received on 3 December 1986 and published in 1987; this is the direct author match, independently of the notebook date discrepancy.}
\GSfield{Present status}{As of November 2026, the corrected formulation is an established result; no unresolved claim is inferred from this occurrence.}
\GSfield{Assessment}{literal statement: correct with missing hypotheses; publication match: exact; verification: read the result; development: none; confidence: high}
\end{GSitem}
\begin{GSitem}{1986-667.2}{A}{Meromorphic lifts and the Kodaira--Spencer quotient}
\GSfield{Kind; source}{method; \GSpage{P001258}}
\GSfield{The notes}{A family is given a lift \(\nabla_t=\partial_t+F(x,t)\partial_x\). One line identifies the tangent space with \(\mathfrak g/\mathfrak k\), but the following displayed Kodaira--Spencer formula uses vector fields extending off \(p\).}
\GSfield{Statement}{A meromorphic lift of a base tangent vector, allowed a pole along the section, gives a Laurent vector field modulo vector fields extending over the punctured curve. For the decorated space its quotient is \(\mathfrak g/\gamma\). Quotienting also by coordinate changes gives the appropriate tangent description for the undecorated moduli problem; \(\mathfrak g/\mathfrak k\) alone omits the curve-dependent stabilizer.}
\GSfield{Reconstruction}{The first quotient is emended using the next formula and the published construction; no change to the diplomatic transcription is proposed.}
\GSfield{Fate}{The tangent identification in \citet[\S1]{kontsevich1987virasoro} agrees with the second notebook formula. The isolated \(\mathfrak g/\mathfrak k\) line is a faithful mathematical error, confirmed by the scan.}
\GSfield{Present status}{As of November 2026, the corrected formulation is an established result; no unresolved claim is inferred from this occurrence.}
\GSfield{Assessment}{literal statement: false as written; publication match: exact; verification: read the result; development: none; confidence: high}
\end{GSitem}
\begin{GSitem}{1986-667.3}{A}{Virasoro action on determinant lines}
\GSfield{Kind; source}{theorem; \GSpage{P001258}}
\GSfield{The notes}{The central extension \(\widetilde{\mathfrak g}\), called the Virasoro algebra, acts on the pullback of the determinant bundle \(d_k\) from \(M_g\).}
\GSfield{Statement}{Let \(\lambda_k=\det H^0(C,K_C^k)\otimes\det H^1(C,K_C^k)^{-1}\). Its pullback to curves with a formal coordinate admits the corresponding lift of the Laurent-vector-field action to a Virasoro central extension. The central character depends on the determinant and cocycle conventions; the notebook does not specify a numerical normalization.}
\GSfield{Reconstruction}{The determinant of cohomology is explicitly defined here; the notes only name \(d_k\).}
\GSfield{Fate}{This is the determinant-line application in \citet[\S2]{kontsevich1987virasoro}, published in 1987. The brief notebook assertion should not be expanded into an unrecorded central-charge formula.}
\GSfield{Present status}{As of November 2026, the corrected formulation is an established result; no unresolved claim is inferred from this occurrence.}
\GSfield{Assessment}{literal statement: correct with missing hypotheses; publication match: exact; verification: read the result; development: none; confidence: high}
\end{GSitem}
\begin{GSitem}{1986-667.4}{E}{Beilinson boundary fragment: constancy and metric normalization}
\GSfield{Kind; source}{remark; \GSpage{P001257}; A. A. Beilinson (boundary fragment)}
\GSfield{The notes}{The preceding Beilinson continuation says that holomorphic functions on \(M_p\), \(p>2\), are constant, and concludes that the zero-curvature metric is unique and can be obtained from a holomorphic section.}
\GSfield{Statement}{This fragment is the continuation already indexed as 1985-666.21. Even when global holomorphic functions on the relevant moduli space are constant, multiplying a flat Hermitian norm by a positive constant preserves its curvature. The conclusion is uniqueness up to that constant after the appropriate global trivialization and monodromy conditions have been supplied; zero curvature alone supplies neither trivialization nor normalization. The Mumford form is a section of a determinant-line expression, not a scalar function on unmarked moduli.}
\GSfield{Reconstruction}{The record binding is editorial provenance. No new problem is created for an already indexed continuation.}
\GSfield{Fate}{The determinant-line interpretation is \citet{beilinsonmanin1986mumford}; the precise earlier occurrence is 1985-666.21, with the metric qualification also indexed as 1985-666.14. It is included here solely because this is the first catalogue record containing P001257, not because Kontsevich spoke this fragment.}
\GSfield{Present status}{As of November 2026, the corrected formulation is an established result; no unresolved claim is inferred from this occurrence.}
\GSfield{Assessment}{literal statement: correct with missing hypotheses; publication match: partial; verification: publisher metadata checked; development: none; confidence: high}
\end{GSitem}

\GStalk{668}{1987-02-24}{A. Yu. Morozov}{P001258-P001263}{Bosonic string measure, determinants, and the conformal anomaly}
The actual talk begins on P001259 and ends on P001261 before Beilinson. It reduces the critical bosonic-string functional integral formally to a moduli integral and explains the determinants, zero modes, and holomorphic factorization underlying that reduction. Gelfand asks about multiplicativity of regularized determinants. The heading reads \(24.2.86\); the catalogue date is preserved with this disagreement recorded.

\begin{GSitem}{1986-668.1}{E}{The genus-g bosonic measure on moduli}
\GSfield{Kind; source}{construction; \GSpage{P001259}}
\GSfield{The notes}{The functional integral over metrics and maps to \(\mathbb R^{26}\) is written as an integral over \(M_p\) of \(|d\mu_p|^2/(\det\operatorname{Im}T)^{13}\) times a vertex-insertion factor.}
\GSfield{Statement}{On the smooth genus-\(g\) moduli space, after fixing conformal gauge, the critical bosonic integrand has a holomorphic determinant-line factor paired with its conjugate and a period-matrix factor \((\det\operatorname{Im}\Omega)^{-13}\). Punctures and vertex insertions alter the corresponding line bundles and factors. This describes the integrand; it does not define a countably additive functional measure or prove convergence at the Deligne--Mumford boundary.}
\GSfield{Reconstruction}{The distinction between an integrand and a convergent amplitude is essential and is not made in the shorthand formula.}
\GSfield{Fate}{\citet[\S\S2--5]{belavinknizhnik1986complex} derive the determinant and holomorphic-factorization expression; \citet{beilinsonmanin1986mumford} identifies the Mumford-form interpretation. Both are 1986 publications, hence precede the catalogue date. Relative to the manuscript date they are contemporaneous.}
\GSfield{Present status}{As of November 2026, the corrected formulation is an established result; no unresolved claim is inferred from this occurrence.}
\GSfield{Assessment}{literal statement: correct with missing hypotheses; publication match: exact; verification: read the result; development: none; confidence: high}
\end{GSitem}
\begin{GSitem}{1986-668.2}{E}{Vertex operators and on-shell insertions}
\GSfield{Kind; source}{construction; \GSpage{P001259}}
\GSfield{The notes}{The scalar vertex is \(\int e^{ikx}\), and a vector vertex is written with one derivative of \(x\).}
\GSfield{Statement}{The free-boson exponential must be normal ordered. An integrated closed-string vertex must have conformal weights \((1,1)\); a boundary open-string vertex has weight \(1\). A single derivative polarization is the open-string vector vertex. Closed-string massless tensor vertices require both holomorphic and antiholomorphic derivatives. Momentum and polarization obey the corresponding on-shell and transversality conditions.}
\GSfield{Reconstruction}{Normal ordering, conformal weights, and the distinction between boundary and bulk vertices are supplied to make the formulas meaningful.}
\GSfield{Fate}{The vertex and Green-function insertion calculation is part of \citet[\S2]{belavinknizhnik1986complex}. The notebook does not keep open and closed sectors separate, so the vector formula cannot be read as a complete closed-string state.}
\GSfield{Present status}{As of November 2026, the corrected formulation is an established result; no unresolved claim is inferred from this occurrence.}
\GSfield{Assessment}{literal statement: correct with missing hypotheses; publication match: partial; verification: read the result; development: none; confidence: high}
\end{GSitem}
\begin{GSitem}{1986-668.3}{E}{Gaussian matter integration and zero modes}
\GSfield{Kind; source}{method; \GSpage{P001259-P001260}}
\GSfield{The notes}{Integration over \(26\) scalar coordinates produces a power \(-13\) of a scalar Laplacian determinant and a Green-function exponential.}
\GSfield{Statement}{For a fixed compact worldsheet metric, the Gaussian integral over nonconstant modes of \(D\) real scalar fields contributes \((\det{}^{\prime}\Delta_0)^{-D/2}\). The prime removes the constant kernel. Integrating the constant modes imposes total target-space momentum zero, while the source-dependent term is the Gaussian quadratic form obtained from the inverse of \(\Delta_0\) on mean-zero functions.}
\GSfield{Reconstruction}{The kernel removal and the momentum-conservation condition are restored.}
\GSfield{Fate}{\citet[\S2]{belavinknizhnik1986complex} explicitly separates these determinant and insertion contributions. At \(D=26\) the exponent is \(-13\); omitting the prime would give a zero determinant.}
\GSfield{Present status}{As of November 2026, the corrected formulation is an established result; no unresolved claim is inferred from this occurrence.}
\GSfield{Assessment}{literal statement: correct with missing hypotheses; publication match: exact; verification: read the result; development: none; confidence: high}
\end{GSitem}
\begin{GSitem}{1986-668.4}{E}{Cancellation of the conformal anomaly at dimension 26}
\GSfield{Kind; source}{theorem; \GSpage{P001260}}
\GSfield{The notes}{The gauge-fixed expression includes the determinant on vector fields and is independent of the conformal factor in dimension \(26\).}
\GSfield{Statement}{For the bosonic Polyakov system of \(D\) free scalar fields and the reparametrization \(bc\) ghost system of weights \((2,-1)\), the total central charge is \(D-26\). Its vanishing at \(D=26\) cancels the local Weyl anomaly in this system; it does not by itself remove tachyonic or degeneration divergences.}
\GSfield{Reconstruction}{The ghost weights and the scope of the cancellation are explicit here.}
\GSfield{Fate}{The critical-dimensional determinant cancellation is established in \citet[\S3]{belavinknizhnik1986complex}. The restricted system and the distinction between local anomaly cancellation and convergence are the needed hypotheses.}
\GSfield{Present status}{As of November 2026, the corrected formulation is an established result; no unresolved claim is inferred from this occurrence.}
\GSfield{Assessment}{literal statement: correct with missing hypotheses; publication match: exact; verification: read the result; development: none; confidence: high}
\end{GSitem}
\begin{GSitem}{1986-668.5}{E}{Gram determinants and determinant of cohomology}
\GSfield{Kind; source}{construction; \GSpage{P001260}}
\GSfield{The notes}{The matrices \(N^{(j)}\) are scalar products of zero modes, and the ratio involving \(\det\Delta_j\), \(\det N^{(j)}\), and \(\det N^{(1-j)}\) is used to express the measure.}
\GSfield{Statement}{For a family of Dolbeault complexes on \(K^j\), the Gram determinants on \(H^0(K^j)\) and its Serre-dual cohomology compensate for changes of zero-mode basis. Together with a zeta determinant of the nonzero spectrum they define a Hermitian norm on the determinant of cohomology. The expression is a line-bundle metric, rather than an invariant numerical determinant without a chosen frame.}
\GSfield{Reconstruction}{The basis transformation and determinant-line interpretation are supplied.}
\GSfield{Fate}{The zero-mode factors occur in \citet[\S3]{belavinknizhnik1986complex}; the determinant-line framework is \citet{quillen1985determinants}. These precede the catalogue date.}
\GSfield{Present status}{As of November 2026, the corrected formulation is an established result; no unresolved claim is inferred from this occurrence.}
\GSfield{Assessment}{literal statement: correct with missing hypotheses; publication match: exact; verification: read the result; development: none; confidence: high}
\end{GSitem}
\begin{GSitem}{1986-668.6}{E}{Beltrami differentials as local deformation coordinates}
\GSfield{Kind; source}{method; \GSpage{P001260}}
\GSfield{The notes}{The complex structure is varied by Beltrami differentials, described as Ahlfors--Bers coordinates.}
\GSfield{Statement}{For a compact Riemann surface of genus \(g\geq2\), infinitesimal deformations are \(H^1(C,T_C)\), dual to holomorphic quadratic differentials. A Beltrami differential modulo infinitesimal diffeomorphisms represents a tangent vector to Teichm\"uller space. Local holomorphic coordinates on that space descend to orbifold charts on moduli, not necessarily to an ordinary global coordinate system.}
\GSfield{Reconstruction}{The genus restriction and the Teichm\"uller/moduli distinction are restored.}
\GSfield{Fate}{\citet[\S3]{belavinknizhnik1986complex} uses this deformation description in the determinant calculation. The method is expository, not a new claim of Morozov.}
\GSfield{Present status}{As of November 2026, the corrected formulation is an established result; no unresolved claim is inferred from this occurrence.}
\GSfield{Assessment}{literal statement: correct with missing hypotheses; publication match: exact; verification: read the result; development: none; confidence: high}
\end{GSitem}
\begin{GSitem}{1986-668.7}{E}{Curvature of determinant metrics and the coefficient 6j squared minus 6j plus 1}
\GSfield{Kind; source}{claim; \GSpage{P001260-P001261}}
\GSfield{The notes}{The Belavin--Knizhnik variation is written as a mixed derivative of a determinant ratio equal to \(c_j=6j^2-6j+1\).}
\GSfield{Statement}{The meaningful local statement concerns the mixed second derivative of the logarithm of the determinant-line norm, with nonzero spectra used. It is a curvature two-form proportional to a universal two-form on moduli, and the coefficient for \(K^j\) is proportional to \(6j^2-6j+1\), with normalization fixed by the curvature convention. A second derivative of an unlogged ratio cannot equal a universal scalar independent of coordinates.}
\GSfield{Reconstruction}{This is a correction of the mathematical shorthand, not a claim that the missing symbols were visible in the manuscript.}
\GSfield{Fate}{Compare the variation formulas in \citet[\S3]{belavinknizhnik1986complex} and the determinant-line formulation of \citet{quillen1985determinants}. The scan confirms that the notebook formula itself suppresses the logarithm and differential-form factor.}
\GSfield{Present status}{As of November 2026, the corrected formulation is an established result; no unresolved claim is inferred from this occurrence.}
\GSfield{Assessment}{literal statement: false as written; publication match: exact; verification: read the result; development: none; confidence: high}
\end{GSitem}
\begin{GSitem}{1986-668.8}{B}{Multiplicativity of regularized elliptic determinants}
\GSfield{Kind; source}{question; \GSpage{P001261}; I. M. Gelfand}
\GSfield{The notes}{Gelfand asks whether \(\det(A_1A_2)=\det A_1\det A_2\) for elliptic operators and says the orders should be arbitrary.}
\GSfield{Statement}{For invertible admissible elliptic pseudodifferential operators \(A,B\) of positive orders on a closed manifold, with specified spectral cuts also admissible for \(AB\), determine \(\log\det_\zeta(AB)-\log\det_\zeta A-\log\det_\zeta B\). In general this multiplicative anomaly is nonzero; determinant multiplicativity is not an unrestricted law.}
\GSfield{Reconstruction}{The admissibility and invertibility hypotheses are added to obtain a well-defined zeta-determinant question; arbitrary-order operators require a different determinant framework.}
\GSfield{Fate}{This repeats GP-1982-40 (1982-587.7). \citet[\S1.1 and the anomaly formulas]{kontsevichvishik1994geometry}, circulated in 1994, gives a systematic treatment, with residue formulas under additional hypotheses and extensions beyond the positive-order setting. The formulation requires spectral cuts, which the notes do not fix.}
\GSfield{Present status}{As of November 2026, the unrestricted multiplicativity assertion is false and the multiplicative anomaly is an established part of elliptic determinant theory. The earlier problem identifier is retained.}
\GSfield{Assessment}{literal statement: correct with missing hypotheses; publication match: exact; verification: read the result; development: none; confidence: high; problem \hyperref[prob:GP-1982-40]{GP-1982-40}}
\end{GSitem}

\GStalk{669}{1987-03-03}{A. A. Beilinson}{P001261-P001265}{Mixed Hodge structures, real splitting, and Green-function periods}
The contribution starts on P001262 and ends at the beginning of P001265. It presents mixed Hodge structures in filtration and split forms, then interprets an extension attached to a punctured curve with two relative marked points through Green functions. The heading reads \(3/\mathrm{III}\text{-}86\), a year earlier than the retained catalogue date.

\begin{GSitem}{1986-669.1}{E}{Mixed Hodge structures over the integers}
\GSfield{Kind; source}{definition; \GSpage{P001262-P001263}}
\GSfield{The notes}{The data are a finitely generated abelian group \(L\), an increasing rational filtration \(W\), and a decreasing complex filtration \(F\), with a pure decomposition on each weight quotient.}
\GSfield{Statement}{An integral mixed Hodge structure consists of a finitely generated abelian group \(L\), a finite increasing filtration \(W\) on \(L_{\mathbb Q}\), and a finite decreasing filtration \(F\) on \(L_{\mathbb C}\), such that \(\operatorname{Gr}^W_m L_{\mathbb Q}\) with the induced filtration is pure of weight \(m\): \(\operatorname{Gr}^W_m L_{\mathbb C}=\bigoplus_{p+q=m}H^{p,q}\) and \(\overline{H^{p,q}}=H^{q,p}\). Tensoring kills torsion, which is retained only in the underlying integral group.}
\GSfield{Fate}{This is Deligne\textquotesingle s definition, published in \citet[\S2.3]{deligne1971hodge}. It is exposition in this talk.}
\GSfield{Present status}{As of November 2026, the corrected formulation is an established result; no unresolved claim is inferred from this occurrence.}
\GSfield{Assessment}{literal statement: correct; publication match: exact; verification: read the result; development: none; confidence: high}
\end{GSitem}
\begin{GSitem}{1986-669.2}{E}{Real splitting and a lowering endomorphism}
\GSfield{Kind; source}{construction; \GSpage{P001263}}
\GSfield{The notes}{The alternative description uses \(L_{\mathbb C}=\bigoplus L^{p,q}\), conjugate symmetry, an endomorphism \(N\) with \(\overline N=-N\), \(N(L^{p,q})\subset\bigoplus_{a<p,b<q}L^{a,b}\), and \(F=e^N F_{\mathrm{split}}\).}
\GSfield{Statement}{For a real mixed Hodge structure there is a canonical real-split mixed Hodge filtration and a unique real endomorphism \(\delta\) of strictly negative bidegrees such that \(F=e^{i\delta}F_{\mathrm{split}}\). Thus the notebook\textquotesingle s \(N\) is \(i\delta\). It is nilpotent by the strict bidegree inequalities. It is not the logarithm of geometric monodromy merely because it is denoted \(N\). Rationality is required of \(W\); conjugation alone gives only a real subspace.}
\GSfield{Reconstruction}{The canonical-splitting interpretation identifies the unnamed lowering operator. The notebook\textquotesingle s apparent inference from real to rational is not valid.}
\GSfield{Fate}{The real-splitting result is stated in \citet[Theorem 1.15]{cattanikaplan1989degenerating}, referring to the 1986 work of Cattani--Kaplan--Schmid. The scan confirms the anti-real condition and the strict double inequality.}
\GSfield{Present status}{As of November 2026, the corrected formulation is an established result; no unresolved claim is inferred from this occurrence.}
\GSfield{Assessment}{literal statement: correct with missing hypotheses; publication match: exact; verification: read the result; development: none; confidence: high}
\end{GSitem}
\begin{GSitem}{1986-669.3}{E}{Purity for smooth complete varieties and the direction of the weight filtration}
\GSfield{Kind; source}{claim; \GSpage{P001263}}
\GSfield{The notes}{Example 1 says that a smooth complex manifold \(C\) has pure \(H^a\) of weight \(a\), and writes \(W_a=H^a\), \(W_{a+1}=0\).}
\GSfield{Statement}{For a smooth projective complex variety (or compact K\"ahler manifold), \(H^a\) carries a pure Hodge structure of weight \(a\). With increasing weights the conditions are \(W_{a-1}=0\) and \(W_a=H^a\). Smooth nonproper algebraic varieties generally have mixed cohomology, and arbitrary complex manifolds need not carry the claimed structure.}
\GSfield{Reconstruction}{The compact K\"ahler or smooth projective hypothesis and the subscript \(a-1\) are editorial repairs.}
\GSfield{Fate}{\citet[\S\S2.1 and 3.2]{deligne1971hodge} distinguishes the pure proper case from the mixed open case. Both the omitted compactness/projectivity hypothesis and the wrong filtration subscript are visible in the notebook.}
\GSfield{Present status}{As of November 2026, the corrected formulation is an established result; no unresolved claim is inferred from this occurrence.}
\GSfield{Assessment}{literal statement: false as written; publication match: exact; verification: read the result; development: none; confidence: high}
\end{GSitem}
\begin{GSitem}{1986-669.4}{E}{The mixed Hodge structure of a punctured curve with relative points}
\GSfield{Kind; source}{construction; \GSpage{P001264}}
\GSfield{The notes}{The example is \(H^1(C\setminus\{x_1,x_2\},\{y_1,y_2\})\), with three weight pieces, and is drawn as a curve punctured at two points and with two relative endpoints.}
\GSfield{Statement}{Let \(C\) be a smooth connected projective complex curve of genus \(g\), and let \(x_1,x_2,y_1,y_2\) be pairwise distinct. For \(H=H^1(C\setminus\{x_1,x_2\},\{y_1,y_2\};\mathbb Q)\), one has \(\operatorname{Gr}^W_0H\simeq\mathbb Q(0)\), \(\operatorname{Gr}^W_1H\simeq H^1(C;\mathbb Q)\), and \(\operatorname{Gr}^W_2H\simeq\mathbb Q(-1)\). Hence \(\dim H=2g+2\). The homological version has dual weights and must not be interchanged silently with this cohomological statement.}
\GSfield{Reconstruction}{Connectedness and four distinct points are imposed; the cohomology/homology convention is fixed explicitly.}
\GSfield{Fate}{This is the relative-curve instance of Deligne\textquotesingle s mixed Hodge structures, \citet[\S3.2]{deligne1971hodge}. The notebook example is explicit; the displayed graded pieces make its intended cohomological convention clear.}
\GSfield{Present status}{As of November 2026, the corrected formulation is an established result; no unresolved claim is inferred from this occurrence.}
\GSfield{Assessment}{literal statement: correct with missing hypotheses; publication match: exact; verification: read the result; development: none; confidence: high}
\end{GSitem}
\begin{GSitem}{1986-669.5}{E}{Relative de Rham representatives and residues}
\GSfield{Kind; source}{method; \GSpage{P001264-P001265}}
\GSfield{The notes}{Differentials on the punctured curve are taken modulo differentials of functions with equal values at \(y_1,y_2\); residues at \(x_1,x_2\) describe the upper weight quotient.}
\GSfield{Statement}{The relative de Rham complex for \((U,Y)\), \(U=C\setminus\{x_1,x_2\}\), \(Y=\{y_1,y_2\}\), computes the relative cohomology. The condition \(f(y_1)=f(y_2)\) removes the endpoint difference of \(df\). The residue theorem leaves one independent residue at the two punctures, giving the one-dimensional \(\mathbb Q(-1)\) weight quotient; the endpoint difference gives the weight-zero quotient.}
\GSfield{Reconstruction}{The relative complex, rather than an unrestricted quotient of all smooth differentials, is intended.}
\GSfield{Fate}{The de Rham and weight-filtration construction is covered by \citet[\S3.2]{deligne1971hodge}. The equal-endpoint condition is the distinctive relative feature recorded here.}
\GSfield{Present status}{As of November 2026, the corrected formulation is an established result; no unresolved claim is inferred from this occurrence.}
\GSfield{Assessment}{literal statement: correct with missing hypotheses; publication match: exact; verification: read the result; development: none; confidence: high}
\end{GSitem}
\begin{GSitem}{1986-669.6}{A}{Green pairings as mixed-Hodge extension data}
\GSfield{Kind; source}{claim; \GSpage{P001264-P001265}}
\GSfield{The notes}{The lowering operator for the relative-curve example is related to an alternating sum of Green functions \(G(x_i,y_j)\).}
\GSfield{Statement}{For disjoint degree-zero divisors \(D=x_1-x_2\) and \(E=y_1-y_2\), a normalized Green kernel gives the real pairing \(G(D,E)=G(x_1,y_1)-G(x_1,y_2)-G(x_2,y_1)+G(x_2,y_2)\). This is the real height-type invariant of the corresponding biextension mixed Hodge structure. Its equality to a specified matrix coefficient of \(N\) requires choices of Tate generator, orientation, and Green normalization; those choices are not recoverable from the brief notes.}
\GSfield{Reconstruction}{The divisors and alternating pairing are unambiguous; no exact scalar multiplier is attributed to the speaker.}
\GSfield{Fate}{The identification belongs to Beilinson\textquotesingle s height-pairing theory, published in \citet{beilinson1987height}. This is a partial match: the precise normalization of the displayed notebook coefficient has not been verified in a primary formula.}
\GSfield{Present status}{As of November 2026, the height/biextension interpretation is established, while the notebook\textquotesingle s normalization remains an editorial uncertainty, not an open mathematical problem.}
\GSfield{Assessment}{literal statement: correct with missing hypotheses; publication match: partial; verification: indirect; development: none; confidence: medium}
\end{GSitem}

\GStalk{670}{1987-03-03}{A. Yu. Morozov}{P001263-P001269}{String amplitudes, Green functions, and the superstring measure}
The continuation headed Morozov starts midway through P001265 and runs through P001267. It connects the string description to gauge fields, discusses explicit bosonic measures and theta functions, and asks how to obtain the superstring measure by eliminating odd variables. The catalogue range includes neighboring contributions; the items below use the actual headings. The manuscript belongs to the \(3/\mathrm{III}\text{-}86\) session.

\begin{GSitem}{1986-670.1}{F}{Ultraviolet motivation and numerical scale estimates}
\GSfield{Kind; source}{remark; \GSpage{P001265}}
\GSfield{The notes}{The notes contrast a string length near \(10^{-33}\,\mathrm{cm}\) with particle scales and describe other theories as inconsistent at short distances.}
\GSfield{Statement}{This is a physical motivation, not a mathematical theorem. A string length is \(\ell_s=\sqrt{\alpha^{\prime}}\) in a stated unit convention; its relation to a gravitational scale depends on coupling and compactification. No universal assertion that every nonstring quantum field theory is ultraviolet inconsistent follows. The written proton size near \(10^{-35}\,\mathrm{cm}\) also conflicts with the scale comparison recorded on P001269.}
\GSfield{Fate}{The critical-string framework is identifiable from \citet[Introduction]{grossharveymartinecrohm1985heterotic}, but that paper does not establish the universal assertion. The specific estimate is retained as notebook error rather than promoted to a research problem.}
\GSfield{Present status}{As of November 2026, this passage remains a historically identifiable motivation with an erroneous numerical entry, not a definite unsolved mathematical claim.}
\GSfield{Assessment}{literal statement: false as written; publication match: thematic only; verification: read the result; development: none; confidence: high}
\end{GSitem}
\begin{GSitem}{1986-670.2}{E}{Ward identities from integrated longitudinal vertices}
\GSfield{Kind; source}{method; \GSpage{P001265-P001266}}
\GSfield{The notes}{A polarization proportional to momentum turns a vector insertion into a total derivative, and the massless string excitation is connected with Yang--Mills theory.}
\GSfield{Statement}{For an on-shell open-string vector vertex on a worldsheet boundary, replacing the polarization by its momentum produces a boundary derivative of the exponential vertex. Decoupling in amplitudes requires the usual BRST or physical-state conditions and control of contact and boundary contributions. Chan--Paton labels supply nonabelian gauge factors. Closed-string massless tensor vertices give a different sector, including gravity.}
\GSfield{Reconstruction}{The boundary, on-shell, and decoupling hypotheses delimit the total-derivative argument.}
\GSfield{Fate}{The distinction between open-string Chan--Paton gauge symmetry and closed-string constructions is already explicit in \citet[Introduction]{grossharveymartinecrohm1985heterotic}; the BRST field-theoretic implementation appears in \citet{witten1986noncommutative}. Both precede the catalogue date.}
\GSfield{Present status}{As of November 2026, the corrected formulation is an established result; no unresolved claim is inferred from this occurrence.}
\GSfield{Assessment}{literal statement: correct with missing hypotheses; publication match: partial; verification: read the result; development: none; confidence: high}
\end{GSitem}
\begin{GSitem}{1986-670.3}{E}{The scalar Green function through the prime form}
\GSfield{Kind; source}{construction; \GSpage{P001266}}
\GSfield{The notes}{The Green function is expressed using a theta function and an Abelian-integral correction involving \((\operatorname{Im}T)^{-1}\). Some quadratic factors are suppressed.}
\GSfield{Statement}{For a marked compact Riemann surface with normalized holomorphic one-forms \(\omega\), period matrix \(\Omega\), and prime form \(E\), the standard scalar Green kernel has the local expression \(-\log|E(z,w)|^2+2\pi\,v^{\mathsf T}(\operatorname{Im}\Omega)^{-1}v\), where \(v=\operatorname{Im}\int_w^z\omega\), up to the metric-dependent normalization needed to fix its additive zero mode. The prime form contains half-differential factors; a theta function alone is not a globally defined scalar replacement.}
\GSfield{Reconstruction}{The complete quadratic form and prime form are editorial reconstructions; no unrecorded normalization is assigned to Morozov.}
\GSfield{Fate}{This construction underlies the vertex-insertion expressions in \citet[\S2]{belavinknizhnik1986complex}. The notebook is a recognizable shorthand, but the missing second vector and half-differentials prevent literal use.}
\GSfield{Present status}{As of November 2026, the corrected formulation is an established result; no unresolved claim is inferred from this occurrence.}
\GSfield{Assessment}{literal statement: correct with missing hypotheses; publication match: partial; verification: read the result; development: none; confidence: high}
\end{GSitem}
\begin{GSitem}{1986-670.4}{E}{Explicit low-genus bosonic measures and Mumford forms}
\GSfield{Kind; source}{claim; \GSpage{P001266-P001267}}
\GSfield{The notes}{The notes list explicit expressions in low genus, including genus four, and mention Manin\textquotesingle s all-genus description.}
\GSfield{Statement}{The critical bosonic measure can be expressed as the norm square of a Mumford form with a period-matrix factor. Explicit theta-constant formulas require distinguishing the full Siegel space from its Jacobian locus and accounting for the genus-dependent modular form and its divisor. A low-genus formula is not a claim that the measure is undefined in higher genus.}
\GSfield{Reconstruction}{The Jacobian-locus and normalization qualifications are supplied. The explicit-formula identification is supported by the authored abstract and bibliographic metadata; the complete formula paper was not read.}
\GSfield{Fate}{\citet{morozov1987explicit} gives explicit one-, two-, three-, and four-loop bosonic vacuum amplitudes; \citet{beilinsonmanin1986mumford} supplies the determinant-line interpretation. The former appeared in 1987, before the catalogue session but after the 1986 manuscript heading; no priority conclusion is drawn from the catalogue alone.}
\GSfield{Present status}{As of November 2026, the corrected formulation is an established result; no unresolved claim is inferred from this occurrence.}
\GSfield{Assessment}{literal statement: correct with missing hypotheses; publication match: exact; verification: indirect; development: none; confidence: high}
\end{GSitem}
\begin{GSitem}{1986-670.5}{B}{Variations of regularized determinants}
\GSfield{Kind; source}{method; \GSpage{P001267}}
\GSfield{The notes}{The proposed calculation varies a determinant using the Green function at coincident points.}
\GSfield{Statement}{For an admissible smooth family of invertible elliptic operators, the variation of the logarithmic zeta determinant is a regularized trace of \(A^{-1}\delta A\), with local correction terms determined by the regularization convention. The unregularized coincident Green kernel is singular and cannot be substituted into an ordinary trace. Zero modes require a reduced determinant and a separate kernel metric.}
\GSfield{Reconstruction}{Admissibility, kernel treatment, and regularization are restored; the incomplete displayed theta expression is not certified.}
\GSfield{Fate}{\citet[\S1]{kontsevichvishik1994geometry} formulates determinant variations in a regularized-trace setting. The notebook\textquotesingle s particular theta-function expression is too incomplete for an exact formula match.}
\GSfield{Present status}{As of November 2026, the corrected formulation is an established result; no unresolved claim is inferred from this occurrence.}
\GSfield{Assessment}{literal statement: correct with missing hypotheses; publication match: partial; verification: read the result; development: none; confidence: medium}
\end{GSitem}
\begin{GSitem}{1986-670.6}{B}{Eliminating odd moduli in the superstring measure}
\GSfield{Kind; source}{question; \GSpage{P001267}}
\GSfield{The notes}{The notes ask how to integrate over the odd variables of supermoduli and obtain an ordinary-moduli measure; the resulting modular ratio is expected to have weight eight.}
\GSfield{Statement}{For a smooth super Riemann surface of fixed genus and spin parity, determine the pushforward of the RNS measure along odd moduli wherever a specified projection exists, and its transformation and boundary behavior. A global holomorphic projection from supermoduli to reduced spin moduli is an additional geometric assertion, not an automatic meaning of Berezin integration. Local odd integration remains meaningful after a splitting is chosen.}
\GSfield{Reconstruction}{The precise global-projection obstruction is a later formulation of one assumption behind the notebook\textquotesingle s proposed calculation.}
\GSfield{Fate}{\citet[Theorem 1.1]{donagiwitten2015supermoduli}, first circulated in April 2013, proves nonprojectedness in genus at least five. Version 1 of the preprint of \citet[Theorem 1.1]{correadonaginoja2026supermoduli}, submitted 30 September 2026, states the nonvanishing of the primary obstruction on both parity components for every genus \(g\geq3\), algebraically and holomorphically. The latter theorem statement was read; its proof has not been independently checked here.}
\GSfield{Present status}{As of 12 November 2026, a global holomorphic projection is obstructed in the stated genera, including genus three according to the September 2026 preprint. This does not assert that local RNS measures or physical amplitudes cease to exist.}
\GSfield{Assessment}{literal statement: correct with missing hypotheses; publication match: partial; verification: read the result; development: none; confidence: high; problem \hyperref[prob:GP-1986-01]{GP-1986-01}}
\end{GSitem}
\begin{GSitem}{1986-670.7}{C}{Determining higher-genus superstring measures}
\GSfield{Kind; source}{problem; \GSpage{P001267}}
\GSfield{The notes}{The computational question remains: find the superstring measure after the odd-variable calculation, rather than only its expected weight.}
\GSfield{Statement}{In the genus-three even-spin case, on the nonhyperelliptic locus where a chosen pushforward is defined, determine the exact coefficients of the RNS top component in the three-function ansatz and prove the claimed normalization and degeneration behavior. This is a modern bounded subproblem of the much broader notebook request, not its verbatim formulation.}
\GSfield{Reconstruction}{The three coefficient problem is a later precise specialization, clearly separated from the all-genus notebook request.}
\GSfield{Fate}{\citet[\S1.7 and \S4.9]{duninbarkowskifedorovsleptsov2025rns}, circulated in May 2025 and published in October 2025, proves membership in a three-dimensional span but labels the coefficient values conjectural. Section 4.9.2 explicitly notes that the numerical computations lack rigorous error bounds. The hyperelliptic polar behavior also prevents identification with an everywhere holomorphic ansatz.}
\GSfield{Present status}{As of 12 November 2026, the cited genus-three coefficient values remain conjectural in the opened publication; no claim of a complete higher-genus solution is made.}
\GSfield{Assessment}{literal statement: correct with missing hypotheses; publication match: partial; verification: read the result; development: none; confidence: high; problem \hyperref[prob:GP-1986-01]{GP-1986-01}}
\GSfield{Worth developing}{Worth developing in a separate run: the three coefficient identities and their normalization are concrete, but a complete RNS derivation and boundary analysis exceed the present bounded verification.}
\end{GSitem}
\begin{GSitem}{1986-670.8}{E}{Weight-eight forms and the genus-four Schottky relation}
\GSfield{Kind; source}{theorem; \GSpage{P001267}}
\GSfield{The notes}{The notes say that cusp forms of weight eight first occur in genus four and identify a relation relevant to the string measure.}
\GSfield{Statement}{For scalar Siegel modular forms of weight eight for the full symplectic modular group, the two rank-sixteen even unimodular theta series agree in genera at most three. Their genus-four difference is a multiple of the Schottky form and vanishes on the Jacobian locus. These full-group statements must not be transferred to level subgroups or to spin-dependent modular forms, which have additional forms.}
\GSfield{Reconstruction}{The full modular group and Jacobian-locus qualifiers are added.}
\GSfield{Fate}{\citet[\S2]{grushevskysalvatimanni2011superstring} records these classical identities and their failure to force the analogous genus-five cosmological constant to vanish. Its preprint first appeared in September 2008; the identities it uses are older than the notebook.}
\GSfield{Present status}{As of November 2026, the corrected formulation is an established result; no unresolved claim is inferred from this occurrence.}
\GSfield{Assessment}{literal statement: correct with missing hypotheses; publication match: partial; verification: read the result; development: none; confidence: high}
\end{GSitem}

\GStalk{671}{1987-03-10}{A. Yu. Morozov}{P001267-P001274}{Compactification, lattice theta series, and heterotic gauge symmetry}
The contribution begins on P001268 and continues through P001273. It combines gauge fields, the string oscillator spectrum, compactification on a sixteen-dimensional torus, modular lattice sums, and root vertex operators. Several formulas mix ordinary bosonic compactification with the chiral heterotic construction. The manuscript date is \(10/\mathrm{III}\text{-}86\), whereas the catalogue says 1987.

\begin{GSitem}{1986-671.1}{E}{Yang--Mills curvature and gauge transformations}
\GSfield{Kind; source}{definition; \GSpage{P001268}}
\GSfield{The notes}{The gauge potential \(A_\mu^a\), its curvature, and infinitesimal transformations are written before the string discussion.}
\GSfield{Statement}{For a Lie-algebra-valued connection one-form \(A\), \(F=dA+A\wedge A\), or \(F_{\mu\nu}=\partial_\mu A_\nu-\partial_\nu A_\mu+[A_\mu,A_\nu]\), and an infinitesimal gauge transformation is \(\delta_\epsilon A=d\epsilon+[A,\epsilon]\). One must retain both antisymmetric curvature indices.}
\GSfield{Reconstruction}{The matrix-valued convention avoids an unrecorded choice of signs for structure constants.}
\GSfield{Fate}{The gauge-field conventions are standard background for the low-energy comparison in \citet{grossharveymartinecrohm1985heterotic}; they are not presented as a new theorem. The missing curvature index is a notebook notational slip.}
\GSfield{Present status}{As of November 2026, the corrected formulation is an established result; no unresolved claim is inferred from this occurrence.}
\GSfield{Assessment}{literal statement: correct with missing hypotheses; publication match: partial; verification: read the result; development: none; confidence: high}
\end{GSitem}
\begin{GSitem}{1986-671.2}{E}{Massless vector states and allowed Chan--Paton groups}
\GSfield{Kind; source}{claim; \GSpage{P001268-P001269}}
\GSfield{The notes}{The first excited string state is interpreted as a vector, and a simple nonabelian gauge group is introduced.}
\GSfield{Statement}{The open bosonic string has a massless vector at oscillator number one in the critical theory. Nonabelian Chan--Paton factors in the usual oriented or unoriented perturbative construction give unitary, orthogonal, or symplectic families subject to consistency conditions. A generic closed bosonic string on flat \(\mathbb R^{26}\) instead has a massless tensor sector and does not produce this vector by the same construction. Arbitrary simple groups are not obtained merely by adding conventional endpoint labels.}
\GSfield{Reconstruction}{The sector and gauge-group qualifications are necessary repairs.}
\GSfield{Fate}{\citet[Introduction]{grossharveymartinecrohm1985heterotic} explicitly motivates the heterotic construction by the limited Chan--Paton possibilities, including the exceptional-group issue. The notebook merges distinct constructions.}
\GSfield{Present status}{As of November 2026, the corrected formulation is an established result; no unresolved claim is inferred from this occurrence.}
\GSfield{Assessment}{literal statement: false as written; publication match: exact; verification: read the result; development: none; confidence: high}
\end{GSitem}
\begin{GSitem}{1986-671.3}{E}{The sixteen-dimensional lattice in the heterotic string}
\GSfield{Kind; source}{construction; \GSpage{P001269-P001270}}
\GSfield{The notes}{The map target is changed from \(\mathbb R^{26}\) to \(T^{16}\times\mathbb R^{10}\), and supersymmetry and the groups \(E_8\times E_8\) and \(SO(32)\) are discussed.}
\GSfield{Statement}{The ten-dimensional heterotic string combines a chiral bosonic matter sector of central charge 26 with a supersymmetric matter sector of central charge 15; the respective ghost systems cancel these matter anomalies. The sixteen extra chiral bosons have momenta in a positive-definite even unimodular rank-sixteen lattice. The two lattices are \(E_8\oplus E_8\) and \(D_{16}^{+}\), giving gauge Lie algebras \(\mathfrak e_8\oplus\mathfrak e_8\) and \(\mathfrak{so}_{32}\). Ordinary nonchiral compactification of the full bosonic theory does not itself supply spacetime supersymmetry.}
\GSfield{Reconstruction}{The chiral construction is made explicit. The global group associated with \(D_{16}^{+}\) is a suitable \(\operatorname{Spin}(32)/\mathbb Z_2\), so the notebook\textquotesingle s \(SO(32)\) is used only at Lie-algebra level.}
\GSfield{Fate}{This is the construction of \citet{grossharveymartinecrohm1985heterotic}, published on 11 February 1985 after receipt in November 1984. It predates either plausible date of this talk.}
\GSfield{Present status}{As of November 2026, the corrected formulation is an established result; no unresolved claim is inferred from this occurrence.}
\GSfield{Assessment}{literal statement: correct with missing hypotheses; publication match: exact; verification: read the result; development: none; confidence: high}
\end{GSitem}
\begin{GSitem}{1986-671.4}{F}{A dynamical mechanism for compactification}
\GSfield{Kind; source}{question; \GSpage{P001269}; I. M. Gelfand}
\GSfield{The notes}{Gelfand asks about obtaining compactification dynamically rather than putting in a torus.}
\GSfield{Statement}{The passage asks for a physical mechanism selecting a lower-dimensional background. It specifies no action including its quantum corrections, no admissible class of backgrounds, and no stability criterion. Those choices change the question, so the notes do not determine a single proposition to prove.}
\GSfield{Fate}{The toroidal construction in \citet{grossharveymartinecrohm1985heterotic} supplies a background, not the general selection mechanism requested. Searches for heterotic compactification and the original construction identify the context but do not justify equating this broad question with a named unsolved conjecture.}
\GSfield{Present status}{As of November 2026, the notebook question remains under-specified; no new problem-index identifier or class D is assigned.}
\GSfield{Assessment}{literal statement: not determinable; publication match: thematic only; verification: read the result; development: none; confidence: high}
\end{GSitem}
\begin{GSitem}{1986-671.5}{E}{Classical winding maps on a worldsheet torus}
\GSfield{Kind; source}{construction; \GSpage{P001270}}
\GSfield{The notes}{The boundary conditions are \(x(z+1)-x(z)=\Lambda_2\) and \(x(z+\tau)-x(z)=\Lambda_1\), but equation (18) uses two plus signs in the affine part.}
\GSfield{Statement}{For real target-space winding vectors and \(\operatorname{Im}\tau>0\), the harmonic affine representative is \[x_{\mathrm{cl}}(z)=\frac{(\Lambda_1-\overline\tau\Lambda_2)z-(\Lambda_1-\tau\Lambda_2)\overline z}{2i\operatorname{Im}\tau}.\] The fluctuation is doubly periodic. This formula satisfies the two stated shifts and is real valued.}
\GSfield{Reconstruction}{\oursComputed\ The displayed corrected affine representative is an editorial algebraic reconstruction from the recorded boundary conditions; the check is logged, not claimed as new research.}
\GSfield{Fate}{The winding-sector decomposition is the compact-boson calculation used in \citet{grossharveymartinecrohm1985heterotic}. The scan confirms the two plus signs in the notebook; direct substitution into its boundary conditions detects the sign error.}
\GSfield{Present status}{As of November 2026, the corrected formulation is an established result; no unresolved claim is inferred from this occurrence.}
\GSfield{Assessment}{literal statement: false as written; publication match: partial; verification: read the result; development: annotated; confidence: high}
\end{GSitem}
\begin{GSitem}{1986-671.6}{E}{Winding sums, Gaussian factors, and Poisson summation}
\GSfield{Kind; source}{method; \GSpage{P001270-P001271}}
\GSfield{The notes}{After Gaussian integration the partition function is a double sum over \(\Lambda_1,\Lambda_2\), with an exponential of a quadratic expression divided by \(\operatorname{Im}\tau\). It is then factorized into theta functions.}
\GSfield{Statement}{For a compact free boson, the classical Euclidean action is a positive multiple of \(\|\Lambda_1-\tau\Lambda_2\|^2/\operatorname{Im}\tau\), so the convergent winding factor has a negative real exponent. Poisson summation rewrites the winding sum as a momentum lattice sum. Holomorphic factorization into a specified finite collection of theta functions needs additional lattice and rationality data; it is not automatic for every real torus metric.}
\GSfield{Reconstruction}{Positivity of the Euclidean action and the conditions for chiral factorization are restored.}
\GSfield{Fate}{The compactification and modular lattice calculation is an essential step of \citet{grossharveymartinecrohm1985heterotic}. The notebook suppresses sign and normalization conventions and also mixes a nonchiral absolute-square expression with a chiral lattice sum.}
\GSfield{Present status}{As of November 2026, the corrected formulation is an established result; no unresolved claim is inferred from this occurrence.}
\GSfield{Assessment}{literal statement: correct with missing hypotheses; publication match: partial; verification: read the result; development: none; confidence: high}
\end{GSitem}
\begin{GSitem}{1986-671.7}{E}{What modularity requires of a positive lattice}
\GSfield{Kind; source}{question; \GSpage{P001271-P001272}; I. M. Gelfand}
\GSfield{The notes}{Gelfand asks whether the even self-dual lattices are the only ones giving the needed modular behavior.}
\GSfield{Statement}{For a positive-definite lattice \(L\) of rank \(r\), the theta series \(\Theta_L(\tau)=\sum_{x\in L}e^{\pi i\tau(x,x)}\) of an even unimodular lattice transforms as a scalar modular form of weight \(r/2\) for the full modular group. If one instead allows a congruence subgroup, a multiplier, or a vector of theta series of discriminant cosets, even unimodularity is not a necessary condition. Thus the word “modular” must include the transformation group and representation before an “only” statement is possible.}
\GSfield{Reconstruction}{The qualified full-group formulation is supplied; the unqualified converse is deliberately not asserted.}
\GSfield{Fate}{The full-group condition is the one imposed in \citet{grossharveymartinecrohm1985heterotic}. Its role is separate from the broader modular theory of nonunimodular lattices. This is an already known modularity criterion, not evidence for an unsolved classification.}
\GSfield{Present status}{As of November 2026, the scalar full-group criterion is classical. The notebook\textquotesingle s unqualified question has different answers for subgroups and vector-valued forms.}
\GSfield{Assessment}{literal statement: not determinable; publication match: partial; verification: read the result; development: none; confidence: high; problem \hyperref[prob:GP-1986-02]{GP-1986-02}}
\end{GSitem}
\begin{GSitem}{1986-671.8}{E}{Two rank-sixteen theta series and their higher-genus difference}
\GSfield{Kind; source}{theorem; \GSpage{P001272}}
\GSfield{The notes}{The two lattice expressions are compared, with an identity involving powers of theta constants and an absolute square.}
\GSfield{Statement}{The lattices \(E_8\oplus E_8\) and \(D_{16}^{+}\) have the same genus-one theta series, the weight-eight form \(E_4^2\). Their Siegel theta series agree through genus three; in genus four their difference is the Schottky form up to scale and vanishes on Jacobians. The holomorphic theta identity is distinct from taking the absolute square in a nonchiral partition function, and the powers of two depend on the theta normalization.}
\GSfield{Reconstruction}{The notebook\textquotesingle s absolute-square line is separated from the holomorphic identity; its numerical normalization is not copied as a theorem.}
\GSfield{Fate}{\citet[\S2]{grushevskysalvatimanni2011superstring} supplies the exact comparison and explains the subsequent genus-five distinction. These low-genus identities are classical, despite the later verification source.}
\GSfield{Present status}{As of November 2026, the corrected formulation is an established result; no unresolved claim is inferred from this occurrence.}
\GSfield{Assessment}{literal statement: correct with missing hypotheses; publication match: exact; verification: read the result; development: none; confidence: high}
\end{GSitem}
\begin{GSitem}{1986-671.9}{E}{Root vertex operators and the affine Lie algebra}
\GSfield{Kind; source}{construction; \GSpage{P001272-P001273}}
\GSfield{The notes}{Exponential operators indexed by lattice roots are said to have commutator given by the operator for the sum of two roots, and zero otherwise.}
\GSfield{Statement}{For an even root lattice, the lattice vertex construction uses cocycles and a Heisenberg algebra. Root operators for \(\alpha+\beta\) a root bracket to a signed root operator. If \(\beta=-\alpha\), their bracket produces Cartan and central terms rather than zero. The oscillator modes realize the corresponding level-one affine Lie algebra; these cocycle and opposite-root cases are indispensable.}
\GSfield{Reconstruction}{The Heisenberg sector, cocycle, and opposite-root term restore the standard construction.}
\GSfield{Fate}{This is the basic-representation construction identified with \citet{frenkelkac1980basic}. Its publisher metadata was opened, but the full original construction was not read in this run; the exact identification is therefore marked indirect. The displayed “zero otherwise” rule is false in the opposite-root case.}
\GSfield{Present status}{As of November 2026, the corrected formulation is an established result; no unresolved claim is inferred from this occurrence.}
\GSfield{Assessment}{literal statement: false as written; publication match: exact; verification: indirect; development: none; confidence: high}
\end{GSitem}

\GStalk{672}{1987-03-24}{D. B. Fuchs}{P001272-P001278}{Topology of real subspace arrangements and complexifiability}
The contribution starts on P001274 and ends on P001276. Fuchs explains the Goresky--MacPherson arrangement theorem through order complexes and stratified Morse theory, including a small example and the special case of \(k\)-arrangements. Gelfand asks about realization, common complex structures, and useful cohomology bases. The manuscript date is \(24.\mathrm{III}.86\), not the catalogue year 1987.

\begin{GSitem}{1986-672.1}{E}{Intersection posets and complements of affine arrangements}
\GSfield{Kind; source}{definition; \GSpage{P001274}}
\GSfield{The notes}{A finite collection of affine subspaces \(L_i\subset\mathbb R^n\) determines its complement and the partially ordered set of intersections.}
\GSfield{Statement}{For a finite arrangement \(\mathcal A\) of proper affine subspaces, put \(M(\mathcal A)=\mathbb R^n\setminus\bigcup\mathcal A\). Its intersection poset consists of nonempty intersections, including the empty intersection \(\widehat0=\mathbb R^n\), ordered by reverse inclusion and labeled by dimension. Parallel affine subspaces can have empty intersection, so the resulting poset need not be the lattice of a central arrangement.}
\GSfield{Reconstruction}{Reverse inclusion, the ambient element, and dimension labels are fixed to avoid sign and interval ambiguities.}
\GSfield{Fate}{The convention is explicit in \citet[\S2]{greenezaslavsky1983whitney} for hyperplanes and underlies the general arrangement treatment of \citet[Part III]{goreskymacpherson1988stratified}.}
\GSfield{Present status}{As of November 2026, the corrected formulation is an established result; no unresolved claim is inferred from this occurrence.}
\GSfield{Assessment}{literal statement: correct with missing hypotheses; publication match: exact; verification: read the result; development: none; confidence: high}
\end{GSitem}
\begin{GSitem}{1986-672.2}{E}{Real hyperplane chambers and the Mobius count}
\GSfield{Kind; source}{theorem; \GSpage{P001274}}
\GSfield{The notes}{For hyperplanes the complement has only zeroth homology, and the number of regions is given combinatorially.}
\GSfield{Statement}{For a finite real affine hyperplane arrangement, every chamber is an open convex set and therefore contractible. If \(\chi_{\mathcal A}(t)=\sum_{X\in L}\mu(\widehat0,X)t^{\dim X}\), then the number of chambers is \((-1)^n\chi_{\mathcal A}(-1)=\sum_{X\in L}|\mu(\widehat0,X)|\). Thus \(H_i(M;\mathbb Z)=0\) for \(i>0\), and \(H_0\) is free of this rank.}
\GSfield{Fate}{\citet{zaslavsky1975facing} first circulated as a 1974 MIT thesis and was published in 1975. The face-count formula is stated again as Theorem A on p.~100 of the opened primary paper \citet{greenezaslavsky1983whitney}.}
\GSfield{Present status}{As of November 2026, the corrected formulation is an established result; no unresolved claim is inferred from this occurrence.}
\GSfield{Assessment}{literal statement: correct; publication match: exact; verification: read the result; development: none; confidence: high}
\end{GSitem}
\begin{GSitem}{1986-672.3}{E}{The Goresky--MacPherson additive formula}
\GSfield{Kind; source}{theorem; \GSpage{P001274-P001275}}
\GSfield{The notes}{The homology of the complement is decomposed into contributions from intersections and the cohomology of their order-complex intervals.}
\GSfield{Statement}{For a finite real affine subspace arrangement and its dimension-labeled intersection poset, the additive formula may be written \[\widetilde H^i(M;\mathbb Z)\simeq\bigoplus_{X\ne\widehat0}\widetilde H_{\operatorname{codim}_{\mathbb R}X-i-2}\bigl(\Delta(\widehat0,X);\mathbb Z\bigr).\] Here \(\Delta(\widehat0,X)\) is the order complex of the open interval. This is a statement about graded groups; it does not identify the ring or assert a canonical splitting.}
\GSfield{Reconstruction}{The cohomological version fixes integral duality conventions. The original 1988 book was bibliographically verified; the full 25 MB file could not be rendered by the web reader, so this entry does not claim a direct reading of that proof.}
\GSfield{Fate}{The full theorem appears in \citet[Part III]{goreskymacpherson1988stratified}; the stratified-Morse method was announced in the 1983 paper \citet{goreskymacpherson1983stratified}. Its group-versus-ring distinction is verified directly in \citet[Introduction and Theorem 1.2]{ziegler1993difference}, circulated in 1991--1992 and published in 1993.}
\GSfield{Present status}{As of November 2026, the corrected formulation is an established result; no unresolved claim is inferred from this occurrence.}
\GSfield{Assessment}{literal statement: correct with missing hypotheses; publication match: exact; verification: indirect; development: none; confidence: high}
\end{GSitem}
\begin{GSitem}{1986-672.4}{E}{Order complexes, reduced degree minus one, and the small example}
\GSfield{Kind; source}{construction; \GSpage{P001275}}
\GSfield{The notes}{Chains of a poset form a simplicial complex; an empty interval contributes in degree \(-1\). The example is three coordinate planes and a generic line in \(\mathbb R^3\).}
\GSfield{Statement}{The order complex has one simplex for each nonempty finite chain. Its augmented reduced homology convention is \(\widetilde H_{-1}(\varnothing;\mathbb Z)=\mathbb Z\), whereas ordinary negative-degree cohomology of a pair is not being used. For the three coordinate planes together with a line through the origin transverse to each plane, the complement has \(b_0=8\), \(b_1=2\), and no higher Betti numbers: the two chambers traversed by the line each lose an open radial ray.}
\GSfield{Reconstruction}{Transversality means that the line is contained in none of the coordinate planes; reduced and ordinary conventions are separated.}
\GSfield{Fate}{The empty-interval convention is required by the arrangement formula in \citet[Part III]{goreskymacpherson1988stratified}. The example is explicitly present in the notebook; its stated ranks are consistent with that formula.}
\GSfield{Present status}{As of November 2026, the corrected formulation is an established result; no unresolved claim is inferred from this occurrence.}
\GSfield{Assessment}{literal statement: correct with missing hypotheses; publication match: partial; verification: indirect; development: none; confidence: high}
\end{GSitem}
\begin{GSitem}{1986-672.5}{F}{Which ordered sets occur as intersection data?}
\GSfield{Kind; source}{question; \GSpage{P001275}; I. M. Gelfand}
\GSfield{The notes}{Gelfand asks whether an arbitrary partially ordered set can occur.}
\GSfield{Statement}{An intersection poset carries more data than an arbitrary order: an ambient element, intersection operations where nonempty, and codimensions constrained by linear algebra. The notes do not specify whether affine or central arrangements, dimension labels, fixed ambient dimension, or fixed codimensions are required. These are different realization problems, and the question cannot be assigned one exact classification theorem as written.}
\GSfield{Fate}{The precise distinction between geometric lattices of hyperplanes and general intersection posets is documented in \citet[\S2]{greenezaslavsky1983whitney} and \citet[Introduction]{ziegler1993difference}. The notebook formulation lacks the constraints necessary for a sharper match.}
\GSfield{Present status}{As of November 2026, no definite new open problem is extracted from this unqualified realization question.}
\GSfield{Assessment}{literal statement: not determinable; publication match: thematic only; verification: read the result; development: none; confidence: high}
\end{GSitem}
\begin{GSitem}{1986-672.6}{E}{Stratified Morse attachment and geometric splitting}
\GSfield{Kind; source}{method; \GSpage{P001275-P001276}}
\GSfield{The notes}{The squared distance to a generic point is used as a Morse function. The notes explicitly warn that relative homology alone does not suffice: a geometric splitting must be checked.}
\GSfield{Statement}{A generic distance function on an arrangement stratification changes sublevel topology at finitely many stratified critical values. Local Morse data determine relative homology, but a direct-sum decomposition of absolute homology needs vanishing of connecting maps and splitting of the resulting extensions. A generic long exact sequence of a pair does not supply this automatically.}
\GSfield{Reconstruction}{The modern statement distinguishes local Morse data from the additional splitting argument.}
\GSfield{Fate}{The method is developed in \citet{goreskymacpherson1983stratified} and \citet[Part III]{goreskymacpherson1988stratified}. The notebook\textquotesingle s explicit caveat is mathematically essential and should be preserved.}
\GSfield{Present status}{As of November 2026, the corrected formulation is an established result; no unresolved claim is inferred from this occurrence.}
\GSfield{Assessment}{literal statement: correct; publication match: exact; verification: read the result; development: none; confidence: high}
\end{GSitem}
\begin{GSitem}{1986-672.7}{E}{Degree concentration for k-arrangements}
\GSfield{Kind; source}{theorem; \GSpage{P001276}}
\GSfield{The notes}{Each \(L_i\) has codimension \(k\), every nonempty intersection has codimension divisible by \(k\), and a contribution of rank \(m\) is asserted to occur in degree \(m(k-1)\), with rank \(|\mu|\).}
\GSfield{Statement}{For a central real \(k\)-arrangement with the stated intersection condition, the intersection lattice is geometric with rank \(r(X)=\operatorname{codim}X/k\). Its open intervals have reduced homology concentrated in degree \(r(X)-2\), free of rank \(|\mu(\widehat0,X)|\). The arrangement formula therefore places the contribution of \(X\) in complement-cohomology degree \(r(X)(k-1)\). For \(k=1\), all contributions belong to reduced degree zero, together with the ordinary base \(H^0\).}
\GSfield{Reconstruction}{The statement is restricted to central arrangements and its rank convention is stated. The notebook has a question mark over the equality.}
\GSfield{Fate}{The specialization is part of \citet[Part III]{goreskymacpherson1988stratified}; \citet[Introduction]{ziegler1993difference} verifies the geometric-lattice property and additive conclusion for \(k=2\). The question mark over the degree in the scan is resolved by the standard degree shift.}
\GSfield{Present status}{As of November 2026, the corrected formulation is an established result; no unresolved claim is inferred from this occurrence.}
\GSfield{Assessment}{literal statement: correct with missing hypotheses; publication match: exact; verification: indirect; development: none; confidence: high}
\end{GSitem}
\begin{GSitem}{1986-672.8}{B}{A common complex structure for a real codimension-two arrangement}
\GSfield{Kind; source}{problem; \GSpage{P001276}; I. M. Gelfand}
\GSfield{The notes}{After the warning that even-dimensional intersections need not give a complex arrangement, Gelfand asks: \ru{при каких усл. сущ. компл. структ., в к-рой все они станут комплексными?}}
\GSfield{Statement}{Given real codimension-two linear subspaces \(L_i\subset\mathbb R^{2d}\), determine whether there is a real linear endomorphism \(J\) with \(J^2=-I\) and \(J(L_i)=L_i\) for every \(i\). The preceding \(2\)-arrangement hypothesis further requires every intersection to have even codimension. Merely requiring codimension two for the individual subspaces is weaker. Linear complexification is also stronger than asking whether the complement is homeomorphic to that of some complex arrangement.}
\GSfield{Reconstruction}{The invariant-subspace equation fixes the meaning of “complex structure.” The notebook does not restrict its question to four transverse planes; that is a bounded published specialization.}
\GSfield{Fate}{\citet[\S2 and Theorem 1.2]{ziegler1993difference} gives explicit distinctions between real and complex arrangements (manuscript dated December 1991; arXiv February 1992; publication 1993). \citet[Theorems 1 and 3]{batzies2007planes} classifies transverse arrangements in \(\mathbb R^4\) by simultaneous conjugation and describes the four-plane complexifiable locus. Thus the four-plane discrimination selected for annotation here is already known, not a new open problem.}
\GSfield{Present status}{As of November 2026, the transverse \(\mathbb R^4\) four-plane case has an explicit published answer. The broad finite-family question admits a simultaneous-invariance formulation; no unproved assertion that it remains open is made.}
\GSfield{Assessment}{literal statement: correct with missing hypotheses; publication match: partial; verification: read the result; development: annotated; confidence: high; problem \hyperref[prob:GP-1986-03]{GP-1986-03}}
\end{GSitem}
\begin{GSitem}{1986-672.9}{F}{The Pappus diagram and an incompletely specified obstruction}
\GSfield{Kind; source}{remark; \GSpage{P001276}}
\GSfield{The notes}{A diagram headed “Goresky\textquotesingle s example” and “Pappus\textquotesingle theorem” is followed by the assertion that complex perturbations preserve the situation and real ones need not.}
\GSfield{Statement}{Pappus\textquotesingle incidence theorem holds over both \(\mathbb R\) and \(\mathbb C\), so the sentence cannot mean that ordinary real projective lines violate Pappus. It must refer to treating real codimension-two subspaces as candidates for complex hyperplanes, where the required common complex structure is additional data. The diagram alone does not specify which real subspaces, incidences, or perturbations are intended.}
\GSfield{Fate}{\citet{ziegler1993difference} gives precise later obstructions in the relevant setting, but does not identify this particular unlabeled diagram. No exact publication match or counterexample is claimed for the fragment.}
\GSfield{Present status}{As of November 2026, the illustrative intention is clear but the particular configuration remains unidentified.}
\GSfield{Assessment}{literal statement: not determinable; publication match: thematic only; verification: read the result; development: none; confidence: medium}
\end{GSitem}
\begin{GSitem}{1986-672.10}{B}{Finding useful cohomology generators and bases}
\GSfield{Kind; source}{question; \GSpage{P001276}; I. M. Gelfand}
\GSfield{The notes}{Gelfand asks how to find the \ru{правильный базис когомологий}.}
\GSfield{Statement}{For a complex hyperplane arrangement with an ordering of its hyperplanes, logarithmic one-forms generate the integral cohomology ring, with circuit-boundary relations; no-broken-circuit monomials give a basis depending on that order. For real \(2\)-arrangements, the analogous relation signs carry geometric information not determined by the intersection lattice. Without a criterion for “right,” no canonical basis is demanded by the notes.}
\GSfield{Reconstruction}{An ordering and the distinction between complex and real arrangements are supplied.}
\GSfield{Fate}{The complex presentation was published in \citet{orliksolomon1980complements}; its exact form and the signed real analogue are reproduced in \citet[Theorems 1.1 and 1.1\textquotesingle]{ziegler1993difference}. The later real result answers a definite part of the request, without manufacturing a canonical-basis conjecture.}
\GSfield{Present status}{As of November 2026, explicit generators, relations, and order-dependent bases are available in the stated setting; the notebook does not specify an additional unresolved canonicality requirement.}
\GSfield{Assessment}{literal statement: correct with missing hypotheses; publication match: partial; verification: read the result; development: none; confidence: high}
\end{GSitem}

\GStalk{673}{1987-03-24}{A. N. Varchenko}{P001276-P001281}{Asymptotic mixed Hodge structures and singularity spectra}
The actual contribution is P001277--P001280. It relates mixed Hodge structures to powers and logarithms in singular integrals, develops the Milnor-fibre setting, and describes the asymptotic filtration and semicontinuity of spectra. Gelfand asks for several functions and for pairs of quadrics. These are separated from the established one-function isolated-singularity theorem.

\begin{GSitem}{1986-673.1}{E}{Pure and mixed Hodge structures in families}
\GSfield{Kind; source}{definition; \GSpage{P001277}}
\GSfield{The notes}{Varchenko recalls pure Hodge decompositions, the weight and Hodge filtrations, and the holomorphic variation of the latter.}
\GSfield{Statement}{For a smooth projective family, the cohomology local system underlies a variation of pure Hodge structure: the Hodge filtration is holomorphic and satisfies Griffiths transversality. For singular or nonproper algebraic varieties the appropriate object is mixed. A locally constant underlying topological system does not make every weight or Hodge datum a topological invariant under arbitrary changes of algebraic structure.}
\GSfield{Reconstruction}{The smooth projective and algebraic hypotheses distinguish the pure and mixed cases.}
\GSfield{Fate}{The foundational structures precede the talk: \citet[\S2]{deligne1971hodge}. The variation statement is background for Varchenko\textquotesingle s asymptotic construction, not a claim first made here.}
\GSfield{Present status}{As of November 2026, the corrected formulation is an established result; no unresolved claim is inferred from this occurrence.}
\GSfield{Assessment}{literal statement: correct with missing hypotheses; publication match: partial; verification: read the result; development: none; confidence: high}
\end{GSitem}
\begin{GSitem}{1986-673.2}{E}{An elementary nodal open-curve mixed Hodge structure}
\GSfield{Kind; source}{construction; \GSpage{P001277}}
\GSfield{The notes}{Two points of \(\mathbb C^*\) are identified, yielding two independent one-dimensional cycles and a nontrivial weight filtration.}
\GSfield{Statement}{Let \(X\) be the algebraic curve obtained from \(\mathbb C^*\) by identifying two distinct points to an ordinary node. The normalization exact sequence gives a weight-zero graph class and a weight-two puncture class: \(\dim H^1(X;\mathbb Q)=2\), \(\operatorname{Gr}^W_0H^1\simeq\mathbb Q(0)\), and \(\operatorname{Gr}^W_2H^1\simeq\mathbb Q(-1)\). The filtration identifies which combination annihilates the normalization cycle; homology carries the dual convention.}
\GSfield{Reconstruction}{The algebraic nodal quotient and cohomological convention are made explicit.}
\GSfield{Fate}{This is an elementary instance of the functorial mixed Hodge construction of \citet{deligne1971hodge}. It illustrates the distinction between topological rank and extension data.}
\GSfield{Present status}{As of November 2026, the corrected formulation is an established result; no unresolved claim is inferred from this occurrence.}
\GSfield{Assessment}{literal statement: correct with missing hypotheses; publication match: partial; verification: read the result; development: none; confidence: high}
\end{GSitem}
\begin{GSitem}{1986-673.3}{E}{Meromorphic continuation of powers of a function}
\GSfield{Kind; source}{theorem; \GSpage{P001277-P001278}}
\GSfield{The notes}{Integrals containing \(f^\lambda\) are continued in \(\lambda\), and rational powers and logarithms are related to filtrations.}
\GSfield{Statement}{For a real polynomial \(f\), use a specified positive region and extend \(f^\lambda\) by zero, or use a specified absolute-value distribution. Initially defined for sufficiently large real part of \(\lambda\), it extends meromorphically as a distribution. A Bernstein functional equation \(P(\lambda)f^{\lambda+1}=b(\lambda)f^\lambda\) gives continuation; resolution controls rational pole locations. A complex-valued \(f^\lambda\) without a chosen branch is not a single-valued distribution.}
\GSfield{Reconstruction}{The sign/branch convention is essential and is supplied.}
\GSfield{Fate}{\citet[Theorems 1 and 1\textquotesingle]{bernstein1972continuation} gives the algebraic continuation argument (received June 1972, published in the 1972 volume). Its proof deliberately avoids resolution, while the stronger pole information comes from the resolution approach cited there.}
\GSfield{Present status}{As of November 2026, the corrected formulation is an established result; no unresolved claim is inferred from this occurrence.}
\GSfield{Assessment}{literal statement: correct with missing hypotheses; publication match: exact; verification: read the result; development: none; confidence: high}
\end{GSitem}
\begin{GSitem}{1986-673.4}{E}{Milnor fibres and the plane-curve Betti-number formula}
\GSfield{Kind; source}{theorem; \GSpage{P001278}}
\GSfield{The notes}{For \(f:(\mathbb C^n,0)\to(\mathbb C,0)\) with an isolated critical point, the fibre is a bouquet of \(\mu\) spheres. In dimension two the notes write \(\mu=2g-(\text{holes})+1\).}
\GSfield{Statement}{For sufficiently small \(0<|t|\ll\varepsilon\), the Milnor fibre \(f^{-1}(t)\cap B_\varepsilon\) is homotopy equivalent to a wedge of \(\mu\) spheres of real dimension \(n-1\), with \(\mu=\dim_{\mathbb C}\mathbb C\{z_1,\ldots,z_n\}/(\partial_1f,\ldots,\partial_nf)\). For a plane-curve fibre of genus \(g\) with \(b\) boundary components, \(\mu=2g+b-1\). The notebook has both signs reversed in this last relation.}
\GSfield{Reconstruction}{The fibre is intersected with a Milnor ball, and “holes” means boundary components. The genus formula is an emendation, not a transcription correction.}
\GSfield{Fate}{This is \citet[Chapters 5--7]{milnor1968singular}; the bouquet statement is also recalled in \citet[\S1]{steenbrink1989spectrum}. The Russian transcription and scan both confirm the erroneous sign pattern.}
\GSfield{Present status}{As of November 2026, the corrected formulation is an established result; no unresolved claim is inferred from this occurrence.}
\GSfield{Assessment}{literal statement: false as written; publication match: exact; verification: read the result; development: none; confidence: high}
\end{GSitem}
\begin{GSitem}{1986-673.5}{B}{Several functions and mixed Hodge structures}
\GSfield{Kind; source}{question; \GSpage{P001278}}
\GSfield{The notes}{Gelfand asks about \(\int f_1^{\lambda_1}\cdots f_k^{\lambda_k}\) and whether a map \(\mathbb C^n\to\mathbb C^k\) gives something like a mixed Hodge structure.}
\GSfield{Statement}{For a specified algebraic map and algebraic fibres, their cohomology carries mixed Hodge structures; variation and degeneration require a stratification and admissibility or nearby-cycle hypotheses. Multivariate powers are governed by systems of functional equations rather than one canonical polynomial in one variable. The notes do not specify whether ordinary fibre cohomology, iterated nearby cycles, or a joint limit along a normal-crossing base is intended, and these objects need not coincide.}
\GSfield{Reconstruction}{The modern alternatives are explicitly distinguished instead of choosing an unrecorded order of limits.}
\GSfield{Fate}{Deligne\textquotesingle s theory already covers algebraic fibre cohomology \citep{deligne1971hodge}; the later mixed-Hodge-module framework supplies functorial direct images and nearby-cycle operations, as stated in \citet[Theorem 1.3]{saito1989introduction}. This is a partial implementation of the request, not an exact solution to an unrecorded multivariable limit problem.}
\GSfield{Present status}{As of November 2026, the specified algebraic cohomology and functorial operations are established. The notebook does not determine a separate precise unresolved multivariable statement.}
\GSfield{Assessment}{literal statement: correct with missing hypotheses; publication match: partial; verification: read the result; development: none; confidence: medium}
\end{GSitem}
\begin{GSitem}{1986-673.6}{F}{A map defined by two quadratic forms}
\GSfield{Kind; source}{question; \GSpage{P001278}; I. M. Gelfand}
\GSfield{The notes}{Gelfand asks what arises for a map given by a pair of quadratic forms.}
\GSfield{Statement}{Given \(q_1,q_2\) on a complex vector space, one can study the affine fibres of \((q_1,q_2)\), a projective intersection of quadrics, or a pencil and its discriminant. The notebook gives no dimension, nondegeneracy, isolated-complete-intersection hypothesis, or choice among these objects. Each choice leads to a different cohomological problem.}
\GSfield{Fate}{The immediately preceding several-functions discussion identifies the motivation, but a search of the mixed-Hodge and intersection-of-quadrics context does not identify a unique assertion. The passage is therefore retained as a question with class F, rather than represented as an overlooked open theorem.}
\GSfield{Present status}{As of November 2026, the intended object and quantifiers cannot be recovered from this sentence.}
\GSfield{Assessment}{literal statement: not determinable; publication match: none found; verification: unresolved lead; development: none; confidence: low}
\end{GSitem}
\begin{GSitem}{1986-673.7}{E}{Gelfand--Leray periods and logarithmic asymptotics}
\GSfield{Kind; source}{construction; \GSpage{P001278-P001279}}
\GSfield{The notes}{A holomorphic top form \(\omega\) gives \(\omega/df\) on a regular fibre, and periods along transported cycles have terms \(t^\alpha(\log t)^k\).}
\GSfield{Statement}{For an isolated hypersurface singularity and a chosen branch over a punctured disc, the period \(\int_{\delta(t)}\omega/df\) has regular-singular asymptotic expansions with rational exponents and finitely bounded logarithmic powers. Monodromy eigenvalues determine exponents modulo integers; a Jordan block of size \(s\) allows logarithmic degree at most \(s-1\), not \(s\). The exact convention for the exponent depends on using \(\omega/df\) rather than another period normalization.}
\GSfield{Reconstruction}{The branch, regular-singular condition, and block-size shift are supplied.}
\GSfield{Fate}{This is the period setting of \citet[\S\S1--2]{varchenko1982asymptotic}, whose Russian paper appeared in 1981 and English translation in 1982 after receipt in March 1980; the announcement \citet{varchenko1980asymptotics} appeared in 1980.}
\GSfield{Present status}{As of November 2026, the corrected formulation is an established result; no unresolved claim is inferred from this occurrence.}
\GSfield{Assessment}{literal statement: correct with missing hypotheses; publication match: exact; verification: read the result; development: none; confidence: high}
\end{GSitem}
\begin{GSitem}{1986-673.8}{E}{The asymptotic Hodge filtration and monodromy weights}
\GSfield{Kind; source}{construction; \GSpage{P001279-P001280}}
\GSfield{The notes}{The Jordan form of monodromy determines a weight filtration, while leading asymptotic terms of forms determine a Hodge filtration.}
\GSfield{Statement}{For the vanishing cohomology \(H^{n-1}\) of an isolated hypersurface singularity, the logarithm of the unipotent monodromy has its monodromy weight filtration centered at \(n-1\) on the nonunit-eigenvalue part and at \(n\) on the eigenvalue-one part. The asymptotic or limiting Hodge filtration contains information beyond the Jordan normal form. Leading exponents with their filtration multiplicities define the spectrum.}
\GSfield{Reconstruction}{The two weight centers and the distinction between a monodromy filtration and the full mixed Hodge structure are restored.}
\GSfield{Fate}{\citet{varchenko1982asymptotic} develops the asymptotic construction; \citet[\S1]{steenbrink1989spectrum} states the Hodge-theoretic definition and relates it to the earlier isolated case. The notebook is reporting established work of the speaker.}
\GSfield{Present status}{As of November 2026, the corrected formulation is an established result; no unresolved claim is inferred from this occurrence.}
\GSfield{Assessment}{literal statement: correct with missing hypotheses; publication match: exact; verification: read the result; development: none; confidence: high}
\end{GSitem}
\begin{GSitem}{1986-673.9}{E}{Bernstein roots, period exponents, and spectral numbers}
\GSfield{Kind; source}{claim; \GSpage{P001279-P001280}}
\GSfield{The notes}{Exponents in the period asymptotics are related to roots of a Bernstein polynomial and to the Hodge spectrum.}
\GSfield{Statement}{Bernstein--Sato roots, monodromy eigenvalues, and spectral numbers are related invariants with different normalizations and multiplicities. Monodromy identifies rational exponents only modulo integers. Spectral numbers also use the Hodge filtration; arbitrary exponents of every term in every asymptotic series are not simply the finite set of Bernstein roots. The precise equality proposed in the abbreviated line cannot be asserted without fixing the polynomial (including reduction), signs, shifts, and choice of leading terms.}
\GSfield{Reconstruction}{The finite leading invariant is separated from the full asymptotic series.}
\GSfield{Fate}{Compare the functional equation in \citet{bernstein1972continuation}, the asymptotic filtration in \citet{varchenko1982asymptotic}, and the spectrum definition in \citet[\S1]{steenbrink1989spectrum}. These sources support the relationship, not an unqualified equality of all three multisets.}
\GSfield{Present status}{As of November 2026, the corrected formulation is an established result; no unresolved claim is inferred from this occurrence.}
\GSfield{Assessment}{literal statement: correct with missing hypotheses; publication match: partial; verification: read the result; development: none; confidence: medium}
\end{GSitem}
\begin{GSitem}{1986-673.10}{E}{Semicontinuity of the singularity spectrum}
\GSfield{Kind; source}{theorem; \GSpage{P001280}}
\GSfield{The notes}{The final statement invokes semicontinuity of the spectrum under perturbation.}
\GSfield{Statement}{For a deformation of an isolated hypersurface singularity, the number of spectral values in a fixed half-open interval of length one, counted with multiplicity, cannot increase when the original singularity splits into nearby isolated singularities with the same critical value: the count for the central singularity is at least the sum of the nearby counts. The interval endpoints and common-critical-value hypothesis are part of the statement.}
\GSfield{Reconstruction}{The interval and splitting hypotheses make the otherwise schematic final sentence precise.}
\GSfield{Fate}{Varchenko announced semicontinuity results in \citet{varchenko1983semicontinuity}; Steenbrink\textquotesingle s 1985 result supplies the general theorem. The precise formulation was read in the author\textquotesingle s later paper \citet[p.~166]{steenbrink1989spectrum}, which cites both works.}
\GSfield{Present status}{As of November 2026, the corrected formulation is an established result; no unresolved claim is inferred from this occurrence.}
\GSfield{Assessment}{literal statement: correct with missing hypotheses; publication match: exact; verification: read the result; development: none; confidence: high}
\end{GSitem}

\GStalk{674}{1987-04-07}{I. V. Cherednik}{P001279-P001283}{RTT algebras, finite-level Yangians, and PBW questions}
Cherednik starts on P001281 and finishes near the top of P001282. The rational RTT equation defines a homogeneous version of a truncated Yangian, with a finite-level PBW claim and a level-one enveloping-algebra example. Gelfand asks for a broader quadratic-algebra criterion. Kogan begins later on P001282. The manuscript heading reads \(7/\mathrm{IV}\text{-}86\).

\begin{GSitem}{1986-674.1}{A}{The rational RTT presentation with a finite level}
\GSfield{Kind; source}{definition; \GSpage{P001281}}
\GSfield{The notes}{With \(R(u)=I+P/u\), the matrix \(L_n(u)=EI+\sum_{q=1}^nu^{-q}\sum e_{ij}\otimes E_{ji}^{q-1}\) satisfies \(R(u-v)L_1(u)L_2(v)=L_2(v)L_1(u)R(u-v)\); \(E\) is central.}
\GSfield{Statement}{After specializing the central homogenizing parameter \(E=1\), the coefficient relations define a level-\(n\) or truncated Yangian. The full Yangian allows infinitely many coefficients. The notebook uses a transposed coefficient matrix and the plus-sign \(R\) convention; both must be adjusted together when comparing with the usual \(R(u)=I-P/u\) convention.}
\GSfield{Reconstruction}{The specialization and convention conversion are made explicit; the central homogeneous presentation is kept distinct from the usual normalized one.}
\GSfield{Fate}{The full RTT algebra is treated in \citet[\S1]{molevnazarovolshanski1996yangians}, first circulated in September 1994; the exact finite-level presentation is recorded in \citet[\S5.2]{futornyovsienko2005kostant}, whose March 2003 preprint credits Cherednik. This identifies the construction without conflating finite and infinite level.}
\GSfield{Present status}{As of November 2026, the corrected formulation is an established result; no unresolved claim is inferred from this occurrence.}
\GSfield{Assessment}{literal statement: correct with missing hypotheses; publication match: exact; verification: read the result; development: none; confidence: high}
\end{GSitem}
\begin{GSitem}{1986-674.2}{A}{PBW and the current-algebra degeneration}
\GSfield{Kind; source}{theorem; \GSpage{P001281-P001282}}
\GSfield{The notes}{The algebra has the size of a polynomial algebra and is a deformation of a matrix current algebra.}
\GSfield{Statement}{For the normalized truncated Yangian over characteristic zero, ordered monomials in its \(nN^2\) coefficients form a PBW basis. With coefficient degree \(q\), the associated graded algebra is a polynomial algebra with weighted generators. With degree \(q-1\), its current-algebra degeneration has the truncated current bracket. These filtrations should not be confused with an assertion that all generators have the same degree in a normalized presentation.}
\GSfield{Reconstruction}{The two filtrations and the normalization \(E=1\) replace the shorthand polynomial-size claim.}
\GSfield{Fate}{\citet[\S5.2, Proposition 3]{futornyovsienko2005kostant} explicitly states the finite-level polynomial associated-graded theorem and credits Cherednik\textquotesingle s 1989 paper. The two standard filtrations for the full Yangian are described in \citet[Theorem 1.22 and \S1.26]{molevnazarovolshanski1996yangians}.}
\GSfield{Present status}{As of November 2026, the corrected formulation is an established result; no unresolved claim is inferred from this occurrence.}
\GSfield{Assessment}{literal statement: correct with missing hypotheses; publication match: partial; verification: read the result; development: none; confidence: high}
\end{GSitem}
\begin{GSitem}{1986-674.3}{E}{Level one recovers the general linear enveloping algebra}
\GSfield{Kind; source}{theorem; \GSpage{P001282}}
\GSfield{The notes}{At level one the algebra is \(U(\mathfrak{gl}_N)\).}
\GSfield{Statement}{For \(L(u)=I+u^{-1}(E_{ji})\) and the notebook\textquotesingle s matching \(R\)-matrix convention, the RTT coefficient relations are exactly \([E_{ij},E_{k\ell}]=\delta_{jk}E_{i\ell}-\delta_{i\ell}E_{kj}\). Thus the normalized level-one algebra is \(U(\mathfrak{gl}_N)\). This is an evaluation quotient of the full Yangian, not an isomorphism from the full Yangian.}
\GSfield{Fate}{\citet[\S1.16]{molevnazarovolshanski1996yangians} states the embedding/evaluation maps; the finite-level identification is explicit in \citet[\S5.2]{futornyovsienko2005kostant}.}
\GSfield{Present status}{As of November 2026, the corrected formulation is an established result; no unresolved claim is inferred from this occurrence.}
\GSfield{Assessment}{literal statement: correct; publication match: exact; verification: read the result; development: none; confidence: high}
\end{GSitem}
\begin{GSitem}{1986-674.4}{B}{PBW criteria beyond the Yangian example}
\GSfield{Kind; source}{question; \GSpage{P001282}; I. M. Gelfand}
\GSfield{The notes}{Gelfand asks how to describe quadratic algebras with polynomial-type size or the PBW property.}
\GSfield{Statement}{One precise later version fixes a homogeneous quadratic Koszul algebra \(A=T(V)/(R)\) and asks which filtered relations with leading part \(R\) yield \(\operatorname{gr}U\simeq A\). For relations \(r-\alpha(r)-\beta(r)\), the necessary overlap conditions in degree three are sufficient when \(A\) is Koszul. This is a deformation criterion for a fixed \(A\), not a classification of every quadratic algebra with a specified Hilbert series.}
\GSfield{Reconstruction}{The fixed quadratic algebra and Koszul hypothesis are later specifications, not words supplied by the notebook.}
\GSfield{Fate}{\citet[Theorem 0.5 and \S3]{bravermangaitsgory1996pbw}, first circulated in November 1994 and published in 1996, establishes the Koszul PBW criterion. It implements a definite part of Gelfand\textquotesingle s broad question; the full unbounded classification is not silently declared solved or open.}
\GSfield{Present status}{As of November 2026, the stated Koszul deformation criterion is established. The notebook contains no fixed generator number or equivalence relation defining a further classification problem.}
\GSfield{Assessment}{literal statement: correct with missing hypotheses; publication match: partial; verification: read the result; development: none; confidence: high}
\end{GSitem}
\begin{GSitem}{1986-674.5}{F}{Yang--Baxter consistency and a PBW converse}
\GSfield{Kind; source}{claim; \GSpage{P001282}}
\GSfield{The notes}{The notes attribute a general connection between exchange relations, PBW, and the Yang--Baxter equation to Zamolodchikov, and mention trigonometric and elliptic alternatives.}
\GSfield{Statement}{For a specified exchange algebra with an ordered reduction system, equality of the two reductions of a triple product leads to the Yang--Baxter equation. To deduce a PBW basis one must also control invertibility, admissible ordering, all ambiguities, and possible higher relations. An arbitrary quadratic quotient need not have PBW merely because a displayed operator satisfies Yang--Baxter; a universal converse needs a specified class of presentations.}
\GSfield{Reconstruction}{The supported triple-overlap mechanism is separated from an unsupported universal equivalence.}
\GSfield{Fate}{The rational RTT case is proved in \citet[Theorem 1.22]{molevnazarovolshanski1996yangians}; the Koszul deformation theorem \citep{bravermangaitsgory1996pbw} shows why an extra structural hypothesis matters. No exact all-\(R\) biconditional corresponding to the unqualified sentence was identified.}
\GSfield{Present status}{As of November 2026, the specific rational construction is established; the universal formulation remains under-specified in these notes.}
\GSfield{Assessment}{literal statement: not determinable; publication match: thematic only; verification: read the result; development: none; confidence: medium}
\end{GSitem}
\begin{GSitem}{1986-674.6}{F}{Associativity and a three-cocycle remark}
\GSfield{Kind; source}{remark; \GSpage{P001282}}
\GSfield{The notes}{A brief remark associates a failure of associativity and a three-cocycle with Jackiw.}
\GSfield{Statement}{A group action associative only up to a phase can lead to a three-cocycle as an associativity obstruction after the relevant objects and multiplication law are specified. Here no group, coefficient module, operators, or cocycle formula is recorded. The sentence cannot be identified with a unique theorem about the RTT algebra.}
\GSfield{Fate}{Searches around the named associativity and cocycle remark locate a broad physics context but no exact notebook formula. It is retained as an unidentified remark rather than folded into the PBW theorem.}
\GSfield{Present status}{As of November 2026, the intended three-cocycle construction is not determined by the fragment.}
\GSfield{Assessment}{literal statement: not determinable; publication match: none found; verification: unresolved lead; development: none; confidence: low}
\end{GSitem}

\GStalk{675}{1987-04-07}{Yan I. Kogan}{P001281-P001287}{Covariant bosonic string states, constraints, and string fields}
The talk begins midway through P001282 and continues through P001286. It derives the oscillator and Virasoro description, explains the mass tower and physical-state constraints, and moves toward gauge symmetry and interacting string fields. Several sign, mode-index, and level labels are faithfully wrong in the notebook and are emended explicitly. The shared heading reads April 1986.

\begin{GSitem}{1986-675.1}{E}{Regge trajectories, tension, and the field-theory limit}
\GSfield{Kind; source}{remark; \GSpage{P001282-P001283}}
\GSfield{The notes}{The string slope controls an approximately linear relation of spin and squared mass, and the limit of small slope recovers a field theory.}
\GSfield{Statement}{For the idealized string spectrum, the leading trajectory has \(J=\alpha^{\prime}M^2+a\), with tension \(T=1/(2\pi\alpha^{\prime})\) in the standard convention. The limit \(\alpha^{\prime}\to0\) at fixed low external energies suppresses massive string states; the surviving effective theory depends on the open/closed sector and gauge labels. This is a controlled low-energy expansion, not equality of complete theories at finite slope.}
\GSfield{Reconstruction}{The sector and low-energy regime are specified.}
\GSfield{Fate}{The low-energy gauge and gravity interpretation is already part of \citet{grossharveymartinecrohm1985heterotic}. The covariant bosonic construction in \citet{witten1986noncommutative} realizes the interacting open-string setting discussed at the end of the talk.}
\GSfield{Present status}{As of November 2026, the corrected formulation is an established result; no unresolved claim is inferred from this occurrence.}
\GSfield{Assessment}{literal statement: correct with missing hypotheses; publication match: partial; verification: read the result; development: none; confidence: high}
\end{GSitem}
\begin{GSitem}{1986-675.2}{E}{The Polyakov action and conformal-gauge constraints}
\GSfield{Kind; source}{construction; \GSpage{P001283}}
\GSfield{The notes}{The action for \(x^\mu\) and the worldsheet metric is reduced to free equations with constraints after conformal gauge fixing.}
\GSfield{Statement}{The Polyakov action is \(S=(4\pi\alpha^{\prime})^{-1}\int\sqrt{|h|}h^{ab}\partial_aX^\mu\partial_bX_\mu\,d^2\sigma\). Variation of \(h\) gives the vanishing stress tensor. Gauge fixing the metric leaves this constraint and the associated ghosts; a target with one timelike coordinate gives an indefinite oscillator inner product before physical-state reduction.}
\GSfield{Reconstruction}{The action normalization and the distinction between gauge fixing and discarding constraints are restored.}
\GSfield{Fate}{This is the critical-string starting point of \citet{polyakov1981quantum}. Its quantized BRST form is written explicitly in \citet[eqs.~(1)--(7)]{furuuchiohta2006noghost}.}
\GSfield{Present status}{As of November 2026, the corrected formulation is an established result; no unresolved claim is inferred from this occurrence.}
\GSfield{Assessment}{literal statement: correct with missing hypotheses; publication match: exact; verification: read the result; development: none; confidence: high}
\end{GSitem}
\begin{GSitem}{1986-675.3}{E}{Oscillator quantization and the stress-tensor modes}
\GSfield{Kind; source}{construction; \GSpage{P001283-P001284}}
\GSfield{The notes}{The modes satisfy a Heisenberg relation with the target metric, and \(L_n\) is written as a quadratic oscillator sum.}
\GSfield{Statement}{For the open bosonic string, \([\alpha_n^\mu,\alpha_m^\nu]=n\delta_{n+m,0}\eta^{\mu\nu}\). The normal-ordered matter stress-tensor modes are \(L_n=\tfrac12\sum_m:\alpha_m\cdot\alpha_{n-m}:\). The two oscillator indices must add to \(n\); the notebook line with a fixed first index \(n\) does not define the intended operator.}
\GSfield{Reconstruction}{The summation index and normal ordering are repaired; this does not change the diplomatic reading.}
\GSfield{Fate}{The exact oscillator and quadratic-mode formulas are given in \citet[eq.~(3)]{furuuchiohta2006noghost}; they are much older expository material, independently visible in the earlier seminar oscillator discussion 1985-662.1--9.}
\GSfield{Present status}{As of November 2026, the corrected formulation is an established result; no unresolved claim is inferred from this occurrence.}
\GSfield{Assessment}{literal statement: false as written; publication match: exact; verification: read the result; development: none; confidence: high}
\end{GSitem}
\begin{GSitem}{1986-675.4}{E}{The open-string mass spectrum and the level-one vector}
\GSfield{Kind; source}{theorem; \GSpage{P001284}}
\GSfield{The notes}{The vacuum is tachyonic, the next state is massless, and higher oscillator levels form a massive tower.}
\GSfield{Statement}{With intercept one, the open critical bosonic string satisfies \(\alpha^{\prime}M^2=N-1\), where \(N=\sum_{n>0}\alpha_{-n}\cdot\alpha_n\) after the usual transverse reduction. Thus \(N=0\) is tachyonic and \(N=1\) is massless. A formula with mass itself proportional to \(N-1\), or without the square-root/sign convention, is not the same statement.}
\GSfield{Reconstruction}{The notebook\textquotesingle s mass normalization is replaced by the conventional squared-mass relation.}
\GSfield{Fate}{The intercept and on-shell operator are explicit in \citet[eq.~(2)]{furuuchiohta2006noghost}; the spectrum is the conventional Veneziano spectrum referred to by \citet{witten1986noncommutative}.}
\GSfield{Present status}{As of November 2026, the corrected formulation is an established result; no unresolved claim is inferred from this occurrence.}
\GSfield{Assessment}{literal statement: correct with missing hypotheses; publication match: partial; verification: read the result; development: none; confidence: high}
\end{GSitem}
\begin{GSitem}{1986-675.5}{E}{The Virasoro commutator and its mode index}
\GSfield{Kind; source}{claim; \GSpage{P001284}}
\GSfield{The notes}{The Virasoro commutator is displayed with \(L_{n-m}\) instead of \(L_{n+m}\).}
\GSfield{Statement}{The matter oscillators of \(D\) free bosons give \([L_n,L_m]=(n-m)L_{n+m}+\tfrac{D}{12}(n^3-n)\delta_{n+m,0}\). The sum of mode indices is forced by the stress-tensor Fourier grading. The ghost system changes the total central charge to \(D-26\).}
\GSfield{Reconstruction}{The subscript \(n+m\) is an editorial correction.}
\GSfield{Fate}{The oscillator realization used in \citet[eqs.~(2)--(4)]{furuuchiohta2006noghost} and the anomaly calculation in \citet{belavinknizhnik1986complex} support the corrected formula. The scan shows the incorrect subscript in the notebook.}
\GSfield{Present status}{As of November 2026, the corrected formulation is an established result; no unresolved claim is inferred from this occurrence.}
\GSfield{Assessment}{literal statement: false as written; publication match: exact; verification: read the result; development: none; confidence: high}
\end{GSitem}
\begin{GSitem}{1986-675.6}{E}{Physical states, null states, and the no-ghost theorem}
\GSfield{Kind; source}{theorem; \GSpage{P001284-P001285}}
\GSfield{The notes}{Physical states satisfy positive-mode constraints and the \(L_0\) equation; negative-norm states are removed in the critical dimension.}
\GSfield{Statement}{At \(D=26\) and intercept one, the physical matter subspace satisfies \(L_n\psi=0\) for \(n>0\) and \((L_0-1)\psi=0\), then is quotiented by null spurious states. In the BRST formulation, one takes the relevant cohomology of a nilpotent BRST operator. The no-ghost theorem identifies the physical space with a nonnegative-norm transverse sector, under the usual momentum hypotheses; it does not say the unreduced covariant Fock space is positive.}
\GSfield{Reconstruction}{The null quotient and the momentum qualification are stated explicitly.}
\GSfield{Fate}{The theorem originates with the early no-ghost work; the opened primary proof \citet{furuuchiohta2006noghost}, first circulated in July 2006, gives the similarity-transformation argument and explicitly assumes nonzero light-cone momentum. That hypothesis is retained here.}
\GSfield{Present status}{As of November 2026, the corrected formulation is an established result; no unresolved claim is inferred from this occurrence.}
\GSfield{Assessment}{literal statement: correct with missing hypotheses; publication match: exact; verification: read the result; development: none; confidence: high}
\end{GSitem}
\begin{GSitem}{1986-675.7}{E}{Component fields and constraints at the first two levels}
\GSfield{Kind; source}{construction; \GSpage{P001285}}
\GSfield{The notes}{The string field is expanded as a scalar, a vector coefficient of \(\alpha_{-1}\), a tensor coefficient of two \(\alpha_{-1}\) modes, and a vector coefficient of \(\alpha_{-2}\). A later constraint uses the earlier vector label.}
\GSfield{Statement}{A string field has independent component fields at each oscillator level. The level-one positive-mode condition gives the massless vector transversality condition. At level two, a state \((B_{\mu\nu}\alpha_{-1}^\mu\alpha_{-1}^\nu+C_\mu\alpha_{-2}^\mu)|p\rangle\) has constraints relating \(B\) and \(C\); it does not involve the separate level-one field \(A_\mu\). The numerical factors depend on the normalization of the symmetric tensor coefficient.}
\GSfield{Reconstruction}{The level-two vector is renamed \(C\) to prevent the notebook\textquotesingle s conflation; unverified numerical factors are not copied.}
\GSfield{Fate}{The level preservation follows from the oscillator and Virasoro definitions, as used in \citet{furuuchiohta2006noghost}. The component-field interpretation is part of the string-field framework \citep{witten1986noncommutative}.}
\GSfield{Present status}{As of November 2026, the corrected formulation is an established result; no unresolved claim is inferred from this occurrence.}
\GSfield{Assessment}{literal statement: false as written; publication match: partial; verification: read the result; development: none; confidence: high}
\end{GSitem}
\begin{GSitem}{1986-675.8}{E}{Null-state gauge transformations and interacting nonabelian terms}
\GSfield{Kind; source}{claim; \GSpage{P001285-P001286}}
\GSfield{The notes}{The vector state changes by a momentum multiple, and the notes connect this with Yang--Mills gauge transformations including a commutator.}
\GSfield{Statement}{At the linearized massless level, a null-state or BRST-exact change induces \(\delta A_\mu=\partial_\mu\epsilon\). The nonabelian term \([A_\mu,\epsilon]\) requires interacting string fields and Chan--Paton algebra data. It cannot be deduced solely from the free positive-mode constraints.}
\GSfield{Reconstruction}{The free and interacting transformations are separated.}
\GSfield{Fate}{\citet{witten1986noncommutative} constructs a nonlinear gauge-invariant open-string field theory. The distinction between its interaction and the free BRST reduction is explicit in the respective settings; the notebook elides that step.}
\GSfield{Present status}{As of November 2026, the corrected formulation is an established result; no unresolved claim is inferred from this occurrence.}
\GSfield{Assessment}{literal statement: correct with missing hypotheses; publication match: partial; verification: publisher metadata checked; development: none; confidence: high}
\end{GSitem}
\begin{GSitem}{1986-675.9}{E}{Associative string products and BRST string-field gauge symmetry}
\GSfield{Kind; source}{construction; \GSpage{P001286}}
\GSfield{The notes}{The talk closes with string fields, BRST ideas, and a connection with noncommutative geometry.}
\GSfield{Statement}{In the cubic open bosonic theory, the graded associative string product, a cyclic pairing, and a nilpotent BRST derivation permit an action of the form \(S=\tfrac12\langle\Psi,Q\Psi\rangle+\tfrac13\langle\Psi,\Psi*\Psi\rangle\), with gauge transformation \(\delta\Psi=Q\Lambda+\Psi*\Lambda-\Lambda*\Psi\) in the appropriate ghost degrees. Analytic definitions and normalization of the string product are additional structure; ordinary multiplication of component fields is not enough.}
\GSfield{Reconstruction}{The standard cubic formulation is supplied as a contextual reconstruction, not an exact formula transcribed from this page.}
\GSfield{Fate}{This identifies the framework of \citet{witten1986noncommutative}, published 5 May 1986. Only its institutional primary metadata and abstract were accessible here; the complete construction is not claimed as directly read. It predates the catalogue date but follows the manuscript April 1986 heading.}
\GSfield{Present status}{As of November 2026, the corrected formulation is an established result; no unresolved claim is inferred from this occurrence.}
\GSfield{Assessment}{literal statement: correct with missing hypotheses; publication match: partial; verification: publisher metadata checked; development: none; confidence: medium}
\end{GSitem}

\GStalk{676}{1987-04-28}{I. V. Cherednik}{P001285-P001290}{Quantum determinants, centers, and fusion}
The actual talk occupies P001287--P001289. It repeats the finite-level RTT construction, defines its determinant by antisymmetrization, states the polynomial-center theorem, and uses tensor products and Young diagrams to build representations. Matrix size \(N\) and level \(n\) are repeatedly interchanged in the notebook. The date visibly reads 28 April 1986.

\begin{GSitem}{1986-676.1}{A}{Finite-level RTT relations and a stray spectral minus sign}
\GSfield{Kind; source}{definition; \GSpage{P001287}}
\GSfield{The notes}{The normalized \(L\) has \(n\) negative powers and size \(N\). In the RTT equation only the left side has \(L^2(-\lambda_2)\).}
\GSfield{Statement}{The level-\(n\) RTT relation uses the same \(L_2(v)\) on both sides. With the single unmatched minus sign, comparison of leading coefficients generally forces unintended relations and does not give the displayed coefficient algebra. The corrected definition is the normalized one from 1986-674.1.}
\GSfield{Reconstruction}{The minus sign is emended from the consistency of the two sides and the immediately following coefficient relations.}
\GSfield{Fate}{\citet[\S5.2]{futornyovsienko2005kostant} gives the exact truncated presentation. The Russian evidence note and scan confirm that the unmatched minus is in the manuscript.}
\GSfield{Present status}{As of November 2026, the corrected formulation is an established result; no unresolved claim is inferred from this occurrence.}
\GSfield{Assessment}{literal statement: false as written; publication match: exact; verification: read the result; development: none; confidence: high}
\end{GSitem}
\begin{GSitem}{1986-676.2}{A}{Quantum determinant by rank-N antisymmetrization}
\GSfield{Kind; source}{construction; \GSpage{P001287}}
\GSfield{The notes}{The determinant is the scalar of the shifted tensor product \(L(\lambda+N-1)\otimes\cdots\otimes L(\lambda)\), but the projector is labeled \(\mathcal P_-^n\).}
\GSfield{Statement}{For \(N\times N\) matrices in the RTT algebra, the top exterior power \(\bigwedge^N\mathbb C^N\) is one-dimensional. With the compatible shift convention, its antisymmetrizer \(A_N\) satisfies \(A_NL_1(\lambda+N-1)\cdots L_N(\lambda)A_N=H(\lambda)A_N\). It is \(N\), the matrix size, that determines the number of factors and the exterior projector, regardless of the level \(n\).}
\GSfield{Reconstruction}{Both occurrences of the projector degree and the tensor-power exponent must be \(N\).}
\GSfield{Fate}{\citet[\S\S2.2--2.7]{molevnazarovolshanski1996yangians} develops the quantum determinant. The finite-level version is explicit in \citet[eq.~(5), \S5.2]{futornyovsienko2005kostant}.}
\GSfield{Present status}{As of November 2026, the corrected formulation is an established result; no unresolved claim is inferred from this occurrence.}
\GSfield{Assessment}{literal statement: false as written; publication match: exact; verification: read the result; development: none; confidence: high}
\end{GSitem}
\begin{GSitem}{1986-676.3}{A}{The center of a truncated Yangian}
\GSfield{Kind; source}{theorem; \GSpage{P001287-P001288}}
\GSfield{The notes}{The principal parts of \(H(\lambda)\) at \(\lambda=-i\) generate the center, with \(q=1,\ldots,n\) and the written range \(i=0,\ldots,n-1\).}
\GSfield{Statement}{For the level-\(n\) Yangian of \(\mathfrak{gl}_N\) over \(\mathbb C\), multiply the rational determinant by \(\prod_{i=0}^{N-1}(\lambda+i)^n\). The result is a monic degree-\(nN\) polynomial whose \(nN\) nonleading coefficients are algebraically independent generators of the center. Equivalently, the \(nN\) principal-part coefficients at \(0,-1,\ldots,-(N-1)\), of orders \(1,\ldots,n\), generate freely. The notebook\textquotesingle s range of pole locations must use \(N-1\).}
\GSfield{Reconstruction}{The partial-fraction and polynomial presentations are distinguished, and matrix size is restored in the pole range.}
\GSfield{Fate}{The finite-level center theorem is stated exactly in \citet[\S5.2, after eq.~(5)]{futornyovsienko2005kostant}, crediting \citet{cherednik1989hidden} and Molev. This gives an author-publication continuation, while the original Cherednik chapter was not directly accessible in this run.}
\GSfield{Present status}{As of November 2026, the corrected formulation is an established result; no unresolved claim is inferred from this occurrence.}
\GSfield{Assessment}{literal statement: false as written; publication match: exact; verification: read the result; development: none; confidence: high}
\end{GSitem}
\begin{GSitem}{1986-676.4}{E}{Evaluation representations and tensor products}
\GSfield{Kind; source}{construction; \GSpage{P001288}}
\GSfield{The notes}{A representation \(V\) of \(\mathfrak{gl}_N\) supplies a matrix \(L_V\); tensor products are formed by multiplying such matrices with spectral shifts.}
\GSfield{Statement}{The evaluation homomorphism from the full Yangian to \(U(\mathfrak{gl}_N)\), followed by a representation and a spectral-translation automorphism, gives an evaluation module. The coproduct \(\Delta T_{ij}(u)=\sum_aT_{ia}(u)\otimes T_{aj}(u)\) realizes tensor products as matrix products. A product of finite-level matrices has degree equal to the sum of their levels; a fixed-level quotient is not automatically a Hopf algebra.}
\GSfield{Reconstruction}{The shift automorphism and full-versus-truncated distinction are supplied.}
\GSfield{Fate}{The evaluation and coproduct statements are \citet[\S1.16 and Theorem 1.28]{molevnazarovolshanski1996yangians}. The final finite-level caution prevents carrying the full Hopf structure to the wrong quotient.}
\GSfield{Present status}{As of November 2026, the corrected formulation is an established result; no unresolved claim is inferred from this occurrence.}
\GSfield{Assessment}{literal statement: correct with missing hypotheses; publication match: exact; verification: read the result; development: none; confidence: high}
\end{GSitem}
\begin{GSitem}{1986-676.5}{A}{Young diagrams and the rational fusion procedure}
\GSfield{Kind; source}{construction; \GSpage{P001288-P001289}}
\GSfield{The notes}{A Young diagram supplies spectral shifts equal to row minus column; a product of elementary \(R\)-matrices and a Young operator produces the associated representation.}
\GSfield{Statement}{For a standard tableau of a partition of \(r\), take its contents with a fixed sign convention and ordered tensor products of fundamental evaluation modules at those contents. Fusion requires regularizing the rational products at coincident contents and applying a properly normalized tableau idempotent. The resulting subquotient realizes the corresponding polynomial \(\mathfrak{gl}_N\) representation when the number of rows is at most \(N\). A bare unnormalized row-column symmetrizer is not interchangeable with every orthogonal tableau idempotent.}
\GSfield{Reconstruction}{Tableau choice, regularization, dimension, and normalization are necessary additions.}
\GSfield{Fate}{This is the line of \citet{cherednik1987gelfand}, published in 1987; its role in Gelfand--Tsetlin bases and Yangian representations is specifically documented in \citet[Introduction]{molevnazarovolshanski1996yangians}. The original Duke text was blocked, so exact fusion normalizations are not certified from it.}
\GSfield{Present status}{As of November 2026, the corrected formulation is an established result; no unresolved claim is inferred from this occurrence.}
\GSfield{Assessment}{literal statement: correct with missing hypotheses; publication match: partial; verification: indirect; development: none; confidence: medium}
\end{GSitem}
\begin{GSitem}{1986-676.6}{E}{Symmetric-group normalization in the Young operator}
\GSfield{Kind; source}{remark; \GSpage{P001289}}
\GSfield{The notes}{The order of the symmetric group is written as \(n\) in a normalization discussion.}
\GSfield{Statement}{The symmetric group on \(n\) letters has order \(n!\). Averaging over a finite group uses its actual order; Young symmetrizers also require their own hook-length normalization to become idempotents. Replacing every such normalization merely by \(n!\) would therefore overcorrect the passage.}
\GSfield{Reconstruction}{Only the group order is emended; the particular projector normalization remains convention dependent.}
\GSfield{Fate}{The tensor and tableau context is the fusion construction of \citet{cherednik1987gelfand}. The group-order slip is elementary and visible in the manuscript; no new theorem is claimed.}
\GSfield{Present status}{As of November 2026, the corrected formulation is an established result; no unresolved claim is inferred from this occurrence.}
\GSfield{Assessment}{literal statement: false as written; publication match: partial; verification: indirect; development: none; confidence: high}
\end{GSitem}

\GStalk{677}{1987-05-05}{G. I. Olshanskii}{P001288-P001292}{A strong-limit enveloping algebra and its Yangian factor}
The talk occupies P001290 and the beginning of P001291. It defines an enlargement of the enveloping algebra through bounded-degree, tail-invariant sequences acting on tensor representations, then identifies its central factor and Yangian factor. The exact 1987 Doklady note was presented by Gelfand on 5 May 1986, matching the manuscript heading and independently corroborating the date discrepancy.

\begin{GSitem}{1986-677.1}{A}{The strong-limit completion on tame tensor representations}
\GSfield{Kind; source}{construction; \GSpage{P001290}}
\GSfield{The notes}{Sequences in \(U(\mathfrak o(n))\) have bounded degrees, commute with a fixed tail orthogonal subgroup, and converge strongly in every tensor representation; sequences are identified by their limits.}
\GSfield{Statement}{Fix \(m\). Consider uniformly bounded PBW-degree sequences in the tail centralizers whose actions have strong limits in every specified tame tensor representation of the infinite classical group. Quotient by equality of all these representation limits; componentwise operations give an algebra \(A_m\), and \(A=\bigcup_mA_m\). The representation family and the bounded-degree condition are part of the definition, not dispensable analytic decoration.}
\GSfield{Reconstruction}{The exact tame representation class from the paper supplies the unspecified universal quantifier over representations.}
\GSfield{Fate}{\citet[\S1, Theorems 1--2]{olshanskii1987extension} gives precisely this strong-limit construction and an equivalent bounded-degree projective-limit construction using \(\mathfrak{gl}_n\) centralizers. It was presented 5 May 1986, received 22 May 1986, published in Russian in 1987, and translated in 1988.}
\GSfield{Present status}{As of November 2026, the corrected formulation is an established result; no unresolved claim is inferred from this occurrence.}
\GSfield{Assessment}{literal statement: correct with missing hypotheses; publication match: exact; verification: read the result; development: none; confidence: high}
\end{GSitem}
\begin{GSitem}{1986-677.2}{A}{The polynomial central factor and the ordinary Yangian}
\GSfield{Kind; source}{theorem; \GSpage{P001290-P001291}}
\GSfield{The notes}{The stable algebra with \(m\) distinguished coordinates is written as a polynomial algebra in central elements tensored with \(Y(\mathfrak{gl}_m)\).}
\GSfield{Statement}{For the preceding strong-limit completion, \(A_0=Z(A)\) is a polynomial algebra in the stable shifted power sums, and \(A_m\simeq A_0\otimes Y(\mathfrak{gl}_m)\). This is the ordinary type-A Yangian, despite the initial orthogonal-group description. It is not the different projective-limit orthogonal centralizer construction that yields twisted Yangians.}
\GSfield{Reconstruction}{The completion is explicitly tied to 1986-677.1 to prevent a false type correction.}
\GSfield{Fate}{\citet[Theorems 3 and 10]{olshanskii1987extension} states both assertions exactly. The 1987 paper, with its 1986 presentation date, settles the identification directly; the later twisted-Yangian framework must not be substituted for its completion.}
\GSfield{Present status}{As of November 2026, the corrected formulation is an established result; no unresolved claim is inferred from this occurrence.}
\GSfield{Assessment}{literal statement: correct; publication match: exact; verification: read the result; development: none; confidence: high}
\end{GSitem}
\begin{GSitem}{1986-677.3}{A}{Orthogonal, symplectic, and unitary versions}
\GSfield{Kind; source}{theorem; \GSpage{P001291}}
\GSfield{The notes}{The same algebra \(A\) arises for orthogonal and symplectic infinite groups, while the unitary case gives \(A\otimes A\).}
\GSfield{Statement}{In the tame strong-limit enlargement described in the cited paper, the \(SO(\infty)\) and \(Sp(\infty)\) constructions give the same abstract algebra \(A\); the corresponding \(U(\infty)\) enlargement has two factors. This comparison concerns that representation-theoretic completion, not an isomorphism between the finite-rank Lie algebras or their ordinary enveloping algebras.}
\GSfield{Reconstruction}{The domain of the comparison is restored.}
\GSfield{Fate}{The assertion appears explicitly in the opening paragraph of \citet{olshanskii1987extension}. The exact primary text confirms what might otherwise look like a transcription error.}
\GSfield{Present status}{As of November 2026, the corrected formulation is an established result; no unresolved claim is inferred from this occurrence.}
\GSfield{Assessment}{literal statement: correct with missing hypotheses; publication match: exact; verification: read the result; development: none; confidence: high}
\end{GSitem}
\begin{GSitem}{1986-677.4}{A}{The small-m cases as a test of the construction}
\GSfield{Kind; source}{method; \GSpage{P001291}; I. M. Gelfand}
\GSfield{The notes}{Gelfand asks first for \(m=0\) and \(m=1\).}
\GSfield{Statement}{The \(m=0\) case isolates the stable central algebra. In the ordinary RTT presentation, \(Y(\mathfrak{gl}_1)\) is commutative, so \(m=1\) tests the first nontrivial centralizer without matrix commutator complications. This is a request for explanatory base cases, not an unresolved classification.}
\GSfield{Reconstruction}{The two base cases are interpreted within the published completion.}
\GSfield{Fate}{The central description and factorization are \citet[Theorems 3 and 10]{olshanskii1987extension}; the rank-one observation follows from the displayed scalar RTT relation.}
\GSfield{Present status}{As of November 2026, the corrected formulation is an established result; no unresolved claim is inferred from this occurrence.}
\GSfield{Assessment}{literal statement: correct with missing hypotheses; publication match: exact; verification: read the result; development: none; confidence: high}
\end{GSitem}

\GStalk{678}{1987-05-05}{I. V. Cherednik}{P001290-P001294}{Highest weights and quantum characteristic polynomials}
The continuation starts midway through P001291 and runs to the end of P001292. It evaluates the quantum determinant in a highest-weight representation and compares it with ordinary determinants, including a two-dimensional antisymmetrizer calculation. The later year heading and September material on the catalogue range are separate records.

\begin{GSitem}{1986-678.1}{E}{Polynomial highest weights and the highest-weight line}
\GSfield{Kind; source}{definition; \GSpage{P001291}}
\GSfield{The notes}{A representation is labeled by \(m_1\geq\cdots\geq m_N\geq0\), and a vector invariant under upper triangular unipotents is described as unique.}
\GSfield{Statement}{Finite-dimensional polynomial irreducible \(GL_N(\mathbb C)\) representations are indexed by partitions with at most \(N\) parts. The upper-unipotent invariants form a one-dimensional highest-weight space. A nonzero vector in that space is unique only up to multiplication by a scalar; a literal unique vector requires a normalization.}
\GSfield{Reconstruction}{“Unique” is replaced by “unique line.”}
\GSfield{Fate}{This is the classical highest-weight setup used in the central-character calculation of \citet[\S2.11]{molevnazarovolshanski1996yangians}.}
\GSfield{Present status}{As of November 2026, the corrected formulation is an established result; no unresolved claim is inferred from this occurrence.}
\GSfield{Assessment}{literal statement: correct with missing hypotheses; publication match: exact; verification: read the result; development: none; confidence: high}
\end{GSitem}
\begin{GSitem}{1986-678.2}{A}{The quantum determinant on a highest-weight module}
\GSfield{Kind; source}{theorem; \GSpage{P001291-P001292}}
\GSfield{The notes}{The central series acts by \(\prod_i(\lambda+N+m_i-i)/(\lambda+N-i)\).}
\GSfield{Statement}{For the level-one evaluation matrix with the notebook\textquotesingle s transpose and positive-shift convention, the normalized quantum determinant acts on highest weight \((m_1,\ldots,m_N)\) by \[\prod_{i=1}^N\frac{\lambda+N-i+m_i}{\lambda+N-i}.\] Removing the denominator gives the shifted characteristic polynomial; the shifts are essential and are not the ordinary determinant of a matrix of commuting entries.}
\GSfield{Reconstruction}{The evaluation, normalization, and weight convention are explicit.}
\GSfield{Fate}{The evaluation of the central quantum determinant and its Harish--Chandra interpretation appear in \citet[\S2.11]{molevnazarovolshanski1996yangians}, with the corresponding conventional change of spectral variable. The use in Gelfand--Tsetlin bases belongs to \citet{cherednik1987gelfand}.}
\GSfield{Present status}{As of November 2026, the corrected formulation is an established result; no unresolved claim is inferred from this occurrence.}
\GSfield{Assessment}{literal statement: correct with missing hypotheses; publication match: exact; verification: read the result; development: none; confidence: high}
\end{GSitem}
\begin{GSitem}{1986-678.3}{E}{The scalar of a two-factor antisymmetrization}
\GSfield{Kind; source}{claim; \GSpage{P001292}}
\GSfield{The notes}{For two \(2\times2\) matrices \(A,B\), the compressed operator \(\mathcal P_-(A\otimes B)\mathcal P_-\) is computed, but the displayed scalar has an inconsistent factor.}
\GSfield{Statement}{With the normalized projector \(\mathcal P_-=(I-P)/2\), the scalar on \(\bigwedge^2\mathbb C^2\) is \[\tfrac12(a_{11}b_{22}+a_{22}b_{11}-a_{12}b_{21}-a_{21}b_{12}).\] In particular \(A=B=I\) gives scalar \(1\). A factor \(1/4\), or no factor at all, cannot satisfy this normalization.}
\GSfield{Reconstruction}{\oursComputed\ The projector normalization determines the factor. The explicit check is recorded as editorial verification.}
\GSfield{Fate}{The compression is the rank-two instance of the determinant construction in \citet[\S2]{molevnazarovolshanski1996yangians}. The scalar normalization is independently checked by the identity-matrix substitution; the notebook formula is emended, not retranscribed.}
\GSfield{Present status}{As of November 2026, the corrected formulation is an established result; no unresolved claim is inferred from this occurrence.}
\GSfield{Assessment}{literal statement: false as written; publication match: partial; verification: read the result; development: annotated; confidence: high}
\end{GSitem}
\begin{GSitem}{1986-678.4}{E}{Capelli-type central coefficients and the classical limit}
\GSfield{Kind; source}{construction; \GSpage{P001292}}
\GSfield{The notes}{The coefficients of the quantum determinant are central, and their relation to familiar determinant or trace expressions is discussed.}
\GSfield{Statement}{In \(U(\mathfrak{gl}_N)\), the coefficients of the shifted quantum characteristic polynomial are central and algebraically independent; their leading PBW symbols are the elementary invariant polynomials on matrices. Other central generators, such as properly interpreted power-sum invariants, are polynomially related to them. Unshifted determinants of noncommuting matrix entries are not automatically central.}
\GSfield{Reconstruction}{The associated-graded meaning of the classical determinant is supplied.}
\GSfield{Fate}{\citet[\S2.11 and the comments in \S2.24]{molevnazarovolshanski1996yangians} identifies the Capelli and Harish--Chandra formulas and their classical antecedents. This is established invariant theory, with the Yangian packaging emphasized in Cherednik\textquotesingle s talk.}
\GSfield{Present status}{As of November 2026, the corrected formulation is an established result; no unresolved claim is inferred from this occurrence.}
\GSfield{Assessment}{literal statement: correct with missing hypotheses; publication match: exact; verification: read the result; development: none; confidence: high}
\end{GSitem}

\GStalk{679}{}{}{P001292-P001294}{1986}
This catalogue record is a year heading, not a separate mathematical talk. The preceding page ends the May quantum-determinant discussion, while P001293 begins the explicitly dated 1 September 1986 session. The overlapping catalogue range does not transfer either contribution to this heading.


\GStalk{680}{1986-09-01}{I. M. Gelfand}{P001292-P001296}{Matroids, Coxeter groups, and quadratic duality}
The actual contribution begins on P001293. Gelfand discusses work with Serganova on matroids and torus strata, explains the Coxeter-group extension, and then presents quadratic algebras in connection with Beilinson and Ginzburg. The agenda also names Cherednik, but it contains no separate Gelfand--Tsetlin result on these pages. Several elementary examples retain useful manuscript errors: intersection-preserving closure, two-generator finiteness, and the distinction between the two groups of order 120.

\begin{GSitem}{1986-680.1}{E}{Matroid closure and the exchange axiom}
\GSfield{Kind; source}{definition; \GSpage{P001293}}
\GSfield{The notes}{The closure operation is extensive and idempotent, and an exchange axiom is inserted above a deleted line. The first axiom instead asserts \(\overline{A\cap B}=\bar A\cap\bar B\).}
\GSfield{Statement}{A finite matroid closure \(\operatorname{cl}:2^E\to2^E\) is extensive, monotone, and idempotent, and satisfies exchange: if \(b\in\operatorname{cl}(A\cup\{a\})\setminus\operatorname{cl}(A)\), then \(a\in\operatorname{cl}(A\cup\{b\})\). Equality between the closure of an intersection and the intersection of closures is not an axiom and fails for ordinary vector matroids. The condition \(\operatorname{cl}(\varnothing)=\varnothing\) means loopless; it does not exclude parallel elements.}
\GSfield{Reconstruction}{The exchange line is preserved, while monotonicity replaces the incorrect intersection identity.}
\GSfield{Fate}{The classical closure and independence formulations are reviewed in \citet{oxleybriefly}. The notebook has substituted a stronger, false identity for monotonicity; its later linear-span example identifies the intended definition.}
\GSfield{Present status}{As of November 2026, the corrected formulation is an established result; no unresolved claim is inferred from this occurrence.}
\GSfield{Assessment}{literal statement: false as written; publication match: exact; verification: read the result; development: none; confidence: high}
\end{GSitem}
\begin{GSitem}{1986-680.2}{E}{Vector matroids and closed singletons}
\GSfield{Kind; source}{construction; \GSpage{P001293}}
\GSfield{The notes}{For vectors \(e_1,\ldots,e_k\), closure means membership in their linear span. The notes also write \(\bar e_i=e_i\).}
\GSfield{Statement}{A finite indexed vector configuration over a field defines a matroid on its index set by linear independence, with \(\operatorname{cl}(A)=\{i:e_i\in\operatorname{span}(e_j:j\in A)\}\). A zero vector is a loop; proportional nonzero vectors give parallel elements. For a loopless representation, every singleton is closed precisely when there are no parallel pairs; together these conditions mean simplicity. The labels, rather than the possibly repeated vectors as an unindexed set, form the ground set.}
\GSfield{Reconstruction}{The extra simplicity hypothesis needed by the singleton assertion is separated from looplessness.}
\GSfield{Fate}{This is the standard representable-matroid construction in \citet{oxleybriefly}, and the starting case of \citet[\S2]{gelfandserganova1987combinatorial}.}
\GSfield{Present status}{As of November 2026, the corrected formulation is an established result; no unresolved claim is inferred from this occurrence.}
\GSfield{Assessment}{literal statement: correct with missing hypotheses; publication match: exact; verification: read the result; development: none; confidence: high}
\end{GSitem}
\begin{GSitem}{1986-680.3}{A}{Grassmannian strata determined by Pluecker support}
\GSfield{Kind; source}{theorem; \GSpage{P001293-P001294}}
\GSfield{The notes}{The Grassmannian is partitioned into strata by grouping similar diagonal-torus orbits, and a matroid determines each stratum.}
\GSfield{Statement}{For \(U\in\operatorname{Gr}(k,n)\) over a fixed field, let \(\mathcal B(U)\) be the \(k\)-subsets whose Pluecker coordinates are nonzero. These are the bases of a rank-\(k\) representable matroid. The locally closed loci with fixed \(\mathcal B(U)\) are the matroid strata. A stratum generally contains many torus orbits; arbitrary abstract matroids need not yield nonempty strata over the chosen field.}
\GSfield{Reconstruction}{Fixed field, nonzero Pluecker support, and the distinction between a stratum and an orbit are explicit.}
\GSfield{Fate}{\citet[\S\S1--2]{gelfandserganova1987combinatorial}, received 30 July 1986 and published in 1987, develops precisely the common-refinement/torus-stratum interpretation. The notebook names Goresky and MacPherson as collaborators on the Grassmannian result; this attribution is retained without treating each stratum as one orbit.}
\GSfield{Present status}{As of November 2026, the corrected formulation is an established result; no unresolved claim is inferred from this occurrence.}
\GSfield{Assessment}{literal statement: correct with missing hypotheses; publication match: exact; verification: read the result; development: none; confidence: high}
\end{GSitem}
\begin{GSitem}{1986-680.4}{E}{Coxeter presentations and the rank-two finiteness question}
\GSfield{Kind; source}{definition; \GSpage{P001294}}
\GSfield{The notes}{The generators satisfy \(s_i^2=1\) and \((s_is_j)^{m_{ij}}=1\). The two-generator example says the group can already be infinite although \(m_{12}\) is specified as an integer.}
\GSfield{Statement}{A Coxeter matrix has \(m_{ii}=1\) and \(m_{ij}=m_{ji}\in\{2,3,\ldots\}\cup\{\infty\}\) for \(i\ne j\); \(m_{ij}=\infty\) means no product-order relation. The rank-two group is dihedral of order \(2m_{12}\) when \(m_{12}<\infty\), and infinite dihedral when \(m_{12}=\infty\). Thus the written finite-integer example cannot be infinite.}
\GSfield{Reconstruction}{The allowed value \(\infty\) restores the intended infinite example.}
\GSfield{Fate}{\citet[Introduction and \S8]{coxeter1934discrete} gives the finite dihedral presentation and the general reflection-group presentation. The elementary finiteness question is settled, not a candidate open problem.}
\GSfield{Present status}{As of November 2026, the corrected formulation is an established result; no unresolved claim is inferred from this occurrence.}
\GSfield{Assessment}{literal statement: false as written; publication match: exact; verification: read the result; development: none; confidence: high}
\end{GSitem}
\begin{GSitem}{1986-680.5}{E}{Icosahedral reflections, the binary group, and the 600-cell}
\GSfield{Kind; source}{remark; \GSpage{P001294}}
\GSfield{The notes}{The notes pass from \(H_3\) of order 120 to the double cover of the order-60 rotation group in \(SU(2)\), then to 120 vertices in \(S^3\) and \(H_4\) of order \(120^2\).}
\GSfield{Statement}{The full icosahedral reflection group \(H_3\) has order 120. The rotational subgroup has order 60; its inverse image in \(SU(2)\) is the binary icosahedral group, also of order 120 but a different group. The resulting quaternionic vertices form the 600-cell, whose 600 facets are tetrahedra and whose full reflection symmetry group \(H_4\) has order 14400. The Coxeter generators are reflections, not rotations about vertices.}
\GSfield{Reconstruction}{The two order-120 groups and the dimension of the tetrahedral facets are distinguished.}
\GSfield{Fate}{The finite reflection groups and their orders occur in \citet[\S9 and the table on p.~618]{coxeter1934discrete}. The notebook changes from a reflection group to a binary rotation group without marking the change; the English word “faces” must be understood as three-dimensional facets here.}
\GSfield{Present status}{As of November 2026, the corrected formulation is an established result; no unresolved claim is inferred from this occurrence.}
\GSfield{Assessment}{literal statement: correct with missing hypotheses; publication match: partial; verification: read the result; development: none; confidence: high}
\end{GSitem}
\begin{GSitem}{1986-680.6}{A}{Bruhat order and Coxeter matroids}
\GSfield{Kind; source}{definition; \GSpage{P001294}}
\GSfield{The notes}{A reduced decomposition of \(a\) occurs inside a reduced decomposition of \(b\), defining Bruhat order; \(a<^w b\) means \(w^{-1}a<w^{-1}b\). A finite subset is a flag matroid when it has a minimum for every such order.}
\GSfield{Statement}{For a finite Coxeter group \(W\), Bruhat order is characterized by a reduced expression occurring as a subsequence, not necessarily a contiguous substring. A nonempty subset of \(W\), or more generally of \(W/W_J\) for a standard parabolic subgroup, is a Coxeter matroid when it has a unique minimum in each translated Bruhat order. In the quotient case use minimal coset representatives and the induced quotient order. A merely minimal element in each finite poset would impose no meaningful condition.}
\GSfield{Reconstruction}{“Minimum” is read as least element; the subword criterion and parabolic quotient are made explicit.}
\GSfield{Fate}{\citet[\S8]{gelfandserganova1987combinatorial} introduces the Weyl/Coxeter-group generalization and its translated-order characterization. The construction was submitted before the September 1986 talk and published in 1987.}
\GSfield{Present status}{As of November 2026, the corrected formulation is an established result; no unresolved claim is inferred from this occurrence.}
\GSfield{Assessment}{literal statement: correct with missing hypotheses; publication match: exact; verification: read the result; development: none; confidence: high}
\end{GSitem}
\begin{GSitem}{1986-680.7}{E}{The bases determine a finite matroid}
\GSfield{Kind; source}{theorem; \GSpage{P001295}}
\GSfield{The notes}{The vector matroid is completely determined by its maximal linearly independent subsets.}
\GSfield{Statement}{For any finite matroid, the independent subsets are exactly the subsets of its bases, so its collection of bases determines its independence, rank, and closure data. Conversely a nonempty family of subsets of one ground set is a matroid basis family precisely when it satisfies the basis-exchange axiom. No realizability hypothesis is needed for this reconstruction.}
\GSfield{Reconstruction}{The extension from vector matroids to arbitrary finite matroids is stated as background, not attributed as a new seminar theorem.}
\GSfield{Fate}{The equivalence of basis and independence axioms is standard and stated in \citet{oxleybriefly}. The notebook explains the representable example, and the same assertion holds abstractly.}
\GSfield{Present status}{As of November 2026, the corrected formulation is an established result; no unresolved claim is inferred from this occurrence.}
\GSfield{Assessment}{literal statement: correct; publication match: exact; verification: read the result; development: none; confidence: high}
\end{GSitem}
\begin{GSitem}{1986-680.8}{E}{The seven-point Fano example and dependence on the field}
\GSfield{Kind; source}{claim; \GSpage{P001295}}
\GSfield{The notes}{A seven-point diagram and the seven triples \(123,165,345,275,174,376,246\) accompany the assertion that not every matroid is vector-representable.}
\GSfield{Statement}{Those triples are the lines of the Fano matroid \(F_7\). It is representable over a field exactly when that field has characteristic two. Thus it supplies a counterexample over \(\mathbb R\) or \(\mathbb C\), the preceding setting, but does not show nonrepresentability over every field. The listed triples identify the matroid more securely than the partly interrupted diagram.}
\GSfield{Reconstruction}{The field is restored and the combinatorial identification uses the written triples.}
\GSfield{Fate}{\citet[Theorem 4.1]{oxleybriefly} states the exact characteristic-two criterion. The diagram is therefore a settled field-dependent example; no universal nonrepresentability claim is assigned to it.}
\GSfield{Present status}{As of November 2026, the corrected formulation is an established result; no unresolved claim is inferred from this occurrence.}
\GSfield{Assessment}{literal statement: correct with missing hypotheses; publication match: exact; verification: read the result; development: none; confidence: high}
\end{GSitem}
\begin{GSitem}{1986-680.9}{E}{Quadratic algebras and their quadratic duals}
\GSfield{Kind; source}{definition; \GSpage{P001295}}
\GSfield{The notes}{The work of Beilinson and Ginzburg is introduced by \(A=T(V)/(W)\), with \(W\subset V\otimes V\), and \(A^!=T(V^*)/(W^\perp)\).}
\GSfield{Statement}{For a finite-dimensional vector space \(V\) over a field \(k\) and \(R\subset V\otimes V\), a connected quadratic algebra is \(A=T_k(V)/(R)\), where \((R)\) is the two-sided homogeneous ideal generated in degree two. Its quadratic dual is \(A^!=T_k(V^*)/(R^\perp)\) under the ordinary tensor pairing. The equality of this algebra with an Ext algebra requires Koszulness and is not part of the definition.}
\GSfield{Reconstruction}{The distinction between quadratic duality and the Koszul Ext identification prevents overattribution.}
\GSfield{Fate}{The construction predates the talk: \citet[\S2, Theorem 2.5]{priddy1970koszul} treats the cohomology relations for Koszul algebras. The exact modern convention is stated in \citet[\S2]{beilinsonginzburgsoergel1996koszul}. The notes contain definitions, not enough data to identify a particular new representation-theoretic duality theorem.}
\GSfield{Present status}{As of November 2026, the corrected formulation is an established result; no unresolved claim is inferred from this occurrence.}
\GSfield{Assessment}{literal statement: correct; publication match: exact; verification: read the result; development: none; confidence: high}
\end{GSitem}
\begin{GSitem}{1986-680.10}{E}{Symmetric and exterior algebras as quadratic duals}
\GSfield{Kind; source}{construction; \GSpage{P001295}}
\GSfield{The notes}{The polynomial algebra has commutator relations and its dual is the Grassmann algebra; one displayed commutator uses mismatched indices.}
\GSfield{Statement}{Over a field of characteristic different from two, the annihilator of the alternating relation space of \(S(V)\) is the symmetric relation space in \(V^*\otimes V^*\), giving \(S(V)^!=\bigwedge V^*\); dually \((\bigwedge V)^!=S(V^*)\). The exterior relations include the squares \(y_i^2=0\), as well as pairwise anticommutation. They cannot be omitted when merely writing the off-diagonal commutator pairing.}
\GSfield{Reconstruction}{The displayed first product should be \(x_ix_k\), and diagonal relations are included. The characteristic restriction avoids a misleading symmetric-tensor shorthand in characteristic two.}
\GSfield{Fate}{\citet[\S2.2(2)]{priddy1970koszul} records the polynomial/exterior dual cohomology calculation. The notebook is an elementary illustration of the preceding definition.}
\GSfield{Present status}{As of November 2026, the corrected formulation is an established result; no unresolved claim is inferred from this occurrence.}
\GSfield{Assessment}{literal statement: correct with missing hypotheses; publication match: exact; verification: read the result; development: none; confidence: high}
\end{GSitem}
\begin{GSitem}{1986-680.11}{E}{A quadratic hypersurface and its homogenized Clifford dual}
\GSfield{Kind; source}{construction; \GSpage{P001296}}
\GSfield{The notes}{For \(A=S(V)/(\sum x_i^2)\), the orthogonal relation space consists of symmetric trace-zero matrices. The resulting relations are \(y_iy_j+y_jy_i=2\delta_{ij}h\), with \(h=y_i^2\) central.}
\GSfield{Statement}{Let \(k\) have characteristic different from two and \(V\) have basis \(x_1,\ldots,x_d\). The quadratic dual of \(S(V)/(x_1^2+\cdots+x_d^2)\) has generators \(y_i\) of degree one, off-diagonal anticommutation, and equal squares. Equivalently it is the graded Clifford algebra with a central degree-two parameter \(h\) and relations \(y_iy_j+y_jy_i=2\delta_{ij}h\). Specializing \(h=1\) gives the usual Clifford algebra; \(h\) is not set to one in the graded dual.}
\GSfield{Reconstruction}{The ground field and degree of \(h\) are specified; the summand \(a_{ij}x_ix_k\) is read as \(a_{ij}x_ix_j\).}
\GSfield{Fate}{This is a direct example of the quadratic-dual definition in \citet[\S2]{beilinsonginzburgsoergel1996koszul}, already computed in the manuscript. No later new theorem is needed to settle it. The claim of centrality is part of the notebook calculation and follows from the equality of all squares.}
\GSfield{Present status}{As of November 2026, the corrected formulation is an established result; no unresolved claim is inferred from this occurrence.}
\GSfield{Assessment}{literal statement: correct with missing hypotheses; publication match: partial; verification: read the result; development: none; confidence: high}
\end{GSitem}

\GStalk{681}{1986-09-29}{V. L. Ginzburg}{P001297-P001299}{Complete symmetric varieties and polynomial quotients of affine Hecke algebras}
Ginzburg presents the De Concini--Procesi compactification of symmetric varieties, then a finite and affine Hecke-algebra construction. The top of P001298 records a problem attributed explicitly to Cherednik: replace the evaluation of the last lattice generator by a polynomial relation. Ovsienko and Kazakov start separate contributions lower on that page and are indexed with their own records. The catalogue initial is preserved although the mathematical attribution to the speaker known as Victor Ginzburg is conventional.

\begin{GSitem}{1986-681.1}{E}{The wonderful compactification of an adjoint symmetric space}
\GSfield{Kind; source}{construction; \GSpage{P001297}}
\GSfield{The notes}{The space \(M=G/H\), with \(H\) a normalizer, has a minimal compactification \(\bar M\). The abbreviation for \(K\) is “compact subgroup”, and \(H/K\) is written as an elementary two-group.}
\GSfield{Statement}{Let \(G\) be an adjoint semisimple complex algebraic group and \(\sigma\) an involution. For the appropriate normalizer of the connected fixed subgroup, the symmetric homogeneous space admits the De Concini--Procesi wonderful projective compactification. The construction uses an algebraic symmetric subgroup; a literal arbitrary compact subgroup of a complex group does not meet these hypotheses. Finite central isogenies and the choice between connected fixed subgroup, full fixed group, and its normalizer must be specified before asserting a quotient formula.}
\GSfield{Reconstruction}{The algebraic symmetric-space hypotheses are reconstructed. The unqualified formula for \(H/K\) is not certified for arbitrary \(K\).}
\GSfield{Fate}{\citet{deconciniprocesi1983complete}, published in 1983 from the 1982 Montecatini setting, is the original construction. Its exact adjoint-group setup is recalled in \citet[\S1]{chirivimaffei2004projective}. The notebook therefore presents established geometry; the isolated abbreviation does not determine an alternative real-form theorem.}
\GSfield{Present status}{As of November 2026, the corrected formulation is an established result; no unresolved claim is inferred from this occurrence.}
\GSfield{Assessment}{literal statement: correct with missing hypotheses; publication match: partial; verification: read the result; development: none; confidence: medium}
\end{GSitem}
\begin{GSitem}{1986-681.2}{E}{Boundary divisors and the orbit stratification}
\GSfield{Kind; source}{theorem; \GSpage{P001297}}
\GSfield{The notes}{The compactification is smooth, has finitely many orbits, its boundary has normal crossings, and there is a unique closed orbit \(G/P\).}
\GSfield{Statement}{For the preceding wonderful compactification of rank \(r\), the boundary consists of \(r\) smooth \(G\)-stable divisors meeting with simple normal crossings. Orbit closures are precisely their partial intersections, including the empty intersection \(X\); the total intersection is the unique closed orbit, a partial flag variety \(G/P\). At a point on exactly \(k\) boundary components, suitable local coordinates give their union by \(z_1\cdots z_k=0\), and their intersection by \(z_1=\cdots=z_k=0\). The codimension of the entire boundary is one, not \(k\).}
\GSfield{Reconstruction}{The union/intersection distinction repairs the notebook parenthesis defining normal crossings.}
\GSfield{Fate}{These are \citet[Theorem 3.1]{deconciniprocesi1983complete}, reproduced explicitly as \citet[Proposition 1.1]{chirivimaffei2004projective}. The latter primary text was read directly.}
\GSfield{Present status}{As of November 2026, the corrected formulation is an established result; no unresolved claim is inferred from this occurrence.}
\GSfield{Assessment}{literal statement: correct with missing hypotheses; publication match: exact; verification: read the result; development: none; confidence: high}
\end{GSitem}
\begin{GSitem}{1986-681.3}{E}{Restriction of line bundles to the closed orbit}
\GSfield{Kind; source}{claim; \GSpage{P001297}}
\GSfield{The notes}{The last numbered property writes an isomorphism between line bundles on \(\bar M\) and those on \(G/P\).}
\GSfield{Statement}{Restriction gives an injection \(\operatorname{Pic}(X)\hookrightarrow\operatorname{Pic}(G/P)\); the image is a specified sublattice of the weight lattice. Surjectivity onto all line bundles of the closed orbit does not hold in general. The corrected assertion determines a line bundle on \(X\) uniquely by its restriction, without claiming that every line bundle on the closed orbit extends.}
\GSfield{Reconstruction}{The arrow is emended from isomorphism to injection.}
\GSfield{Fate}{The precise theorem is \citet[Proposition 8.1]{deconciniprocesi1983complete}, restated as \citet[Propositions 1.2--1.4]{chirivimaffei2004projective}; these propositions distinguish the injective restriction from its image. Thus the notebook isomorphism is an overstatement.}
\GSfield{Present status}{As of November 2026, the corrected formulation is an established result; no unresolved claim is inferred from this occurrence.}
\GSfield{Assessment}{literal statement: false as written; publication match: exact; verification: read the result; development: none; confidence: high}
\end{GSitem}
\begin{GSitem}{1986-681.4}{E}{The finite Hecke algebra and its commuting generators}
\GSfield{Kind; source}{definition; \GSpage{P001297}}
\GSfield{The notes}{The generators satisfy the quadratic and braid relations, with commuting condition \(j\ne i,i+1\).}
\GSfield{Statement}{For \(q\ne0\), the type-A Hecke algebra has generators \(T_1,\ldots,T_{n-1}\), relations \((T_i+1)(T_i-q)=0\), the adjacent braid relation, and \(T_iT_j=T_jT_i\) for \(|i-j|>1\). It has basis \(T_w\) indexed by \(S_n\), and at \(q=1\) specializes to the symmetric-group algebra. The written condition incorrectly permits commuting with \(T_{i-1}\).}
\GSfield{Reconstruction}{The generator normalization is explicitly matched rather than silently changing the quadratic relation.}
\GSfield{Fate}{The finite type-A case is included in \citet[\S2]{mathas2004matrix}, whose generators use the alternate convention \(T_i\mapsto-T_i\). The notebook relation requires the absolute-distance condition.}
\GSfield{Present status}{As of November 2026, the corrected formulation is an established result; no unresolved claim is inferred from this occurrence.}
\GSfield{Assessment}{literal statement: false as written; publication match: exact; verification: read the result; development: none; confidence: high}
\end{GSitem}
\begin{GSitem}{1986-681.5}{E}{Affine Hecke generators and the evaluation map}
\GSfield{Kind; source}{construction; \GSpage{P001297-P001298}}
\GSfield{The notes}{Commuting lattice generators \(x_i\) are adjoined, with \(x_iT_i=T_ix_{i+1}-(q-1)x_{i+1}\). The proposed evaluation has \(x_n\mapsto1\), \(x_{n-1}\mapsto qT_{n-1}^2\).}
\GSfield{Statement}{For the affine Hecke algebra of \(GL_n\), use commuting invertible \(X_i\), the finite Hecke generators, and the Bernstein relations. With the written convention they imply \(T_iX_iT_i=qX_{i+1}\), together with \(T_iX_j=X_jT_i\) for \(j\ne i,i+1\). Thus evaluation \(X_n=1\) recursively gives \(X_i=qT_i^{-1}X_{i+1}T_i^{-1}\), in particular \(X_{n-1}=qT_{n-1}^{-2}\). The displayed positive square is incompatible with the displayed exchange relation.}
\GSfield{Reconstruction}{\oursComputed\ Invertibility and the nonadjacent Bernstein relation are supplied; the evaluation recursion is stated in the notebook convention.}
\GSfield{Fate}{The Jucys--Murphy elements in \citet[\S2]{mathas2004matrix} give the equivalent finite-algebra realization after reversing the indices and matching generator conventions. The inverse-square correction is an editorial consistency check of the notebook relations, not a new representation-theoretic theorem.}
\GSfield{Present status}{As of November 2026, the corrected formulation is an established result; no unresolved claim is inferred from this occurrence.}
\GSfield{Assessment}{literal statement: false as written; publication match: partial; verification: read the result; development: annotated; confidence: high}
\end{GSitem}
\begin{GSitem}{1986-681.6}{B}{Polynomial quotients of the affine Hecke algebra}
\GSfield{Kind; source}{problem; \GSpage{P001298, upper section}; I. V. Cherednik}
\GSfield{The notes}{Cherednik asks for the target algebra and homomorphism when \(x_n^2\mapsto1\), or more generally when a polynomial relation in \(x_n\) is imposed. The problem is marked unsolved.}
\GSfield{Statement}{Over \(\mathbb C\), let \(q\ne0\) and let \(p(z)=\prod_{a=1}^r(z-Q_a)\) be monic with every \(Q_a\ne0\). The quotient of the affine type-A Hecke algebra by the two-sided ideal generated by \(p(X_n)\) is a cyclotomic Hecke algebra of type \(G(r,1,n)\), after reversing the affine convention. It has dimension \(r^n n!\), with a PBW basis of bounded powers of its commuting Murphy elements times \(T_w\). Repeated roots and nonsemisimple parameters do not invalidate this basis theorem. For \(r=2\) it is the type-B Hecke case, including \(p(z)=z^2-1\).}
\GSfield{Reconstruction}{The polynomial is interpreted as a universal quotient relation, which gives a definite later solution to the question of its target. Nonzero roots match the invertible affine generators.}
\GSfield{Fate}{\citet{arikikoike1994hecke} and the independent cyclotomic construction of Broue--Malle supplied this later answer. The exact presentation and Ariki--Koike basis theorem are read in \citet[\S2, (2.2)]{mathas2004matrix}, first circulated in August 2001 and published in 2004. This identifies the universal quotient; it does not claim that every such quotient admits an evaluation back into the original \(n!\)-dimensional finite Hecke algebra.}
\GSfield{Present status}{As of 12 November 2026, the universal polynomial-quotient problem is settled by cyclotomic Hecke algebras and their basis theorem. No generic semisimplicity assumption is needed for this answer.}
\GSfield{Assessment}{literal statement: correct; publication match: exact; verification: read the result; development: none; confidence: high; problem \hyperref[prob:GP-1986-04]{GP-1986-04}}
\end{GSitem}

\GStalk{683}{1986-09-29}{V. Yu. Ovsienko}{P001297-P001304}{KdV and super-KdV as Euler equations}
The mathematical report occupies only the middle of P001298, between Ginzburg\textquotesingle s Hecke-algebra discussion and Kazakov\textquotesingle s random surfaces. It identifies KdV and its supersymmetric extensions with Euler equations for Virasoro-type groups; the catalogue\textquotesingle s wide page range does not transfer adjacent talks to Ovsienko.

\begin{GSitem}{1986-683.1}{A}{KdV as the Euler equation on the Virasoro group}
\GSfield{Kind; source}{theorem; \GSpage{P001298}}
\GSfield{The notes}{KdV is described as an Euler equation on the Virasoro group, extending the familiar Euler equation for a Lie group.}
\GSfield{Statement}{For the Virasoro central extension of vector fields on the circle, the Euler--Arnold equation of the right-invariant \(L^2\) metric is, up to choices of signs and scale, \(u_t+3uu_x+c u_{xxx}=0\). The fixed dual central coordinate produces dispersion; the metric and coadjoint convention are indispensable.}
\GSfield{Reconstruction}{The metric, fixed central coordinate, and coefficient convention are supplied rather than read from the one-line note.}
\GSfield{Fate}{\citet[Prop. 1]{ovsienkokhesin1987super} proves the Euler interpretation of KdV and then constructs the super-extension in Proposition 2. The manuscript was received 8 August 1986 and published in 1987; this is the directly matching near-contemporary publication.}
\GSfield{Present status}{As of 12 November 2026 this is an established result under the hypotheses in the modern statement.}
\GSfield{Assessment}{literal statement: correct with missing hypotheses; publication match: exact; verification: read the result; development: none; confidence: high}
\end{GSitem}
\begin{GSitem}{1986-683.2}{A}{Euler equations on the Neveu--Schwarz and Ramond superalgebras}
\GSfield{Kind; source}{construction; \GSpage{P001298}}
\GSfield{The notes}{The report adds the supersymmetric KdV equations and names the Schwarz and Ramond algebras.}
\GSfield{Statement}{The super Euler construction uses an inertia pairing with an inverse derivative on the odd component. In the Neveu--Schwarz case the antiperiodic odd modes have no zero mode, so this pairing is nondegenerate. In the Ramond case the periodic odd zero mode gives a degenerate pairing and must be treated accordingly. The component fields have their stated even and odd Grassmann parities; replacing odd fields by commuting scalar functions changes the equation.}
\GSfield{Reconstruction}{The name Neveu--Schwarz, the inverse-derivative inertia, and the different zero-mode behavior are supplied from the publication. The one-line notebook report preserves no explicit component system.}
\GSfield{Fate}{\citet[Prop. 2 and the subsequent Ramond remark]{ovsienkokhesin1987super}, received 8 August 1986 and published in 1987, gives the directly matching construction, including the degeneracy qualification.}
\GSfield{Present status}{As of 12 November 2026 this is an established result under the hypotheses in the modern statement.}
\GSfield{Assessment}{literal statement: correct with missing hypotheses; publication match: partial; verification: read the result; development: none; confidence: high}
\end{GSitem}

\GStalk{682}{1986-09-29}{V. A. Kazakov}{P001297-P001299}{Planar maps, matrix integrals, and two-dimensional gravity}
The catalogue range overlaps Ginzburg and Ovsienko. Kazakov actually begins in the lower part of P001298 and continues through P001302: triangulated spheres, the cubic Hermitian matrix model, its planar expansion, eigenvalue saddle, and critical exponent. The formal combinatorial integral must be distinguished from a convergent integral over real Hermitian matrices. Record 695 revisits the counting conventions.

\begin{GSitem}{1986-682.1}{E}{Triangulations as a discretization of the gravity sum}
\GSfield{Kind; source}{construction; \GSpage{P001298-P001299}}
\GSfield{The notes}{The notes replace an integral over metrics on the sphere by \(Z(\beta)=\sum_n\sum_{\Gamma_n}e^{-\beta n}\), allowing generalized triangulations, and say forbidding exceptional gluings does not change the answer.}
\GSfield{Statement}{For a specified finite class of connected oriented maps dual to cubic ribbon graphs, define a formal area-generating series with symmetry weights \(1/|\operatorname{Aut}\Gamma|\). Changes between unrestricted maps and simple triangulations alter coefficients and the critical coupling; in the pure-gravity universality class their leading susceptibility exponent agrees. This does not construct a measure on all smooth metrics.}
\GSfield{Reconstruction}{Orientation, connectedness, symmetry weights, and the distinction between exact coefficients and a universal exponent are supplied editorially.}
\GSfield{Fate}{\citet[Secs. 2 and 4]{brezinitzyksonparisizuber1978planar} already derive the planar generating series. The continuum gravity interpretation is developed in \citet{polyakov1981quantum}; no equality of measures is proved in these notes.}
\GSfield{Present status}{As of 12 November 2026 this is an established result under the hypotheses in the modern statement.}
\GSfield{Assessment}{literal statement: correct with missing hypotheses; publication match: partial; verification: read the result; development: none; confidence: high}
\end{GSitem}
\begin{GSitem}{1986-682.2}{E}{Formal cubic Hermitian matrix integral}
\GSfield{Kind; source}{construction; \GSpage{P001299}}
\GSfield{The notes}{\(I(g)=\int e^{-\operatorname{tr}(U^2+gU^3/\sqrt N)}dU\), followed by \(Z=\lim N^{-2}\log I(e^{-\beta/2})\).}
\GSfield{Statement}{Expand the cubic exponential formally against a normalized Gaussian measure and use \(N^{-2}\log[I(g)/I(0)]\) for connected vacuum graphs. For real nonzero \(g\), the displayed real integral diverges along one scalar direction. A convergent complex-contour definition requires additional data and is not implied by the formal power series.}
\GSfield{Reconstruction}{The formal interpretation and division by \(I(0)\) repair the displayed expression; they are not extra claims attributed to Kazakov.}
\GSfield{Fate}{\citet[Sec. 4, Eq. (42)]{brezinitzyksonparisizuber1978planar} explicitly states the instability of the cubic integral and uses its formal series. The notebook suppresses both this qualification and Gaussian normalization.}
\GSfield{Present status}{As of 12 November 2026 this is an established result under the hypotheses in the modern statement.}
\GSfield{Assessment}{literal statement: false as written; publication match: exact; verification: read the result; development: annotated; confidence: high}
\end{GSitem}
\begin{GSitem}{1986-682.3}{E}{Wick contractions and the genus power of N}
\GSfield{Kind; source}{method; \GSpage{P001299-P001300}}
\GSfield{The notes}{A six-factor Gaussian moment is illustrated by one contraction; double lines make closed index loops contribute \(N\), and a graph is assigned \(N^{F-V/2}=N^{2-2\chi}\), with \(\chi\) called its Euler number.}
\GSfield{Statement}{For a Gaussian Hermitian matrix with density proportional to \(e^{-\operatorname{tr}U^2/2}\), \(\mathbb E(U_{ij}U_{kl})=\delta_{il}\delta_{jk}\); every even moment sums over all pairings. A connected cubic ribbon graph has \(3V=2E\), so its power is \(N^{F-E+V}=N^{2-2g}\), where \(g\) is genus. With density \(e^{-\operatorname{tr}U^2}\), the covariance has an additional factor \(1/2\).}
\GSfield{Reconstruction}{The missing pairing sum, index order, Gaussian scaling, and genus/Euler-characteristic distinction are explicit editorial repairs.}
\GSfield{Fate}{The topological expansion is the pre-existing large-\(N\) method of \citet{thooft1974planar}; \citet[Sec. 2]{brezinitzyksonparisizuber1978planar} specializes it to matrix counting.}
\GSfield{Present status}{As of 12 November 2026 this is an established result under the hypotheses in the modern statement.}
\GSfield{Assessment}{literal statement: false as written; publication match: exact; verification: read the result; development: annotated; confidence: high}
\end{GSitem}
\begin{GSitem}{1986-682.4}{E}{Eigenvalue reduction and the Vandermonde square}
\GSfield{Kind; source}{method; \GSpage{P001300}}
\GSfield{The notes}{Writing \(U=\Omega\Lambda\Omega^{-1}\) changes the matrix integral into an eigenvalue integral with \(\prod_{i<j}(\lambda_i-\lambda_j)^2\).}
\GSfield{Statement}{Lebesgue measure on complex Hermitian matrices decomposes, away from repeated eigenvalues, as \(C_N\Delta(\lambda)^2\prod_i d\lambda_i\,d\bar\Omega\), with angular measure on \(U(N)/U(1)^N\); ordering and permutation factors are absorbed in \(C_N\). For invariant integrands the angular factor cancels in normalized ratios.}
\GSfield{Fate}{\citet[Sec. 3]{brezinitzyksonparisizuber1978planar} uses precisely this reduction. It is standard Weyl integration rather than a new matrix-gravity theorem.}
\GSfield{Present status}{As of 12 November 2026 this is an established result under the hypotheses in the modern statement.}
\GSfield{Assessment}{literal statement: correct with missing hypotheses; publication match: exact; verification: read the result; development: none; confidence: high}
\end{GSitem}
\begin{GSitem}{1986-682.5}{E}{Uniform equilibrium density in the circular ensemble}
\GSfield{Kind; source}{remark; \GSpage{P001300}}
\GSfield{The notes}{Dyson\textquotesingle s circular ensemble is mentioned as a case in which the limiting density is constant.}
\GSfield{Statement}{For Haar-distributed unitary matrices, the normalized one-point eigenangle density is constant on the circle; the corresponding large-\(N\) logarithmic-gas equilibrium measure is normalized Lebesgue measure. This concerns the circular ensemble with no external potential; the Hermitian cubic model has a nonconstant equilibrium density.}
\GSfield{Fate}{\citet[Sec. IV, Thm. 6; Sec. V, Thm. 8 and Eq. (79)]{dyson1962statistical} defines the Haar unitary ensemble and its circular eigenangle density. The original paper, received 15 September 1961 and published in 1962, was opened and these results read. The exact source used by the speaker is not identified.}
\GSfield{Present status}{As of 12 November 2026 this is an established result under the hypotheses in the modern statement.}
\GSfield{Assessment}{literal statement: correct with missing hypotheses; publication match: partial; verification: read the result; development: none; confidence: medium}
\end{GSitem}
\begin{GSitem}{1986-682.6}{E}{Planar saddle equation and one-cut resolvent}
\GSfield{Kind; source}{method; \GSpage{P001300-P001301}}
\GSfield{The notes}{The rescaled eigenvalues satisfy \(2\lambda_i+3g\lambda_i^2=2\sum_{j\ne i}(\lambda_i-\lambda_j)^{-1}\). A density of total mass two and an analytic resolvent are introduced.}
\GSfield{Statement}{After \(U=\sqrt N M\), the saddle equation is \(V^{\prime}(\lambda_i)=\frac2N\sum_{j\ne i}(\lambda_i-\lambda_j)^{-1}\). If \(\nu=2\rho\), then \(\int\nu=2\), \(V^{\prime}(x)=\operatorname{PV}\int\nu(y)/(x-y)\,dy\), and \(R(z)=\int\nu(y)/(z-y)\,dy\sim2/z\). On a one-cut branch the polynomial-times-square-root ansatz is fixed by these conditions. Boundary values satisfy \(R_+-R_-=-2\pi i\nu\).}
\GSfield{Reconstruction}{The \(1/N\), principal value, boundary-value factor \(i\), and specified formal branch are supplied.}
\GSfield{Fate}{\citet[Sec. 4, Eqs. (43)--(46)]{brezinitzyksonparisizuber1978planar} gives the cubic solution. The asymptotic condition selects the formal branch analytic near zero coupling; it does not justify an unstable real equilibrium problem.}
\GSfield{Present status}{As of 12 November 2026 this is an established result under the hypotheses in the modern statement.}
\GSfield{Assessment}{literal statement: false as written; publication match: exact; verification: read the result; development: annotated; confidence: high}
\end{GSitem}
\begin{GSitem}{1986-682.7}{E}{Exact cubic free-energy series and critical radius}
\GSfield{Kind; source}{theorem; \GSpage{P001302}}
\GSfield{The notes}{The notes give \(-\frac12\sum_{k\ge1}(72g^2)^k\Gamma(3k/2)/[\Gamma(k/2+1)(k+2)!]\), then write \(g_c=1/(108\sqrt3)\).}
\GSfield{Statement}{In the source normalization \(\operatorname{tr}(U^2/2+gU^3/\sqrt N)\), the negative planar logarithm has the displayed series and its nearest singularity in \(g^2\) is \(g_c^2=1/(108\sqrt3)\). The notebook\textquotesingle s earlier quadratic coefficient is different, so one must either restore \(U^2/2\) or rescale \(g\). The coefficients are weighted connected maps with \(2k\) cubic vertices.}
\GSfield{Reconstruction}{Normalization is compared to the original paper; the series cannot simultaneously use the notebook\textquotesingle s unrescaled quadratic term and the source\textquotesingle s critical constant.}
\GSfield{Fate}{\citet[Sec. 4, Eqs. (49)--(52)]{brezinitzyksonparisizuber1978planar} is an exact pre-1986 match, including the squared critical coupling.}
\GSfield{Present status}{As of 12 November 2026 this is an established result under the hypotheses in the modern statement.}
\GSfield{Assessment}{literal statement: false as written; publication match: exact; verification: read the result; development: annotated; confidence: high}
\end{GSitem}
\begin{GSitem}{1986-682.8}{E}{Pure-gravity singular exponent and map asymptotics}
\GSfield{Kind; source}{claim; \GSpage{P001302}}
\GSfield{The notes}{\(Z(\beta)\sim(\beta-\beta_c)^{5/2}\) and map coefficients have an \(n^{-7/2}\) factor.}
\GSfield{Statement}{The nonanalytic part of the unrooted connected planar free energy has exponent \(5/2\). Transfer to coefficients gives a \(k^{-7/2}\) power multiplying exponential growth, with vertex parity and coupling conventions retained. Analytic background terms are present; rooting a map changes the power by a factor proportional to its number of edges.}
\GSfield{Reconstruction}{The words nonanalytic part, weighted unrooted count, and parity qualification are editorially necessary.}
\GSfield{Fate}{The cubic solution in \citet[Sec. 4]{brezinitzyksonparisizuber1978planar} supplies this singularity. Record 695 clarifies why unlabelled integer counts are not the coefficients of the raw free energy.}
\GSfield{Present status}{As of 12 November 2026 this is an established result under the hypotheses in the modern statement.}
\GSfield{Assessment}{literal statement: correct with missing hypotheses; publication match: partial; verification: read the result; development: none; confidence: high}
\end{GSitem}
\begin{GSitem}{1986-682.9}{A}{Adding spins to the random surface}
\GSfield{Kind; source}{remark; \GSpage{P001302}}
\GSfield{The notes}{The final paragraph proposes putting a second kind of variable, including spins, on the triangulation.}
\GSfield{Statement}{The annealed partition function sums both over maps and over matter configurations on each map. An Ising field on a cubic dual graph can be encoded by two matrix species with a covariance coupling different spins.}
\GSfield{Fate}{Kazakov develops this construction in record 685 and in \citet{kazakov1986ising}; no independently specified conjecture occurs in this closing remark.}
\GSfield{Present status}{As of 12 November 2026 the Ising-on-random-map model is solved in the planar limit; this does not identify every possible added matter system.}
\GSfield{Assessment}{literal statement: correct with missing hypotheses; publication match: partial; verification: publisher metadata checked; development: none; confidence: high}
\end{GSitem}

\GStalk{684}{1986-09-29}{V. L. Ginzburg}{P001302-P001305}{Generalized Lie algebras: manuscript attribution to D. Gurevich}
The catalogue identifies Ginzburg on 29 September. The original heading on P001303 instead reads D. Gurevich, 13 October 1986, followed by generalized Lie algebras and groups. We retain the catalogue record but attribute the surviving mathematical items explicitly to D. I. Gurevich. They concern involutive Yang--Baxter operators, braided Lie identities, super and color examples, and braided commutators. Kazakov begins later on P001304.

\begin{GSitem}{1986-684.1}{E}{An involutive Yang--Baxter symmetry}
\GSfield{Kind; source}{definition; \GSpage{P001303}; D. I. Gurevich}
\GSfield{The notes}{A linear map \(S:L\otimes L\to L\otimes L\) satisfies \(S^2=1\) and \(S_1S_2S_1=S_2S_1S_2\).}
\GSfield{Statement}{Let \(L\) be finite-dimensional over a characteristic-zero field. An involution \(S\) satisfying the braid equation determines representations of every symmetric group on \(L^{\otimes n}\), by acting in adjacent tensor factors; dropping involutivity gives braid-group representations.}
\GSfield{Fate}{\citet[pp. 797--799]{gurevich1986yangbaxter} introduced this generalized formal Lie-theoretic framework before the manuscript talk: the note was received 10 April 1985 and published in 1986.}
\GSfield{Present status}{As of 12 November 2026 this is an established result under the hypotheses in the modern statement.}
\GSfield{Assessment}{literal statement: correct; publication match: exact; verification: read the result; development: none; confidence: high}
\end{GSitem}
\begin{GSitem}{1986-684.2}{E}{Braided antisymmetry, Jacobi identity, and bracket covariance}
\GSfield{Kind; source}{definition; \GSpage{P001303}; D. I. Gurevich}
\GSfield{The notes}{Three identities for the bracket are given in right-acting operator notation: antisymmetry, a three-term Jacobi identity, and a compatibility relation with \(S\).}
\GSfield{Statement}{In a symmetric monoidal category with symmetry \(S\), a Lie object has a morphism \(b:L\otimes L\to L\) satisfying \(bS=-b\) and the categorical Jacobi identity. On ordinary vector spaces the further compatibility that \(b\) is a morphism for the specified tensor symmetry is essential; a bare pair consisting of any Yang--Baxter operator and any antisymmetric bracket is insufficient.}
\GSfield{Reconstruction}{The categorical formulation avoids silently reversing the notebook\textquotesingle s right-acting products; the bracket covariance is retained as a separate hypothesis.}
\GSfield{Fate}{\citet[pp. 798--800]{gurevich1986yangbaxter} states the corresponding generalized Lie identities. The notebook is a compressed exposition of that theory, not an independent new axiom system.}
\GSfield{Present status}{As of 12 November 2026 this is an established result under the hypotheses in the modern statement.}
\GSfield{Assessment}{literal statement: correct with missing hypotheses; publication match: exact; verification: read the result; development: none; confidence: high}
\end{GSitem}
\begin{GSitem}{1986-684.3}{E}{Ordinary and super Lie algebras as examples}
\GSfield{Kind; source}{construction; \GSpage{P001303-P001304}; D. I. Gurevich}
\GSfield{The notes}{The flip gives ordinary Lie algebras; a signed flip determined by degrees gives superalgebras.}
\GSfield{Statement}{For the ordinary flip \(S(x\otimes y)=y\otimes x\) the axioms recover Lie algebras. For a \(\mathbb Z/2\)-graded space, \(S(x\otimes y)=(-1)^{|x||y|}y\otimes x\), with a parity-preserving bracket, gives super antisymmetry and super Jacobi. The degree of a nonzero homogeneous bracket is \(|x|+|y|\) modulo two.}
\GSfield{Reconstruction}{Additivity of parity is a condition on homogeneous brackets, not a definition of the bracket of arbitrary nonhomogeneous vectors.}
\GSfield{Fate}{These are the foundational examples recalled in \citet[pp. 798--799]{gurevich1986yangbaxter}; Lie superalgebras were already established well before 1986.}
\GSfield{Present status}{As of 12 November 2026 this is an established result under the hypotheses in the modern statement.}
\GSfield{Assessment}{literal statement: correct with missing hypotheses; publication match: exact; verification: read the result; development: none; confidence: high}
\end{GSitem}
\begin{GSitem}{1986-684.4}{E}{Color symmetries and color brackets}
\GSfield{Kind; source}{construction; \GSpage{P001304}; D. I. Gurevich}
\GSfield{The notes}{\(S(t_i\otimes t_j)=a(i,j)t_j\otimes t_i\), with conditions on \(a\) said to follow from the axioms; the notes cite Trostel, Journal of Mathematical Physics 25(11), 1984.}
\GSfield{Statement}{For an abelian grading group \(G\), a bicharacter \(\varepsilon:G\times G\to k^\times\) satisfying \(\varepsilon(g,h)\varepsilon(h,g)=1\) defines such an involutive symmetry. A homogeneous bracket of degree zero then satisfies \(\varepsilon\)-skew symmetry and the corresponding color Jacobi identity. Arbitrary coefficients \(a(i,j)\) need compatibility with the degrees and nonzero structure constants.}
\GSfield{Reconstruction}{The grading group and bicharacter supply one standard precise realization of the displayed coefficients; the notes do not specify them.}
\GSfield{Fate}{The notebook citation is identifiable as \citet{trostel1984color}, pp. 3183--3189. The publisher metadata matches the year, volume, issue, and subject; it does not establish the uncertain attribution to Mosolov elsewhere in the line.}
\GSfield{Present status}{As of 12 November 2026 this is an established result under the hypotheses in the modern statement.}
\GSfield{Assessment}{literal statement: correct with missing hypotheses; publication match: partial; verification: publisher metadata checked; development: none; confidence: medium}
\end{GSitem}
\begin{GSitem}{1986-684.5}{E}{The braided commutator of an associative algebra}
\GSfield{Kind; source}{construction; \GSpage{P001304}; D. I. Gurevich}
\GSfield{The notes}{An associative algebra produces a generalized Lie bracket by subtracting the product with reversed factors through \(S\).}
\GSfield{Statement}{If \((A,\mu)\) is an associative algebra object in the symmetric tensor category defined by an involutive braiding \(S\), then \(b=\mu-\mu S\) is a Lie object. Associativity alone, combined with an unrelated solution of the braid equation, does not suffice: multiplication must be compatible with the braiding.}
\GSfield{Reconstruction}{The compatibility hypothesis is explicit in the modern statement; the brief notebook sentence leaves it implicit.}
\GSfield{Fate}{\citet[pp. 799--800]{gurevich1986yangbaxter} treats associative and Lie objects in the same compatible setting.}
\GSfield{Present status}{As of 12 November 2026 this is an established result under the hypotheses in the modern statement.}
\GSfield{Assessment}{literal statement: correct with missing hypotheses; publication match: exact; verification: read the result; development: none; confidence: high}
\end{GSitem}
\begin{GSitem}{1986-684.6}{F}{Historical attribution of braiding formalism to Drinfeld}
\GSfield{Kind; source}{remark; \GSpage{P001304}; D. I. Gurevich}
\GSfield{The notes}{The notes say that Drinfeld put braids in their final form.}
\GSfield{Statement}{This sentence records a historical assessment; it does not specify a theorem, construction, or priority claim precise enough to adjudicate. The mathematical identities actually present are indexed separately.}
\GSfield{Fate}{Gurevich\textquotesingle s 1986 paper provides a concrete contemporaneous source for the adjacent mathematics, but no exact publication is identified with this informal phrase.}
\GSfield{Present status}{As of 12 November 2026 the phrase remains an insufficiently specified historical attribution; no open mathematical problem is inferred.}
\GSfield{Assessment}{literal statement: not determinable; publication match: thematic only; verification: read the result; development: none; confidence: low}
\end{GSitem}

\GStalk{685}{1986-09-29}{V. A. Kazakov}{P001303-P001308}{Ising spins on random triangulations and Wick expansion}
The catalogue date is preserved, but this exposition begins after the Gurevich heading dated 13 October and occupies P001304--P001306. It introduces annealed Ising gravity, a two-matrix representation, a critical temperature, and a Wick-theorem intervention by A. N. Varchenko. Ginzburg\textquotesingle s braid discussion starts with the next dated heading and is not part of this record.

\begin{GSitem}{1986-685.1}{A}{Annealed Ising model on cubic planar maps}
\GSfield{Kind; source}{construction; \GSpage{P001304-P001305}}
\GSfield{The notes}{Spins \(\sigma_i=\pm1\) are summed on the vertices of a random dual graph, with edge interaction and magnetic field.}
\GSfield{Statement}{For each connected cubic map \(\Gamma\), use \(Z_\Gamma(\beta,H)=\sum_{\sigma}\exp[\beta\sum_{\{i,j\}\in E}\sigma_i\sigma_j+H\sum_i\sigma_i]\), then sum over maps with a specified area and automorphism weight. Each unoriented edge is counted once. This is an annealed ensemble: summing over geometry is part of the partition function.}
\GSfield{Reconstruction}{Edge counting and symmetry weights remove otherwise ambiguous factors in the notebook expression.}
\GSfield{Fate}{\citet{kazakov1986ising}, published December 1986, exactly solves a closely related planar Ising matrix model. \citet{boulatovkazakov1987ising} develops the magnetic-field critical behavior; its preprint circulated in 1986.}
\GSfield{Present status}{As of 12 November 2026 this is an established result under the hypotheses in the modern statement.}
\GSfield{Assessment}{literal statement: correct with missing hypotheses; publication match: partial; verification: publisher metadata checked; development: none; confidence: high}
\end{GSitem}
\begin{GSitem}{1986-685.2}{A}{Continuum Majorana field coupled to gravity}
\GSfield{Kind; source}{claim; \GSpage{P001305}}
\GSfield{The notes}{The model is associated with a fermionic field and gravity; the notes emphasize that the continuum identification is difficult.}
\GSfield{Statement}{The critical Ising matter theory is the Majorana conformal field theory with central charge \(1/2\). Its coupling to random planar maps has continuum gravitational critical exponents. A formal identification of a discrete annealed partition sum with a path integral over metrics and Grassmann fields is not, by itself, an equality of rigorously defined measures.}
\GSfield{Reconstruction}{The central charge and the distinction between discrete exact solvability and continuum measure construction are editorial context.}
\GSfield{Fate}{The matrix solution of \citet{kazakov1986ising} and its extension in \citet{boulatovkazakov1987ising} establish exact discrete planar results; their interpretation as matter coupled to gravity is a further continuum argument.}
\GSfield{Present status}{As of 12 November 2026 the Ising-gravity universality class and its exponents are established; the notebook supplies no precise measure-convergence theorem to classify separately.}
\GSfield{Assessment}{literal statement: correct with missing hypotheses; publication match: partial; verification: publisher metadata checked; development: none; confidence: high}
\end{GSitem}
\begin{GSitem}{1986-685.3}{A}{Two-matrix encoding of Ising edge weights}
\GSfield{Kind; source}{construction; \GSpage{P001305}}
\GSfield{The notes}{The integral has a mixed quadratic term \(-2cUV\) and cubic terms \(g(e^HU^3+e^{-H}V^3)/\sqrt N\), with \(c=e^{-2\beta}\).}
\GSfield{Statement}{Use two Hermitian Gaussian matrices with covariance proportional to \(\begin{pmatrix}1&c\\c&1\end{pmatrix}\), for \(|c|<1\), and expand the cubic terms formally. A propagator between different species carries weight \(c\) relative to equal species. Multiplication by a common edge factor converts this to \(e^{\beta\sigma_i\sigma_j}\). Cubic integrals on a real contour remain formal unless contours are specified.}
\GSfield{Reconstruction}{The covariance, positivity restriction, and common edge rescaling are supplied; the notebook uses an unnormalized quadratic form.}
\GSfield{Fate}{This is the matrix-model mechanism of \citet{kazakov1986ising} and \citet{boulatovkazakov1987ising}. A quartic model enumerates maps of a different valency and has different critical couplings.}
\GSfield{Present status}{As of 12 November 2026 this is an established result under the hypotheses in the modern statement.}
\GSfield{Assessment}{literal statement: correct with missing hypotheses; publication match: partial; verification: publisher metadata checked; development: none; confidence: high}
\end{GSitem}
\begin{GSitem}{1986-685.4}{A}{Critical Ising coupling on cubic random maps}
\GSfield{Kind; source}{theorem; \GSpage{P001305}}
\GSfield{The notes}{The transition is at \(\beta_*=\frac12\log(2\sqrt7+1)\).}
\GSfield{Statement}{For the standard cubic two-matrix Ising model at zero magnetic field, the matter critical point has \(c_*=e^{-2\beta_*}=(2\sqrt7-1)/27=1/(2\sqrt7+1)\). This is a joint critical point of matter and the map-generating function, not the fixed square-lattice Ising temperature.}
\GSfield{Reconstruction}{The zero-field cubic-map normalization is necessary to identify the constant.}
\GSfield{Fate}{The cubic critical constant is checked against the explicit Boulatov--Kazakov model in \citet[Eq. (2.9)]{sato2020ising}; its antecedent is \citet{boulatovkazakov1987ising}. The 1986 quartic-lattice paper alone is not an exact match for this numerical constant.}
\GSfield{Present status}{As of 12 November 2026 this is an established result under the hypotheses in the modern statement.}
\GSfield{Assessment}{literal statement: correct with missing hypotheses; publication match: exact; verification: read the result; development: none; confidence: high}
\end{GSitem}
\begin{GSitem}{1986-685.5}{E}{Wick theorem as an all-pairings identity}
\GSfield{Kind; source}{theorem; \GSpage{P001305-P001306}; A. N. Varchenko}
\GSfield{The notes}{Varchenko explains that the Gaussian moment is the sum of products of pairwise covariances; odd moments vanish.}
\GSfield{Statement}{For centered jointly Gaussian variables \(X_1,\ldots,X_{2m}\), \[\mathbb E\prod_iX_i=\sum_{\pi}\prod_{\{a,b\}\in\pi}\mathbb E(X_aX_b)\], where \(\pi\) runs over all perfect matchings; moments of odd total degree are zero. Repeated variables are allowed. The matrix version is obtained by substituting its covariance.}
\GSfield{Fate}{This is the classical Gaussian pairing theorem used by \citet[Sec. 2]{brezinitzyksonparisizuber1978planar}. The intervention explains the mechanism rather than proposes a new theorem.}
\GSfield{Present status}{As of 12 November 2026 this is an established result under the hypotheses in the modern statement.}
\GSfield{Assessment}{literal statement: correct; publication match: exact; verification: read the result; development: none; confidence: high}
\end{GSitem}
\begin{GSitem}{1986-685.6}{E}{Expansion coefficients and Gaussian index contraction}
\GSfield{Kind; source}{method; \GSpage{P001306}}
\GSfield{The notes}{The expansion is written with Gaussian moments of powers of traces; the displayed contractions and Taylor coefficients are compressed.}
\GSfield{Statement}{The formal identity is \(\langle e^{-gP}\rangle_0=\sum_{n\ge0}(-g)^n\langle P^n\rangle_0/n!\). Hermitian entry contractions cross the indices, \(\langle U_{ij}U_{kl}\rangle_0\propto\delta_{il}\delta_{jk}\). Logarithms select connected diagrams, with automorphism weights remaining after the factorials are included.}
\GSfield{Reconstruction}{Taylor factorials and covariance orientation are restored explicitly; these are repairs to the displayed shorthand.}
\GSfield{Fate}{\citet[Secs. 2--4]{brezinitzyksonparisizuber1978planar} supplies the relevant expansion and graph conventions. These details explain the rational, rather than unlabelled integer, coefficients stressed later in record 695.}
\GSfield{Present status}{As of 12 November 2026 this is an established result under the hypotheses in the modern statement.}
\GSfield{Assessment}{literal statement: false as written; publication match: exact; verification: read the result; development: annotated; confidence: high}
\end{GSitem}

\GStalk{686}{1986-10-20}{V. L. Ginzburg}{P001306-P001310}{Polynomial monodromy and the Burau representation}
The actual mathematical section is P001307--P001308. It constructs braid representations from cohomology of a punctured plane, identifies a Hecke relation, and asks faithfulness of Burau. The same letter is used for an exponent and a multiplicative monodromy parameter; the hyperelliptic example makes the ambiguity substantive.

\begin{GSitem}{1986-686.1}{E}{Discriminant complements and the braid group}
\GSfield{Kind; source}{construction; \GSpage{P001307-P001308}}
\GSfield{The notes}{Roots are replaced by coefficients \(\sigma_1,\ldots,\sigma_n\); the discriminant is removed, and the fundamental group is identified with the group of braids.}
\GSfield{Statement}{The space of monic complex polynomials of degree \(n\) with simple roots is the unordered configuration space of \(n\) points in \(\mathbb C\), hence has fundamental group \(B_n\). The ordered root space is its \(S_n\)-cover and has pure braid group as fundamental group. The universal punctured-plane fibration has fiber \(\mathbb C\setminus\{z_1,\ldots,z_n\}\).}
\GSfield{Fate}{The topological construction underlying the Burau representation is recalled in \citet[Sec. 2]{bigelow1999burau}; it was classical before this report.}
\GSfield{Present status}{As of 12 November 2026 this is an established result under the hypotheses in the modern statement.}
\GSfield{Assessment}{literal statement: correct; publication match: exact; verification: read the result; development: none; confidence: high}
\end{GSitem}
\begin{GSitem}{1986-686.2}{E}{Rank-one twisted cohomology and reduced Burau}
\GSfield{Kind; source}{construction; \GSpage{P001308}}
\GSfield{The notes}{The multivalued function \(f^\alpha\), where \(f(z)=\prod(z-z_i)\), defines coefficients for first cohomology; moving the roots gives a braid representation.}
\GSfield{Statement}{Let \(t=e^{2\pi i\alpha}\ne1\) be the common local monodromy about each puncture. Then \(H^1(\mathbb C\setminus\{z_i\},\mathcal L_t)\) has dimension \(n-1\). Monodromy in the configuration space gives the reduced Burau representation, up to duality and basis conventions. The universal form is over \(\mathbb Z[t,t^{-1}]\), and specializing \(t\) is a separate operation.}
\GSfield{Reconstruction}{The nontrivial local-monodromy hypothesis and separation of formal parameter from exponent are supplied.}
\GSfield{Fate}{\citet[Sec. 2]{bigelow1999burau} uses the equivalent infinite-cyclic-cover construction. It verifies the representation being discussed, but not the notebook\textquotesingle s attribution of every step to Givental.}
\GSfield{Present status}{As of 12 November 2026 this is an established result under the hypotheses in the modern statement.}
\GSfield{Assessment}{literal statement: correct with missing hypotheses; publication match: partial; verification: read the result; development: none; confidence: high}
\end{GSitem}
\begin{GSitem}{1986-686.3}{E}{The Hecke quadratic relation and parameter conversion}
\GSfield{Kind; source}{claim; \GSpage{P001308}}
\GSfield{The notes}{The braid generator is said to satisfy \((T_i+1)(T_i-\alpha)=0\), although \(\alpha\) has just denoted the exponent of \(f\).}
\GSfield{Statement}{For the conventional Burau generator \(B_i\), \((B_i-1)(B_i+t)=0\). With \(T_i=-B_i\), this becomes \((T_i+1)(T_i-t)=0\). Thus the Hecke parameter is the multiplicative monodromy \(t\), not the additive exponent \(\alpha\); inverse conventions can replace \(t\) by \(t^{-1}\).}
\GSfield{Reconstruction}{The displayed polynomial is reconciled by an explicit normalization, not silently copied.}
\GSfield{Fate}{The universal representation and its matrices are stated in \citet[Sec. 1]{bharathrambirmanbrendle2026burau}; the notebook\textquotesingle s formula is consistent only after a parameter renaming and sign convention.}
\GSfield{Present status}{As of 12 November 2026 this is an established result under the hypotheses in the modern statement.}
\GSfield{Assessment}{literal statement: false as written; publication match: exact; verification: read the result; development: annotated; confidence: high}
\end{GSitem}
\begin{GSitem}{1986-686.4}{B}{Faithfulness of the Burau representation}
\GSfield{Kind; source}{problem; \GSpage{P001308}}
\GSfield{The notes}{The stated goal is to prove faithfulness; the case \(n=3\) is already known.}
\GSfield{Statement}{For which \(n\) is the universal reduced Burau representation \(B_n\to\operatorname{GL}_{n-1}(\mathbb Z[t,t^{-1}])\) injective? This concerns the indeterminate-parameter representation and does not assert faithfulness at every complex specialization.}
\GSfield{Reconstruction}{The original all-rank question is separated from the exceptional rank-four case and from specialized representations.}
\GSfield{Fate}{\citet[Theorem 1.1]{bigelow1999burau}, first posted 20 April 1999 and published in 1999, proves nonfaithfulness for \(n\ge5\). \citet[Main Theorem, Sec. 1]{bharathrambirmanbrendle2026burau}, posted 6 July and revised 14 September 2026, claims faithfulness for \(n=4\). Its actual theorem and PDF were read; the proof has not been independently certified in this project.}
\GSfield{Present status}{As of 12 November 2026, faithfulness for \(n\le3\) and nonfaithfulness for \(n\ge5\) are established; \(n=4\) has a recent affirmative preprint claim, which must be identified as such. The unrestricted conjecture in the notes is false.}
\GSfield{Assessment}{literal statement: correct; publication match: exact; verification: read the result; development: none; confidence: high; problem \hyperref[prob:GP-1986-30]{GP-1986-30}}
\end{GSitem}
\begin{GSitem}{1986-686.5}{E}{Hyperelliptic monodromy specialization}
\GSfield{Kind; source}{remark; \GSpage{P001308}}
\GSfield{The notes}{The final line says \(\alpha=-1\) yields monodromy of elliptic or hyperelliptic integrals.}
\GSfield{Statement}{The double-cover interpretation corresponds to local monodromy \(t=-1\), hence exponent \(\alpha\in\frac12+\mathbb Z\) if \(f^\alpha\) is used. Its periods lie in the appropriate anti-invariant part of homology of the branched double cover; compactification and parity of the number of branch points affect the dimension.}
\GSfield{Reconstruction}{The distinction between exponent \(-1\), which has trivial monodromy, and multiplicative parameter \(-1\) resolves the apparent contradiction.}
\GSfield{Fate}{The Burau-cover description in \citet[Sec. 2]{bigelow1999burau} supports the specialization, but the notebook conflates the exponent and its exponential.}
\GSfield{Present status}{As of 12 November 2026 this is an established result under the hypotheses in the modern statement.}
\GSfield{Assessment}{literal statement: false as written; publication match: partial; verification: read the result; development: annotated; confidence: high}
\end{GSitem}

\GStalk{687}{1986-10-20}{A. A. Beilinson}{P001308-P001310}{The Deligne--Illusie proof of Hodge degeneration}
The short report at the top of P001309 announces an algebraic proof of degeneration of the Hodge--de Rham spectral sequence, referring to an Orsay preprint by Deligne and Illusie. The theorem over characteristic zero and the new positive-characteristic method are separate items.

\begin{GSitem}{1986-687.1}{E}{Hodge--de Rham degeneration in characteristic zero}
\GSfield{Kind; source}{theorem; \GSpage{P001309}}
\GSfield{The notes}{For a smooth projective complex variety, the spectral sequence from coherent differential forms to de Rham cohomology degenerates.}
\GSfield{Statement}{For a smooth proper variety \(X\) over a characteristic-zero field, \(E_1^{p,q}=H^q(X,\Omega_X^p)\Rightarrow H_{\mathrm{dR}}^{p+q}(X)\) degenerates at \(E_1\). Over \(\mathbb C\), algebraic de Rham cohomology identifies with singular cohomology with complex coefficients; projective smoothness also supplies the classical Hodge decomposition.}
\GSfield{Reconstruction}{The indexing is standardized and smooth proper is distinguished from the more restrictive projective hypothesis in the note.}
\GSfield{Fate}{The statement predates the seminar. \citet[Introduction and Cor. 2.4]{deligneillusie1987relevements} gives a new algebraic proof through reduction to positive characteristic, published in 1987.}
\GSfield{Present status}{As of 12 November 2026 this is an established result under the hypotheses in the modern statement.}
\GSfield{Assessment}{literal statement: correct; publication match: exact; verification: read the result; development: none; confidence: high}
\end{GSitem}
\begin{GSitem}{1986-687.2}{B}{Lifting modulo p squared and the truncated Cartier splitting}
\GSfield{Kind; source}{method; \GSpage{P001309}}
\GSfield{The notes}{The notes credit a new algebraic proof to Deligne and Illusie, with a preprint at Orsay.}
\GSfield{Statement}{If \(X\) is smooth over a perfect field of characteristic \(p>0\) and lifts smoothly to \(W_2(k)\), then the truncated Frobenius pushforward \(\tau_{<p}F_*\Omega_X^\bullet\) splits in the derived category as \(\bigoplus_{i<p}\Omega^i_{X^{(1)}}[-i]\). For proper \(X\) with \(\dim X<p\), this yields full Hodge--de Rham degeneration. Characteristic-zero degeneration follows by spreading out and choosing suitable large primes.}
\GSfield{Reconstruction}{The lifting, perfectness, truncation, and dimension hypotheses are not present in the brief report and are reconstructed from the theorem.}
\GSfield{Fate}{\citet[Thm. 2.1 and Cor. 2.4]{deligneillusie1987relevements} is the precise matching result. The notebook witnesses circulation by October 1986; the article appeared in Inventiones 89 (1987), 247--270.}
\GSfield{Present status}{As of 12 November 2026 this is an established result under the hypotheses in the modern statement.}
\GSfield{Assessment}{literal statement: correct with missing hypotheses; publication match: exact; verification: read the result; development: none; confidence: high}
\end{GSitem}

\GStalk{688}{1986-10-20}{V. A. Kazakov}{P001308-P001312}{Unitary orbital integrals and reduced gauge models}
The section from the lower part of P001309 through P001311 develops the Harish-Chandra--Itzykson--Zuber integral, recalls the principal-minor spherical-function formula, and explains the radial heat operator. It closes with a two-source integral and the Eguchi--Kawai matrix reduction. The final assertion that the reduced integral contains all gauge theory needs a center-symmetry qualification.

\begin{GSitem}{1986-688.1}{E}{The normalized Harish-Chandra--Itzykson--Zuber integral}
\GSfield{Kind; source}{theorem; \GSpage{P001309-P001310}}
\GSfield{The notes}{The unitary integral of \(\exp\operatorname{tr}(\Lambda U M U^*)\) is a determinant of exponentials divided by two Vandermondes.}
\GSfield{Statement}{For Hermitian \(A,B\) with eigenvalues \(a_i,b_j\) and normalized Haar measure, \[\int_{U(N)}e^{t\operatorname{tr}(AUBU^*)}\,dU=\left(\prod_{k=1}^{N-1}k!\right)t^{-N(N-1)/2}\frac{\det[e^{ta_i b_j}]}{\Delta(a)\Delta(b)}.\] Initially take distinct eigenvalues and \(t\ne0\); all apparent singularities extend by continuity. The same Vandermonde ordering is used in numerator conventions for both spectra.}
\GSfield{Fate}{\citet[Eq. (3.12)]{itzyksonzuber1980planar} gives the directly matching matrix formula, building on Harish-Chandra\textquotesingle s earlier orbital-integral work. The author-hosted original PDF and formula were inspected.}
\GSfield{Present status}{As of 12 November 2026 this is an established result under the hypotheses in the modern statement.}
\GSfield{Assessment}{literal statement: correct with missing hypotheses; publication match: exact; verification: read the result; development: none; confidence: high}
\end{GSitem}
\begin{GSitem}{1986-688.2}{E}{Principal-minor integral and shifted spherical parameters}
\GSfield{Kind; source}{theorem; \GSpage{P001310}; I. M. Gelfand}
\GSfield{The notes}{Gelfand recalls an integral of complex powers of successive principal minors and attributes it to Gelfand--Naimark, 1951.}
\GSfield{Statement}{For \(R>0\), principal minors \(\Delta_k\), \(s\in\mathbb C^N\), and \(s_{N+1}=0\), average \(\prod_k\Delta_k(URU^*)^{s_k-s_{k+1}}\). If \(a_j=s_j+N-j\), the value is \(\bigl(\prod_{k=1}^{N-1}k!\bigr)\det[r_i^{a_j}]/[\Delta(r)\Delta(a)]\). Positive minors fix the logarithms; colliding parameters are handled by limits.}
\GSfield{Reconstruction}{The exact principal-minor exponents and Weyl shifts repair the compressed notebook display.}
\GSfield{Fate}{The old formula and its original book locator are reproduced with proof in \citet[Prop. 2.1]{faraut2014rayleigh}. Record 692 gives a fuller representation-theoretic explanation. The shifts \(N-j\) must not be dropped.}
\GSfield{Present status}{As of 12 November 2026 this is an established result under the hypotheses in the modern statement.}
\GSfield{Assessment}{literal statement: correct with missing hypotheses; publication match: partial; verification: read the result; development: none; confidence: high}
\end{GSitem}
\begin{GSitem}{1986-688.3}{E}{Radial heat operator on Hermitian matrices}
\GSfield{Kind; source}{method; \GSpage{P001310-P001311}}
\GSfield{The notes}{The heat equation on the real vector space of Hermitian matrices becomes a radial equation in eigenvalues; multiplication by the Vandermonde makes it free.}
\GSfield{Statement}{With the Hilbert--Schmidt Euclidean metric on \(\operatorname{Herm}_N\), a conjugation-invariant smooth function has radial Laplacian \(\Delta(\lambda)^{-2}\sum_i\partial_i[\Delta(\lambda)^2\partial_i]\). Conjugation by \(\Delta\) gives \(\Delta^{-1}\sum_i\partial_i^2\Delta\). The radial heat kernel is therefore an antisymmetrized product of scalar Gaussian kernels after accounting for the Jacobian.}
\GSfield{Reconstruction}{The metric fixes heat-time and Gaussian normalization; smooth continuation through collisions is understood.}
\GSfield{Fate}{\citet[Sec. 3, Eqs. (3.13)--(3.17)]{itzyksonzuber1980planar} gives this derivation. It is specific to the complex Hermitian, or Dyson-index-two, radial operator; analogous real and quaternionic operators retain interaction terms.}
\GSfield{Present status}{As of 12 November 2026 this is an established result under the hypotheses in the modern statement.}
\GSfield{Assessment}{literal statement: correct with missing hypotheses; publication match: exact; verification: read the result; development: none; confidence: high}
\end{GSitem}
\begin{GSitem}{1986-688.4}{E}{Berezin and Harish-Chandra antecedents of the radial method}
\GSfield{Kind; source}{remark; \GSpage{P001311}; I. M. Gelfand}
\GSfield{The notes}{Gelfand says Berezin had obtained the result and Harish-Chandra used it repeatedly.}
\GSfield{Statement}{This intervention links the matrix heat calculation to the radial-part method in harmonic analysis. It does not identify a unique dissertation proposition; the exact unitary integral and its radial proof are indexed separately.}
\GSfield{Fate}{\citet{berezingelfand1956spherical} is a verified pre-1986 primary source for spherical-function harmonic analysis. The notebook\textquotesingle s particular dissertation attribution is retained as a historical recollection, not upgraded to a proved priority assertion.}
\GSfield{Present status}{As of 12 November 2026 the mathematical antecedent is established, but the precise dissertation passage intended by this intervention is not identified.}
\GSfield{Assessment}{literal statement: correct with missing hypotheses; publication match: thematic only; verification: publisher metadata checked; development: none; confidence: medium}
\end{GSitem}
\begin{GSitem}{1986-688.5}{B}{A unitary integral with two coupled external-source pairs}
\GSfield{Kind; source}{question; \GSpage{P001311}}
\GSfield{The notes}{The question is to evaluate \(\int e^{\operatorname{tr}(A_1UB_1U^*)+\operatorname{tr}(A_2UB_2U^*)}\,dU\), related in the notes to a chiral field.}
\GSfield{Statement}{For fixed finite matrices, the exponential has an absolutely convergent Taylor expansion on the compact group. Each coefficient is a finite sum of unitary moments computable by Weingarten calculus. In general the two-source exponent does not reduce to one HCIZ pair; an all-coupling planar asymptotic formula is a stronger question than finite-dimensional evaluability.}
\GSfield{Reconstruction}{The convergent finite-\(N\) method is an editorial application of the integration theorem; it is not asserted to be the speaker\textquotesingle s intended closed formula.}
\GSfield{Fate}{\citet[Sec. 2, Cor. 2.4]{collinssniady2006integration}, first circulated February 2004 and published 2006, gives all polynomial Haar moments, hence an exact coefficientwise method. The notebook does not specify the source scaling needed to identify a unique large-\(N\) chiral-field problem. This is the same integral family as 1986-695.11 and 1986-698.1. \ours\ The chapter gives an exact \(N=2\) formula, a convergent finite-\(N\) moment expansion with a remainder bound, and a three-source obstruction to treating relative angles as independent.}
\GSfield{Present status}{As of 12 November 2026 finite-\(N\) evaluation by a convergent moment series is available. The vague chiral-field identification is not treated as a separate solved or open theorem.}
\GSfield{Assessment}{literal statement: correct with missing hypotheses; publication match: partial; verification: read the result; development: developed in this chapter; confidence: high; problem \hyperref[prob:GP-1986-60]{GP-1986-60}}
\end{GSitem}
\begin{GSitem}{1986-688.6}{E}{Eguchi--Kawai reduction and its symmetry hypothesis}
\GSfield{Kind; source}{construction; \GSpage{P001311}}
\GSfield{The notes}{The integral over \(d\) unitary matrices with plaquette terms is called the lattice Yang--Mills partition function, containing all \(d\)-dimensional physics.}
\GSfield{Statement}{The one-site action is \(N\beta\sum_{\mu<\nu}\operatorname{Re}\operatorname{tr}(U_\mu U_\nu U_\mu^*U_\nu^*)\), up to normalization. Equality of planar Wilson-loop equations with an extended lattice requires large-\(N\) factorization, unbroken center symmetry, and translation invariance in the unreduced model. The untwisted bosonic one-site model breaks center symmetry at weak coupling in four dimensions, so the literal universal equivalence fails.}
\GSfield{Reconstruction}{The factor of \(N\), real part, and symmetry conditions are supplied. Oppositely oriented terms in the notebook can themselves make the action real.}
\GSfield{Fate}{\citet{eguchikawai1982reduction} proposes the reduction. \citet[Sec. 2]{garciaperez2020largeN} states the exact symmetry conditions and explains the original failure; twisted, deformed, and adjoint-fermion variants are distinct models.}
\GSfield{Present status}{As of 12 November 2026 this is an established result under the hypotheses in the modern statement.}
\GSfield{Assessment}{literal statement: false as written; publication match: partial; verification: read the result; development: annotated; confidence: high}
\end{GSitem}
\begin{GSitem}{1986-688.7}{C}{Analytic solution of reduced planar gauge integrals in dimensions three and four}
\GSfield{Kind; source}{problem; \GSpage{P001311}}
\GSfield{The notes}{The integral is underlined as unevaluated, and gauge theory as unsolved.}
\GSfield{Statement}{Determine the large-\(N\) free energy and Wilson-loop observables of a precisely normalized \(d\)-matrix reduced gauge model for \(d=3,4\), through its phases, and establish the regime in which these observables reproduce the extended-lattice theory. Solving a finite compact integral by Taylor expansion is not this nonperturbative planar problem.}
\GSfield{Reconstruction}{The free-energy scaling, observables, and distinction between the unrepaired original model and symmetry-preserving variants make the question precise.}
\GSfield{Fate}{The reduction of \citet{eguchikawai1982reduction} is a reformulation, not a general analytic solution. \citet[Secs. 1--2]{garciaperez2020largeN} describes continuing numerical work; \citet{neuberger2020quenched} distinguishes competing quenching prescriptions. Neither proves an all-coupling solution of the notebook integral.}
\GSfield{Present status}{As of 12 November 2026 exact two-dimensional results and controlled reductions/numerics do not provide the requested general three- or four-dimensional planar solution. The original untwisted weak-coupling equivalence is false, and must be repaired before a continuum-physics claim.}
\GSfield{Assessment}{literal statement: correct with missing hypotheses; publication match: partial; verification: read the result; development: none; confidence: high; problem \hyperref[prob:GP-1986-31]{GP-1986-31}}
\GSfield{Worth developing}{Worth developing only after fixing a model and phase; a bounded strong-coupling coefficient calculation or symmetry test is feasible, whereas an all-coupling solution is not a bounded project.}
\end{GSitem}

\GStalk{689}{1986-10-27}{I. S. Zakharevich}{P001310-P001313}{Orbital integration, interlacing, and concentration remarks}
The new dated heading begins at P001312. The report revisits the orbital-integral heat proof and the spectral distribution of products of positive matrices, then uses sections of ellipsoids and interlacing arrays. Gelfand adds two high-dimensional intuition questions. Their printed numerical and dynamical assertions require separate qualifications.

\begin{GSitem}{1986-689.1}{E}{Heat-kernel proof of the unitary orbital formula}
\GSfield{Kind; source}{method; \GSpage{P001312}}
\GSfield{The notes}{The report starts from \(g_t=\Delta g\) on Hermitian matrices and states the determinant formula for the orbital integral.}
\GSfield{Statement}{Compare the full Euclidean heat kernel on \(\operatorname{Herm}_N\) with its radial kernel after multiplication by \(\Delta(\lambda)\). Antisymmetrization over eigenvalue permutations yields the determinant of scalar heat kernels and therefore HCIZ, normalized by the value at a scalar matrix.}
\GSfield{Fate}{This is the method already given in \citet[Sec. 3]{itzyksonzuber1980planar}; it repeats 1986-688.3 with a new exposition.}
\GSfield{Present status}{As of 12 November 2026 this is an established result under the hypotheses in the modern statement.}
\GSfield{Assessment}{literal statement: correct; publication match: exact; verification: read the result; development: none; confidence: high}
\end{GSitem}
\begin{GSitem}{1986-689.2}{B}{Spectrum of a product of positive Hermitian matrices}
\GSfield{Kind; source}{question; \GSpage{P001312}}
\GSfield{The notes}{Given the spectra of \(H_1,H_2>0\), the notes ask how the spectrum of \(H_1H_2\) is distributed.}
\GSfield{Statement}{The product is similar to \(H_1^{1/2}H_2H_1^{1/2}>0\). Two different questions must be distinguished: which eigenvalue lists occur, and the pushforward distribution when relative eigenvectors have Haar law. The first is a multiplicative spectral-support problem; the second needs a specified measure and cannot be answered by support inequalities alone.}
\GSfield{Reconstruction}{Haar-relative-eigenbasis sampling is supplied as the natural precise probabilistic interpretation; the note does not explicitly choose it.}
\GSfield{Fate}{\citet{klyachko2000random} relates product singular-value support to additive Hermitian spectral inequalities. Principal-minor spherical transforms, as in \citet[Sec. 2]{faraut2014rayleigh}, are relevant to the distribution question. Neither reference is claimed to be the exact unstated density formula intended in the note.}
\GSfield{Present status}{As of 12 November 2026 finite-dimensional spectral support has established solutions; the notebook\textquotesingle s unspecified distribution question is not identified with one universal density theorem.}
\GSfield{Assessment}{literal statement: correct with missing hypotheses; publication match: partial; verification: publisher metadata checked; development: none; confidence: medium}
\end{GSitem}
\begin{GSitem}{1986-689.3}{E}{Ellipsoid sections and interlacing arrays}
\GSfield{Kind; source}{theorem; \GSpage{P001312}}
\GSfield{The notes}{A plane section of an ellipsoid with axes \(a,b,c\) is described by \(a\le a_1\le b\le b_1\le c\); repeated sections form a triangular array.}
\GSfield{Statement}{For a Hermitian operator restricted orthogonally to a codimension-one subspace, sorted eigenvalues interlace. Repeated principal compressions produce Gelfand--Tsetlin patterns. For ellipsoids, the quadratic-form eigenvalues are inverse squared semiaxes; the inequality must be translated consistently if the symbols denote geometric lengths.}
\GSfield{Reconstruction}{Compression, axis convention, and ordering are supplied.}
\GSfield{Fate}{\citet[Thm. 1.1 and Prop. 2.3]{faraut2014rayleigh} states the interlacing theorem and orbital-measure pushforward. The geometric interlacing is classical; the modern probability formulation is context, not a 1986 novelty attribution.}
\GSfield{Present status}{As of 12 November 2026 this is an established result under the hypotheses in the modern statement.}
\GSfield{Assessment}{literal statement: correct with missing hypotheses; publication match: exact; verification: read the result; development: none; confidence: high}
\end{GSitem}
\begin{GSitem}{1986-689.4}{E}{Principal-minor powers as spherical coordinates for integration}
\GSfield{Kind; source}{method; \GSpage{P001312}}
\GSfield{The notes}{The next integral averages powers of \(a_{11}\), a leading \(2\times2\) determinant, and successive minors of \(U^*\mu U\).}
\GSfield{Statement}{For positive definite \(\mu\), these complex powers are unambiguously defined using positive real logarithms. Their unitary average is the spherical function of \(\operatorname{GL}_N(\mathbb C)/U(N)\); the exponent differences and \(N-j\) shifts are as in 1986-688.2 and 1986-692.1.}
\GSfield{Reconstruction}{Positivity of \(\mu\) is needed for the elementary complex-power interpretation.}
\GSfield{Fate}{\citet[Sec. 2, Prop. 2.1]{faraut2014rayleigh} reproduces the Gelfand--Naimark integration scheme. This occurrence motivates the representation-theoretic explanation given on 3 November.}
\GSfield{Present status}{As of 12 November 2026 this is an established result under the hypotheses in the modern statement.}
\GSfield{Assessment}{literal statement: correct with missing hypotheses; publication match: exact; verification: read the result; development: none; confidence: high}
\end{GSitem}
\begin{GSitem}{1986-689.5}{E}{The numerical equatorial-slab assertion in dimension 1000}
\GSfield{Kind; source}{claim; \GSpage{P001312}}
\GSfield{The notes}{The question asks what fraction of a 1000-dimensional ball lies in a layer of thickness \(R/100\) about an equator; the answer written is almost all.}
\GSfield{Statement}{For uniform volume in the radius-\(R\) ball of dimension \(d\), \(\Pr(|x_1|\le aR)=I_{a^2}(1/2,(d+1)/2)\). At \(d=1000\), even the generous half-width interpretation \(a=0.01\) gives approximately \(0.248\); interpreting \(R/100\) as total thickness gives approximately \(0.126\). Thus the literal numerical example does not illustrate almost-full concentration.}
\GSfield{Reconstruction}{Two possible meanings of thickness are tested rather than choosing one silently.}
\GSfield{Fate}{The formula follows from the cross-sectional volume of a ball. \oursComputed\ The numerical values are independently checked by quadrature and the regularized beta function in the accompanying code; this is an editorial correction, not a new concentration theorem.}
\GSfield{Present status}{As of 12 November 2026 the general concentration heuristic is standard, but this specific dimension-to-width ratio is quantitatively false.}
\GSfield{Assessment}{literal statement: false as written; publication match: none found; verification: read the result; development: annotated; confidence: high}
\end{GSitem}
\begin{GSitem}{1986-689.6}{F}{Time spent near faces of a high-dimensional simplex billiard}
\GSfield{Kind; source}{claim; \GSpage{P001313}}
\GSfield{The notes}{A particle reflects at faces of a 1000-dimensional simplex; the notes say it spends almost all time within \(1/100\) of the faces.}
\GSfield{Statement}{The statement does not specify units for the boundary layer, the shape or inradius of the simplex, initial-condition measure, or whether the assertion is an ensemble or almost-every-trajectory limit. Concentration of volume near a boundary alone cannot establish a time-average claim without a suitable dynamical hypothesis.}
\GSfield{Reconstruction}{The required data are identified without replacing the notebook sentence by an invented general billiard conjecture.}
\GSfield{Fate}{No definite billiard theorem can be matched to these missing quantifiers. The neighboring ball question explains the heuristic motivation, but supplies neither ergodicity nor a trajectorywise argument.}
\GSfield{Present status}{As of 12 November 2026 this occurrence is underdetermined; no open problem or class-D inference is made.}
\GSfield{Assessment}{literal statement: not determinable; publication match: none found; verification: unresolved lead; development: none; confidence: low}
\end{GSitem}

\GStalk{690}{1986-10-27}{B. V. Lidskii}{P001311-P001314}{The diagonal moment map}
The short intervention occupies the middle of P001313. It identifies the diagonal projection on a unitary orbit as a moment map and relates its Fourier transform to the orbital integral. The uncertain handwritten initial is not used to override the catalogue identification. Gelfand\textquotesingle s question about a super analogue of the symmetric group follows.

\begin{GSitem}{1986-690.1}{E}{Diagonal projection as a moment map and its Fourier transform}
\GSfield{Kind; source}{construction; \GSpage{P001313}}
\GSfield{The notes}{On Hermitian matrices with fixed eigenvalues, take \(H\mapsto\operatorname{diag}H\); the Fourier transform is written down.}
\GSfield{Statement}{On the coadjoint orbit of \(U(N)\) with its Kirillov--Kostant--Souriau form, diagonal projection is the moment map for the diagonal torus. Its image is the convex hull of permutations of the fixed eigenvalue list. The Fourier transform of the pushed-forward normalized orbital measure is \(\int e^{i\operatorname{tr}(D U\Lambda U^*)}\,dU\), hence HCIZ with \(t=i\).}
\GSfield{Reconstruction}{The symplectic form, measure normalization, and Schur--Horn image describe the standard precise construction; the note only states the map and transform.}
\GSfield{Fate}{\citet{duistermaatheckman1982variation} supplies the general moment-map measure context; \citet[Sec. 3]{itzyksonzuber1980planar} supplies the explicit Fourier integral. These were pre-existing results.}
\GSfield{Present status}{As of 12 November 2026 this is an established result under the hypotheses in the modern statement.}
\GSfield{Assessment}{literal statement: correct with missing hypotheses; publication match: partial; verification: read the result; development: none; confidence: high}
\end{GSitem}
\begin{GSitem}{1986-690.2}{F}{What is the super analogue of the symmetric group?}
\GSfield{Kind; source}{question; \GSpage{P001313}; I. M. Gelfand}
\GSfield{The notes}{Gelfand asks \ru{Что такое супераналог симметрической группы?}}
\GSfield{Statement}{No criterion for analogue is specified: permutation action on super tensor powers, a Schur--Weyl centralizer, a spin double cover, and a Lie-superalgebra Weyl object give different answers. The signed permutation action is already present in 1986-684.3, but does not uniquely answer this broader question.}
\GSfield{Fate}{The neighboring color-symmetry exposition and \citet{gurevich1986yangbaxter} provide relevant context only. The phrase does not determine a precise research problem or a unique later theorem.}
\GSfield{Present status}{As of 12 November 2026 the intended notion remains unidentified from this occurrence; no claim of an undeveloped idea is made.}
\GSfield{Assessment}{literal statement: not determinable; publication match: thematic only; verification: read the result; development: none; confidence: low}
\end{GSitem}

\GStalk{691}{1986-10-27}{V. L. Ginzburg}{P001312-P001315}{Exact stationary phase and localization}
The lower part of P001313 gives an approximate Duistermaat--Heckman formula, its coadjoint-orbit example, and Gelfand\textquotesingle s questions about higher asymptotic terms and characteristic numbers. A final conjectural sentence is explicitly marked mine by the note-taker; it must not be attributed to Gelfand or Ginzburg.

\begin{GSitem}{1986-691.1}{E}{Exact stationary phase for a Hamiltonian circle action}
\GSfield{Kind; source}{theorem; \GSpage{P001313}}
\GSfield{The notes}{For a closed symplectic manifold and periodic Hamiltonian flow, the first stationary-phase term is said to give the exact integral. The written prefactor is \((2\pi)^{-n}\), without a signature phase.}
\GSfield{Statement}{Let \((M^{2n},\omega)\) be compact without boundary and let \(H\) generate a circle action with isolated fixed points. With Liouville measure \(\omega^n/n!\), the oscillatory integral is exactly its stationary-phase sum. In local Darboux coordinates, for real \(t>0\), each point contributes \((2\pi/t)^n e^{itH(p)}e^{i\pi\operatorname{sgn}(\operatorname{Hess}_pH)/4}/\sqrt{|\det\operatorname{Hess}_pH|}\). Consistent orientation and moment-map conventions fix signs.}
\GSfield{Reconstruction}{Isolated nondegenerate fixed points, the Liouville factorial, absolute determinant, signature phase, and \(t\)-scaling are supplied.}
\GSfield{Fate}{\citet{duistermaatheckman1982variation} proves exactness; \citet[Thm. (1.5)]{duistermaatheckman1983addendum} treats nonisolated fixed sets. The notebook calls its own formula approximate, appropriately, but its displayed constant and phases cannot be retained literally.}
\GSfield{Present status}{As of 12 November 2026 this is an established result under the hypotheses in the modern statement.}
\GSfield{Assessment}{literal statement: false as written; publication match: exact; verification: read the result; development: annotated; confidence: high}
\end{GSitem}
\begin{GSitem}{1986-691.2}{E}{Coadjoint-orbit symplectic form and the unitary integral}
\GSfield{Kind; source}{construction; \GSpage{P001313}}
\GSfield{The notes}{For \(X=U^*\mu U\), the form is written \(\omega(X)(\xi,\eta)=\operatorname{tr}(X[\xi,\eta])\), with \(\xi,\eta\) called tangent vectors.}
\GSfield{Statement}{If tangent vectors are represented by infinitesimal generators \(\xi,\eta\) as \(\operatorname{ad}^*\xi X,\operatorname{ad}^*\eta X\), the orbit form is \(\omega_X(\operatorname{ad}^*\xi X,\operatorname{ad}^*\eta X)=\langle X,[\xi,\eta]\rangle\), up to a uniform sign. It descends modulo the stabilizer; generators are not literally arbitrary tangent matrices. A regular unitary orbit is \(U(N)/U(1)^N\); repeated eigenvalues enlarge the stabilizer.}
\GSfield{Reconstruction}{The generator/tangent distinction and regular-orbit qualification repair the notation.}
\GSfield{Fate}{This is the orbit example of exact localization in \citet{duistermaatheckman1982variation}; the resulting alternating exponential sum agrees with HCIZ in \citet{itzyksonzuber1980planar}.}
\GSfield{Present status}{As of 12 November 2026 this is an established result under the hypotheses in the modern statement.}
\GSfield{Assessment}{literal statement: correct with missing hypotheses; publication match: exact; verification: read the result; development: none; confidence: high}
\end{GSitem}
\begin{GSitem}{1986-691.3}{E}{Higher stationary-phase terms and characteristic numbers}
\GSfield{Kind; source}{question; \GSpage{P001313}; I. M. Gelfand}
\GSfield{The notes}{Gelfand asks when the second term vanishes, when the first term is exact, and how characteristic numbers can arise.}
\GSfield{Statement}{Within compact Hamiltonian torus geometry, localization on a fixed component integrates the exponential of the restricted symplectic form divided by the equivariant Euler class of its normal bundle. Curvature expansions produce characteristic classes and a finite-degree contribution on each component. Outside a specified Hamiltonian phase-amplitude class, vanishing of one coefficient does not characterize exactness.}
\GSfield{Reconstruction}{The compact Hamiltonian subclass is distinguished from the unbounded scope of the informal question.}
\GSfield{Fate}{\citet[Thm. (1.5) and its expansion on pp. 154--155]{duistermaatheckman1983addendum} already gives the fixed-component and characteristic-class mechanism. It is a precise partial answer to the intervention, not a classification of all oscillatory integrals.}
\GSfield{Present status}{As of 12 November 2026 the localization answer in the specified subclass is established. The note does not define a general classification problem beyond it.}
\GSfield{Assessment}{literal statement: correct with missing hypotheses; publication match: partial; verification: read the result; development: none; confidence: high}
\end{GSitem}
\begin{GSitem}{1986-691.4}{F}{The note-taker proposes replacing the circle by a nilpotent group}
\GSfield{Kind; source}{conjecture; \GSpage{P001313}; Notebook author (explicitly: mine)}
\GSfield{The notes}{The sentence begins \ru{Гипотеза (моя)} and suggests a nilpotent Lie-group action, with a two-dimensional example expected to give two asymptotic terms.}
\GSfield{Statement}{No Hamiltonian, phase, measure, compactness condition, fixed-point behavior, or asymptotic parameter is specified. A group action by itself does not determine the number of nonzero stationary-phase terms. The dimension-two nilpotent case is abelian, so even the proposed mechanism is not determined by a nonabelian example.}
\GSfield{Reconstruction}{We retain the conjectural fragment and attribution, but do not invent missing hypotheses or a class-D open problem.}
\GSfield{Fate}{The fixed-component theorem of \citet{duistermaatheckman1983addendum} is contextual, not a result matching this conjecture. The original scan confirms that the parenthetical attribution is the note-taker\textquotesingle s own.}
\GSfield{Present status}{As of 12 November 2026 the mathematical proposition intended by this sentence is not identifiable.}
\GSfield{Assessment}{literal statement: not determinable; publication match: thematic only; verification: read the result; development: none; confidence: low}
\end{GSitem}

\GStalk{692}{1986-11-03}{G. I. Olshanskii}{P001313-P001316}{Spherical functions behind principal-minor integrals}
The actual exposition begins on P001314 and ends in the upper part of P001315. Olshanskii explains the Gelfand--Naimark formula through highest-weight characters and spherical functions, with HCIZ as a limiting form. Gelfand asks which other integrals admit such an interpretation.

\begin{GSitem}{1986-692.1}{E}{The principal-minor spherical integral in full normalization}
\GSfield{Kind; source}{theorem; \GSpage{P001314}}
\GSfield{The notes}{The integrand is a product of powers of leading principal minors of \(URU^*\); the answer has alternating powers of eigenvalues and parameter differences.}
\GSfield{Statement}{Let \(r_i>0\), \(R=\operatorname{diag}(r_i)\), and \(a_j=s_j+N-j\). For normalized Haar measure, \[\int_{U(N)}\prod_{k=1}^N\Delta_k(URU^*)^{s_k-s_{k+1}}dU=\frac{\bigl(\prod_{k=1}^{N-1}k!\bigr)\det[r_i^{a_j}]}{\Delta(r)\prod_{i<j}(a_i-a_j)},\quad s_{N+1}=0.\] Both eigenvalue and parameter collisions have removable limits; all complex powers use positive bases.}
\GSfield{Reconstruction}{The exact normalization, exponent differences, and Weyl shifts reconstruct the abbreviated formula.}
\GSfield{Fate}{\citet[Prop. 2.1, citing Gelfand--Naimark, II.9.3]{faraut2014rayleigh} gives this formula with proof. The original theorem predates 1986 by decades.}
\GSfield{Present status}{As of 12 November 2026 this is an established result under the hypotheses in the modern statement.}
\GSfield{Assessment}{literal statement: correct with missing hypotheses; publication match: exact; verification: read the result; development: none; confidence: high}
\end{GSitem}
\begin{GSitem}{1986-692.2}{E}{Contraction from spherical functions to HCIZ}
\GSfield{Kind; source}{method; \GSpage{P001314-P001315}}
\GSfield{The notes}{The notes say formula (2), the exponential orbital integral, follows from the first formula.}
\GSfield{Statement}{With fixed Hermitian \(X\), let \(R_\tau=I+X/\tau>0\) for sufficiently large \(\tau\), and replace the spherical parameters by \(\tau s\). The principal-minor powers tend to an exponential whose exponent telescopes to \(\operatorname{tr}(\operatorname{diag}(s)UXU^*)\). Compactness permits passage under the integral; the determinant side has the corresponding confluent limit.}
\GSfield{Reconstruction}{The scaling and positivity range are editorial reconstruction of the indicated implication.}
\GSfield{Fate}{\citet[Sec. 2, after Prop. 2.2]{faraut2014rayleigh} states precisely this Mehler--Heine-type limit. It explains the relation asserted in the notebook without claiming that the limit calculation was written there.}
\GSfield{Present status}{As of 12 November 2026 this is an established result under the hypotheses in the modern statement.}
\GSfield{Assessment}{literal statement: correct with missing hypotheses; publication match: partial; verification: read the result; development: none; confidence: high}
\end{GSitem}
\begin{GSitem}{1986-692.3}{E}{Characters, highest weights, and spherical matrix coefficients}
\GSfield{Kind; source}{construction; \GSpage{P001315}}
\GSfield{The notes}{For integral parameters the formula is related to Weyl\textquotesingle s character formula; arbitrary complex parameters are interpreted through representations of \(\operatorname{GL}_N(\mathbb C)\).}
\GSfield{Statement}{Dominant integral highest weights give Schur characters, expressed by the alternant divided by a Vandermonde. Averaging a highest-weight coefficient gives a normalized spherical function on \(\operatorname{GL}_N(\mathbb C)/U(N)\). Complex continuation defines the spherical family, but arbitrary complex inducing parameters do not give unitary representations.}
\GSfield{Reconstruction}{The distinction between polynomial highest weights, analytic continuation, and the unitary parameter range is supplied.}
\GSfield{Fate}{\citet[Sec. 2]{faraut2014rayleigh} verifies the spherical formula; its shifts explain the character relation. Unitarity must be checked separately and is discussed for \(\operatorname{GL}_N\) in record 697.}
\GSfield{Present status}{As of 12 November 2026 this is an established result under the hypotheses in the modern statement.}
\GSfield{Assessment}{literal statement: correct with missing hypotheses; publication match: partial; verification: read the result; development: none; confidence: high}
\end{GSitem}
\begin{GSitem}{1986-692.4}{F}{Which integrals are evaluations of spherical functions?}
\GSfield{Kind; source}{question; \GSpage{P001315}; I. M. Gelfand}
\GSfield{The notes}{Gelfand asks which other integrals can be evaluated in this representation-theoretic way.}
\GSfield{Statement}{The note specifies neither a class of integrands nor an admissible group and representation construction. A precise finite family is given by principal-minor averages and their orbital-integral limits, but a classification of all integrals having some representation-theoretic realization is not a defined question.}
\GSfield{Fate}{The adjacent formulas are established examples; \citet[Secs. 2--3]{faraut2014rayleigh} explains their common harmonic-analytic origin. No exact later theorem is matched to the unrestricted wording.}
\GSfield{Present status}{As of 12 November 2026 this is a research-program remark, not an identifiable unsolved problem.}
\GSfield{Assessment}{literal statement: not determinable; publication match: thematic only; verification: read the result; development: none; confidence: low}
\end{GSitem}

\GStalk{693}{1986-11-03}{V. G. Drinfeld}{P001314-P001320}{Quantizing KdV and the second Gelfand--Dickey bracket}
The section actually runs from the lower part of P001315 through P001318. It distinguishes the compatible KdV brackets, their free-field and Miura realizations, higher-order operators, and quantum commuting integrals. Gelfand and Feigin intervene. Three final questions concern quantization of KdV integrals, their spectrum in irreducible Virasoro representations, and general second-bracket quantization. The spectral question receives a bounded level-two development in this chapter; P001319 instead begins Beilinson on 10 November.

\begin{GSitem}{1986-693.1}{E}{The two compatible Poisson structures of KdV}
\GSfield{Kind; source}{construction; \GSpage{P001315}}
\GSfield{The notes}{The notes write \(P_1=-2\partial\) and \(P_2=-\frac12\partial^3-u\partial-\partial u\), then ask about quantization.}
\GSfield{Statement}{On local differential-polynomial functionals of a smooth periodic field \(u\), or with boundary conditions eliminating integration-by-parts terms, \(P_1,P_2\) define compatible Hamiltonian operators. Every constant linear combination is Poisson. Appropriate Hamiltonians and rescaling of \(u,x,t\) give the KdV hierarchy; coefficients cannot be compared across conventions without the rescaling.}
\GSfield{Reconstruction}{The class of functionals and boundary conditions are supplied; the notebook does not fix a functional-analytic phase space.}
\GSfield{Fate}{\citet[Secs. 2--3]{drinfeldsokolov1985lie} surveys the classical structures with proofs, following the Gelfand--Dickey theory. The Russian original appeared in 1984, the English translation in 1985.}
\GSfield{Present status}{As of 12 November 2026 this is an established result under the hypotheses in the modern statement.}
\GSfield{Assessment}{literal statement: correct with missing hypotheses; publication match: exact; verification: read the result; development: none; confidence: high}
\end{GSitem}
\begin{GSitem}{1986-693.2}{F}{A simultaneous quantization with two Planck constants}
\GSfield{Kind; source}{question; \GSpage{P001315}}
\GSfield{The notes}{Two compatible brackets suggest two Planck constants.}
\GSfield{Statement}{The notebook does not specify a flat two-parameter associative algebra, a representation, or which classical pencil it should recover. In the displayed classical pencil, shifting \(u\) by a constant changes \(P_2\) by a multiple of \(\partial\), so independence of two quantum parameters is not automatic.}
\GSfield{Reconstruction}{The obstruction is a missing definition, not evidence of an undeveloped mathematical conjecture.}
\GSfield{Fate}{The classical pencil in \citet{drinfeldsokolov1985lie} and the second-structure quantum KdV construction in \citet{bazhanovlukyanovzamolodchikov1996integrable} are relevant, but neither is identified with this unspecified simultaneous-quantization proposal.}
\GSfield{Present status}{As of 12 November 2026 no definite two-parameter assertion is recoverable from this sentence.}
\GSfield{Assessment}{literal statement: not determinable; publication match: thematic only; verification: read the result; development: none; confidence: medium}
\end{GSitem}
\begin{GSitem}{1986-693.3}{F}{Nonlinear quantization and perturbing a measure}
\GSfield{Kind; source}{remark; \GSpage{P001315}; I. M. Gelfand and V. G. Drinfeld}
\GSfield{The notes}{Gelfand says that genuinely nonlinear quantization is not understood, mentions quantum groups and a perturbation of a measure; Drinfeld contrasts approaches associated with Faddeev and Albeverio--Høegh-Krohn.}
\GSfield{Statement}{No measure space, perturbing functional, regularization, or proposed equivalence is stated. The passage records methodological preferences about quantization, not a theorem comparing the named constructions.}
\GSfield{Fate}{The surrounding quantum-KdV problem acquired concrete algebraic constructions, but the broad methodological comparison has no exact mathematical formulation here.}
\GSfield{Present status}{As of 12 November 2026 this remains an under-specified programmatic remark, not a separate open problem.}
\GSfield{Assessment}{literal statement: not determinable; publication match: thematic only; verification: read the result; development: none; confidence: low}
\end{GSitem}
\begin{GSitem}{1986-693.4}{E}{The second bracket as the Virasoro Lie--Poisson structure}
\GSfield{Kind; source}{construction; \GSpage{P001316}}
\GSfield{The notes}{Fourier expansion of the second bracket yields the Virasoro algebra and Kirillov\textquotesingle s orbit picture.}
\GSfield{Statement}{The affine hyperplane of fixed central coordinate in the dual Virasoro algebra has a Lie--Poisson bracket of the form \(a\partial^3+2u\partial+u^{\prime}\), up to sign and scale. Quantization replaces the Fourier-mode Poisson brackets by Virasoro commutators. This identifies the algebra of fields; it does not by itself order the nonlinear Hamiltonians.}
\GSfield{Reconstruction}{The chapter fixes the standard convention \([L_m,L_n]=(m-n)L_{m+n}+c(m^3-m)\delta_{m+n,0}/12\); the notebook uses a different index/sign convention.}
\GSfield{Fate}{\citet[Sec. 2]{ovsienkokhesin1987super} recalls the already known Virasoro Lie--Poisson identification; \citet[Eq. (2), (11)--(14)]{bazhanovlukyanovzamolodchikov1996integrable} separates it from the quantum charges.}
\GSfield{Present status}{As of 12 November 2026 this is an established result under the hypotheses in the modern statement.}
\GSfield{Assessment}{literal statement: correct with missing hypotheses; publication match: exact; verification: read the result; development: none; confidence: high}
\end{GSitem}
\begin{GSitem}{1986-693.5}{E}{Classical local commuting KdV Hamiltonians}
\GSfield{Kind; source}{theorem; \GSpage{P001316}}
\GSfield{The notes}{The integrals \(F_n\) are said to commute, and a quantum analogue is attributed to Zamolodchikov.}
\GSfield{Statement}{For a scalar second-order Lax operator, residues of appropriate odd half-integer powers produce local functionals forming a Poisson-commutative hierarchy. They are defined modulo total derivatives; periodicity or decay makes their integrals well defined. Classical commutativity does not imply commutativity after an arbitrary ordering of modes.}
\GSfield{Reconstruction}{The formal pseudodifferential convention and total-derivative quotient are supplied.}
\GSfield{Fate}{\citet[Sec. 2]{drinfeldsokolov1985lie} gives the classical residue construction. The quantum question is separated as GP-1986-32.}
\GSfield{Present status}{As of 12 November 2026 this is an established result under the hypotheses in the modern statement.}
\GSfield{Assessment}{literal statement: correct with missing hypotheses; publication match: exact; verification: read the result; development: none; confidence: high}
\end{GSitem}
\begin{GSitem}{1986-693.6}{B}{Quantization of the full family of local KdV integrals}
\GSfield{Kind; source}{problem; \GSpage{P001316; P001318}}
\GSfield{The notes}{The closing question 1(a) asks to quantize the KdV integrals; an earlier line says Zamolodchikov has quantum analogues.}
\GSfield{Statement}{Construct local normally ordered Virasoro charges \(I_{2k-1}\) deforming the classical KdV Hamiltonians, with \([I_{2k-1},I_{2l-1}]=0\), rather than merely quantizing the underlying bracket. The construction is naturally formal or generic in the deformation parameter; specialized representations and convergence need separate treatment.}
\GSfield{Reconstruction}{The normalization and formal/generic meaning of quantization are supplied. No assertion that an arbitrary ordering works is attributed to the notes.}
\GSfield{Fate}{\citet[Thm. 4.7.3, with the generic-parameter hypotheses of Sec. 4.7.1]{feiginfrenkel1996integrals} proves deformation and commutativity in the relevant affine \(\mathfrak{sl}_2\) case (1993 circulation; final arXiv version September 1995; publication 1996). \citet[Eq. (11) and quantum monodromy construction]{bazhanovlukyanovzamolodchikov1996integrable} gives explicit initial charges and transfer operators (1994 circulation; 1996 publication).}
\GSfield{Present status}{As of 12 November 2026 the algebraic construction of commuting quantum KdV charges is established. Their simultaneous spectral classification is a separate, only partly resolved problem.}
\GSfield{Assessment}{literal statement: correct; publication match: exact; verification: read the result; development: none; confidence: high; problem \hyperref[prob:GP-1986-32]{GP-1986-32}}
\end{GSitem}
\begin{GSitem}{1986-693.7}{E}{Miura transformation and the Heisenberg Poisson field}
\GSfield{Kind; source}{construction; \GSpage{P001316}}
\GSfield{The notes}{The substitution \(u=v^{\prime}-v^2\) is called Poisson for the second Gelfand--Dickey bracket, with a derivative-of-delta bracket on \(v\).}
\GSfield{Statement}{For the displayed normalization \(P_2=-\frac12\partial^3-u\partial-\partial u\), take \(P_v=\frac12\partial\). The differential of the Miura map is \(D=\partial-2v\), and \(DP_vD^*=P_2\). If \(P_v=\partial\), the pushforward is \(2P_2\). This is an identity of formal Hamiltonian differential operators; global surjectivity on smooth periodic fields is not asserted.}
\GSfield{Reconstruction}{The explicit factor \(1/2\) is an editorial normalization check, not a claim that the handwritten Heisenberg coefficient had that value. \ours\ The displayed Hamiltonian-operator identity and its factor \(1/2\) are checked directly in the saved symbolic calculation.}
\GSfield{Fate}{\citet[Sec. 3.7]{drinfeldsokolov1985lie} treats the Miura Poisson realization. The factor depends on the bracket normalization, which the notebook does not keep consistently fixed.}
\GSfield{Present status}{As of 12 November 2026 this is an established result under the hypotheses in the modern statement.}
\GSfield{Assessment}{literal statement: correct with missing hypotheses; publication match: exact; verification: read the result; development: annotated; confidence: high}
\end{GSitem}
\begin{GSitem}{1986-693.8}{E}{Normal-ordered free-boson realization of Virasoro}
\GSfield{Kind; source}{construction; \GSpage{P001316}}
\GSfield{The notes}{The modes of the Heisenberg field are used to form quadratic \(L_n\); the displayed sum has only \(p>q\).}
\GSfield{Statement}{With \([a_m,a_n]=m\delta_{m+n,0}\), the normally ordered modes \(L_n=\frac12\sum_{k\in\mathbb Z}:a_{n-k}a_k:\) act on a Heisenberg Fock module and satisfy Virasoro at \(c=1\). A background-charge derivative changes the central charge. Omitting the diagonal term when \(n=2k\), or omitting normal ordering and zero-mode data, does not give the same operator.}
\GSfield{Reconstruction}{The conventional \(c=1\) realization is stated exactly; the notebook\textquotesingle s unrestricted sum is replaced by its required completion and ordering.}
\GSfield{Fate}{\citet[Secs. 4.2 and 4.4]{feiginfrenkel1996integrals} develops the Heisenberg vertex-algebra realization; \citet[Eq. (31)]{bazhanovlukyanovzamolodchikov1996integrable} uses the corresponding quantum Miura field. The construction itself predates the report.}
\GSfield{Present status}{As of 12 November 2026 this is an established result under the hypotheses in the modern statement.}
\GSfield{Assessment}{literal statement: false as written; publication match: exact; verification: read the result; development: annotated; confidence: high}
\end{GSitem}
\begin{GSitem}{1986-693.9}{E}{Higher-order scalar operators and Miura factorization}
\GSfield{Kind; source}{construction; \GSpage{P001317}}
\GSfield{The notes}{\(L=\partial^n+u_{n-2}\partial^{n-2}+\cdots+u_0=(\partial+v_1)\cdots(\partial+v_n)\), with \(\sum v_i=0\), is associated with \(\mathfrak{sl}_n\), Wilson, Drinfeld, and Sokolov.}
\GSfield{Statement}{Expanding the ordered product defines a differential-polynomial map from \(n\) first-order fields with sum zero to the coefficients of a monic scalar operator without its \(\partial^{n-1}\) term. The ordering matters because \(\partial f=f\partial+f^{\prime}\). Reversing the map may require a differential-field extension or local solutions; it is not an unconditional global factorization over the original smooth periodic coefficient ring.}
\GSfield{Reconstruction}{The forward construction is distinguished from the stronger, generally obstructed global inverse statement.}
\GSfield{Fate}{\citet[Sec. 3.7, Prop. 3.26]{drinfeldsokolov1985lie} supplies the generalized Miura setting. The reference to Picard--Vessiot theory in the notes is consistent with the field-extension qualification.}
\GSfield{Present status}{As of 12 November 2026 this is an established result under the hypotheses in the modern statement.}
\GSfield{Assessment}{literal statement: correct with missing hypotheses; publication match: exact; verification: read the result; development: none; confidence: high}
\end{GSitem}
\begin{GSitem}{1986-693.10}{B}{The nonlinear W3 algebra and higher quantum Miura fields}
\GSfield{Kind; source}{claim; \GSpage{P001317}}
\GSfield{The notes}{The \(n=3\) second bracket gives a nonlinear algebra; Fateev, Zamolodchikov, and a Feigin report in April 1986 are mentioned.}
\GSfield{Statement}{The principal quantum \(W_3\) algebra has generating fields of spins two and three whose operator products close nonlinearly through normally ordered composites. It is the \(\mathfrak{sl}_3\) instance of quantum Drinfeld--Sokolov reduction. This is different from a finite-dimensional Lie algebra with bilinear structure constants.}
\GSfield{Reconstruction}{The \(W_3\) identification is strong, but the notes do not contain its full operator-product coefficients.}
\GSfield{Fate}{\citet{fateevzamolodchikov1987z3} publishes the matching spin-three symmetry; \citet{feiginfrenkel1990quantization} gives the general \(\mathfrak{sl}_n\) BRST interpretation. The precise content of Feigin\textquotesingle s April conference remark is not separately recoverable.}
\GSfield{Present status}{As of 12 November 2026 this is an established result under the hypotheses in the modern statement.}
\GSfield{Assessment}{literal statement: correct with missing hypotheses; publication match: partial; verification: publisher metadata checked; development: none; confidence: high}
\end{GSitem}
\begin{GSitem}{1986-693.11}{F}{Sklyanin and quantum KdV on a finite interval}
\GSfield{Kind; source}{remark; \GSpage{P001317}}
\GSfield{The notes}{One line mentions Sklyanin\textquotesingle s quantum KdV on a finite interval with boundary conditions.}
\GSfield{Statement}{No endpoint conditions, Hamiltonian, reflection matrix, or representation are specified. Periodic quantum KdV cannot simply be substituted for this boundary problem.}
\GSfield{Fate}{The quantum inverse-scattering context is documented by \citet{bazhanovlukyanovzamolodchikov1996integrable}, but it is not an exact publication match for the unrecorded boundary conditions. The occurrence remains unidentified.}
\GSfield{Present status}{As of 12 November 2026 no precise interval problem is determined by the surviving sentence.}
\GSfield{Assessment}{literal statement: not determinable; publication match: thematic only; verification: read the result; development: none; confidence: low}
\end{GSitem}
\begin{GSitem}{1986-693.12}{E}{The sweep method as a Riccati or Miura transformation}
\GSfield{Kind; source}{method; \GSpage{P001317-P001318}; I. M. Gelfand}
\GSfield{The notes}{Gelfand considers \(y^{\prime\prime}-a(x)y=0\), propagates \(y^{\prime}+\alpha y=0\), and writes \(\alpha^{\prime}-\alpha^2=0\) with an exclamation mark.}
\GSfield{Statement}{On an interval where a chosen solution does not vanish, \(\alpha=-y^{\prime}/y\) satisfies \(\alpha^{\prime}-\alpha^2=-a\). Conversely this Riccati equation yields the factorization \(\partial^2-a=(\partial-\alpha)(\partial+\alpha)\). Zeros of \(y\) create Riccati poles; propagation must then change charts or use a matrix first-order system.}
\GSfield{Reconstruction}{The omitted forcing \(-a\), zero-free interval, and chart limitation are supplied.}
\GSfield{Fate}{The factorization mechanism is the scalar Miura construction discussed in \citet[Sec. 3.7]{drinfeldsokolov1985lie}. \ours\ The corrected sign and nonzero right-hand side follow directly by differentiating \(-y^{\prime}/y\); the scan confirms that zero is the actual notebook reading.}
\GSfield{Present status}{As of 12 November 2026 this is an established result under the hypotheses in the modern statement.}
\GSfield{Assessment}{literal statement: false as written; publication match: exact; verification: read the result; development: annotated; confidence: high}
\end{GSitem}
\begin{GSitem}{1986-693.13}{F}{Higher-order and elliptic boundary propagation}
\GSfield{Kind; source}{remark; \GSpage{P001318}; I. M. Gelfand}
\GSfield{The notes}{Gelfand says the sweep method also works in higher order and sketches shrinking a contour for \(\Delta u=0\), mentioning Keldysh.}
\GSfield{Statement}{Higher-order scalar ODEs can be converted to first-order systems and boundary subspaces propagated by their fundamental solution. For elliptic PDEs, contour propagation leads to boundary operator families and is not the same finite-dimensional Riccati equation. The note gives no domain, regularity, solvability, or operator formula for the PDE version.}
\GSfield{Reconstruction}{The finite-dimensional ODE content is separated from the unidentified PDE claim.}
\GSfield{Fate}{The ODE part is standard and compatible with the factorization context of \citet{drinfeldsokolov1985lie}; no exact Keldysh theorem is identified with the elliptic sketch.}
\GSfield{Present status}{As of 12 November 2026 the ODE propagation mechanism is established; this specific elliptic-boundary assertion lacks sufficient data for assessment.}
\GSfield{Assessment}{literal statement: not determinable; publication match: thematic only; verification: read the result; development: none; confidence: medium}
\end{GSitem}
\begin{GSitem}{1986-693.14}{E}{Feigin: invariant forms, Casimirs, and loop algebras}
\GSfield{Kind; source}{construction; \GSpage{P001318}; B. L. Feigin}
\GSfield{The notes}{Feigin writes the quadratic Casimir \(\sum e_i e^i\), loop \(\mathfrak{sl}_n\), and scalar products obtained from coefficients of \(\operatorname{tr}(uv)\); higher Casimirs are linked to Gelfand--Dickey Hamiltonians.}
\GSfield{Statement}{For a finite-dimensional Lie algebra with nondegenerate invariant symmetric form, dual bases define a central quadratic Casimir in its enveloping algebra. On the loop algebra, coefficient or residue pairings \(\operatorname{Res}\operatorname{tr}(u(z)v(z))z^m dz\) are invariant. Higher invariant polynomials, after Hamiltonian reduction, give classical \(W\)-fields; quantum versions require normal ordering and corrections.}
\GSfield{Reconstruction}{Nondegeneracy, invariance, and the distinction between central finite-dimensional Casimirs and reduced local fields are supplied.}
\GSfield{Fate}{\citet[Secs. 2.4 and 4.6]{feiginfrenkel1996integrals} explains the classical and quantum reductions. The bare Casimir construction is old; the exact unpublished formula allegedly given in April is not filled in.}
\GSfield{Present status}{As of 12 November 2026 this is an established result under the hypotheses in the modern statement.}
\GSfield{Assessment}{literal statement: correct with missing hypotheses; publication match: partial; verification: read the result; development: none; confidence: high}
\end{GSitem}
\begin{GSitem}{1986-693.15}{C}{Spectrum of quantum KdV charges on irreducible Virasoro modules}
\GSfield{Kind; source}{problem; \GSpage{P001318}}
\GSfield{The notes}{Question 1(b) asks to calculate the spectrum in irreducible representations of the Virasoro algebra, immediately after quantizing the KdV integrals.}
\GSfield{Statement}{Fix \(c,h\in\mathbb C\), the irreducible highest-weight module \(L(c,h)\), and normalized commuting local KdV charges. Determine their joint generalized eigenspaces and nilpotent actions grade by grade, and a complete all-grade parameterization where available. An ordinary eigenbasis requires semisimplicity, which is not automatic outside a positive unitary setting.}
\GSfield{Reconstruction}{The spectrum is read as the spectrum of the quantum KdV family, not merely \(L_0\). The highest-weight restriction and generalized-eigenspace target are modern specifications absent from the notebook.}
\GSfield{Fate}{\citet[Eq. (11)]{bazhanovlukyanovzamolodchikov1996integrable} supplies explicit charges and a spectral program. The opened v1 of \citet[Defs. 6.5--6.6, Prop. 6.7, Thm. 6.8, Sec. 7]{masoeroruzza2026odeim} gives a 2026 preprint completeness theorem at \(c=-2\), \(h=(\beta^2-1)/8\), \(\beta\notin\mathbb Z\), for admissible QQ solutions modulo entire-function gauge, and diagonalizes the first three charges in its free-field realization. This does not settle arbitrary \(c,h\).}
\GSfield{Present status}{As of 12 November 2026 the general formulation remains only partly resolved. \oursComputed\ This chapter computes the level-two charge and proves a nonunitary Jordan block at \(c=1,h=-1/32\) survives in the irreducible quotient; the result is explicitly a bounded calculation, with no historical novelty claim.}
\GSfield{Assessment}{literal statement: correct; publication match: partial; verification: read the result; development: developed in this chapter; confidence: high; problem \hyperref[prob:GP-1986-33]{GP-1986-33}}
\GSfield{Worth developing}{Worth a bounded development: exact low-grade matrices and singular-vector quotient checks reveal semisimplicity obstructions without claiming all-grade completeness.}
\end{GSitem}
\begin{GSitem}{1986-693.16}{B}{Quantization of the general second Gelfand--Dickey bracket}
\GSfield{Kind; source}{problem; \GSpage{P001318}}
\GSfield{The notes}{Question 2 asks for quantization of the second Gelfand--Dickey bracket in the general case.}
\GSfield{Statement}{For monic scalar differential operators of order \(n\), with the \(\partial^{n-1}\) coefficient removed, construct a vertex or completed mode algebra deforming the second Adler--Gelfand--Dickey Poisson algebra. In the principal \(\mathfrak{sl}_n\) setting this is the quantum \(W_n\) algebra; the precise reduction and deformation parameter must be specified.}
\GSfield{Reconstruction}{The scalar principal-reduction setting is the strongest explicit setting supported by the preceding \(\mathfrak{sl}_n\) discussion.}
\GSfield{Fate}{\citet{feiginfrenkel1990quantization} identifies quantum Drinfeld--Sokolov reduction of affine \(\mathfrak{sl}_n\) with \(W_n\). \citet[Thm. 4.6.9 and Sec. 4.6.10]{feiginfrenkel1996integrals} proves the more general simple-Lie-algebra construction and explains its BRST relation.}
\GSfield{Present status}{As of 12 November 2026 the algebraic quantization in the principal scalar-operator setting is established. This statement does not claim that every global analytic boundary-value problem is thereby solved.}
\GSfield{Assessment}{literal statement: correct; publication match: exact; verification: read the result; development: none; confidence: high; problem \hyperref[prob:GP-1986-34]{GP-1986-34}}
\end{GSitem}

\GStalk{694}{1986-11-10}{A. A. Beilinson}{}{Virasoro uniformization and differential equations on moduli of curves}
The catalogue gives no page range, but the original P001319 heading clearly names Beilinson and 10 November 1986. The coherent continuation P001319--P001321 supplies surviving notes: formal-coordinate uniformization of pointed curves, determinant-line twists, and differential equations obtained from singular Virasoro vectors. This attribution is an explicit editorial recovery from the manuscript; the empty catalogue page field is retained. The note credits Kontsevich and Beilinson--Schechtman independently and refers to Belavin through Drinfeld.

\begin{GSitem}{1986-694.1}{E}{The universal pointed curve and its dimensions}
\GSfield{Kind; source}{construction; \GSpage{P001319}}
\GSfield{The notes}{The diagram \(X\to M\) has a genus-\(g\) curve as fiber; the displayed dimension is \(6g-6\).}
\GSfield{Statement}{For \(g\ge2\), the smooth moduli stack \(\mathcal M_g\) has complex dimension \(3g-3\), hence real dimension \(6g-6\). Its universal curve is \(\mathcal M_{g,1}\to\mathcal M_g\), with total complex dimension \(3g-2\). On a coarse moduli space a universal family need not exist without a level structure or stack interpretation.}
\GSfield{Reconstruction}{The genus restriction, complex versus real dimension, and stack interpretation are supplied.}
\GSfield{Fate}{This is classical moduli theory, used explicitly in \citet[Sec. 1.1 and Sec. 4.2]{benzvifrenkel2004geometric}. The notebook dimension is correct only as a real dimension of the base.}
\GSfield{Present status}{As of 12 November 2026 this is an established result under the hypotheses in the modern statement.}
\GSfield{Assessment}{literal statement: correct with missing hypotheses; publication match: exact; verification: read the result; development: annotated; confidence: high}
\end{GSitem}
\begin{GSitem}{1986-694.2}{E}{Cotangent and Hodge determinant line twists}
\GSfield{Kind; source}{construction; \GSpage{P001319}}
\GSfield{The notes}{For integers \(j,c\), the fiber at \((K,x)\) is \((\Omega^1_{K,x})^{\otimes j}\otimes[\det\Omega^1(K)]^c\).}
\GSfield{Statement}{On \(\mathcal M_{g,1}\), write \(\psi\) for the cotangent line at the marked point and \(\lambda=\det\pi_*\omega_\pi\) for the Hodge line pulled back from \(\mathcal M_g\). The displayed bundle is \(\psi^{\otimes j}\otimes\lambda^{\otimes c}\), with negative integral powers interpreted as duals. A complex central parameter generally defines a twist of differential operators, not an ordinary fractional tensor power of a line bundle.}
\GSfield{Reconstruction}{The notation \(\Omega^1(K)\) is interpreted as the finite-dimensional space of global holomorphic differentials. The exponent \(c\) is not identified numerically with a Virasoro central charge without fixing the central-extension convention.}
\GSfield{Fate}{\citet{beilinsonschechtman1988determinant}, received 12 February 1988 and published that December, develops determinant bundles and Virasoro action. \citet[Sec. 4.2]{benzvifrenkel2004geometric} treats both actual line bundles and arbitrary complex twists.}
\GSfield{Present status}{As of 12 November 2026 this is an established result under the hypotheses in the modern statement.}
\GSfield{Assessment}{literal statement: correct with missing hypotheses; publication match: partial; verification: read the result; development: none; confidence: high}
\end{GSitem}
\begin{GSitem}{1986-694.3}{A}{Formal-coordinate uniformization and the Virasoro lift}
\GSfield{Kind; source}{theorem; \GSpage{P001319-P001320}}
\GSfield{The notes}{The space \(\widehat X\ni(K,x,q)\) carries an action of formal vector fields; \(t\mathbb C[[t]]\partial_t\) changes the coordinate and the remaining modes vary the curve.}
\GSfield{Statement}{The formal-coordinate torsor over \(\mathcal M_{g,1}\) has a transitive infinitesimal action of \(\operatorname{Der}\mathbb C((t))\), compatible with the coordinate-change group \(\operatorname{Aut}\mathbb C[[t]]\). The central generator acts trivially on the base; the lift to determinant-line differential operators is governed by the Virasoro central extension. This is an infinitesimal algebraic statement, not a global action of an unspecified infinite-dimensional Lie group.}
\GSfield{Reconstruction}{The distinction between the centerless action on coordinates and the central lift on a line twist makes the notebook shorthand precise.}
\GSfield{Fate}{The directly corresponding publication is \citet{beilinsonschechtman1988determinant}. The action and its determinant-line lift are stated and used in \citet[Sec. 1.1 and Prop. 4.2.3]{benzvifrenkel2004geometric}, first circulated January 2003. The precise independent Kontsevich priority is recorded from the notebook, without an identified separate manuscript. \citet[Thm. 1.1]{esnaulttsai2000determinant}, circulated October 1999 and published in 2000, corrects the original paper\textquotesingle s Theorem 2.3.1(ii), a divisor-twist reduction step; this later repair does not invalidate the uniformization principle.}
\GSfield{Present status}{As of 12 November 2026 this is an established result under the hypotheses in the modern statement.}
\GSfield{Assessment}{literal statement: correct with missing hypotheses; publication match: exact; verification: read the result; development: none; confidence: high}
\end{GSitem}
\begin{GSitem}{1986-694.4}{E}{Kodaira--Spencer tangent space through punctured-curve vector fields}
\GSfield{Kind; source}{method; \GSpage{P001320}}
\GSfield{The notes}{The tangent quotient is written as \(\mathit{Vir}/\Phi_q(\mathcal T(K\backslash X))\), then with the point \(x\) in the denominator; two horizontal lifts are compared.}
\GSfield{Statement}{For a smooth pointed curve \(C,p\) with formal coordinate \(q\), the tangent space to the coordinate torsor is the quotient \(\operatorname{Der}\mathbb C((q))/H^0(C\setminus\{p\},T_C)\), where algebraic vector fields have finite-order poles at \(p\) and are expanded in \(q\). Comparing trivializations of an infinitesimal deformation on the formal disc and punctured complement gives the quotient; changing the complement lift adds a global punctured-curve vector field.}
\GSfield{Reconstruction}{The correct point removal, algebraic rather than arbitrary holomorphic vector fields, and finite Laurent pole condition are supplied.}
\GSfield{Fate}{\citet[Sec. 4.2]{benzvifrenkel2004geometric} gives the uniformization and stabilizer formulation. The notebook itself explains the two-lift mechanism. The capital \(X\) in the first denominator is incompatible with the adjoining definitions.}
\GSfield{Present status}{As of 12 November 2026 this is an established result under the hypotheses in the modern statement.}
\GSfield{Assessment}{literal statement: false as written; publication match: exact; verification: read the result; development: annotated; confidence: high}
\end{GSitem}
\begin{GSitem}{1986-694.5}{A}{Singular Virasoro vectors produce canonical differential operators}
\GSfield{Kind; source}{method; \GSpage{P001321}}
\GSfield{The notes}{A reducible Verma module \(M_{(c,g)}\) contains a highest-weight vector \(v^{\prime}=L(v)\); applying \(L\in U(\mathit{Vir})\) to a lifted section gives an equation between line twists.}
\GSfield{Statement}{Given a singular vector of definite grade in a Virasoro highest-weight module, localization on the formal-coordinate torsor, followed by equivariant descent, produces a differential operator between the appropriate twisted sheaves. Contragredient conventions determine the direction and weight shift. The construction requires compatibility with the coordinate-change action; a section is represented by an equivariant function, not literally by a single vector of a fixed Verma module.}
\GSfield{Reconstruction}{The highest weight written \(g\) is renamed \(h\) when needed to avoid confusing it with curve genus. No unrecorded central-charge-to-Hodge-power formula is imposed.}
\GSfield{Fate}{The genus-zero null-vector mechanism occurs in \citet{belavinpolyakovzamolodchikov1984infinite}. Its moduli-space globalization belongs to \citet{beilinsonschechtman1988determinant}; \citet[Secs. 2.2--2.4 and 4.2]{benzvifrenkel2004geometric} gives a precise localization framework.}
\GSfield{Present status}{As of 12 November 2026 this is an established result under the hypotheses in the modern statement.}
\GSfield{Assessment}{literal statement: correct with missing hypotheses; publication match: partial; verification: read the result; development: none; confidence: high}
\end{GSitem}
\begin{GSitem}{1986-694.6}{F}{Finite-rank solution local systems require a specified conformal-block sector}
\GSfield{Kind; source}{claim; \GSpage{P001321}}
\GSfield{The notes}{The closing sentence says that solutions of the singular-vector equations form a finite-rank local system.}
\GSfield{Statement}{A system of differential equations gives a finite-rank local system on a smooth locus when its corresponding differential module is an integrable connection of finite rank there. Mere reducibility of one arbitrary Virasoro Verma module does not establish this property. To identify the intended conformal-block theorem one needs the representations, insertions, quotient relations, and finiteness hypotheses, none of which are recorded.}
\GSfield{Reconstruction}{The missing finiteness condition is exposed rather than supplied by silently choosing a minimal model.}
\GSfield{Fate}{The null-vector equations of \citet{belavinpolyakovzamolodchikov1984infinite} and localization of \citet{beilinsonschechtman1988determinant} explain the context. They are not evidence for an unconditional finite-rank theorem for every reducible Verma module.}
\GSfield{Present status}{As of 12 November 2026 finite-rank conformal-block systems exist under substantial additional hypotheses, but this particular unrestricted closing assertion is insufficiently specified.}
\GSfield{Assessment}{literal statement: not determinable; publication match: thematic only; verification: indirect; development: none; confidence: medium}
\end{GSitem}

\GStalk{695}{1986-11-17}{V. A. Kazakov}{P001321-P001326}{Random surfaces, Gaussian embeddings, and coupled matrix integrals}
The actual section is P001322--P001325, jointly headed Kazakov and Zakharevich; interventions by Gelfand are marked. It compares the bosonic-string action with triangulated Gaussian surfaces, revisits weighted planar-map counting, and presents a two-source unitary integral and three-state matrix model. The final three- and four-dimensional gauge question repeats record 688. The lattice-theory report beginning P001326 belongs to record 696.

\begin{GSitem}{1986-695.1}{A}{Bosonic action and Gaussian surfaces on triangulations}
\GSfield{Kind; source}{construction; \GSpage{P001322}}
\GSfield{The notes}{The continuum expression integrates \(X:S\to\mathbb R^n\) and a metric \(g\), then replaces it by a sum over triangulations with an edge energy \(\sum_{i\sim j}|x_i-x_j|^2\).}
\GSfield{Statement}{For a fixed finite connected triangulation, assign target vectors to vertices and a positive quadratic graph energy. A random-surface model sums the resulting Gaussian integrals over a specified map class with area and symmetry weights. The formal continuum Polyakov integral additionally needs a measure, gauge fixing, and treatment of zero modes; the displayed discretization is not a proved equality of continuum measures.}
\GSfield{Reconstruction}{The graph class, weights, and distinction between a finite Gaussian model and a formal metric integral are supplied.}
\GSfield{Fate}{\citet{polyakov1981quantum} gives the continuum motivation. \citet{boulatovkazakovkostovmigdal1986surfaces}, circulated as PR-1123 in September 1986 and published 29 December 1986, analyzes this dynamical-triangulation model.}
\GSfield{Present status}{As of 12 November 2026 this is an established result under the hypotheses in the modern statement.}
\GSfield{Assessment}{literal statement: correct with missing hypotheses; publication match: partial; verification: publisher metadata checked; development: none; confidence: high}
\end{GSitem}
\begin{GSitem}{1986-695.2}{E}{Vertex area weights are not curvature defects}
\GSfield{Kind; source}{claim; \GSpage{P001322-P001323}}
\GSfield{The notes}{The notes set \(\sigma_i=n_i/3\), call it a defect, and explain it by dividing each regular triangle among its three vertices.}
\GSfield{Statement}{If each equilateral triangle has unit area and \(n_i\) triangles meet at vertex \(i\), its barycentric area allocation is \(n_i/3\), and summing these allocations counts faces. The angular curvature defect is instead \(2\pi-n_i\pi/3\) at an interior vertex. The two quantities encode different geometric data.}
\GSfield{Reconstruction}{The elementary incidence and angle calculation is an editorial correction, explicitly marked as ours.}
\GSfield{Fate}{The discrete area and curvature observables are distinguished in \citet{boulatovkazakovkostovmigdal1986surfaces}. \ours\ Counting triangle--vertex incidences gives the area identity and directly explains the notebook diagram; it does not turn \(n_i/3\) into an angle defect.}
\GSfield{Present status}{As of 12 November 2026 this is an established result under the hypotheses in the modern statement.}
\GSfield{Assessment}{literal statement: false as written; publication match: partial; verification: indirect; development: annotated; confidence: high}
\end{GSitem}
\begin{GSitem}{1986-695.3}{E}{Piecewise affine Dirichlet energy and the equilateral restriction}
\GSfield{Kind; source}{claim; \GSpage{P001322-P001323}}
\GSfield{The notes}{The integrands are said to agree for regular embeddings, and a fixed metric with a linear map on each triangle is proposed.}
\GSfield{Statement}{For a piecewise affine map on a Euclidean triangulated surface, its Dirichlet energy is a sum of edge squared differences weighted by cotangents of the opposite angles, up to the chosen factor \(1/2\). Equilateral triangles make these weights constant. For an arbitrary intrinsic triangulation the unweighted graph energy is a different discretization; target images being equilateral is not the necessary intrinsic-metric condition.}
\GSfield{Reconstruction}{The intrinsic domain metric is separated from the shape of its target image; the edge-energy constant must be fixed before claiming equality.}
\GSfield{Fate}{The dynamical-triangulation model of \citet{boulatovkazakovkostovmigdal1986surfaces} fixes regular intrinsic triangles. The statement in the notes is a normalization-level discretization observation, not a theorem of convergence to the continuum path integral.}
\GSfield{Present status}{As of 12 November 2026 this is an established result under the hypotheses in the modern statement.}
\GSfield{Assessment}{literal statement: correct with missing hypotheses; publication match: partial; verification: indirect; development: none; confidence: high}
\end{GSitem}
\begin{GSitem}{1986-695.4}{E}{Gauge topology does not require abandoning cubic lattices}
\GSfield{Kind; source}{question; \GSpage{P001323}; I. M. Gelfand}
\GSfield{The notes}{Gelfand asks how to write Yang--Mills when characteristic classes appear, saying that one cannot restrict to a cubic lattice.}
\GSfield{Statement}{The relevant precise issue is a geometric definition of topological charge for lattice gauge fields with a controlled continuum limit. A hypercubic cellulation does not preclude such a definition: admissibility or small-curvature restrictions can separate topological sectors. The notebook does not specify a different proposed characteristic-class discretization.}
\GSfield{Reconstruction}{The topological-charge formulation is the narrowest concrete interpretation supported by the intervention; it is not an invented general open gauge problem.}
\GSfield{Fate}{\citet{luscher1982topology}, circulated June 1981 and published in 1982, constructs an integer topological charge for sufficiently small-action-density \(SU(2)\) fields on a four-dimensional hypercubic lattice, invariant under continuous admissible deformations and converging to the continuum Chern number. Thus the claimed obstruction to cubic lattices is too strong.}
\GSfield{Present status}{As of 12 November 2026 this is an established result under the hypotheses in the modern statement.}
\GSfield{Assessment}{literal statement: false as written; publication match: partial; verification: publisher metadata checked; development: none; confidence: high}
\end{GSitem}
\begin{GSitem}{1986-695.5}{F}{High-dimensional concentration does not make every graph triangle equilateral}
\GSfield{Kind; source}{claim; \GSpage{P001323}; I. S. Zakharevich}
\GSfield{The notes}{Zakharevich says a random embedding of any triangulated sphere makes all its triangles regular on average, appealing to concentration near a shell.}
\GSfield{Statement}{Independent identically distributed isotropic vertex vectors can have pair distances concentrated at a common scale in large target dimension. Coordinates distributed by a graph Gaussian, however, have edge variances determined by the graph covariance and effective resistances; those variances need not be equal. A shell-concentration slogan alone therefore does not justify the asserted equilateral replacement.}
\GSfield{Reconstruction}{The two different random-embedding laws are distinguished editorially; no new asymptotic result is claimed.}
\GSfield{Fate}{\citet{boulatovkazakovkostovmigdal1986surfaces} studies limiting target dimensions, but its verified abstract does not supply this universal equilateral theorem. No probability law, limit, or error estimate is stated in the notebook.}
\GSfield{Present status}{As of 12 November 2026 concentration statements are available for specified distributions; this unrestricted heuristic is not an identified theorem.}
\GSfield{Assessment}{literal statement: not determinable; publication match: thematic only; verification: indirect; development: none; confidence: medium}
\end{GSitem}
\begin{GSitem}{1986-695.6}{E}{Target-dimension factorization and removal of translation modes}
\GSfield{Kind; source}{method; \GSpage{P001323}; I. S. Zakharevich}
\GSfield{The notes}{The \(n\)-dimensional embedding integral is written as the \(n\)-th power of a one-dimensional Gaussian integral; the parameter \(n\) is then continued analytically.}
\GSfield{Statement}{Fix one vertex of a connected graph and write its scalar energy as \(x^{\mathsf T}L_*x\), with positive reduced Laplacian \(L_*\). The scalar integral is \(\pi^{(v-1)/2}(\det L_*)^{-1/2}\), and \(n\) independent coordinates give its \(n\)-th power for integer \(n\ge0\). The positive determinant permits a specified entire exponential in complex \(n\) for each fixed graph. Analyticity of an infinite map sum requires separate convergence control.}
\GSfield{Reconstruction}{The pinned vertex and measure normalization are explicit editorial repairs; the unpinned integral is infinite.}
\GSfield{Fate}{The Gaussian-embedding model is the one analyzed by \citet{boulatovkazakovkostovmigdal1986surfaces}. \ours\ Diagonalizing \(L_*\) gives the displayed finite Gaussian factor and exposes the missing translational zero-mode removal.}
\GSfield{Present status}{As of 12 November 2026 this is an established result under the hypotheses in the modern statement.}
\GSfield{Assessment}{literal statement: false as written; publication match: partial; verification: indirect; development: annotated; confidence: high}
\end{GSitem}
\begin{GSitem}{1986-695.7}{A}{The solvable dimensions zero and minus two}
\GSfield{Kind; source}{claim; \GSpage{P001323}}
\GSfield{The notes}{Kazakov says the sum can be calculated at \(n=0\) and \(n=-2\).}
\GSfield{Statement}{At \(n=0\) the matter determinant disappears, leaving weighted pure planar maps. At \(n=-2\), after fixing the Gaussian normalization, the weight contains \(\det L_*\); the matrix-tree theorem identifies this determinant with the number of spanning trees of the graph. Negative target dimension is determinant analytic continuation, not a Euclidean embedding in a negative-dimensional space.}
\GSfield{Reconstruction}{The determinant interpretation supplies the precise meaning of the formal negative dimension; the word only is retained as a 1986 contextual remark, not an eternal no-go theorem.}
\GSfield{Fate}{Pure-map solvability is explicit in \citet[Sec. 4]{brezinitzyksonparisizuber1978planar}. \citet{boulatovkazakovkostovmigdal1986surfaces} identifies \(D=-2\) as exactly solvable and reports its logarithmic mean-square-size behavior. The article appeared in December 1986 after September preprint circulation.}
\GSfield{Present status}{As of 12 November 2026 this is an established result under the hypotheses in the modern statement.}
\GSfield{Assessment}{literal statement: correct with missing hypotheses; publication match: partial; verification: publisher metadata checked; development: none; confidence: high}
\end{GSitem}
\begin{GSitem}{1986-695.8}{E}{Rational map weights and the apparent noninteger counting formula}
\GSfield{Kind; source}{question; \GSpage{P001323-P001324}; V. A. Kazakov and I. S. Zakharevich}
\GSfield{The notes}{The cubic series is first called an unlabelled count. At \(k=100\), \(72^{100}149!/(102!50!)\) is observed not to be an integer. Zakharevich replaces the count by a labelled count divided by \(n!\).}
\GSfield{Statement}{Connected matrix free energies enumerate maps with inverse-automorphism weights, so rational coefficients are expected. For a labelling set on which the combinatorial automorphism action is faithful, the orbit--stabilizer relation gives \[\#\{\text{labelled maps}\}/n!=\sum_{[T]}1/|\operatorname{Aut}T|\]. The exact labels must include the flag or cyclic-order data used by the Feynman model. The negative sign in the repeated series belongs to the source convention for negative free energy, not a positive counting function.}
\GSfield{Reconstruction}{The weighting and sign conventions repair the count; triangle labels alone are not asserted to remove every possible residual automorphism.}
\GSfield{Fate}{\citet[Secs. 2 and 4]{brezinitzyksonparisizuber1978planar} gives the graph weights and cubic series. \oursComputed\ The 5-adic valuation of the displayed \(k=100\) rational number is \(-1\), confirming its nonintegrality; the exact check is saved with the chapter evidence.}
\GSfield{Present status}{As of 12 November 2026 this is an established result under the hypotheses in the modern statement.}
\GSfield{Assessment}{literal statement: false as written; publication match: exact; verification: read the result; development: annotated; confidence: high}
\end{GSitem}
\begin{GSitem}{1986-695.9}{E}{Critical exponents survive only regular changes of coupling}
\GSfield{Kind; source}{claim; \GSpage{P001324}}
\GSfield{The notes}{The notes transfer a \(5/2\) singularity in \(b\) to one in \(\beta\), although \(b=b(\beta)\) is not evaluated; the coefficient power is \(-7/2\).}
\GSfield{Statement}{The singular exponent transfers unchanged if \(b(\beta)\) is analytic near criticality and \(b^{\prime}(\beta_*)\ne0\). If its first nonzero derivative has order \(m\), the composed singularity instead has power \(5m/2\). The coefficient asymptotic \(k^{-7/2}\) pertains to the weighted unrooted pure-map series, with its parity convention, and includes an exponential growth factor.}
\GSfield{Reconstruction}{Analytic background terms and the nonvanishing derivative condition are supplied; the note does not determine \(b(\beta)\).}
\GSfield{Fate}{\citet[Sec. 4]{brezinitzyksonparisizuber1978planar} establishes the cubic singularity. \ours\ Substituting the first nonzero term of the Taylor expansion of \(b-b_*\) gives the necessary change-of-variable qualification.}
\GSfield{Present status}{As of 12 November 2026 this is an established result under the hypotheses in the modern statement.}
\GSfield{Assessment}{literal statement: correct with missing hypotheses; publication match: partial; verification: read the result; development: annotated; confidence: high}
\end{GSitem}
\begin{GSitem}{1986-695.10}{E}{Normalized planar logarithm and higher-valency map potentials}
\GSfield{Kind; source}{method; \GSpage{P001324}}
\GSfield{The notes}{The displayed generating function is \(Z=\lim_{N\to\infty}\log I_N\); the potential may be extended from cubic to cubic plus quartic.}
\GSfield{Statement}{The connected planar free energy is the formal limit \[\lim_{N\to\infty}N^{-2}\log[I_N(g,\lambda)/I_N(0,0)]\], with a fixed Gaussian normalization. A term \(\operatorname{tr}M^k\) inserts a vertex of valency \(k\) in the dual ribbon graph. Cubic and quartic terms therefore mix triangle and quadrilateral faces in the dual map. A positive stabilizing leading coefficient gives a convergent real integral; an unstable cubic potential alone remains formal.}
\GSfield{Reconstruction}{The factor \(N^{-2}\), Gaussian denominator, and potential sign conditions repair the displayed limit.}
\GSfield{Fate}{\citet[Secs. 2--4]{brezinitzyksonparisizuber1978planar} treats these matrix integrals and the planar expansion. This repeats 1986-682.2--.3, with an additional quartic coupling.}
\GSfield{Present status}{As of 12 November 2026 this is an established result under the hypotheses in the modern statement.}
\GSfield{Assessment}{literal statement: false as written; publication match: exact; verification: read the result; development: annotated; confidence: high}
\end{GSitem}
\begin{GSitem}{1986-695.11}{B}{Two diagonal source terms and the principal chiral model}
\GSfield{Kind; source}{question; \GSpage{P001325}}
\GSfield{The notes}{Equation (12) has \(\operatorname{tr}(p_1U^{-1}p_1U)+\operatorname{tr}(p_2U^{-1}p_2U)\), with diagonal \(p_1,p_2\), and is identified with a two-dimensional principal chiral field at large \(N\).}
\GSfield{Statement}{For finite \(N\) and specified source matrices the compact integral is entire in the couplings, and its Taylor coefficients are exactly given by unitary Weingarten calculus. A planar chiral-field equivalence additionally needs a prescribed source eigenvalue distribution, scaling, and a justified reduction or quenching procedure. Exact field-theory information does not automatically evaluate the unspecified finite integral.}
\GSfield{Reconstruction}{The distinction between an exact compact integral and a continuum planar equivalence is editorial; the latter is not assigned an artificial solved status.}
\GSfield{Fate}{\citet[Cor. 2.4]{collinssniady2006integration} supplies a convergent finite-\(N\) coefficient method. This is the same two-source integral family as 1986-688.5 and GP-1986-60; the chapter development gives explicit bounded coefficients and a remainder estimate.}
\GSfield{Present status}{As of 12 November 2026 exact finite-dimensional coefficient evaluation is established; the notebook does not specify a unique continuum reduction theorem.}
\GSfield{Assessment}{literal statement: correct with missing hypotheses; publication match: partial; verification: read the result; development: developed in this chapter; confidence: high; problem \hyperref[prob:GP-1986-60]{GP-1986-60}}
\end{GSitem}
\begin{GSitem}{1986-695.12}{B}{Three-state spins and a three-matrix potential}
\GSfield{Kind; source}{construction; \GSpage{P001325}}
\GSfield{The notes}{Equation (13) uses three Hermitian matrices, symmetric pair couplings, and quartic self-interactions, with \(Z_3\) spins said to replace Ising spins.}
\GSfield{Statement}{Three matrix species with equal unlike-species covariance encode a three-state Potts interaction on the corresponding ribbon graphs. A quartic vertex generates four-valent dual graphs, so calling the same model a sum of purely triangular dual faces needs correction or a duality convention. With the printed sign, a positive quartic coupling destabilizes the real integral; the Feynman expansion is formal.}
\GSfield{Reconstruction}{The Potts interpretation, graph valency, and formal-integral qualification prevent a false exact publication match.}
\GSfield{Fate}{\citet[Sec. 1]{daul1995potts}, written November 1994 and posted February 1995, treats the \(q\)-state random-planar Potts matrix model and computes three-state scaling information. Its displayed potential is cubic; it is not an exact all-coupling evaluation of the notebook quartic integral.}
\GSfield{Present status}{As of 12 November 2026 this is an established result under the hypotheses in the modern statement.}
\GSfield{Assessment}{literal statement: false as written; publication match: partial; verification: read the result; development: none; confidence: high}
\end{GSitem}
\begin{GSitem}{1986-695.13}{E}{Angular variables around a three-matrix cycle are constrained}
\GSfield{Kind; source}{claim; \GSpage{P001325}}
\GSfield{The notes}{After diagonalization, equation (14) integrates three pairwise angular matrices independently and its final factor is again \(\mu\).}
\GSfield{Statement}{If \(M_i=V_i\Lambda_iV_i^*\), set \(U_{12}=V_1^*V_2\), \(U_{23}=V_2^*V_3\), \(U_{31}=V_3^*V_1\). Then \(U_{12}U_{23}U_{31}=I\); only two are independent after removing common conjugation. The third spectral pair is \(\nu,\lambda\), not a second \(\nu,\mu\) pair. A product of three independent HCIZ integrals therefore gives a different model.}
\GSfield{Reconstruction}{The shared-unitary constraint and final eigenvalue factor are repaired; the scan confirms the repeated \(\mu\) is the literal reading.}
\GSfield{Fate}{\citet[Sec. 1, discussion preceding Eq. (3)]{daul1995potts} explicitly identifies the lack of independence of cyclic relative angles and introduces an auxiliary Gaussian matrix to remove the coupling obstruction. \ours\ The displayed closure identity follows by multiplying the three relative matrices.}
\GSfield{Present status}{As of 12 November 2026 this is an established result under the hypotheses in the modern statement.}
\GSfield{Assessment}{literal statement: false as written; publication match: exact; verification: read the result; development: annotated; confidence: high}
\end{GSitem}
\begin{GSitem}{1986-695.14}{C}{The higher-dimensional gauge integral recurs}
\GSfield{Kind; source}{problem; \GSpage{P001325}}
\GSfield{The notes}{The one-site unitary plaquette integral is repeated; \(d=2\) is called solved and \(d=3,4\) are requested.}
\GSfield{Statement}{The problem is the nonperturbative planar free energy and Wilson-loop observables of a fixed \(d\)-matrix plaquette model in dimensions three and four, together with a justified center-symmetric volume-reduction regime. The printed action also needs a real part, or an explicit sum over both plaquette orientations, for the usual real Euclidean gauge weight.}
\GSfield{Reconstruction}{The phase, orientation convention, and symmetry conditions are needed to formulate the literal target accurately.}
\GSfield{Fate}{\citet{eguchikawai1982reduction} is the original reduction; \citet[Secs. 1--2]{garciaperez2020largeN} explains the center-symmetry obstruction and modern variants. This is a repeat of GP-1986-31, not a new independent problem.}
\GSfield{Present status}{As of 12 November 2026 the general three- and four-dimensional analytic planar solution is not supplied by the established reduction methods; the original weak-coupling four-dimensional untwisted equivalence fails.}
\GSfield{Assessment}{literal statement: correct with missing hypotheses; publication match: partial; verification: read the result; development: none; confidence: high; problem \hyperref[prob:GP-1986-31]{GP-1986-31}}
\GSfield{Worth developing}{Worth a bounded coefficient or symmetry calculation for a specified model; the unrestricted all-coupling higher-dimensional planar solution is a much larger target.}
\end{GSitem}

\GStalk{696}{1986-11-17}{I. M. Gelfand}{P001324-P001329}{Subspace configurations and preprojective reductions of modular lattices}
The manuscript section P001326--P001327 names a joint project with B. V. Lidskii and V. A. Ponomarev. It reviews subspace configurations and the free modular lattice, distinguishes word equality from linear equivalence, and announces a preprojective quotient and perfect elements. The printed 1987 paper supplies necessary restrictions omitted in the notebook. Finite axiomatizability, decidability, and classification of representations are different questions.

\begin{GSitem}{1986-696.1}{E}{Projective geometries and continuous geometries as lattice antecedents}
\GSfield{Kind; source}{remark; \GSpage{P001326}}
\GSfield{The notes}{Gelfand lists Dedekind, Weyl on projective axioms, and von Neumann on factors; dimension is described by lengths of chains, with Desargues in dimension at least three.}
\GSfield{Statement}{The lattice of subspaces of an \(n\)-dimensional vector space is complemented and modular with maximal strict-chain length \(n\). Projective dimension is \(n-1\). Incidence projective spaces of projective dimension at least three are Desarguesian; arbitrary projective planes need not be. Von Neumann continuous geometries replace finite rank by a continuous dimension theory under additional lattice-completeness conditions.}
\GSfield{Reconstruction}{The vector-space/projective-dimension shift and the distinction between finite projective geometry and continuous geometry are made explicit.}
\GSfield{Fate}{The coordinatization framework is developed and precisely recalled in \citet[Introduction and Sec. 2]{herrmann1984arithmetic}. The notebook compresses several classical theories and does not specify the separate work attributed to the uncertain name Disona.}
\GSfield{Present status}{As of 12 November 2026 this is an established result under the hypotheses in the modern statement.}
\GSfield{Assessment}{literal statement: correct with missing hypotheses; publication match: partial; verification: read the result; development: none; confidence: high}
\end{GSitem}
\begin{GSitem}{1986-696.2}{E}{Two and three subspaces: finite indecomposable types}
\GSfield{Kind; source}{theorem; \GSpage{P001326}}
\GSfield{The notes}{Two subspaces are called trivial; three decompose into one-dimensional pieces and a configuration of three lines. The associated free modular lattice has 28 elements.}
\GSfield{Statement}{An ordered triple of subspaces of a finite-dimensional vector space over a field decomposes into one-dimensional blocks, each subspace being zero or the entire line, and copies of the two-dimensional block of three distinct lines. The decomposition is unique up to isomorphism and order of blocks. Separately, the abstract free modular lattice on three generators has 28 elements; this counts lattice terms modulo modular identities, not indecomposable representation types.}
\GSfield{Reconstruction}{The field, ordering, finite dimension, and difference between lattice-element and representation counts are supplied.}
\GSfield{Fate}{\citet[Opening paragraph]{gelfandponomarev1971quadruples} recalls the elementary two- and three-subspace cases. \citet[p. 513]{herrmann1983four} confirms the 28-element free modular lattice. These facts long precede the talk.}
\GSfield{Present status}{As of 12 November 2026 this is an established result under the hypotheses in the modern statement.}
\GSfield{Assessment}{literal statement: correct with missing hypotheses; publication match: exact; verification: read the result; development: none; confidence: high}
\end{GSitem}
\begin{GSitem}{1986-696.3}{E}{Four-subspace classification and its parameter families}
\GSfield{Kind; source}{theorem; \GSpage{P001326}}
\GSfield{The notes}{The four-subspace problem is attributed to long work of Gelfand and Ponomarev, with a reference to Manin.}
\GSfield{Statement}{Classify ordered quadruples of subspaces of a finite-dimensional vector space up to a common linear isomorphism. The category admits direct-sum decomposition into indecomposables and an explicit classification containing discrete families and regular one-parameter data. Over a general field, irreducible polynomial and rational-canonical data replace eigenvalue parameters; the classification is not a finite list of configurations.}
\GSfield{Reconstruction}{Field-dependent regular parameters are retained, rather than assuming an algebraically closed field without saying so.}
\GSfield{Fate}{\citet{gelfandponomarev1971quadruples}, received 21 October 1970, gives the matching canonical classification; it was already reported in earlier seminar record 1970-198. The 1986 passage is retrospective.}
\GSfield{Present status}{As of 12 November 2026 this is an established result under the hypotheses in the modern statement.}
\GSfield{Assessment}{literal statement: correct with missing hypotheses; publication match: exact; verification: read the result; development: none; confidence: high}
\end{GSitem}
\begin{GSitem}{1986-696.4}{E}{The free modular lattice and its linear representations}
\GSfield{Kind; source}{definition; \GSpage{P001326}}
\GSfield{The notes}{Terms in \(e_1,\ldots,e_r\) use sum and intersection, with equality by axioms; a representation preserves both operations.}
\GSfield{Statement}{Let \(D^r\) be the free lattice in the variety of modular lattices on \(r\) generators. A finite-dimensional linear representation over \(k\) is a lattice homomorphism \(\rho:D^r\to L(V)\), with \(L(V)\) the subspace lattice. Thus a representation is exactly a configuration of \(r\) ordered subspaces of \(V\), extended to all lattice terms. Abstract equality in \(D^r\) need not be the same relation as equality in every such representation.}
\GSfield{Reconstruction}{The field and finite-dimensional representation category are specified.}
\GSfield{Fate}{\citet[Sec. 1, pp. 524--525]{gelfandlidskiiponomarev1987preprojective} states these definitions. The project builds on \citet{gelfandponomarev1974modular}.}
\GSfield{Present status}{As of 12 November 2026 this is an established result under the hypotheses in the modern statement.}
\GSfield{Assessment}{literal statement: correct; publication match: exact; verification: read the result; development: none; confidence: high}
\end{GSitem}
\begin{GSitem}{1986-696.5}{E}{Undecidable equality in free modular lattices}
\GSfield{Kind; source}{theorem; \GSpage{P001326}}
\GSfield{The notes}{The notes say the word-identity problem is undecidable for \(r\ge5\).}
\GSfield{Statement}{There is no algorithm deciding whether two lattice terms are equal in the free modular lattice on \(r\) generators when \(r\ge4\). The result concerns equality modulo all modular identities, not equality in all finite-dimensional subspace lattices over a fixed field.}
\GSfield{Reconstruction}{The distinct decision problems are separated; Ryzhikova\textquotesingle s theorem is not a contradiction of the free-modular-lattice result.}
\GSfield{Fate}{\citet[Thm. 1, p. 513]{herrmann1983four} proves the four-generator case and recalls Freese\textquotesingle s earlier \(r\ge5\) theorem. The note is correct but gives a nonoptimal threshold already in 1986. \citet[Theorem]{ryzhikova1987decidability} establishes the contrasting decidability of subspace-lattice identities, submitted 16 May 1986 and published in 1987.}
\GSfield{Present status}{As of 12 November 2026 this is an established result under the hypotheses in the modern statement.}
\GSfield{Assessment}{literal statement: correct; publication match: exact; verification: read the result; development: none; confidence: high}
\end{GSitem}
\begin{GSitem}{1986-696.6}{E}{Finite equational axiomatization of linear equivalence}
\GSfield{Kind; source}{problem; \GSpage{P001326}}
\GSfield{The notes}{Ryzhikova, Semenov\textquotesingle s diploma student, is associated with the question whether linear equivalence can be formulated by finitely many axioms.}
\GSfield{Statement}{In the equational interpretation, ask whether all identities valid in subspace lattices of finite-dimensional rational vector spaces follow from finitely many lattice identities. Equivalently, is the lattice variety generated by these subspace lattices finitely based? The analogous class including all fields contains the rational case and is also generated by finite-dimensional lattices. The notebook does not fix a field or say whether axioms must be equations.}
\GSfield{Reconstruction}{The equational and field quantifiers are explicit editorial specifications. This is not the earlier decidability problem GP-1985-20.}
\GSfield{Fate}{\citet[Thm. 1.4]{herrmann1984arithmetic}, received February 1982 and revised August 1983, proves that a modular variety containing all rational projective geometries cannot both be finitely based and be generated by its finite-dimensional members. This gives a negative answer to the stated standard interpretation. \citet{ryzhikova1987decidability} answers a different question, algorithmic recognition, affirmatively.}
\GSfield{Present status}{As of 12 November 2026 the finite equational basis question under the rational/all-field interpretation has a pre-1986 negative answer. Fixed positive-characteristic variants or unrestricted first-order axioms are not silently identified with this formulation.}
\GSfield{Assessment}{literal statement: correct with missing hypotheses; publication match: partial; verification: read the result; development: none; confidence: medium; problem \hyperref[prob:GP-1986-35]{GP-1986-35}}
\end{GSitem}
\begin{GSitem}{1986-696.7}{E}{Morphisms and indecomposability of lattice representations}
\GSfield{Kind; source}{definition; \GSpage{P001326}}
\GSfield{The notes}{A morphism \(V_1\to V_2\) takes every represented subspace into its counterpart; indecomposability is introduced through direct sums.}
\GSfield{Statement}{For representations \(\rho_i:D^r\to L(V_i)\), define \(\operatorname{Hom}(\rho_1,\rho_2)=\{f\in\operatorname{Hom}_k(V_1,V_2):f(\rho_1(x))\subseteq\rho_2(x)\ \forall x\in D^r\}\). It suffices to impose the inclusions on the generators. A direct-sum splitting must split all represented subspaces simultaneously; indecomposable means no nontrivial such splitting.}
\GSfield{Reconstruction}{The requirement of simultaneous splitting prevents confusion with arbitrary decompositions of the ambient vector space.}
\GSfield{Fate}{\citet[Sec. 1, p. 525]{gelfandlidskiiponomarev1987preprojective} gives exactly this category. It is the standard category underlying the earlier subspace-configuration work.}
\GSfield{Present status}{As of 12 November 2026 this is an established result under the hypotheses in the modern statement.}
\GSfield{Assessment}{literal statement: correct; publication match: exact; verification: read the result; development: none; confidence: high}
\end{GSitem}
\begin{GSitem}{1986-696.8}{E}{Preprojective means finiteness of source types, not finiteness of morphisms}
\GSfield{Kind; source}{definition; \GSpage{P001326-P001327}}
\GSfield{The notes}{The page break gives a definition ending with only finitely many morphisms \(\tau:\rho\to\rho_0\).}
\GSfield{Statement}{A representation \(\rho_0\) is preprojective in this finite-dimensional subspace category when only finitely many isomorphism classes of indecomposable \(\rho\) have \(\operatorname{Hom}(\rho,\rho_0)\ne0\). Over an infinite field every nonzero Hom space itself has infinitely many scalar multiples; counting individual morphisms would therefore destroy the intended definition. Arnold\textquotesingle s phrase without parameters is an intuition, not an equivalent categorical definition.}
\GSfield{Reconstruction}{The fixed target, indecomposable source, nonzero Hom condition, and isomorphism-class quotient are supplied.}
\GSfield{Fate}{\citet[Sec. 1, definition and Prop. 1]{gelfandlidskiiponomarev1987preprojective} supplies the intended representation-type finiteness. The actual scan confirms the notebook says morphisms, so the repair is mathematical rather than paleographic.}
\GSfield{Present status}{As of 12 November 2026 this is an established result under the hypotheses in the modern statement.}
\GSfield{Assessment}{literal statement: false as written; publication match: exact; verification: read the result; development: annotated; confidence: high}
\end{GSitem}
\begin{GSitem}{1986-696.9}{A}{Explicit classification of preprojective representations}
\GSfield{Kind; source}{theorem; \GSpage{P001327}}
\GSfield{The notes}{Theorem 1 says all preprojective representations of a free lattice on an uncertain \(r\) or (4) generators are described.}
\GSfield{Statement}{For the free modular lattice \(D^r\) in the range considered by the published paper, the indecomposable preprojective representations consist of the zero-valued one-dimensional representation \(\rho_{0,0}\) and explicitly indexed representations \(\rho_{t,l}\) constructed from homogeneous pieces of a graded algebra \(A_k^r\). The algebraic construction works over a specified field \(k\); it does not classify all indecomposable \(r\)-subspace systems.}
\GSfield{Reconstruction}{The uncertain handwritten generator count is not forcibly read as four. The modern statement specifies the preprojective subcategory and cites the explicit algebraic index rather than claiming a full wild-type classification.}
\GSfield{Fate}{\citet[Sec. 2, Prop. 6]{gelfandlidskiiponomarev1987preprojective}, received 9 December 1986 and published in 1987, gives the exact classification and points to earlier Gelfand--Ponomarev constructions. This is a direct publication match to the announced description.}
\GSfield{Present status}{As of 12 November 2026 this is an established result under the hypotheses in the modern statement.}
\GSfield{Assessment}{literal statement: correct with missing hypotheses; publication match: partial; verification: read the result; development: none; confidence: high}
\end{GSitem}
\begin{GSitem}{1986-696.10}{A}{The preprojective quotient and its representation category}
\GSfield{Kind; source}{construction; \GSpage{P001327}}
\GSfield{The notes}{Identify lattice elements when every preprojective representation has the same value. Theorem 2 says the preprojective representations of the lattice and its reduction coincide.}
\GSfield{Statement}{Over \(\mathbb Q\), define \(x\sim_{\mathrm{pp}}y\) when \(\rho(x)=\rho(y)\) for every indecomposable preprojective representation, and let \(P_{\mathbb Q}D^r=D^r/\sim_{\mathrm{pp}}\). The equivalence is a lattice congruence. The published classification of all indecomposable representations of this quotient consists of the original preprojective representations together with the dual all-full one-dimensional representation \(\rho_{0,0}^{\prime}\). Thus a statement about all quotient representations needs this extra object.}
\GSfield{Reconstruction}{The rational ground field and the distinction between preprojective representations and all quotient representations are supplied.}
\GSfield{Fate}{\citet[Sec. 1, definition and Prop. 5]{gelfandlidskiiponomarev1987preprojective} gives this precise quotient statement. The notebook\textquotesingle s narrower assertion about preprojectives is compatible with it, but does not itself state the all-representation classification.}
\GSfield{Present status}{As of 12 November 2026 this is an established result under the hypotheses in the modern statement.}
\GSfield{Assessment}{literal statement: correct with missing hypotheses; publication match: exact; verification: read the result; development: none; confidence: high}
\end{GSitem}
\begin{GSitem}{1986-696.11}{A}{Describe the lattice obtained by preprojective reduction}
\GSfield{Kind; source}{problem; \GSpage{P001327}}
\GSfield{The notes}{The problem of finding the reduction is marked will be done and then actually already done.}
\GSfield{Statement}{For \(D^r\) over \(\mathbb Q\), describe exactly which families \(U_\rho\subset V_\rho\), indexed by nontrivial indecomposable preprojectives, arise as \(U_\rho=\rho(x)\) for one lattice element \(x\). In the large-rank range of the published theorem, the criterion is compatibility \(\sum_{\tau<\rho}U_\tau\operatorname{Hom}(\tau,\rho)\subseteq U_\rho\), with equality for all but finitely many \(\rho\). Equivalently the quotient is described by finitely generated positive-degree homogeneous right ideals of the explicit graded algebra.}
\GSfield{Reconstruction}{The family-compatibility criterion is reconstructed from the published theorem; the notebook contains only the announcement. The conservative rank subrange avoids silently guessing an OCR inequality.}
\GSfield{Fate}{\citet[Thms. 1A--1C]{gelfandlidskiiponomarev1987preprojective} publishes the announced description in 1987. The rank condition is essential: its following remark says the assertion is false at \(r=4\) and had not been proved there for \(r=5\). We use the criterion only for \(r\ge7\), a safe subrange of the OCR-rendered condition \(r>6\); the precise endpoint requires a cleaner scan of the paper. The endpoint \(r=6\) is unresolved in our source verification only; no assertion that it is mathematically open is made.}
\GSfield{Present status}{As of 12 November 2026 the published large-rank rational description is established. The 1987 rank-five gap is a historical limitation of that theorem, not evidence that a rank-five problem remains open today.}
\GSfield{Assessment}{literal statement: correct with missing hypotheses; publication match: partial; verification: read the result; development: none; confidence: medium; problem \hyperref[prob:GP-1986-36]{GP-1986-36}}
\end{GSitem}
\begin{GSitem}{1986-696.12}{E}{Perfect elements are tested on indecomposable representations}
\GSfield{Kind; source}{definition; \GSpage{P001327}}
\GSfield{The notes}{A perfect element is defined by \(\rho(a)=0\) or \(V\) for every representation, and compared to the center of a group ring.}
\GSfield{Statement}{The standard definition requires \(\rho(a)\in\{0,V\}\) for every indecomposable finite-dimensional representation, over the fields under consideration. On a decomposable representation different summands may take different alternatives, so the image can be a proper nonzero direct summand. The analogy with central idempotents concerns such block selection; it is not an isomorphism of lattice and ring centers.}
\GSfield{Reconstruction}{The omitted indecomposable restriction is a substantive correction confirmed against the original P001327 scan.}
\GSfield{Fate}{\citet{gelfandponomarev1974modular} introduced these perfect elements. \citet[Sec. 1, Prop. 3 and following discussion]{gelfandlidskiiponomarev1987preprojective} restates the indecomposable qualifier and the construction from suitable upward-closed subsets of preprojectives.}
\GSfield{Present status}{As of 12 November 2026 this is an established result under the hypotheses in the modern statement.}
\GSfield{Assessment}{literal statement: false as written; publication match: exact; verification: read the result; development: annotated; confidence: high}
\end{GSitem}

\GStalk{697}{1986-11-24}{A. V. Zelevinsky}{P001327-P001331}{The unitary dual of general linear groups}
The talk begins on P001328 and ends at the top of P001330, before Kazakov. It reports the unitary-dual classification of Tadić and Vogan through determinant characters, complementary series, and normalized parabolic induction, then discusses square-integrable building blocks over local fields. Gelfand adds a kernel model for invariant Hermitian forms. The main formulas are recoverable, but the absolute-value normalization, modular character, center quotient, and segment exponents need correction or qualification.

\begin{GSitem}{1986-697.1}{E}{Unitary characters of the complex multiplicative group}
\GSfield{Kind; source}{construction; \GSpage{P001328}}
\GSfield{The notes}{\(\delta(z)=(z/|z|)^k|z|^{i\rho}\), \(k\in\mathbb Z\), \(\rho\in\mathbb R\), and \(u(\delta,n)=\delta\circ\det\).}
\GSfield{Statement}{Every continuous unitary character of \(\mathbb C^\times\) has the displayed form, using the ordinary complex modulus. Composing with determinant gives a one-dimensional unitary representation of \(\operatorname{GL}_n(\mathbb C)\). These are the \(u(\delta,n)\) building blocks in the complex-field classification because square-integrable representations modulo center occur here only for \(\operatorname{GL}_1(\mathbb C)\).}
\GSfield{Reconstruction}{Continuity and the ordinary-modulus convention are specified.}
\GSfield{Fate}{\citet[Secs. 1.2 and 9.1]{badulescurenard2010unitary} explicitly gives these blocks and places them in the established Tadić--Vogan classification. The characters themselves are classical.}
\GSfield{Present status}{As of 12 November 2026 this is an established result under the hypotheses in the modern statement.}
\GSfield{Assessment}{literal statement: correct; publication match: exact; verification: read the result; development: none; confidence: high}
\end{GSitem}
\begin{GSitem}{1986-697.2}{E}{Stein complementary-series blocks and the complex absolute value}
\GSfield{Kind; source}{construction; \GSpage{P001328}}
\GSfield{The notes}{The second family is induced from \(|\det|^\alpha u(\delta,n)\otimes|\det|^{-\alpha}u(\delta,n)\), with \(0<\alpha<1/2\).}
\GSfield{Statement}{Let \(\nu(g)=|\det g|_{\mathbb C}=|\det g|^2\), where the rightmost modulus is ordinary. With normalized parabolic induction, \(\pi(u,\alpha)=\nu^\alpha u\times\nu^{-\alpha}u\) is irreducible and unitarizable for \(0<\alpha<1/2\). If ordinary modulus is used directly, the corresponding exponent interval is \(0<a<1\). The completion uses an invariant positive form; the naively twisted inducing characters are not themselves unitary.}
\GSfield{Reconstruction}{The normalized local absolute value and the modular factor in induction are indispensable; the notebook uses the same bars as in its ordinary-modulus character formula without defining the change.}
\GSfield{Fate}{\citet[Sec. 1.2 and Thm. 7.2]{badulescurenard2010unitary} states the precise complex normalization. \citet{vogan1986unitary}, published October 1986, establishes the surrounding Archimedean classification before this talk.}
\GSfield{Present status}{As of 12 November 2026 this is an established result under the hypotheses in the modern statement.}
\GSfield{Assessment}{literal statement: correct with missing hypotheses; publication match: exact; verification: read the result; development: annotated; confidence: high}
\end{GSitem}
\begin{GSitem}{1986-697.3}{E}{Induction, modular normalization, and the relevant adjunction}
\GSfield{Kind; source}{definition; \GSpage{P001328-P001329}; I. M. Gelfand}
\GSfield{The notes}{Gelfand defines induced functions by \(f(hg)=\pi(h)f(g)\), then says restriction and induction are dual functors.}
\GSfield{Statement}{For parabolic \(P=MN\) in a reductive local group, normalized smooth induction uses an equivariance factor \(\delta_P^{1/2}\) in addition to the inducing \(M\)-representation, with a consistent left/right convention. In the non-Archimedean smooth category its Frobenius adjunction involves the normalized Jacquet functor, not merely ordinary restriction to \(M\). For finite groups ordinary induction and restriction satisfy the familiar Frobenius reciprocity. Function spaces and completions matter for locally compact groups.}
\GSfield{Reconstruction}{The categorical and analytic qualifications distinguish an informal induction definition from the precise parabolic functor.}
\GSfield{Fate}{\citet[Sec. 1]{tadic1986classification} and \citet[Secs. 2.4--3]{badulescurenard2010unitary} fix the normalized parabolic conventions used in the classification. The notebook gives only the underlying equivariant-function picture.}
\GSfield{Present status}{As of 12 November 2026 this is an established result under the hypotheses in the modern statement.}
\GSfield{Assessment}{literal statement: correct with missing hypotheses; publication match: partial; verification: read the result; development: none; confidence: high}
\end{GSitem}
\begin{GSitem}{1986-697.4}{E}{Irreducibility, exhaustiveness, and uniqueness of the complex unitary building blocks}
\GSfield{Kind; source}{theorem; \GSpage{P001329}}
\GSfield{The notes}{Theorem (5) says all products of the two block types are irreducible and unitary, every irreducible unitary representation occurs, and equivalence means permutation of blocks.}
\GSfield{Statement}{For \(\operatorname{GL}_n(\mathbb C)\), take a finite multiset of determinant-character blocks \(u(\delta,m)\) and complementary blocks \(\pi(u(\delta,m),\alpha)\), \(0<\alpha<1/2\), whose degrees sum to \(n\). Their normalized parabolic product is irreducible and unitarizable. Every irreducible unitary representation occurs, and the multiset of blocks is unique. This classifies irreducibles; general unitary representations require direct-integral multiplicity data.}
\GSfield{Reconstruction}{Normalized induction, the local absolute value, and irreducibility rather than arbitrary unitary decompositions are specified.}
\GSfield{Fate}{\citet{vogan1986unitary} was published in October 1986. The Tadić approach circulated as MPI/SFB 85--22 in 1985; \citet[Secs. 1.2, 7.2 and 12]{badulescurenard2010unitary} gives the precise theorem and later proof history, including the separate ingredients needed for the Archimedean approach.}
\GSfield{Present status}{As of 12 November 2026 this is an established result under the hypotheses in the modern statement.}
\GSfield{Assessment}{literal statement: correct with missing hypotheses; publication match: exact; verification: read the result; development: none; confidence: high}
\end{GSitem}
\begin{GSitem}{1986-697.5}{E}{Square integrability must be taken modulo the center}
\GSfield{Kind; source}{definition; \GSpage{P001329}}
\GSfield{The notes}{\(D^u\) is the set of square-integrable representations of \(\operatorname{GL}_k(F)\), with square-integrable matrix coefficients.}
\GSfield{Statement}{The relevant building blocks are irreducible unitary representations whose matrix coefficients are square-integrable on \(\operatorname{GL}_k(F)/Z\), with unitary central character. They need not be square-integrable on the full group, whose center is noncompact. For \(F=\mathbb C\), only \(k=1\) occurs; for \(F=\mathbb R\), \(k=1,2\) occur; non-Archimedean local fields have higher-degree discrete-series blocks.}
\GSfield{Reconstruction}{The omitted center quotient is substantive and is restored explicitly.}
\GSfield{Fate}{\citet[Secs. 7 and 9]{badulescurenard2010unitary} uses square integrability modulo center throughout. \citet[Sec. 5]{tadic1986classification} gives the non-Archimedean Langlands ingredients.}
\GSfield{Present status}{As of 12 November 2026 this is an established result under the hypotheses in the modern statement.}
\GSfield{Assessment}{literal statement: false as written; publication match: exact; verification: read the result; development: annotated; confidence: high}
\end{GSitem}
\begin{GSitem}{1986-697.6}{E}{Speh blocks as Langlands quotients and the correct segment}
\GSfield{Kind; source}{construction; \GSpage{P001329}}
\GSfield{The notes}{\(u(\delta,n)\) is the unique irreducible quotient of an induction with first exponent \((n-1)/2\), second \(n/2-1\), and final exponent \(-(n-1)/2\) without a final \(\delta\).}
\GSfield{Statement}{For an irreducible unitary square-integrable-mod-center representation \(\delta\) of \(\operatorname{GL}_k(F)\), with \(F\) a local field, set \(u(\delta,n)\) to be the unique Langlands quotient of \(\nu^{(n-1)/2}\delta\times\nu^{(n-3)/2}\delta\times\cdots\times\nu^{-(n-1)/2}\delta\). Here \(\nu=|\det|_F\) and the exponents decrease by one. The standard module is generally reducible; its quotient is the block used in the unitary classification.}
\GSfield{Reconstruction}{The scan confirms the second exponent \(n/2-1\), so replacing it by \((n-3)/2\) is a faithful-error correction, not a transcription claim; the last \(\delta\) is restored.}
\GSfield{Fate}{\citet[Sec. 7.1, Eqs. (7.1)--(7.3)]{badulescurenard2010unitary} gives the exact formula, with \(\nu_\delta\) replacing \(\nu\) for general division-algebra inner forms. \citet[Sec. 5]{tadic1986classification} supplies the split non-Archimedean version.}
\GSfield{Present status}{As of 12 November 2026 this is an established result under the hypotheses in the modern statement.}
\GSfield{Assessment}{literal statement: false as written; publication match: exact; verification: read the result; development: annotated; confidence: high}
\end{GSitem}
\begin{GSitem}{1986-697.7}{E}{The field-uniform classification is a theorem for local fields}
\GSfield{Kind; source}{theorem; \GSpage{P001329}}
\GSfield{The notes}{The notes say the same theorem remains valid for every field after replacing characters by \(D^u\).}
\GSfield{Statement}{For a nondiscrete locally compact field \(F\), the split general linear unitary dual is parametrized by normalized products of \(u(\delta,n)\) and \(\nu^\alpha u(\delta,n)\times\nu^{-\alpha}u(\delta,n)\), \(0<\alpha<1/2\), with unique block multisets. The topology, Haar measure, and admissible smooth or Harish-Chandra category are part of the theorem. It is not a theorem about arbitrary abstract fields without topology.}
\GSfield{Reconstruction}{The local-field qualification repairs the phrase every field; finite groups of rational points are a separate representation problem.}
\GSfield{Fate}{\citet[Theorem A and Sec. 7]{tadic1986classification} proves the non-Archimedean local-field result in 1986. \citet{vogan1986unitary} treats Archimedean groups; \citet[Thm. 7.2 and its hypotheses]{badulescurenard2010unitary} unifies the block formulation.}
\GSfield{Present status}{As of 12 November 2026 this is an established result under the hypotheses in the modern statement.}
\GSfield{Assessment}{literal statement: false as written; publication match: exact; verification: read the result; development: annotated; confidence: high}
\end{GSitem}
\begin{GSitem}{1986-697.8}{F}{Claims about other groups lack a fixed real form}
\GSfield{Kind; source}{question; \GSpage{P001329}; I. M. Gelfand and A. A. Kirillov}
\GSfield{The notes}{Gelfand asks about other groups. Kirillov lists \(B_2=C_2\), \(G_2\), \(SU(2,2)\), \(SU(2,n)\), and names Knapp and Duflo; \(SO(5)\sim Sp(4)\) is written.}
\GSfield{Statement}{A unitary-dual claim must specify the locally compact group, its real or complex form, covering group, and central character conventions. Root-system symbols \(B_2,C_2,G_2\) do not supply these data. The relation between orthogonal and symplectic rank-two groups is a Lie-algebra or isogeny relation, not an unqualified isomorphism of the displayed global groups. The note contains no parameter list or theorem for the named non-general-linear groups.}
\GSfield{Reconstruction}{The absence of a specified real form is retained as a source limitation instead of silently choosing complex groups.}
\GSfield{Fate}{\citet[References, Knapp--Speh 1983 and Duflo 1980 entries]{vogan1986unitary} identifies relevant pre-existing classification literature, but this is a thematic lead rather than an exact match to every group in the notebook list. No separate open conjecture is inferred.}
\GSfield{Present status}{As of 12 November 2026 the notebook list is too compressed to classify each intended real-form theorem responsibly.}
\GSfield{Assessment}{literal statement: not determinable; publication match: thematic only; verification: publisher metadata checked; development: none; confidence: medium}
\end{GSitem}
\begin{GSitem}{1986-697.9}{E}{Quaternionic general linear groups and reducibility normalization}
\GSfield{Kind; source}{remark; \GSpage{P001329}; A. A. Kirillov}
\GSfield{The notes}{The final attribution says Vogan did the quaternionic ones.}
\GSfield{Statement}{The unitary dual of \(\operatorname{GL}_n(\mathbb H)\) belongs to the Archimedean general-linear classification. Its analogous block construction uses the reduced norm and a reducibility step \(\nu_\delta\) attached to each discrete-series block. One cannot copy the split-field exponents with an undefined quaternionic determinant; in particular the reducibility step changes for one-dimensional \(\mathbb H^\times\) blocks.}
\GSfield{Reconstruction}{The group is inferred from the preceding general-linear context; no claim about all quaternionic classical groups is attributed to the phrase.}
\GSfield{Fate}{\citet{vogan1986unitary} is the established 1986 classification. \citet[Prop. 9.10 and Secs. 10--12]{badulescurenard2010unitary} proves the detailed compatibility with the Tadić block description and explains the exceptional step \(\nu_\delta=\nu^2\) for one-dimensional \(\mathbb H^\times\) blocks.}
\GSfield{Present status}{As of 12 November 2026 this is an established result under the hypotheses in the modern statement.}
\GSfield{Assessment}{literal statement: correct with missing hypotheses; publication match: partial; verification: read the result; development: none; confidence: high}
\end{GSitem}
\begin{GSitem}{1986-697.10}{E}{Riesz-kernel Hermitian forms and the positive range}
\GSfield{Kind; source}{method; \GSpage{P001330}; I. M. Gelfand}
\GSfield{The notes}{Gelfand writes the double integral with kernel \(|x_1-x_2|^{-\rho}\) and says it converges for \(\rho<2\).}
\GSfield{Statement}{On \(\mathbb R^2\simeq\mathbb C\), for a Schwartz function \(f\) and \(0<\rho<2\), the quadratic form \(\iint |x-y|^{-\rho}f(x)\overline{f(y)}\,dx\,dy\) is finite and positive. In the Fourier convention \(\widehat f(\xi)=\int e^{-ix\cdot\xi}f(x)dx\), it equals a positive constant times \(\int |\xi|^{\rho-2}|\widehat f(\xi)|^2d\xi\). Local kernel integrability requires \(\rho<2\); positivity and decay assumptions are separate issues.}
\GSfield{Reconstruction}{The ambient real dimension two, Schwartz test space, and lower positivity bound are reconstructed from the threshold; the note does not specify its test-function space.}
\GSfield{Fate}{This is the classical homogeneous-distribution mechanism treated in \citet[Ch. II, Sec. 3]{gelfandshilov1964generalized}; its representation-theoretic application belongs to intertwining-operator theory such as \citet{knappstein1971intertwining}. \ours\ Positivity in the stated range also follows directly by writing \(|x|^{-\rho}=\Gamma(\rho/2)^{-1}\int_0^\infty t^{\rho/2-1}e^{-t|x|^2}dt\) and using the positive Fourier transform of each Gaussian.}
\GSfield{Present status}{As of 12 November 2026 this is an established result under the hypotheses in the modern statement.}
\GSfield{Assessment}{literal statement: correct with missing hypotheses; publication match: partial; verification: publisher metadata checked; development: annotated; confidence: high}
\end{GSitem}
\begin{GSitem}{1986-697.11}{E}{Analytic continuation alone does not preserve positivity}
\GSfield{Kind; source}{remark; \GSpage{P001330}; I. M. Gelfand}
\GSfield{The notes}{For \(\rho>2\) the notes demand analytic continuation and say the form will no longer be positive definite; a Fourier-side weighted norm is sketched.}
\GSfield{Statement}{The homogeneous kernel admits meromorphic continuation as a tempered distribution, with poles at even positive values of \(\rho\) in real dimension two. Away from poles, its Fourier coefficient is proportional to \(2^{2-\rho}\pi\Gamma(1-\rho/2)/\Gamma(\rho/2)\); its sign depends on the parameter interval and any scalar normalization. Thus continuation is not automatically positive, but the blanket claim of nonpositivity for every \(\rho>2\) is also too strong. Invariant subquotients and moment-vanishing conditions require their own analysis.}
\GSfield{Reconstruction}{The unnormalized continuation, its poles, and scalar sign are distinguished from positivity of a specified representation-theoretic quotient.}
\GSfield{Fate}{\citet[Chs. II--III]{gelfandshilov1964generalized} provides the homogeneous-distribution framework; \citet{knappstein1971intertwining} develops the intertwining and complementary-series context. The notebook gives no specific representation whose positivity range can be fully reconstructed.}
\GSfield{Present status}{As of 12 November 2026 this is an established result under the hypotheses in the modern statement.}
\GSfield{Assessment}{literal statement: false as written; publication match: partial; verification: publisher metadata checked; development: annotated; confidence: high}
\end{GSitem}

\GStalk{698}{1986-11-24}{V. A. Kazakov}{P001329-P001334}{Unitary integrals and the Ising model on random planar maps}
The identifiable talk begins on P001330. It moves from a unitary integral with several diagonal sources and reduced gauge models to the two-matrix representation of an Ising model on fluctuating planar maps. The ordinary Harish-Chandra--Itzykson--Zuber integral and biorthogonal-polynomial reduction are the calculational bridge. P001329 belongs to the preceding discussion; the February lecture starts on P001333.

\begin{GSitem}{1986-698.1}{C}{The unitary integral with two independent diagonal sources}
\GSfield{Kind; source}{question; \GSpage{P001330}}
\GSfield{The notes}{Kazakov asks to calculate \(\int d\Omega\,\exp\sum_{\mu=1}^{2}\operatorname{Tr}(\Omega^\dagger p_\mu\Omega q_\mu)\), with diagonal \(p_\mu,q_\mu\), and says that the two-dimensional case has not been calculated.}
\GSfield{Statement}{For prescribed diagonal complex \(N\times N\) matrices \(P_1,P_2,Q_1,Q_2\), determine \(Z_N=\int_{U(N)}\exp\{\sum_{\mu=1}^{2}\operatorname{Tr}(U^*P_\mu UQ_\mu)\}\,dU\), where Haar measure has mass one. Distinguish an exact convergent series, an efficient finite formula, and a useful large-\(N\) analysis.}
\GSfield{Reconstruction}{Haar normalization, complex parameters and the distinction among meanings of evaluation are editorial specifications. The notebook states no particular target form of answer.}
\GSfield{Fate}{\citet[Eq.~(3), Sections~2.2--2.3 and Appendix]{holguin2024multimatrix} study this same integral, obtaining an explicit \(U(2)\) expression and a \(U(3)\) expansion; their preprint appeared on 6 July 2023 and the article on 4 April 2024. They also identify the obstruction to a naive extension of the one-source localization formula. The chapter supplies an explicitly labelled rederivation and general convergent Haar expansion, not a claim of a new closed solution. Their angular convention in equation~(2.16) has total mass one quarter; comparison with this chapter\textquoteright s probability Haar measure therefore includes an overall normalization, not a correction of the paper.}
\GSfield{Present status}{As of 12 November 2026, exact series evaluation is available and the two-dimensional unitary rank is explicit. A compact formula comparable to the one-source determinant, together with the desired general large-\(N\) control, is not supplied by the cited work; these stronger targets must not be conflated with the existence of an exact evaluation series.}
\GSfield{Assessment}{literal statement: correct with missing hypotheses; publication match: exact; verification: read the result; development: developed in this chapter; confidence: high; problem \hyperref[prob:GP-1986-60]{GP-1986-60}}
\GSfield{Worth developing}{Worth developing in a bounded form: a checked finite-rank formula and a convergent expansion with an error bound make the original computational question concrete without claiming a general closed-form solution.}
\end{GSitem}
\begin{GSitem}{1986-698.2}{E}{Expansion near the identity and planar perturbation theory}
\GSfield{Kind; source}{method; \GSpage{P001330}}
\GSfield{The notes}{The notes put \(\Omega=e^{iA}\), isolate a Gaussian part, expand the rest of the exponential and the Haar measure, and retain planar diagrams.}
\GSfield{Statement}{An expansion about a specified nondegenerate saddle of a unitary matrix integral produces Gaussian Wick contractions and interaction vertices, including the Jacobian of exponential coordinates. At fixed formal order, double-line contractions are sorted by the Euler characteristic of their ribbon graph; planar terms are the leading powers of \(N\) under the usual matrix scaling.}
\GSfield{Reconstruction}{The saddle and normalization conditions are necessary additions. The notebook does not justify a globally valid identity-centred expansion.}
\GSfield{Fate}{This is the established matrix-model diagrammatic method of \citet[Sections~2--3]{bessisitzyksonzuber1980graphical}. It does not itself establish convergence or select the dominant saddle of the multi-source integral.}
\GSfield{Present status}{As of 12 November 2026, the precise statement above is established; the qualifications distinguish it from the abbreviated notebook formulation.}
\GSfield{Assessment}{literal statement: correct with missing hypotheses; publication match: partial; verification: read the result; development: annotated; confidence: high}
\end{GSitem}
\begin{GSitem}{1986-698.3}{E}{One-site reduction of lattice gauge theory}
\GSfield{Kind; source}{claim; \GSpage{P001331}}
\GSfield{The notes}{A product of \(d\) unitary integrals with plaquette traces is asserted to be equivalent to the ordinary gauge theory as \(N\to\infty\).}
\GSfield{Statement}{The Eguchi--Kawai loop-equation reduction requires, in addition to large-\(N\) factorization, the unbroken centre symmetry that annihilates open reduced loops. The unrestricted equivalence of the original one-site model at every coupling is false; in the relevant weak-coupling regime the centre symmetry can break.}
\GSfield{Fate}{\citet{eguchikawai1982reduction} proposed the reduction in 1982. \citet[abstract and model construction]{bhanothellerneuberger1982quenched} identified the symmetry-breaking obstruction and proposed quenching in the same year. This also connects with occurrences 1982-565.3--1982-565.6, but is not identical to the earlier specific free-energy problem.}
\GSfield{Present status}{As of 12 November 2026, the precise statement above is established; the qualifications distinguish it from the abbreviated notebook formulation.}
\GSfield{Assessment}{literal statement: false as written; publication match: partial; verification: publisher metadata checked; development: annotated; confidence: high}
\end{GSitem}
\begin{GSitem}{1986-698.4}{A}{Two-matrix Ising model on random quadrangulations}
\GSfield{Kind; source}{construction; \GSpage{P001331}}
\GSfield{The notes}{The integral has Gaussian coupling \(2cuv\) and separate quartic vertices \(u^4+v^4\); it is said to count two-coloured quadrangular decompositions with a weight for interfaces.}
\GSfield{Statement}{Treated as a formal expansion of a normalized two-Hermitian-matrix Gaussian integral with positive-definite quadratic covariance, the connected free energy sums connected orientable ribbon graphs. Quartic vertices are assigned one of two colours; dualization gives quadrangulations with an Ising spin on each dual face or original vertex, according to convention. Graphs carry automorphism factors, and the mixed propagator fixes the domain-wall weight.}
\GSfield{Reconstruction}{Formal expansion, connectedness, automorphism weights and the distinction between the Gaussian parameter \(c\) and an exponential Ising coupling have been supplied. The displayed \(e^{2cm}\) cannot simply use the same \(c\) as the quadratic matrix coupling.}
\GSfield{Fate}{\citet{kazakov1986ising} gives the exact planar Ising solution on a fluctuating lattice. The journal article appeared on 8 December 1986, shortly after the recorded 24 November talk; publication metadata verifies the identification, while the precise notebook parametrization needs the original conventions. The general graphical framework is already in \citet{bessisitzyksonzuber1980graphical}.}
\GSfield{Present status}{As of 12 November 2026, the precise statement above is established; the qualifications distinguish it from the abbreviated notebook formulation.}
\GSfield{Assessment}{literal statement: correct with missing hypotheses; publication match: partial; verification: publisher metadata checked; development: annotated; confidence: high}
\end{GSitem}
\begin{GSitem}{1986-698.5}{E}{Hermitian diagonalization and the Vandermonde square}
\GSfield{Kind; source}{theorem; \GSpage{P001331}}
\GSfield{The notes}{Diagonalization of each matrix produces \(\Delta(x)^2\Delta(y)^2\) and one relative unitary integration.}
\GSfield{Statement}{For an integrable conjugation-invariant function on Hermitian \(N\times N\) matrices, Lebesgue measure disintegrates into a convention-dependent constant times \(\prod_i dx_i\,\Delta(x)^2\) and normalized orbit measure. For two matrices the relative eigenbasis is Haar distributed; stabilizers and eigenvalue permutations are included in the constant.}
\GSfield{Fate}{\citet[Hermitian eigenvalue reduction]{bessisitzyksonzuber1980graphical} use precisely this reduction in matrix enumeration. The square, rather than a single Vandermonde, is essential before the angular integral.}
\GSfield{Present status}{As of 12 November 2026, the precise statement above is established; the qualifications distinguish it from the abbreviated notebook formulation.}
\GSfield{Assessment}{literal statement: correct with missing hypotheses; publication match: exact; verification: read the result; development: annotated; confidence: high}
\end{GSitem}
\begin{GSitem}{1986-698.6}{E}{The one-source unitary orbital integral}
\GSfield{Kind; source}{theorem; \GSpage{P001331}}
\GSfield{The notes}{The angular integral is written as \(\det(e^{2cx_i y_j})/[\Delta(x)\Delta(y)]\), without a normalizing factor.}
\GSfield{Statement}{For normalized Haar measure and pairwise distinct eigenvalues, \[\int_{U(N)} e^{t\operatorname{Tr}(XUYU^*)}\,dU=\left(\prod_{j=1}^{N-1}j!\right)t^{-N(N-1)/2}\frac{\det(e^{tx_i y_j})}{\Delta(x)\Delta(y)}.\] Coincident eigenvalues and \(t=0\) are interpreted by analytic continuation of the removable singularities.}
\GSfield{Fate}{The matrix-integral form is the established formula in \citet{itzyksonzuber1980planar}. The omitted \(t\)-dependent factor matters if the notebook equality is literal; some eigenvalue-independent constants can instead be absorbed into a partition-function normalization.}
\GSfield{Present status}{As of 12 November 2026, the precise statement above is established; the qualifications distinguish it from the abbreviated notebook formulation.}
\GSfield{Assessment}{literal statement: false as written; publication match: exact; verification: read the result; development: annotated; confidence: high}
\end{GSitem}
\begin{GSitem}{1986-698.7}{E}{Biorthogonal polynomials and the partition-function product}
\GSfield{Kind; source}{construction; \GSpage{P001332}}
\GSfield{The notes}{Polynomials are orthogonal for a two-variable weight. The notebook gives a moment determinant and a product of norms as the matrix integral.}
\GSfield{Statement}{For a bilinear moment functional \(\langle f,g\rangle=\iint f(x)g(y)w(x,y)\,dx\,dy\), assume all leading moment minors are nonzero. Unique monic families \(P_n,Q_n\) satisfy \(\langle P_n,Q_m\rangle=h_n\delta_{nm}\). A determinant formula for a monic polynomial must be divided by its leading moment minor. For symmetric weight one may take \(P_n=Q_n\). Andr\'eief expansion gives the labelled eigenvalue integral \(N!\prod_{n=0}^{N-1}h_n\), before any angular normalization.}
\GSfield{Fate}{The orthogonal and biorthogonal reduction belongs to the established matrix method; see \citet{bessisitzyksonzuber1980graphical} and the two-matrix application in \citet{kazakov1986ising}. The factorial and index shifts depend on whether the eigenvalues are ordered; the notebook does not state that convention.}
\GSfield{Present status}{As of 12 November 2026, the precise statement above is established; the qualifications distinguish it from the abbreviated notebook formulation.}
\GSfield{Assessment}{literal statement: false as written; publication match: partial; verification: read the result; development: annotated; confidence: high}
\end{GSitem}
\begin{GSitem}{1986-698.8}{A}{Parametric planar free energy and its critical points}
\GSfield{Kind; source}{claim; \GSpage{P001332}}
\GSfield{The notes}{The notes give \(g(z)=z/(1-3z)^2-c^2z+3c^2z^3\), an integral expression for \(F\), and locate singularities by \(g^{\prime}(z)=0\).}
\GSfield{Statement}{On the branch of the stated parametric solution analytic at the Gaussian point, critical points of the inverse relation \(z=z(g)\) occur where \(g^{\prime}(z)=0\), unless a singularity or cancellation changes the local analysis. Extracting asymptotics of map coefficients additionally requires identifying the dominant singularity and the order of vanishing. The displayed free-energy integral depends on the partition-function normalization.}
\GSfield{Fate}{The explicit two-matrix Ising solution is identified by \citet{kazakov1986ising}. Its publication metadata is verified, but a line-by-line match of the notebook free-energy normalization is not asserted. The inverse-function criterion explains the stated mechanism; the notebook alone does not justify the claimed global asymptotics.}
\GSfield{Present status}{As of 12 November 2026, the planar Ising matrix model is exactly solved. The present identification is of the model and critical-point mechanism; the precise additive and multiplicative normalization of the notebook expression remains an editorial qualification.}
\GSfield{Assessment}{literal statement: correct with missing hypotheses; publication match: partial; verification: publisher metadata checked; development: annotated; confidence: medium}
\end{GSitem}

\GStalk{699}{}{-}{P001332-P001334}{YEAR 1987}
This is a catalogue heading marking the calendar-year transition, not a second mathematical lecture. The last Kazakov formulas occupy P001332; the next explicit manuscript date is 9 February 1987 on P001333. No independent mathematical occurrence is assigned to this structural record.


\GStalk{700}{1987-02-09}{Don Zagier}{P001332-P001344}{Hecke traces, Eisenstein series, and the Riemann hypothesis}
After a question about polygon gluings, Zagier develops level-one modular forms, Hecke operators and a kernel proof of the Eichler--Selberg trace formula. The second half concerns Eisenstein series, horocycle averages and toroidal automorphic forms. Several normalizing constants in the notebook require correction. The talk ends at the dated heading on P001343; the wide catalogue interval overlaps both neighbours.

\begin{GSitem}{1986-700.1}{F}{Faltings and p-adic string theory}
\GSfield{Kind; source}{remark; \GSpage{P001333}; I. M. Gelfand}
\GSfield{The notes}{Gelfand mentions Faltings in connection with p-adic strings, without a mathematical assertion or bibliographic identification.}
\GSfield{Statement}{The record preserves a thematic pointer to p-adic approaches to string theory. It specifies neither a p-adic curve, a moduli problem, nor an amplitude to calculate.}
\GSfield{Fate}{No exact theorem or problem can be recovered from this sentence. It is not legitimate to identify it with an arbitrary later paper on p-adic strings.}
\GSfield{Present status}{As of 12 November 2026, the mathematical referent of this particular remark remains unidentified; this is a source limitation, not a claim that the subject was undeveloped.}
\GSfield{Assessment}{literal statement: not determinable; publication match: none found; verification: unresolved lead; development: none; confidence: low}
\end{GSitem}
\begin{GSitem}{1986-700.2}{E}{Genus enumeration of oriented polygon gluings}
\GSfield{Kind; source}{question; \GSpage{P001333}}
\GSfield{The notes}{Pair the sides of an oriented \(2n\)-gon with opposite orientation; cyclic rotations remain distinguishable. Count the resulting surfaces by genus.}
\GSfield{Statement}{Let \(\epsilon_g(n)\) count perfect matchings of the labelled sides of a fixed oriented \(2n\)-gon, glued orientation-reversingly, whose quotient is a connected closed oriented surface of genus \(g\). The number is zero for \(2g>n\). Labels fix a boundary root, so no quotient by rotation is taken.}
\GSfield{Fate}{\citet[Theorem~2 and its recurrence]{harerzagier1986euler} had already solved this enumeration in 1986. The later Kazakov lecture 1986-702 gives the matrix-integral derivation. The connection is to orbifold Euler characteristics of moduli spaces; Teichm\"uller space itself is contractible.}
\GSfield{Present status}{As of 12 November 2026, the precise statement above is established; the qualifications distinguish it from the abbreviated notebook formulation.}
\GSfield{Assessment}{literal statement: correct with missing hypotheses; publication match: exact; verification: read the result; development: annotated; confidence: high; problem \hyperref[prob:GP-1986-62]{GP-1986-62}}
\end{GSitem}
\begin{GSitem}{1986-700.3}{E}{Level-one holomorphic modular and cusp forms}
\GSfield{Kind; source}{definition; \GSpage{P001333-P001334}}
\GSfield{The notes}{The notes impose the weight transformation law, Fourier expansion and vanishing at infinity, but write determinant \(\pm1\) and a negative Fourier exponent.}
\GSfield{Statement}{For even \(k\), a level-one modular form is holomorphic on \(\mathbb H\), satisfies \(f((az+b)/(cz+d))=(cz+d)^kf(z)\) for \(\left(\begin{smallmatrix}a&b\\c&d\end{smallmatrix}\right)\in SL_2(\mathbb Z)\), and is holomorphic at the cusp. Thus \(f(z)=\sum_{n\ge0}a_nq^n\), \(q=e^{2\pi iz}\); a cusp form has \(a_0=0\). Matrices of determinant \(-1\) do not act holomorphically on \(\mathbb H\) by this formula.}
\GSfield{Fate}{This is the classical setting fixed explicitly in \citet[Section~1]{zagier1977quadratic} and \citet[Section~1]{kohnenzagier1984periods}. The manuscript signs are faithfully transcribed but mathematically inconsistent with its later definition of \(q\).}
\GSfield{Present status}{As of 12 November 2026, the precise statement above is established; the qualifications distinguish it from the abbreviated notebook formulation.}
\GSfield{Assessment}{literal statement: false as written; publication match: exact; verification: read the result; development: annotated; confidence: high}
\end{GSitem}
\begin{GSitem}{1986-700.4}{E}{Integral bases, small weights, and the discriminant form}
\GSfield{Kind; source}{theorem; \GSpage{P001334}}
\GSfield{The notes}{The notes list the vanishing of several cusp-form spaces, integral Fourier coefficients, the Eisenstein series and \(\Delta=q-24q^2+252q^3-\cdots\).}
\GSfield{Statement}{For level one, \(M_*\) over \(\mathbb C\) is generated by \(E_4,E_6\), and \(S_k=\Delta M_{k-12}\), where \(\Delta=(E_4^3-E_6^2)/1728\). In particular \(S_k=0\) for \(k=2,4,6,8,10,14\) and all odd \(k\). A basis with integral Fourier coefficients can be chosen; this does not mean every complex linear combination has integral coefficients.}
\GSfield{Fate}{\citet[Section~1 and examples]{kohnenzagier1984periods} use the rational Fourier structure and the discriminant form. The displayed coefficients and vanishing statements are standard pre-1987 material.}
\GSfield{Present status}{As of 12 November 2026, the precise statement above is established; the qualifications distinguish it from the abbreviated notebook formulation.}
\GSfield{Assessment}{literal statement: correct with missing hypotheses; publication match: exact; verification: read the result; development: annotated; confidence: high}
\end{GSitem}
\begin{GSitem}{1986-700.5}{E}{Normalized Eisenstein series and Fourier coefficients}
\GSfield{Kind; source}{definition; \GSpage{P001334}}
\GSfield{The notes}{The primitive lattice sum defining \(E_k\) has denominator power 2, whereas the transformation weight and subsequent examples use \(k\).}
\GSfield{Statement}{For even \(k\ge4\), \(E_k(z)=\tfrac12\sum_{(c,d)=1}(cz+d)^{-k}=1-(2k/B_k)\sum_{n\ge1}\sigma_{k-1}(n)q^n\). Absolute convergence justifies rearrangement. In particular \(E_4=1+240(q+9q^2+28q^3+73q^4+\cdots)\). Weight 2 requires quasimodular or nonholomorphic corrections.}
\GSfield{Fate}{The Eisenstein normalizations and their products are treated in \citet[Sections~1--2]{kohnenzagier1984periods}. The denominator exponent in the notebook is an error, not a different normalization.}
\GSfield{Present status}{As of 12 November 2026, the precise statement above is established; the qualifications distinguish it from the abbreviated notebook formulation.}
\GSfield{Assessment}{literal statement: false as written; publication match: exact; verification: read the result; development: annotated; confidence: high}
\end{GSitem}
\begin{GSitem}{1986-700.6}{E}{Hecke eigenforms and coefficient multiplicativity}
\GSfield{Kind; source}{theorem; \GSpage{P001334-P001335}}
\GSfield{The notes}{An Euler form is characterized by multiplicative coefficients, and the Hecke operators have a simultaneous eigenbasis.}
\GSfield{Statement}{On \(S_k(SL_2(\mathbb Z))\), the usual Hecke operators commute and are self-adjoint for the Petersson product. A normalized eigenform \(f=\sum a_nq^n\), \(a_1=1\), satisfies \(T_nf=a_nf\) and \(a_ma_n=\sum_{d\mid(m,n)}d^{k-1}a_{mn/d^2}\). Coprime multiplicativity alone omits the prime-power relation needed for the full Hecke characterization.}
\GSfield{Fate}{\citet[Section~1]{zagier1977quadratic} fixes the normalized eigenbasis and its eigenvalues; \citet[Section~1]{kohnenzagier1984periods} uses the same Hecke structure in the period rationality theorem.}
\GSfield{Present status}{As of 12 November 2026, the precise statement above is established; the qualifications distinguish it from the abbreviated notebook formulation.}
\GSfield{Assessment}{literal statement: correct with missing hypotheses; publication match: exact; verification: read the result; development: annotated; confidence: high}
\end{GSitem}
\begin{GSitem}{1986-700.7}{E}{Normalization of the slash action and Hecke operators}
\GSfield{Kind; source}{construction; \GSpage{P001335}}
\GSfield{The notes}{The \(GL_2^+\) slash action includes \(\det(A)^{k/2}\). The displayed coset definition of \(T_n\) also has a factor \(1/2\), while the examples have \(T_1=1\).}
\GSfield{Statement}{Use \((f|_kA)(z)=\det(A)^{k/2}(cz+d)^{-k}f(Az)\) and \(T_nf=n^{k/2-1}\sum_{A\in SL_2(\mathbb Z)\backslash M_n}f|_kA\), with each coset counted once. Then \(T_1=1\) and \(T_2f=2^{k-1}f(2z)+\tfrac12f(z/2)+\tfrac12f((z+1)/2)\). A separate factor \(1/2\) is appropriate only if both signs are deliberately counted.}
\GSfield{Fate}{The normalization is the one underlying \citet[Section~1]{zagier1977quadratic}. The notebook combines the quotient convention with the double-counting factor; its concrete \(T_2\) example identifies the intended correction.}
\GSfield{Present status}{As of 12 November 2026, the precise statement above is established; the qualifications distinguish it from the abbreviated notebook formulation.}
\GSfield{Assessment}{literal statement: false as written; publication match: exact; verification: read the result; development: annotated; confidence: high}
\end{GSitem}
\begin{GSitem}{1986-700.8}{E}{Weighted Bergman reproducing identity}
\GSfield{Kind; source}{theorem; \GSpage{P001336}}
\GSfield{The notes}{A half-plane integral reproduces a holomorphic function with constant \((-1)^{k/2}\pi/[2^{k-2}(k-1)]\).}
\GSfield{Statement}{For even \(k\ge2\), \(w\in\mathbb H\), and a holomorphic function in the weighted Bergman space for measure \(y^{k-2}dx\,dy\), \[\int_{\mathbb H}\frac{f(z)y^{k-2}}{(\bar z-w)^k}\,dx\,dy=\frac{(-1)^{k/2}\pi}{2^{k-2}(k-1)}f(w).\] A growth or integrability hypothesis is essential; an arbitrary holomorphic function need not have the displayed integral.}
\GSfield{Fate}{This kernel calculation is the analytic ingredient in \citet[Section~2]{zagier1977quadratic}. Its use for automorphic forms follows by unfolding a convergent Poincar\'e kernel, not by assuming every modular form is square-integrable over the entire half-plane.}
\GSfield{Present status}{As of 12 November 2026, the precise statement above is established; the qualifications distinguish it from the abbreviated notebook formulation.}
\GSfield{Assessment}{literal statement: correct with missing hypotheses; publication match: exact; verification: read the result; development: annotated; confidence: high}
\end{GSitem}
\begin{GSitem}{1986-700.9}{E}{The Hecke kernel and its spectral trace}
\GSfield{Kind; source}{method; \GSpage{P001336-P001337}}
\GSfield{The notes}{A determinant-\(n\) matrix sum is expanded in a Hecke eigenbasis; setting the two variables on the appropriate diagonal yields \(\operatorname{Tr}T_n\).}
\GSfield{Statement}{For even \(k>2\), the convergent holomorphic kernel obtained by summing the weighted Bergman kernel over integer matrices of determinant \(n\) represents the normalized Hecke operator on the finite-dimensional Petersson space \(S_k\). Its spectral expansion runs over \(\dim S_k\) eigenforms, with reciprocal Petersson norms and eigenvalues \(a_j(n)\). Integrating the diagonal gives \(\sum_{j=1}^{\dim S_k}a_j(n)\).}
\GSfield{Fate}{\citet[Section~2 and equations~(14)--(19)]{zagier1977quadratic} develops this kernel method. The notebook upper summation bound \(n\) in the eigenform sum is a transcription-faithful mathematical slip: the number of eigenforms does not depend on the Hecke index.}
\GSfield{Present status}{As of 12 November 2026, the precise statement above is established; the qualifications distinguish it from the abbreviated notebook formulation.}
\GSfield{Assessment}{literal statement: false as written; publication match: exact; verification: read the result; development: annotated; confidence: high}
\end{GSitem}
\begin{GSitem}{1986-700.10}{E}{Binary quadratic forms in the trace calculation}
\GSfield{Kind; source}{construction; \GSpage{P001337-P001338}}
\GSfield{The notes}{Matrices of determinant \(n\) are reindexed by \(t=a+d\) and \(Q=[c,d-a,-b]\), with discriminant \(t^2-4n\), and the associated real-analytic kernel is integrated.}
\GSfield{Statement}{The matrix-to-form substitution has \(\operatorname{disc}(Q)=t^2-4n\). The elliptic contribution comes from negative discriminants with stabilizer weights. Positive nonsquare discriminants cancel in the specialized trace calculation, whereas positive squares contribute the divisor term; the scalar discriminant-zero contribution must also be retained. Sign pairs must be counted consistently with the initial factor \(1/2\).}
\GSfield{Fate}{\citet[Section~1, equations~(17)--(23), and Sections~3--4]{zagier1977quadratic} explicitly separates these cases. The notebook sentence about discriminants cannot be read as saying that all nonnegative contributions vanish.}
\GSfield{Present status}{As of 12 November 2026, the precise statement above is established; the qualifications distinguish it from the abbreviated notebook formulation.}
\GSfield{Assessment}{literal statement: false as written; publication match: exact; verification: read the result; development: annotated; confidence: high}
\end{GSitem}
\begin{GSitem}{1986-700.11}{E}{The Eichler--Selberg trace formula at level one}
\GSfield{Kind; source}{theorem; \GSpage{P001338-P001339}}
\GSfield{The notes}{The final trace is a finite sum of weighted class numbers and a divisor correction.}
\GSfield{Statement}{For even \(k\ge4\) and \(n\ge1\), put \(p_k(t,n)=[X^{k-2}](1-tX+nX^2)^{-1}\). Write \(H(D)\) for the Hurwitz weighted class number of negative discriminant \(D\), with \(H(0)=-1/12\). Then \[\operatorname{Tr}(T_n|S_k)=-\frac12\sum_{t^2\le4n}p_k(t,n)H(t^2-4n)-\frac12\sum_{d\mid n}\min(d,n/d)^{k-1}.\] The \(D=0\) convention incorporates the scalar term when \(n\) is a square.}
\GSfield{Fate}{This is the pre-existing Eichler--Selberg formula, recovered in \citet[Section~1, equation~(23)]{zagier1977quadratic}. The notebook gives its kernel derivation, not a new trace formula. Weight 2 is excluded because its correction terms differ.}
\GSfield{Present status}{As of 12 November 2026, the precise statement above is established; the qualifications distinguish it from the abbreviated notebook formulation.}
\GSfield{Assessment}{literal statement: correct with missing hypotheses; publication match: exact; verification: read the result; development: annotated; confidence: high}
\end{GSitem}
\begin{GSitem}{1986-700.12}{E}{Hurwitz weights and the trace polynomials}
\GSfield{Kind; source}{remark; \GSpage{P001339}}
\GSfield{The notes}{A table lists class numbers and \(p_2=1,p_3=t,p_4=t^2-n\), with a tentative identification as Legendre polynomials.}
\GSfield{Statement}{The polynomial recursion is \(p_{k+1}=tp_k-np_{k-1}\), so \[p_k(t,n)=n^{(k-2)/2}U_{k-2}(t/(2\sqrt n))\] for \(n>0\), where \(U_j\) is the Chebyshev polynomial of the second kind. Hurwitz numbers include imprimitive forms: \(H(-3)=1/3\), \(H(-4)=1/2\), \(H(-16)=3/2\), and \(H(-32)=3\).}
\GSfield{Fate}{\citet[equations~(18), (20)]{zagier1977quadratic} specifies both the generating polynomial and the imprimitive class-number weights. The notebook table entry 2 for discriminant \(-32\) counts only the primitive order and is not the required Hurwitz number.}
\GSfield{Present status}{As of 12 November 2026, the precise statement above is established; the qualifications distinguish it from the abbreviated notebook formulation.}
\GSfield{Assessment}{literal statement: false as written; publication match: exact; verification: read the result; development: annotated; confidence: high}
\end{GSitem}
\begin{GSitem}{1986-700.13}{E}{Nonholomorphic Eisenstein series and Rankin unfolding}
\GSfield{Kind; source}{definition; \GSpage{P001340}}
\GSfield{The notes}{The primitive Eisenstein sum is integrated against an invariant function, leaving the Mellin transform of its constant Fourier term.}
\GSfield{Statement}{For \(\operatorname{Re}s>1\), set \(E(z,s)=\tfrac12\sum_{(c,d)=1}y^s|cz+d|^{-2s}\). For a smooth automorphic \(F\) with sufficient decay, \[\int_{SL_2(\mathbb Z)\backslash\mathbb H}F(z)E(z,s)\,\frac{dx\,dy}{y^2}=\int_0^\infty a_0(y)y^{s-2}\,dy,\quad a_0(y)=\int_0^1F(x+iy)\,dx.\] Complex conjugation may be inserted consistently on both sides, not on one side only.}
\GSfield{Fate}{\citet[Section~1]{zagier1981eisenstein} gives this transform, and \citet[Sections~1--2]{zagier1982rankin} treats functions of nonrapid decay by explicit regularization. The latter paper was received in June 1981 and published in volume 28 (1982).}
\GSfield{Present status}{As of 12 November 2026, the precise statement above is established; the qualifications distinguish it from the abbreviated notebook formulation.}
\GSfield{Assessment}{literal statement: correct with missing hypotheses; publication match: exact; verification: read the result; development: annotated; confidence: high}
\end{GSitem}
\begin{GSitem}{1986-700.14}{E}{Completion, continuation, and poles of the Eisenstein series}
\GSfield{Kind; source}{theorem; \GSpage{P001340-P001341}}
\GSfield{The notes}{The notes assert a functional equation using \((2\pi)^s\Gamma(s)E(z,s)\), and list poles at 1 and half of zeta zeros.}
\GSfield{Statement}{The correct completion for the primitive sum is \[E^*(z,s)=\pi^{-s}\Gamma(s)\zeta(2s)E(z,s)\]. It is invariant under \(s\mapsto1-s\), with its standard simple poles at 0 and 1. The uncompleted series has residue \(3/\pi\) at 1; zeros of \(\zeta(2s)\) give possible poles after division, subject to cancellation. Multiplying by \(\Gamma(s)\) alone does not give the stated symmetry.}
\GSfield{Fate}{\citet[Section~1]{zagier1981eisenstein} supplies the exact completion and residue. The missing zeta factor and the power of \(\pi\) are material errors in the recorded formula.}
\GSfield{Present status}{As of 12 November 2026, the precise statement above is established; the qualifications distinguish it from the abbreviated notebook formulation.}
\GSfield{Assessment}{literal statement: false as written; publication match: exact; verification: read the result; development: annotated; confidence: high}
\end{GSitem}
\begin{GSitem}{1986-700.15}{C}{Horocycle averages and the Riemann hypothesis}
\GSfield{Kind; source}{theorem; \GSpage{P001341-P001342}}
\GSfield{The notes}{The constant term is said to tend to the average over the fundamental domain, with oscillations of size \(y^{3/4}\). The normalization \(3/\pi\) disappears in one displayed formula.}
\GSfield{Statement}{For suitable smooth compactly supported automorphic \(F\), \(a_0(y)\to(3/\pi)\int_{\Gamma\backslash\mathbb H}F\,d\mu\) as \(y\downarrow0\). Mellin poles associated with zeta zeros control finer terms. A bound of the form \(O_F(y^{3/4-\varepsilon})\) for all such tests is a Riemann-hypothesis strength assertion; it is not an unconditional rate for every measurable invariant function.}
\GSfield{Fate}{\citet[Section~1 and the horocycle discussion]{zagier1981eisenstein} explains the connection between equidistribution remainders and zeta zeros. The current unresolved status of RH is independently recorded by \citet{clay2026riemann}.}
\GSfield{Present status}{As of 12 November 2026, qualitative equidistribution under the stated test-function hypotheses is established; the universal RH-strength remainder is conditional. The notebook should not be read as an announced proof of RH.}
\GSfield{Assessment}{literal statement: false as written; publication match: exact; verification: read the result; development: annotated; confidence: high; problem \hyperref[prob:GP-1986-61]{GP-1986-61}}
\GSfield{Worth developing}{The statement is useful as an exact conditional implication and a normalization check, but proving its full Riemann-hypothesis strength is not a bounded chapter project.}
\end{GSitem}
\begin{GSitem}{1986-700.16}{E}{Quadratic-field periods and toroidal automorphic forms}
\GSfield{Kind; source}{construction; \GSpage{P001342}}
\GSfield{The notes}{Toroidal integrals of the Eisenstein series contain a quadratic-field zeta function, and therefore vanish at suitable zeros of \(\zeta(s)\).}
\GSfield{Statement}{For a nonsplit rational torus attached to a quadratic field, the appropriately normalized toroidal period of an Eisenstein series factors into a Hecke L-function product. The toroidal subspace consists of automorphic forms whose periods over every such torus vanish. Analytic continuation and the exact normalization are required at zeros and poles; one cannot simply substitute a singular quotient.}
\GSfield{Fate}{\citet[Sections~2--4]{zagier1981eisenstein} develops the proposal already before the talk. \citet[main theorem]{cornelissenlorscheid2011toroidal}, first circulated in June 2009 and revised in August 2011, identifies the Eisenstein-derivative part through zero multiplicities and the cuspidal part through central L-values.}
\GSfield{Present status}{As of 12 November 2026, the later structural theorem clarifies the toroidal space, including derivatives at multiple zeros. It does not prove that all zeta zeros lie on the critical line.}
\GSfield{Assessment}{literal statement: correct with missing hypotheses; publication match: exact; verification: read the result; development: annotated; confidence: high}
\end{GSitem}
\begin{GSitem}{1986-700.17}{C}{The toroidal spectral strategy for the Riemann hypothesis}
\GSfield{Kind; source}{question; \GSpage{P001342-P001343}}
\GSfield{The notes}{The notes form a span \(W\) of quadratic-form transforms and an orthogonal complement \(V\), note Laplacian invariance, and suggest that unitarity of the corresponding representation would imply RH.}
\GSfield{Statement}{The recognizable strategy seeks a spectral realization of the toroidal automorphic space that forces the relevant Eisenstein parameters onto the critical line. The necessary conclusion is temperedness of the relevant constituents; unitarity alone allows complementary-series representations and is insufficient. One must specify the topology, invariant pairing, spectral identification and exclusion of non-tempered constituents. Formal invariance under the Laplacian proves none of these requirements.}
\GSfield{Fate}{\citet[final sections]{zagier1981eisenstein} gives the earlier strategy and its difficulties. \citet[Introduction and main results]{cornelissenlorscheid2011toroidal} identifies the toroidal spectrum and explicitly states the RH connection in terms of temperedness. It does not prove that needed restriction. \citet{clay2026riemann} confirms that RH remains open.}
\GSfield{Present status}{As of 12 November 2026, the algebraic toroidal spectrum is substantially understood, but the required critical-line spectral restriction is unproved. The notebook\textquoteright s unqualified appeal to unitarity does not establish RH.}
\GSfield{Assessment}{literal statement: false as written; publication match: partial; verification: read the result; development: annotated; confidence: medium; problem \hyperref[prob:GP-1986-61]{GP-1986-61}}
\GSfield{Worth developing}{Worth preserving as a precisely conditional spectral strategy; the notebook leaves the functional-analytic object too unspecified for a new proof attempt here.}
\end{GSitem}

\GStalk{701}{1987-02-16}{Don Zagier}{P001342-P001350}{Critical values and rational periods of modular forms}
The lecture connects even zeta-values, Eisenstein series and Mellin transforms of cusp forms, then explains the two rational period structures of Hecke eigenforms using Rankin unfolding. It explicitly points to Kohnen and Zagier\textquoteright s already published work on rational periods. The manuscript heading is on P001343 and the next lecture begins on P001349.

\begin{GSitem}{1986-701.1}{E}{Two-colour map enumeration with a prescribed interface length}
\GSfield{Kind; source}{question; \GSpage{P001343}}
\GSfield{The notes}{At the opening, maps with fixed numbers of black and white triangles and a prescribed boundary between colours are proposed.}
\GSfield{Statement}{Count rooted orientable planar maps of the specified face degrees with two face colours, tracking separately the numbers of both colours and the number of edges separating colours. Equivalently, use the corresponding formal two-matrix generating function, with separate vertex variables and mixed-propagator weight; unrooted counts require automorphism weights.}
\GSfield{Fate}{The two-matrix Ising solution of \citet{kazakov1986ising} treats this enumeration framework. The notebook does not fix rooting, degenerate faces, or a requested coefficient formula, so identification with the exact quartic parameterization is partial. The same variant returns at the end of record 702.}
\GSfield{Present status}{As of 12 November 2026, the precise statement above is established; the qualifications distinguish it from the abbreviated notebook formulation.}
\GSfield{Assessment}{literal statement: correct with missing hypotheses; publication match: partial; verification: publisher metadata checked; development: annotated; confidence: medium; problem \hyperref[prob:GP-1986-63]{GP-1986-63}}
\end{GSitem}
\begin{GSitem}{1986-701.2}{E}{Even and negative zeta-values and moduli Euler characteristics}
\GSfield{Kind; source}{theorem; \GSpage{P001343}}
\GSfield{The notes}{Bernoulli formulas for \(\zeta(2n)\) and \(\zeta(1-2g)\) are related to the Euler characteristic of the moduli space; the list gives \(\zeta(-3)=1/252\).}
\GSfield{Statement}{For \(n\ge1\), \(\zeta(2n)=(-1)^{n+1}B_{2n}(2\pi)^{2n}/[2(2n)!]\) and \(\zeta(1-2n)=-B_{2n}/(2n)\). For \(g\ge2\), \(\chi_{\mathrm{orb}}(\mathcal M_{g,1})=\zeta(1-2g)\) and \(\chi_{\mathrm{orb}}(\mathcal M_g)=\zeta(1-2g)/(2-2g)\). Thus \(\zeta(-3)=1/120\), while \(\zeta(-5)=-1/252\).}
\GSfield{Fate}{The moduli-space formula is the main theorem of \citet{harerzagier1986euler}, published before this lecture. Marking one point and using the orbifold Euler characteristic are indispensable; replacing these by the ordinary Euler characteristic of unpointed Teichm\"uller space is false.}
\GSfield{Present status}{As of 12 November 2026, the precise statement above is established; the qualifications distinguish it from the abbreviated notebook formulation.}
\GSfield{Assessment}{literal statement: false as written; publication match: exact; verification: read the result; development: annotated; confidence: high}
\end{GSitem}
\begin{GSitem}{1986-701.3}{E}{Eisenstein products and recurrences for even zeta-values}
\GSfield{Kind; source}{method; \GSpage{P001343-P001344}}
\GSfield{The notes}{Identities such as \(E_4^2=E_8\) and \(E_4E_6=E_{10}\) yield recurrences for even zeta-values.}
\GSfield{Statement}{Level-one dimensions give \(E_4^2=E_8\) and \(E_4E_6=E_{10}\) for normalized Eisenstein series. Comparing their Fourier coefficients and the Bernoulli normalization produces convolution identities for divisor sums and even zeta-values. At weights with cusp forms, the Eisenstein product generally has a nonzero cuspidal component, which must be retained.}
\GSfield{Fate}{\citet[initial sections]{kohnenzagier1984periods} supplies the Eisenstein-product framework used later in the lecture. The assertion is an established technique, with the small-weight identities representing special dimension-one cases.}
\GSfield{Present status}{As of 12 November 2026, the precise statement above is established; the qualifications distinguish it from the abbreviated notebook formulation.}
\GSfield{Assessment}{literal statement: correct with missing hypotheses; publication match: exact; verification: read the result; development: annotated; confidence: high}
\end{GSitem}
\begin{GSitem}{1986-701.4}{E}{Mellin transform and functional equation of a cusp-form L-function}
\GSfield{Kind; source}{theorem; \GSpage{P001344-P001345}}
\GSfield{The notes}{The completed Dirichlet series is expressed as an integral of \(f(iy)\), and inversion gives the reflection formula.}
\GSfield{Statement}{If \(f(z)=\sum_{n\ge1}a_nq^n\in S_k(SL_2(\mathbb Z))\), then initially in a right half-plane \(\Lambda(f,s)=(2\pi)^{-s}\Gamma(s)\sum a_nn^{-s}=\int_0^\infty f(iy)y^{s-1}\,dy\). Cuspidal decay at both ends makes the integral entire, and \(\Lambda(f,s)=(-1)^{k/2}\Lambda(f,k-s)\). This applies to each cusp form, without an eigenform hypothesis; Euler products need that additional hypothesis.}
\GSfield{Fate}{\citet[Section~1]{kohnenzagier1984periods} uses precisely these completed values and their functional equation in defining period polynomials.}
\GSfield{Present status}{As of 12 November 2026, the precise statement above is established; the qualifications distinguish it from the abbreviated notebook formulation.}
\GSfield{Assessment}{literal statement: correct with missing hypotheses; publication match: exact; verification: read the result; development: annotated; confidence: high}
\end{GSitem}
\begin{GSitem}{1986-701.5}{E}{The two periods of a normalized Hecke eigenform}
\GSfield{Kind; source}{theorem; \GSpage{P001345}}
\GSfield{The notes}{The critical values split into even and odd groups, each rationally proportional to a period; the two periods are said probably to be transcendental.}
\GSfield{Statement}{For a normalized level-one Hecke eigenform \(f\) with coefficient field \(K_f\), there are nonzero periods \(\Omega_f^\pm\), defined up to \(K_f^\times\), such that the appropriately parity-normalized critical values \(L(f,m)/( (2\pi i)^m\Omega_f^{(-1)^m})\), \(1\le m\le k-1\), lie in \(K_f\). For rational Fourier coefficients, the normalized period ratios are rational. Absolute transcendence is a separate assertion and depends on specifying the normalization.}
\GSfield{Fate}{The period theorem predates the lecture; \citet[Theorem~1 and period-polynomial discussion]{kohnenzagier1984periods} gives the rational structures used here. The notebook\textquoteright s transcendence aside is not a proved part of that theorem.}
\GSfield{Present status}{As of 12 November 2026, the precise statement above is established; the qualifications distinguish it from the abbreviated notebook formulation.}
\GSfield{Assessment}{literal statement: correct with missing hypotheses; publication match: exact; verification: read the result; development: annotated; confidence: high}
\end{GSitem}
\begin{GSitem}{1986-701.6}{E}{Critical periods of the discriminant form}
\GSfield{Kind; source}{remark; \GSpage{P001345}}
\GSfield{The notes}{A table expresses the integrals \(r_n=\int_0^\infty\Delta(iy)y^n\,dy\) in two periods, including a denominator 691.}
\GSfield{Statement}{For the real-coefficient weight-12 discriminant form, \(r_n=(2\pi)^{-n-1}\Gamma(n+1)L(\Delta,n+1)\), \(0\le n\le10\), and \(r_n=r_{10-n}\). The even-indexed and odd-indexed period vectors each span a rational line. Numerical values of the generators have no invariant meaning until the rational scaling convention is fixed.}
\GSfield{Fate}{\citet[examples and tables for weight 12]{kohnenzagier1984periods} gives the exact rational period relations. The occurrence of 691 belongs to these exact arithmetic ratios, not to a floating-point inference.}
\GSfield{Present status}{As of 12 November 2026, the precise statement above is established; the qualifications distinguish it from the abbreviated notebook formulation.}
\GSfield{Assessment}{literal statement: correct with missing hypotheses; publication match: partial; verification: read the result; development: annotated; confidence: high}
\end{GSitem}
\begin{GSitem}{1986-701.7}{E}{Rankin unfolding of products of Eisenstein series}
\GSfield{Kind; source}{method; \GSpage{P001345-P001347}}
\GSfield{The notes}{The Petersson product of \(f\) with \(G_{2n}G_{k-2n}\) is unfolded to a Dirichlet series involving \(a_m\sigma_{k-1-2n}(m)\). The written range makes one weight negative.}
\GSfield{Statement}{For even \(k\) and integers \(n\) with \(2n\ge4\), \(k-2n\ge4\), pair the holomorphic Eisenstein product with a weight-\(k\) cusp form. Absolute convergence in an auxiliary parameter, followed where necessary by analytic continuation, justifies unfolding. Fourier integration yields a Gamma factor \(\Gamma(k-1)=(k-2)!\), with conjugations consistent with the Petersson convention.}
\GSfield{Fate}{\citet[Rankin method and products of Eisenstein series]{kohnenzagier1984periods} gives this argument. The notebook\textquoteright s inequality \(k/2<n\le k-2\) is incompatible with its chosen weights; an admissible basic range is \(2\le n\le k/2-2\).}
\GSfield{Present status}{As of 12 November 2026, the precise statement above is established; the qualifications distinguish it from the abbreviated notebook formulation.}
\GSfield{Assessment}{literal statement: false as written; publication match: exact; verification: read the result; development: annotated; confidence: high}
\end{GSitem}
\begin{GSitem}{1986-701.8}{E}{Hecke factorization of the Rankin Dirichlet series}
\GSfield{Kind; source}{theorem; \GSpage{P001347}}
\GSfield{The notes}{The divisor-sum Dirichlet series factors into two L-values divided by an even zeta-value.}
\GSfield{Statement}{For a normalized level-one eigenform and a complex parameter \(a\), the identity \[\sum_{m\ge1}\frac{a_m\sigma_a(m)}{m^s}=\frac{L(f,s)L(f,s-a)}{\zeta(2s-a-k+1)}\] holds where the series and Euler products converge absolutely and elsewhere by meromorphic continuation. Setting \(s=k-1\), \(a=k-1-2n\) gives the notebook\textquoteright s quotient \(L(f,k-1)L(f,2n)/\zeta(2n)\), interpreted by continuation if needed.}
\GSfield{Fate}{\citet[Rankin product calculation]{kohnenzagier1984periods} uses the Hecke relation to obtain this factorization and derive period comparisons. The identity does not say that \(L(f,2n)/\zeta(2n)\) is independent of \(n\).}
\GSfield{Present status}{As of 12 November 2026, the precise statement above is established; the qualifications distinguish it from the abbreviated notebook formulation.}
\GSfield{Assessment}{literal statement: correct with missing hypotheses; publication match: exact; verification: read the result; development: annotated; confidence: high}
\end{GSitem}
\begin{GSitem}{1986-701.9}{E}{Three rational structures on the cusp-form space}
\GSfield{Kind; source}{construction; \GSpage{P001347-P001348}}
\GSfield{The notes}{The lecture distinguishes rational Fourier coefficients and rational even and odd periods, and cites Kohnen--Zagier under the title ``Modular forms with rational coefficients.}
\GSfield{Statement}{The Fourier expansion defines a rational structure on level-one cusp forms. Period polynomials, decomposed by parity, give two further rational structures after choosing the appropriate period scalings; Hecke equivariance relates them. The critical-period length for a weight-\(w\) form is \(w-2\), so writing \(S_{2k}\) while keeping the range \(0,\ldots,k-2\) mixes conventions.}
\GSfield{Fate}{The explicitly cited work is \citet{kohnenzagier1984periods}, whose actual title is \emph{Modular forms with rational periods}, published in 1984, not a new 1987 result. It contains the rational-structure comparison rather than the claimed equality of all normalized critical values.}
\GSfield{Present status}{As of 12 November 2026, the precise statement above is established; the qualifications distinguish it from the abbreviated notebook formulation.}
\GSfield{Assessment}{literal statement: false as written; publication match: exact; verification: read the result; development: annotated; confidence: high}
\end{GSitem}

\GStalk{702}{1987-03-02}{V. A. Kazakov}{P001348-P001356}{Gaussian matrix integrals and polygon gluings}
The talk gives the matrix-integral derivation of the Harer--Zagier enumeration of orientable polygon gluings, including Wick pairings, Hermitian eigenvalues, Hermite polynomials and generating functions. It closes with variants having triangular faces and colours. The identifiable text occupies P001349--P001354.

\begin{GSitem}{1986-702.1}{E}{The rooted orientable gluing numbers}
\GSfield{Kind; source}{definition; \GSpage{P001349}}
\GSfield{The notes}{The lecture returns to pairings of \(2n\) polygon sides and counts them by the genus of the glued surface.}
\GSfield{Statement}{The same numbers \(\epsilon_g(n)\) of 1986-700.2 are used: all sides retain their cyclic labels and each pair is glued with opposite induced boundary orientation. Euler\textquoteright s relation is \(V-n+1=2-2g\), hence \(V=n+1-2g\). Pictures suggesting prohibitions on special side pairings must be separated from the unrestricted Harer--Zagier count.}
\GSfield{Fate}{\citet[Theorem~2]{harerzagier1986euler} gives the already established formula. The matrix computation below uses the unrestricted pairing ensemble, including adjacent pairings; imposing the pictured restrictions would define a different enumeration.}
\GSfield{Present status}{As of 12 November 2026, the precise statement above is established; the qualifications distinguish it from the abbreviated notebook formulation.}
\GSfield{Assessment}{literal statement: correct with missing hypotheses; publication match: exact; verification: read the result; development: annotated; confidence: high; problem \hyperref[prob:GP-1986-62]{GP-1986-62}}
\end{GSitem}
\begin{GSitem}{1986-702.2}{E}{Gaussian integration by parts and Wick contractions}
\GSfield{Kind; source}{method; \GSpage{P001350-P001351}}
\GSfield{The notes}{A Gaussian Hermitian matrix integral and an integration-by-parts identity produce pairings of the factors in a trace.}
\GSfield{Statement}{For a centred Hermitian Gaussian matrix with covariance \(\mathbb E M_{ij}M_{kl}=N^{-1}\delta_{il}\delta_{jk}\), Wick\textquoteright s rule pairs entries and each contraction contributes that covariance. Equivalently, Gaussian integration by parts differentiates a polynomial observable against its covariance. The measure is normalized to expectation one; treating all complex off-diagonal entries as independent real entries would give wrong constants.}
\GSfield{Fate}{\citet[matrix-integral proof of Theorem~2]{harerzagier1986euler} and \citet[Gaussian diagrammatics]{bessisitzyksonzuber1980graphical} give this exact combinatorial method, established before the talk.}
\GSfield{Present status}{As of 12 November 2026, the precise statement above is established; the qualifications distinguish it from the abbreviated notebook formulation.}
\GSfield{Assessment}{literal statement: correct with missing hypotheses; publication match: exact; verification: read the result; development: annotated; confidence: high}
\end{GSitem}
\begin{GSitem}{1986-702.3}{E}{Topological expansion of a Gaussian trace moment}
\GSfield{Kind; source}{theorem; \GSpage{P001351}}
\GSfield{The notes}{Index cycles are identified with vertices of the glued polygon and the power of \(N\) records genus.}
\GSfield{Statement}{With the covariance just specified, \(\mathbb E[N^{-1}\operatorname{Tr}M^{2n}]=\sum_{g=0}^{\lfloor n/2\rfloor}\epsilon_g(n)N^{-2g}\). For a pairing, the number of free index loops is \(V=n+1-2g\), while the \(n\) propagators and normalized trace contribute \(N^{-n-1}\). This is a finite identity for every positive integer \(N\), not merely an asymptotic expansion.}
\GSfield{Fate}{This is the exact matrix interpretation in \citet[proof of Theorem~2]{harerzagier1986euler}. It connects to the earlier planar-enumeration occurrence 1982-565.1.}
\GSfield{Present status}{As of 12 November 2026, the precise statement above is established; the qualifications distinguish it from the abbreviated notebook formulation.}
\GSfield{Assessment}{literal statement: correct with missing hypotheses; publication match: exact; verification: read the result; development: annotated; confidence: high}
\end{GSitem}
\begin{GSitem}{1986-702.4}{E}{Eigenvalue density and Hermite polynomials}
\GSfield{Kind; source}{method; \GSpage{P001351-P001352}}
\GSfield{The notes}{Diagonalization, Vandermonde determinants and orthogonal polynomials reduce the matrix moment to one-dimensional Gaussian integrals.}
\GSfield{Statement}{For a Gaussian Hermitian ensemble, write the Vandermonde as a determinant of monic Hermite polynomials orthogonal for the chosen Gaussian weight. Integrating all but one eigenvalue gives the one-point density as the Gaussian weight times \(\sum_{j=0}^{N-1}P_j(x)^2/h_j\). Its \(2n\)-th moment equals \(\mathbb E\operatorname{Tr}M^{2n}\). The choice \(e^{-x^2}\) versus \(e^{-x^2/2}\) rescales both the polynomials and the moments.}
\GSfield{Fate}{\citet[Gaussian integral computation]{harerzagier1986euler} gives this reduction. The notebook shifts Gaussian conventions and indices; these cannot be combined without rescaling.}
\GSfield{Present status}{As of 12 November 2026, the precise statement above is established; the qualifications distinguish it from the abbreviated notebook formulation.}
\GSfield{Assessment}{literal statement: correct with missing hypotheses; publication match: exact; verification: read the result; development: annotated; confidence: high}
\end{GSitem}
\begin{GSitem}{1986-702.5}{E}{Stochastic quantization as an alternative matrix method}
\GSfield{Kind; source}{remark; \GSpage{P001352}}
\GSfield{The notes}{A remark attributes a Markov-process approach to Parisi.}
\GSfield{Statement}{The identifiable method realizes a Euclidean field or matrix distribution as an invariant law of a Langevin process, with noise and drift satisfying the appropriate fluctuation-dissipation relation. For a matrix model, existence of the invariant probability and convergence need hypotheses on the potential. This remark does not specify a new stochastic dynamics.}
\GSfield{Fate}{The seminar had already recorded stochastic quantization in 1982-565.2. Its role here is an alternative computational viewpoint, while the actual enumeration proceeds by the Gaussian integral; see the established graphical framework in \citet{bessisitzyksonzuber1980graphical}.}
\GSfield{Present status}{As of 12 November 2026, the precise statement above is established; the qualifications distinguish it from the abbreviated notebook formulation.}
\GSfield{Assessment}{literal statement: correct with missing hypotheses; publication match: thematic only; verification: read the result; development: annotated; confidence: medium}
\end{GSitem}
\begin{GSitem}{1986-702.6}{E}{The Harer--Zagier moment generating function}
\GSfield{Kind; source}{theorem; \GSpage{P001353}}
\GSfield{The notes}{Two generating formulas are written with double-factorial and normalization shifts that conflict with the zeroth moment.}
\GSfield{Statement}{Let \(C_n(N)=\sum_g\epsilon_g(n)N^{n+1-2g}\), with \(C_0(N)=N\). Then, as a formal power series, \[1+2\sum_{n\ge0}\frac{C_n(N)}{(2n-1)!!}x^{n+1}=\left(\frac{1+x}{1-x}\right)^N,\quad (-1)!!=1.\] Equivalently \(\sum_{N\ge0}C_n(N)t^N=(2n-1)!!\,t(1+t)^n/(1-t)^{n+2}\). These are unnormalized polynomial moments; dividing by powers of \(N\) gives the normalized trace moments.}
\GSfield{Fate}{\citet[Theorem~2 and generating-function proof]{harerzagier1986euler} gives the exact generating identity. The constant and first coefficients are decisive normalization checks; the recorded formula cannot use \((2n)!!\) with the same moments.}
\GSfield{Present status}{As of 12 November 2026, the precise statement above is established; the qualifications distinguish it from the abbreviated notebook formulation.}
\GSfield{Assessment}{literal statement: false as written; publication match: exact; verification: read the result; development: annotated; confidence: high}
\end{GSitem}
\begin{GSitem}{1986-702.7}{E}{Closed genus formula and the Harer--Zagier recursion}
\GSfield{Kind; source}{theorem; \GSpage{P001353-P001354}}
\GSfield{The notes}{The lecture extracts a coefficient of a power of a hyperbolic-tangent expression to obtain the genus count.}
\GSfield{Statement}{For \(0\le2g\le n\), \[\epsilon_g(n)=\frac{(2n)!}{(n+1)!(n-2g)!}[u^{2g}]\left(\frac{u/2}{\tanh(u/2)}\right)^{n+1}.\] Equivalently, \((n+1)\epsilon_g(n)=(4n-2)\epsilon_g(n-1)+(2n-1)(n-1)(2n-3)\epsilon_{g-1}(n-2)\), with the usual zero conventions and \(\epsilon_0(0)=1\). In particular genus zero gives Catalan numbers.}
\GSfield{Fate}{These are \citet[Theorem~2 and displayed recurrence]{harerzagier1986euler}. The theorem was published in 1986, so the March 1987 presentation is expository. The square-root substitution in the notebook needs a factor-of-two correction to agree with the \(u/2\) convention.}
\GSfield{Present status}{As of 12 November 2026, the precise statement above is established; the qualifications distinguish it from the abbreviated notebook formulation.}
\GSfield{Assessment}{literal statement: false as written; publication match: exact; verification: read the result; development: annotated; confidence: high; problem \hyperref[prob:GP-1986-62]{GP-1986-62}}
\end{GSitem}
\begin{GSitem}{1986-702.8}{E}{Maps with triangular faces and a marked polygonal boundary}
\GSfield{Kind; source}{question; \GSpage{P001354}}
\GSfield{The notes}{Variants add triangular faces, fix a \(2n\)-sided boundary and genus, or introduce two colours and an interface length.}
\GSfield{Statement}{A formal matrix model with a cubic interaction and an insertion \(\operatorname{Tr}M^{2n}\) generates orientable maps with triangular interaction vertices and one marked boundary, sorted by genus. More colours use several matrices and mixed covariances. Exact coefficient extraction depends on rooting and whether the specified polygons are faces of the primal or dual map.}
\GSfield{Fate}{The graphical construction is established in \citet{bessisitzyksonzuber1980graphical}; the Ising specialization is \citet{kazakov1986ising}. These references verify a method for the variants, not one unique numerical formula, because the notebook leaves the counting conventions unstated.}
\GSfield{Present status}{As of 12 November 2026, the precise statement above is established; the qualifications distinguish it from the abbreviated notebook formulation.}
\GSfield{Assessment}{literal statement: correct with missing hypotheses; publication match: partial; verification: read the result; development: annotated; confidence: medium; problem \hyperref[prob:GP-1986-63]{GP-1986-63}}
\end{GSitem}

\GStalk{704}{1987-03-09}{I. M. Gelfand}{P001354-P001358}{Fractal measures, factors, and proposed connections}
The short Gelfand preamble on P001355 proposes connections among fractional-dimensional measures, representations, factors and chaotic dynamics. It includes a recognizable reference to the uniqueness of the hyperfinite type III\(_1\) factor, recorded with a problematic type label. The following explicit Beilinson heading controls attribution of the zeta-function material.

\begin{GSitem}{1986-704.1}{F}{Fractional-dimensional measures and representation parameters}
\GSfield{Kind; source}{remark; \GSpage{P001355}}
\GSfield{The notes}{Gelfand mentions a work on integral fractional-linear transformations of the upper half-plane and sets carrying measures of intermediate Hausdorff dimension, in connection with a series of representations.}
\GSfield{Statement}{The notes identify neither the work nor the acting subgroup, invariant or quasi-invariant measure, and representation functor. A relation between dimension and a principal or complementary-series parameter is a plausible theme, but there is no unique theorem or conjecture to reconstruct from the sentence.}
\GSfield{Fate}{Searches for the 1986--1987 combinations of Gelfand, Hausdorff measures, fractional-linear actions and representation series did not recover an exact source. The nearby multifractal literature is only thematic and is deliberately not substituted for the missing reference.}
\GSfield{Present status}{As of 12 November 2026, the exact referent remains unidentified. No claim of mathematical nondevelopment follows from this negative identification result.}
\GSfield{Assessment}{literal statement: not determinable; publication match: none found; verification: unresolved lead; development: none; confidence: low}
\end{GSitem}
\begin{GSitem}{1986-704.2}{F}{Factors arising from representations of infinite-dimensional groups}
\GSfield{Kind; source}{conjecture; \GSpage{P001355}}
\GSfield{The notes}{The word ``hypothesis introduces the claim that factors arise from representations of infinite-dimensional groups.}
\GSfield{Statement}{To turn this into a mathematical conjecture one must specify the class of groups, representations, generated von Neumann algebras, factor types and any universality assertion. Without those restrictions the statement is either a broad observation or a tautological realization using unitary groups; the notebook gives no discriminating condition.}
\GSfield{Fate}{No precise publication match is asserted. The statement is preserved as a research-program remark adjacent to factor classification, not converted into an invented universal classification problem.}
\GSfield{Present status}{As of 12 November 2026, this formulation is insufficient to identify a specific open question or a theorem with a determinate fate.}
\GSfield{Assessment}{literal statement: not determinable; publication match: none found; verification: unresolved lead; development: none; confidence: low}
\end{GSitem}
\begin{GSitem}{1986-704.3}{E}{Uniqueness of the hyperfinite type III1 factor}
\GSfield{Kind; source}{theorem; \GSpage{P001355}}
\GSfield{The notes}{The notes say physics produces factors of type \(III_\infty\), of which exactly one is hyperfinite.}
\GSfield{Statement}{The recognizable theorem is uniqueness, up to isomorphism, of the injective (equivalently, in the separable setting, hyperfinite) type \(III_1\) factor with separable predual. There is no standard Connes type \(III_\infty\); the subscript \(\infty\) also occurs in the conventional name \(R_\infty\) for the unique factor, which may explain the notation.}
\GSfield{Reconstruction}{Type \(III_1\), separability and injectivity are the editorial reconstruction of the recognizable uniqueness statement. The manuscript\textquoteright s type label is retained in notes rather than silently corrected.}
\GSfield{Fate}{\citet[Section~2]{haagerup1987bicentralizer} completes the uniqueness proof; the paper was received on 21 January 1985 and printed in April 1987. \citet[publisher description of the 1985--1986 lectures]{wright1989uniqueness} confirms circulation before the March 1987 meeting. This is therefore E, not a new conjecture by Gelfand.}
\GSfield{Present status}{As of 12 November 2026, the precise statement above is established; the qualifications distinguish it from the abbreviated notebook formulation.}
\GSfield{Assessment}{literal statement: false as written; publication match: exact; verification: read the result; development: annotated; confidence: high}
\end{GSitem}
\begin{GSitem}{1986-704.4}{F}{Factors, attractors, chaos, and characteristic functions}
\GSfield{Kind; source}{remark; \GSpage{P001355}}
\GSfield{The notes}{A list juxtaposes factors, attractors, chaos, Hausdorff measures and characteristic functions.}
\GSfield{Statement}{The list records proposed connections among operator algebras, dynamics and fractal measures. It specifies no map, equality, invariant or question; each named theory has many inequivalent ways of interacting with the others.}
\GSfield{Fate}{The local context supports a thematic reading only. Neither the type-\(III_1\) theorem nor a later general theory of multifractals can serve as an exact solution of an unspecified list.}
\GSfield{Present status}{As of 12 November 2026, no definite mathematical assertion can be extracted from the list without inventing substantial additional content.}
\GSfield{Assessment}{literal statement: not determinable; publication match: none found; verification: unresolved lead; development: none; confidence: low}
\end{GSitem}

\GStalk{703}{1987-03-09}{A. A. Beilinson}{P001354-P001356}{Zeta functions, regulators, and Artin L-functions}
Beilinson\textquoteright s explicit heading is on P001355 after Gelfand\textquoteright s factor discussion. His exposition continues through P001358, beyond the catalogue endpoint: archimedean and nonarchimedean local factors, special values, regulators, p-adic interpolation and Artin L-functions. The attribution follows the manuscript headings rather than the overlapping catalogue spans.

\begin{GSitem}{1986-703.1}{E}{The Euler product and completed Riemann zeta function}
\GSfield{Kind; source}{theorem; \GSpage{P001355}}
\GSfield{The notes}{Beilinson lists the Euler product, meromorphic continuation with a single pole at 1, and the symmetric completion \(\pi^{-s/2}\Gamma(s/2)\zeta(s)\).}
\GSfield{Statement}{The Euler product converges absolutely for \(\operatorname{Re}s>1\). The zeta function continues meromorphically to \(\mathbb C\), with one simple pole at 1, and \(\Lambda(s)=\pi^{-s/2}\Gamma(s/2)\zeta(s)\) satisfies \(\Lambda(s)=\Lambda(1-s)\). The completed function itself has poles at both 0 and 1; only \(\tfrac12s(s-1)\Lambda(s)\) is entire.}
\GSfield{Fate}{These are classical facts reviewed through the local-global integral method in \citet[Sections~2--4]{kudla2003tate}. Tate\textquoteright s thesis circulated in 1950 and was printed in 1967; the source distinguishes those dates.}
\GSfield{Present status}{As of 12 November 2026, the precise statement above is established; the qualifications distinguish it from the abbreviated notebook formulation.}
\GSfield{Assessment}{literal statement: correct; publication match: exact; verification: read the result; development: annotated; confidence: high}
\end{GSitem}
\begin{GSitem}{1986-703.2}{E}{Archimedean and p-adic local zeta integrals}
\GSfield{Kind; source}{construction; \GSpage{P001355-P001356}}
\GSfield{The notes}{The Gamma factor is described as a Mellin transform of a Gaussian; the nonarchimedean analogue uses the characteristic function of \(\mathbb Z_p\setminus\{0\}\), called an open subgroup.}
\GSfield{Statement}{With additive normalization \(e^{-\pi x^2}\) and multiplicative measure \(dx/|x|\), the real local integral is \(\int_{\mathbb R^\times}e^{-\pi x^2}|x|^s\,dx/|x|=\pi^{-s/2}\Gamma(s/2)\). With \(\operatorname{vol}(\mathbb Z_p^\times)=1\), \(\int_{\mathbb Q_p^\times}1_{\mathbb Z_p}(x)|x|_p^s\,d^\times x=(1-p^{-s})^{-1}\) for \(\operatorname{Re}s>0\). The set \(\mathbb Z_p\setminus\{0\}\) is not a multiplicative subgroup.}
\GSfield{Fate}{\citet[local zeta integrals, Sections~2--3]{kudla2003tate} records these factors and the Haar conventions. Kirillov\textquoteright s volume analogy is heuristic: the volume of an \(s\)-ball has a different Gamma expression, even under analytic continuation in dimension.}
\GSfield{Present status}{As of 12 November 2026, the precise statement above is established; the qualifications distinguish it from the abbreviated notebook formulation.}
\GSfield{Assessment}{literal statement: false as written; publication match: exact; verification: read the result; development: annotated; confidence: high}
\end{GSitem}
\begin{GSitem}{1986-703.3}{E}{Trivial zeros and Bernoulli special values}
\GSfield{Kind; source}{theorem; \GSpage{P001356}}
\GSfield{The notes}{The functional equation gives simple zeros at negative even integers, and \(\zeta(1-n)=-B_n/n\) is stated at negative arguments.}
\GSfield{Statement}{For \(n\ge1\), \(\zeta(-2n)=0\) and the zero is simple. For integers \(m\ge2\), \(\zeta(1-m)=-B_m/m\), using \(B_1=-1/2\); the omitted \(m=1\) case must be handled separately because \(\zeta(0)=-1/2\). The simplicity at negative even integers uses the functional equation together with nonvanishing of the positive real zeta value.}
\GSfield{Fate}{The analytic continuation in \citet{kudla2003tate} and the Bernoulli calculation reproduced in \citet[Section~2]{gelbart2012reciprocal} support the corrected conventions. These facts predate the lecture.}
\GSfield{Present status}{As of 12 November 2026, the precise statement above is established; the qualifications distinguish it from the abbreviated notebook formulation.}
\GSfield{Assessment}{literal statement: correct with missing hypotheses; publication match: exact; verification: read the result; development: annotated; confidence: high}
\end{GSitem}
\begin{GSitem}{1986-703.4}{E}{Regularized power sums and the choice of local parameter}
\GSfield{Kind; source}{method; \GSpage{P001356}}
\GSfield{The notes}{The notes write \(\zeta(-n)=\sum k^n=(t\,d/dt)^n(1-t)^{-1}|_{t=1}\), clarified as the first regular term.}
\GSfield{Statement}{The divergent power sum cannot be evaluated by ordinary substitution. A compatible finite-part prescription uses \(t=e^{-u}\) and the constant term at \(u=0\) of \(\sum_{k\ge1}k^n e^{-ku}\); that constant is \(\zeta(-n)\). Taking instead the constant Laurent coefficient in \(t-1\) is a different operation and generally gives another answer. The regularization parameter must be part of the statement.}
\GSfield{Fate}{The Bernoulli generating-function method underlying the special-value formula is explicit in \citet[Section~2]{gelbart2012reciprocal}. The notebook does not specify a parameter; its abbreviated equality is therefore not literally correct as an ordinary series or a coordinate-free finite part.}
\GSfield{Present status}{As of 12 November 2026, the precise statement above is established; the qualifications distinguish it from the abbreviated notebook formulation.}
\GSfield{Assessment}{literal statement: false as written; publication match: partial; verification: read the result; development: annotated; confidence: high}
\end{GSitem}
\begin{GSitem}{1986-703.5}{E}{Borel regulators and derivatives at negative even integers}
\GSfield{Kind; source}{theorem; \GSpage{P001356}}
\GSfield{The notes}{The corrected line relates \(\zeta^{\prime}(-2n)\) to the covolume of the image of \(K_{4n+1}(\mathbb Z)\) in \(\mathbb R\).}
\GSfield{Statement}{For \(n\ge1\), the torsion-free part of \(K_{4n+1}(\mathbb Z)\) has rank one. Under the Borel regulator with the corresponding Tate normalization, its image is a lattice in \(\mathbb R\), and the nonzero leading zeta-value \(\zeta^{\prime}(-2n)\) agrees with its regulator covolume up to a nonzero rational factor and the declared orientation convention. Torsion is killed by the real regulator.}
\GSfield{Fate}{\citet[Section~6, especially~6.2--6.4]{borel1977zeta} proves the commensurability statement; the paper was received in November 1976 and published in 1977. The notebook\textquoteright s replacement of \(\zeta(1-2n)\) by a derivative at \(-2n\) is mathematically meaningful and should be preserved.}
\GSfield{Present status}{As of 12 November 2026, the precise statement above is established; the qualifications distinguish it from the abbreviated notebook formulation.}
\GSfield{Assessment}{literal statement: correct with missing hypotheses; publication match: exact; verification: read the result; development: annotated; confidence: high}
\end{GSitem}
\begin{GSitem}{1986-703.6}{E}{p-adic interpolation and the missing measure regularization}
\GSfield{Kind; source}{construction; \GSpage{P001356}}
\GSfield{The notes}{The displayed identity makes \((1-p^{n-1})\zeta(1-n)\) a moment of one \(\mathbb Z_p\)-valued measure on \(\mathbb Z_p^\times\), with no auxiliary factor.}
\GSfield{Statement}{The Kubota--Leopoldt construction interpolates Euler-factor-removed zeta or Dirichlet L-values on the appropriate components of p-adic weight space. The zeta branch has a pole and is naturally a pseudomeasure; a bounded measure requires a regularizing factor. For example, in the convention of the cited source and even \(n\ge2\), Mazur\textquoteright s bounded measure satisfies \(\int x^{n-1}d\mu_{1,c}=-(1-c^{-n})(1-p^{n-1})\zeta(1-n)\), where \(c>1\) is prime to \(p\).}
\GSfield{Fate}{The original interpolation is \citet{kubotaleopoldt1964}; the explicit bounded-measure formula and its distinction from the unbounded Bernoulli distribution are read in \citet[equations~(2.1), (2.12)--(2.13)]{gelbart2012reciprocal}. The notebook omits both the regularization and the character-component conventions.}
\GSfield{Present status}{As of 12 November 2026, the precise statement above is established; the qualifications distinguish it from the abbreviated notebook formulation.}
\GSfield{Assessment}{literal statement: false as written; publication match: exact; verification: read the result; development: annotated; confidence: high}
\end{GSitem}
\begin{GSitem}{1986-703.7}{E}{Components of p-adic weight space}
\GSfield{Kind; source}{definition; \GSpage{P001356}}
\GSfield{The notes}{The units are split into \((1+p\mathbb Z_p)\) and the \((p-1)\)-st roots of unity, giving \(p-1\) components of the character space.}
\GSfield{Statement}{For odd \(p\), \(\mathbb Z_p^\times=\mu_{p-1}\times(1+p\mathbb Z_p)\). Continuous p-adic characters split according to their finite character on \(\mu_{p-1}\); each component is parametrized analytically by the character on the pro-p factor. For \(p=2\), the corresponding decomposition is \(\{\pm1\}\times(1+4\mathbb Z_2)\), so the written odd-prime formula must not be used unchanged.}
\GSfield{Fate}{\citet[discussion of weight space and the Iwasawa algebra]{gelbart2012reciprocal} makes the character decomposition explicit in the interpolation construction. The prime restriction is missing from the notes.}
\GSfield{Present status}{As of 12 November 2026, the precise statement above is established; the qualifications distinguish it from the abbreviated notebook formulation.}
\GSfield{Assessment}{literal statement: correct with missing hypotheses; publication match: exact; verification: read the result; development: annotated; confidence: high}
\end{GSitem}
\begin{GSitem}{1986-703.8}{C}{Riemann hypothesis and prime distribution}
\GSfield{Kind; source}{conjecture; \GSpage{P001356}}
\GSfield{The notes}{The final two numbered topics are RH and its connection with the distribution of primes, without an attempted proof.}
\GSfield{Statement}{The precise named conjecture is \(\operatorname{Re}\rho=1/2\) for every nontrivial zero \(\rho\) of \(\zeta\). Explicit-formula estimates relate the zero locations to the error in prime-counting functions; the prime number theorem by itself is much weaker than RH.}
\GSfield{Fate}{This repeats the conjecture behind record 700. \citet{clay2026riemann} lists RH as unsolved; no resolution is inferred from the adjacent exposition of known analytic continuation.}
\GSfield{Present status}{Open as of 12 November 2026; the notebook gives a topic heading, not a new equivalent criterion or claimed solution.}
\GSfield{Assessment}{literal statement: correct with missing hypotheses; publication match: exact; verification: read the result; development: annotated; confidence: high; problem \hyperref[prob:GP-1986-61]{GP-1986-61}}
\GSfield{Worth developing}{Its role is to identify the common unsolved target of the two lectures; an attempted proof is outside a bounded critical annotation.}
\end{GSitem}
\begin{GSitem}{1986-703.9}{E}{Finite-image Artin representations and ramified Euler factors}
\GSfield{Kind; source}{definition; \GSpage{P001357}}
\GSfield{The notes}{A continuous complex representation of the absolute Galois group is localized at primes; inertia and Frobenius are introduced, followed by the unramified determinant factor.}
\GSfield{Statement}{A continuous finite-dimensional representation \(\rho:G_{\mathbb Q}\to GL(V)\) with the usual complex topology has finite image. It is unramified outside finitely many primes. At every prime, the Artin factor is \(L_p(\rho,s)=\det(1-p^{-s}\rho(\mathrm{Fr}_p)|V^{I_p})^{-1}\), with a consistent arithmetic-Frobenius convention. The choice of decomposition group changes the data by conjugacy. The notebook formula on all of \(V\) applies at unramified primes.}
\GSfield{Fate}{\citet[Introduction]{gunhazrasahu2025artin} gives the full inertia-invariant definition and finite-Galois setup. The formulation is classical; the 2025 paper is used to verify the exact general convention and current status, not as the origin of Artin theory.}
\GSfield{Present status}{As of 12 November 2026, the precise statement above is established; the qualifications distinguish it from the abbreviated notebook formulation.}
\GSfield{Assessment}{literal statement: correct with missing hypotheses; publication match: exact; verification: read the result; development: annotated; confidence: high}
\end{GSitem}
\begin{GSitem}{1986-703.10}{E}{Dirichlet and Dedekind examples of Artin formalism}
\GSfield{Kind; source}{theorem; \GSpage{P001357-P001358}}
\GSfield{The notes}{One-dimensional characters give Dirichlet functions; inducing the trivial character from \(G_K\) gives \(\zeta_K\); direct sums multiply L-functions.}
\GSfield{Statement}{For a finite-order one-dimensional character of \(G_{\mathbb Q}\), the Artin L-function is the primitive Dirichlet L-function associated with its conductor. For any number field \(K\), \(L(\operatorname{Ind}_{G_K}^{G_{\mathbb Q}}1,s)=\zeta_K(s)\). Moreover \(L(\rho_1\oplus\rho_2,s)=L(\rho_1,s)L(\rho_2,s)\). Using an imprimitive character modulo a larger integer removes extra Euler factors and must be distinguished from the primitive Artin function.}
\GSfield{Fate}{\citet[Section~1]{kudla2003tate} explains primitive and imprimitive characters, while \citet[Section~2]{mantillasoler2019equivalence} gives the induced Galois representation for the Dedekind zeta function.}
\GSfield{Present status}{As of 12 November 2026, the precise statement above is established; the qualifications distinguish it from the abbreviated notebook formulation.}
\GSfield{Assessment}{literal statement: correct with missing hypotheses; publication match: exact; verification: read the result; development: annotated; confidence: high}
\end{GSitem}
\begin{GSitem}{1986-703.11}{E}{An Artin L-function determines the representation but not the field}
\GSfield{Kind; source}{theorem; \GSpage{P001358}}
\GSfield{The notes}{The notes explicitly distinguish recovering \(\rho\) from \(L(\rho,s)\) and the failure to recover \(K\) from \(\zeta_K(s)\).}
\GSfield{Statement}{Finite-image complex representations of \(G_{\mathbb Q}\) with identical Artin L-functions are isomorphic: equality of Dirichlet coefficients determines unramified Frobenius traces, Chebotarev determines the finite-group character, and complex finite-group representations are semisimple. However nonconjugate subgroups can induce isomorphic permutation representations, yielding nonisomorphic arithmetically equivalent number fields.}
\GSfield{Fate}{\citet[Theorem~2.5, Corollary~2.8 and Section~3 examples]{mantillasoler2019equivalence} gives this argument and explicit field phenomena, tracing the arithmetic-equivalence criterion to Perlis (1977) and Gassmann. The later paper confirms rather than originates the notebook\textquoteright s distinction.}
\GSfield{Present status}{As of 12 November 2026, the precise statement above is established; the qualifications distinguish it from the abbreviated notebook formulation.}
\GSfield{Assessment}{literal statement: correct; publication match: exact; verification: read the result; development: annotated; confidence: high}
\end{GSitem}
\begin{GSitem}{1986-703.12}{C}{Artin holomorphy}
\GSfield{Kind; source}{conjecture; \GSpage{P001358}}
\GSfield{The notes}{All complex Artin functions are said to be meromorphic; for nontrivial irreducible \(\rho\), holomorphy everywhere is conjectured.}
\GSfield{Statement}{Brauer induction gives meromorphic continuation to \(\mathbb C\) of every finite-image complex Artin L-function. Artin\textquoteright s conjecture asserts that \(L(\rho,s)\) is entire when \(\rho\) is nontrivial and irreducible. Trivial summands must be excluded because they contribute a Dedekind-zeta pole. Monomial representations and many further classes are known, but the general conjecture remains open.}
\GSfield{Fate}{\citet[Introduction and Theorems in Section~1]{gunhazrasahu2025artin}, published 17 February 2025, explicitly distinguishes the classical meromorphic-continuation theorem from the still-open holomorphy assertion and gives new partial criteria.}
\GSfield{Present status}{As of 12 November 2026, meromorphic continuation is established for all complex Artin representations, while the general nontrivial irreducible holomorphy conjecture is unresolved.}
\GSfield{Assessment}{literal statement: correct; publication match: exact; verification: read the result; development: annotated; confidence: high; problem \hyperref[prob:GP-1986-64]{GP-1986-64}}
\GSfield{Worth developing}{The exact finite-image formulation and the boundary between Brauer meromorphy and entire continuation deserve preservation; no general proof is attempted here.}
\end{GSitem}
\begin{GSitem}{1986-703.13}{F}{The scope of the l-adic continuation question}
\GSfield{Kind; source}{remark; \GSpage{P001358}}
\GSfield{The notes}{For l-adic representations, the notes say that even meromorphic continuation is unknown.}
\GSfield{Statement}{This sentence cannot apply uniformly to arbitrary continuous \(\ell\)-adic representations without additional conditions. To formulate the intended arithmetic problem one needs, at a minimum, a representation unramified outside finitely many places, suitable algebraic Frobenius eigenvalues and a complex realization; the geometric or motivic setting supplies the usual conjectural analytic continuation problem. These hypotheses are not stated in the notebook.}
\GSfield{Fate}{The adjoining Artin discussion, verified in \citet{gunhazrasahu2025artin}, concerns finite-image complex representations and cannot settle this broader question. No specific motivic representation or exact conjecture is identified here.}
\GSfield{Present status}{As of 12 November 2026, the unqualified sentence is not a well-defined universal theorem or conjecture. Important geometric cases are known, but none is specified sufficiently here for a precise fate assignment.}
\GSfield{Assessment}{literal statement: not determinable; publication match: thematic only; verification: unresolved lead; development: none; confidence: medium}
\end{GSitem}
\begin{GSitem}{1986-703.14}{E}{The completed Artin functional equation}
\GSfield{Kind; source}{theorem; \GSpage{P001358}}
\GSfield{The notes}{The last formula relates \(L^*(\rho,s)\) to \(L^*(\rho^*,1-s)\), with an epsilon factor containing Gauss sums.}
\GSfield{Statement}{For a finite-image complex Artin representation over \(\mathbb Q\), let \(a_\pm\) be the dimensions of the \(\pm1\)-eigenspaces of complex conjugation and \(A(\rho)\) its Artin conductor. With \(\Gamma_{\mathbb R}(s)=\pi^{-s/2}\Gamma(s/2)\), set \(\Lambda(\rho,s)=A(\rho)^{s/2}\Gamma_{\mathbb R}(s)^{a_+}\Gamma_{\mathbb R}(s+1)^{a_-}L(\rho,s)\). Then \(\Lambda(\rho,s)=W(\rho)\Lambda(\rho^*,1-s)\) meromorphically, where \(|W(\rho)|=1\). If the conductor is omitted from the completion, the epsilon factor depends on \(s\).}
\GSfield{Fate}{The Artin analytic framework in \citet{gunhazrasahu2025artin} distinguishes known functional equations from conjectural holomorphy; the local-global normalization is illustrated in \citet[Sections~3--4]{kudla2003tate}. The notebook has not fixed which conductor factor is absorbed in \(L^*\).}
\GSfield{Present status}{As of 12 November 2026, the precise statement above is established; the qualifications distinguish it from the abbreviated notebook formulation.}
\GSfield{Assessment}{literal statement: correct with missing hypotheses; publication match: partial; verification: read the result; development: annotated; confidence: high}
\end{GSitem}

\GStalk{705}{1987-03-23}{V. E. Zakharov}{P001359-P001364}{Dressing methods and spectral differential operators}
Zakharov reviews the Marchenko and nonlocal d-bar dressing constructions and then considers commuting operators that differentiate in the spectral parameter. The novelty claim concerns this last extension, not inverse scattering itself. Several displayed Lax and matrix identities have sign or hypothesis problems. The Migdal heading appears partway down P001363.

\begin{GSitem}{1986-705.1}{E}{KdV as a Lax compatibility condition}
\GSfield{Kind; source}{method; \GSpage{P001359}}
\GSfield{The notes}{The displayed equation is \(u_t+6uu_x+u_{xxx}=0\), while the spectral equations use \(L=\partial_x^2+u\) and \(A=4\partial_x^3+6u\partial_x+3u_x\). The free wave is \(e^{ikx+4ik^3t}\).}
\GSfield{Statement}{For this pair, compatibility \(L_t=[A,L]\) gives \(u_t=6uu_x+u_{xxx}\). A free solution with wave number \(k\) has spectral eigenvalue \(-k^2\) and time factor \(e^{-4ik^3t}\). The standard KdV convention in the first line is recovered by reversing time and making the same change throughout the pair. The potential and eigenfunction must be smooth enough for the differentiations.}
\GSfield{Fate}{The inverse-scattering and dressing framework is classical and is reviewed in \citet[Section~2, equations~(25)--(31)]{zakharovshabat1974scheme}. The notebook combines two opposite time conventions. This is a correction of a faithfully recorded formula, not a new integrability claim.}
\GSfield{Present status}{As of 12 November 2026, the precise statement above is established; the qualifications distinguish it from the abbreviated notebook formulation.}
\GSfield{Assessment}{literal statement: false as written; publication match: partial; verification: read the result; development: annotated; confidence: high}
\end{GSitem}
\begin{GSitem}{1986-705.2}{E}{Marchenko dressing and its solvability conditions}
\GSfield{Kind; source}{construction; \GSpage{P001359-P001360}}
\GSfield{The notes}{The integral equation is \(K(x,z)+F(x+z)+\int K(x,s)F(s+z)\,ds=0\), with three proposed integration conventions, \(F_t=8F_{\xi\xi\xi}\), and \(u=2(d/dx)K(x,x)\).}
\GSfield{Statement}{In the right Marchenko convention the integral is over \([x,\infty)\). Assume sufficient smoothness and decay to differentiate and integrate by parts, and invertibility of the corresponding integral operator so that the normalized equation has a unique solution. The diagonal formula then produces the potential for the consistently signed Lax system. The left and half-difference prescriptions require their own boundary or formal-integral conventions; they are not unconditional solutions for every function \(F\).}
\GSfield{Fate}{\citet[Section~1, equations~(4)--(6) and Theorems~1--2]{zakharovshabat1974scheme} develops the general dressing scheme cited explicitly in the lecture. The notebook omits the unique-solvability and boundary assumptions and carries over the KdV time-sign inconsistency. Its extension to growing exponentials is a formal prescription, not a convergent integral on an infinite interval.}
\GSfield{Present status}{As of 12 November 2026, the precise statement above is established; the qualifications distinguish it from the abbreviated notebook formulation.}
\GSfield{Assessment}{literal statement: correct with missing hypotheses; publication match: partial; verification: read the result; development: annotated; confidence: high}
\end{GSitem}
\begin{GSitem}{1986-705.3}{E}{The KP extension needs a nonlocal coefficient in its Lax pair}
\GSfield{Kind; source}{construction; \GSpage{P001360}}
\GSfield{The notes}{The kernel condition becomes \(F_{xx}-F_{zz}=\alpha F_y\), but the displayed wave equation has coefficient 2 and retains the one-dimensional time operator unchanged.}
\GSfield{Statement}{A consistent KP pair is \((\alpha\partial_y-\partial_x^2-u)\psi=0\) and \((\partial_t-4\partial_x^3-6u\partial_x-3u_x-3\alpha w)\psi=0\), with \(w_x=u_y\). Compatibility yields \((u_t-6uu_x-u_{xxx})_x=3\alpha^2u_{yy}\). Boundary data fix the integration constant in \(w\). Omitting \(w\) imposes extra restrictions and does not give generic KP solutions.}
\GSfield{Reconstruction}{The auxiliary potential \(w\), its boundary convention and a consistent \(\alpha\)-normalization are supplied editorially.}
\GSfield{Fate}{The multidimensional dressing construction in \citet[Section~2, equations~(36)--(38)]{zakharovshabat1974scheme} is the precise antecedent acknowledged by the notes. The reconstructed pair fixes a declared convention; it is not a verbatim restoration of an absent notebook term.}
\GSfield{Present status}{As of 12 November 2026, the precise statement above is established; the qualifications distinguish it from the abbreviated notebook formulation.}
\GSfield{Assessment}{literal statement: false as written; publication match: partial; verification: read the result; development: annotated; confidence: high}
\end{GSitem}
\begin{GSitem}{1986-705.4}{E}{The normalized nonlocal d-bar problem}
\GSfield{Kind; source}{method; \GSpage{P001360-P001361}}
\GSfield{The notes}{The Fourier-transformed kernel leads to \(\bar\partial_\lambda\chi=\int\chi(\eta)T(\eta,\lambda)\,d\eta\,d\bar\eta\), with \(\chi(\infty)=1\); the coefficient of \(\lambda^{-1}\) determines \(u\).}
\GSfield{Statement}{Choose an area-form convention and a kernel for which the normalized nonlocal \(\bar\partial\)-problem is uniquely solvable, for example in a Fredholm setting with no normalized homogeneous solution. Its Cauchy transform must use the matching distributional normalization \(\bar\partial(1/z)=\pi\delta_0\). Differentiating a uniquely solvable problem with compatible parameter dependence supplies the auxiliary linear equations. The potential is extracted from the first coefficient at infinity after fixing the oscillatory convention.}
\GSfield{Fate}{\citet[Section~1]{zakharovmanakov1984multidimensional} explicitly makes uniqueness the central dressing hypothesis and announces the nonlocal d-bar extension. \citet{fokaszakharov1992dressing} later treats this method directly. The original lecture is identifiable as this construction; an arbitrary kernel does not automatically yield a nonsingular solution.}
\GSfield{Present status}{As of 12 November 2026, the precise statement above is established; the qualifications distinguish it from the abbreviated notebook formulation.}
\GSfield{Assessment}{literal statement: correct with missing hypotheses; publication match: partial; verification: read the result; development: annotated; confidence: high}
\end{GSitem}
\begin{GSitem}{1986-705.5}{E}{Commuting rational flows and Riemann-surface dressing data}
\GSfield{Kind; source}{construction; \GSpage{P001361-P001362}}
\GSfield{The notes}{The kernel evolves by \(T_{x_i}=I_i(\eta)T-TI_i(\lambda)\), with mutually commuting rational matrix functions; the notes also allow functions on a Riemann surface.}
\GSfield{Statement}{For compatible commuting multipliers, the kernel flows commute. A normalized uniquely solvable dressing problem then generates compatible dressed operators. On a compact Riemann surface the meromorphic functions, allowed poles, normalization divisor and any vector bundle must be specified. The polynomial choices for KP must agree with the same spectral and time conventions used to define the Fourier transform.}
\GSfield{Fate}{\citet{zakharovmanakov1984multidimensional} supplies the established multidimensional dressing framework. The recorded cubic flow \(4\lambda^3\), without the imaginary factor dictated by the earlier oscillatory wave, cannot be matched to that earlier convention without a change of time variable. The unrestricted Riemann-surface sentence is a programmatic extension rather than a complete existence theorem.}
\GSfield{Present status}{As of 12 November 2026, the precise statement above is established; the qualifications distinguish it from the abbreviated notebook formulation.}
\GSfield{Assessment}{literal statement: correct with missing hypotheses; publication match: partial; verification: read the result; development: annotated; confidence: high}
\end{GSitem}
\begin{GSitem}{1986-705.6}{E}{First-order differentiation in a variable spectral parameter}
\GSfield{Kind; source}{method; \GSpage{P001362}}
\GSfield{The notes}{The new proposal is \(L_i=\partial_{x_i}+K_i(\lambda)+f_i(\lambda)\partial_\lambda\), with rational coefficients and moving poles; the gravitational example is attributed to Belinsky and Zakharov.}
\GSfield{Statement}{The compatibility equation for such operators contains both matrix curvature and the commutator of the scalar spectral vector fields. Once an ansatz for poles and residues is fixed, its coefficients give nonlinear equations in the deformation variables. The count of equations and unknowns and solvability depend on that ansatz, so they are not universal consequences of writing \([L_i,L_j]=0\).}
\GSfield{Fate}{\citet[Introduction, equations~(1.9)--(1.10), and subsequent constructions]{burtsevzakharovmikhailov1987variable} treats inverse scattering with a variable spectral parameter and includes the gravitational motivation. The paper was received on 6 February 1986 and published in March 1987, so this first-order part was circulating before the lecture, even though the notebook labels the wider proposal new.}
\GSfield{Present status}{As of 12 November 2026, the precise statement above is established; the qualifications distinguish it from the abbreviated notebook formulation.}
\GSfield{Assessment}{literal statement: correct with missing hypotheses; publication match: partial; verification: read the result; development: annotated; confidence: high}
\end{GSitem}
\begin{GSitem}{1986-705.7}{E}{The zero-curvature condition for spectral differential operators}
\GSfield{Kind; source}{claim; \GSpage{P001362}}
\GSfield{The notes}{For \(L_i=\partial_{x_i}+M_i(\lambda,\partial_\lambda)\), the notes write \([L_i,L_j]=0\) if and only if \([M_i,M_j]=0\).}
\GSfield{Statement}{The full condition is \(\partial_{x_i}M_j-\partial_{x_j}M_i+[M_i,M_j]=0\). It reduces to the written commutativity condition only when the \(M_i\) are independent of all deformation variables, or when their cross derivatives cancel for a specified reason. Rationality in \(\lambda\) alone does not remove these terms.}
\GSfield{Fate}{This is the elementary compatibility identity used in the variable-spectral-parameter framework of \citet{burtsevzakharovmikhailov1987variable}. The sentence becomes correct for the constant-coefficient-in-\(x\) example that follows, but not for the preceding claim about freely moving poles.}
\GSfield{Present status}{As of 12 November 2026, the precise statement above is established; the qualifications distinguish it from the abbreviated notebook formulation.}
\GSfield{Assessment}{literal statement: false as written; publication match: partial; verification: read the result; development: annotated; confidence: high}
\end{GSitem}
\begin{GSitem}{1986-705.8}{E}{A corrected commuting pair of matrix differential operators}
\GSfield{Kind; source}{construction; \GSpage{P001362}}
\GSfield{The notes}{The notes require \(I^2=1\), \(A^2=-1\), \(IA+AI=0\), set \(C=-(I+A)/2\), \(M_1=I\partial_\lambda+A/\lambda\), \(M_2=M_1\partial_\lambda+C/\lambda^2\), but display \(A=\left(\begin{smallmatrix}0&1\\1&0\end{smallmatrix}\right)\).}
\GSfield{Statement}{With \(I=\operatorname{diag}(1,-1)\) and \(A=\left(\begin{smallmatrix}0&1\\-1&0\end{smallmatrix}\right)\), the required matrix relations hold and \([M_1,M_2]=0\) on vector-valued meromorphic functions away from \(\lambda=0\). With the notebook\textquoteright s plus/plus matrix the relations fail and \([M_1,M_2]=2\operatorname{Id}/\lambda^3\).}
\GSfield{Fate}{\oursComputed{} The exact operator composition and a second calculation by direct action on polynomial vectors independently verify both conclusions. The second check uses two basis vectors, seven degrees and four nonzero rational spectral values for both matrices, giving 112 exact comparisons. The supplied code and output record these checks. The scan confirms that both displayed off-diagonal signs are positive; this is a faithful mathematical error. The general compatibility method predates the talk, but no exact published priority claim is made for this particular illustrative pair.}
\GSfield{Present status}{As of 12 November 2026, the bounded algebraic issue is completely checked in this companion. This verification repairs an example; it does not establish the lecture\textquoteright s proposed general classification of nonlinear systems.}
\GSfield{Assessment}{literal statement: false as written; publication match: partial; verification: read the result; development: annotated; confidence: high}
\end{GSitem}
\begin{GSitem}{1986-705.9}{F}{Solving the common system for higher spectral differential operators}
\GSfield{Kind; source}{question; \GSpage{P001362}}
\GSfield{The notes}{After the commuting example the notes ask to find solutions of \(L_i y=0\), and say that a great deal of computation is required.}
\GSfield{Statement}{For the proposed operators containing higher derivatives in the spectral variable, a common-solution problem needs specified coefficient functions, domains, singularities, initial or asymptotic data, and a function class for \(y\). The notebook supplies an algebraic example but no such analytic problem. Formal compatibility alone does not specify a unique solution or an inverse-scattering theorem.}
\GSfield{Fate}{\citet{burtsevzakharovmikhailov1987variable} identifies the first-order spectral-derivative background. A bibliographic lead to Zakharov\textquoteright s 1988 conference contribution on commuting operators and a nonlocal d-bar problem was located, but its original text was not recovered; it is not used to certify the unrestricted higher-order proposal.}
\GSfield{Present status}{As of 12 November 2026, the analytic problem intended here is underdetermined. The verified algebraic example is separated in item 1986-705.8; no general open problem is inferred from the missing formulation.}
\GSfield{Assessment}{literal statement: not determinable; publication match: thematic only; verification: unresolved lead; development: none; confidence: medium}
\end{GSitem}
\begin{GSitem}{1986-705.10}{F}{Degenerating soliton kernels to rational solutions}
\GSfield{Kind; source}{method; \GSpage{P001363}}
\GSfield{The notes}{The prescription begins with Gelfand--Levitan--Marchenko equations in the spectral variable, takes \(n\)-soliton solutions, and degenerates them to rational solutions; the next step is to construct the kernels \(K(\lambda,\lambda^{\prime},x_i)\) and \(F(\lambda,\lambda^{\prime},x_i)\).}
\GSfield{Statement}{A definite degeneration theorem must specify the nonlinear system, a family of scattering or kernel data, the rescaling and coalescence parameters, and the topology in which the potentials or wave functions converge. None of these limiting data is fixed in this passage. The assertion is therefore a proposed method, rather than a stated theorem for all rational solutions.}
\GSfield{Fate}{\citet[Section~3]{zakharovshabat1974scheme} gives explicit finite-rank soliton kernels in the established dressing method. That construction is an antecedent, not an exact publication match for the unspecified spectral-variable rational limit suggested here.}
\GSfield{Present status}{As of 12 November 2026, rational limits of soliton constructions are a recognizable topic, but this notebook passage does not select one equation or degeneration theorem whose precise fate can be assigned.}
\GSfield{Assessment}{literal statement: not determinable; publication match: thematic only; verification: read the result; development: none; confidence: medium}
\end{GSitem}
\begin{GSitem}{1986-705.11}{F}{The proposed scheme is said to include dissipative systems}
\GSfield{Kind; source}{claim; \GSpage{P001363}}
\GSfield{The notes}{The single sentence says that dissipative systems are included too.}
\GSfield{Statement}{To test this assertion one must define the evolution, the real form, the phase space and a dissipated functional or contraction property. A non-self-adjoint spectral operator by itself is not a definition of dissipative dynamics. The passage gives none of these data.}
\GSfield{Fate}{The spectral-parameter extension of \citet{burtsevzakharovmikhailov1987variable} provides related operator machinery, but it does not identify which dissipative system this sentence intends. No exact theorem is assigned on the strength of the shared vocabulary.}
\GSfield{Present status}{As of 12 November 2026, the specific claim remains unidentifiable from the notes. This is a limitation of the recorded statement, not evidence that dissipative integrable models have not been developed.}
\GSfield{Assessment}{literal statement: not determinable; publication match: thematic only; verification: unresolved lead; development: none; confidence: low}
\end{GSitem}
\begin{GSitem}{1986-705.12}{F}{Spectral differential operators multiplying the evolution derivatives}
\GSfield{Kind; source}{construction; \GSpage{P001363}}
\GSfield{The notes}{The next generalization is \(L_i=R_i(\lambda,\partial_\lambda)\partial_{x_i}-M_i(\lambda,\partial_\lambda)\).}
\GSfield{Statement}{The displayed operator pencil is well-defined formally once coefficient regularity and an operator ordering are chosen. A compatibility or dressing assertion additionally requires the orders, dependence on the evolution variables, domains, and any invertibility assumptions on \(R_i\). The notebook states the ansatz but supplies no resulting equation or theorem.}
\GSfield{Fate}{The inspected first-order predecessor \citet{burtsevzakharovmikhailov1987variable} is a thematic antecedent. The missing conference text cannot be used to claim an exact later treatment of arbitrary \(R_i\) and \(M_i\).}
\GSfield{Present status}{As of 12 November 2026, the ansatz is clear but its intended mathematical assertion is not. No universal solvability theorem or undeveloped open problem is asserted.}
\GSfield{Assessment}{literal statement: not determinable; publication match: thematic only; verification: unresolved lead; development: none; confidence: medium}
\end{GSitem}

\GStalk{706}{1987-03-23}{A. A. Migdal}{P001362-P001369}{Matrix loop equations and Wilson loops}
Migdal derives Schwinger--Dyson identities for matrix traces, planar factorization and the resolvent equation, before connecting them with gauge-theory loop equations. Gelfand interrupts with the cube-plane problem and its answer 96. The next dated meeting begins on P001368; P001362 and the beginning of P001363 remain Zakharov\textquoteright s.

\begin{GSitem}{1986-706.1}{E}{Matrix expectation and the integration contour}
\GSfield{Kind; source}{definition; \GSpage{P001363}}
\GSfield{The notes}{The matrix integral has weight \(\exp[-(N/\lambda)\operatorname{Tr}V(A)]\), while \(A^*=-A\); the loop observable is \(\psi(t)=N^{-1}\operatorname{Tr}e^{tA}\).}
\GSfield{Statement}{A rigorous one-matrix model can be defined on Hermitian matrices with real confining polynomial \(V\) and \(\lambda>0\), normalized by its finite partition function. A skew-Hermitian convention instead needs a compatible contour and potential obtained by \(A=iH\); it is not the same real-eigenvalue integral without changing these data. Formal power-series expectations provide another explicitly distinct interpretation.}
\GSfield{Reconstruction}{The confining Hermitian model is the precise setting used for the analytical statements below; it is an editorial choice, not a silent change to the recorded anti-Hermitian line.}
\GSfield{Fate}{\citet[Sections~2--3]{brezinitzyksonparisizuber1978planar} defines the Hermitian planar model and its eigenvalue reduction. The notebook switches between anti-Hermitian gauge fields and real-eigenvalue matrix formulas without declaring the contour change.}
\GSfield{Present status}{As of 12 November 2026, the precise statement above is established; the qualifications distinguish it from the abbreviated notebook formulation.}
\GSfield{Assessment}{literal statement: correct with missing hypotheses; publication match: exact; verification: read the result; development: annotated; confidence: high}
\end{GSitem}
\begin{GSitem}{1986-706.2}{E}{The finite-matrix Schwinger--Dyson identity}
\GSfield{Kind; source}{theorem; \GSpage{P001363-P001364}}
\GSfield{The notes}{An integral of a total matrix derivative yields a convolution equation. The first formula omits the expectation signs, which the notebook itself restores on the next page.}
\GSfield{Statement}{Under the confining assumptions of the preceding item, integration by parts and the Duhamel identity \(D(e^{tA})[H]=\int_0^t e^{sA}He^{(t-s)A}\,ds\) give \[0=\int_0^t\mathbb E[\psi(s)\psi(t-s)]\,ds-\lambda^{-1}V^{\prime}(\partial_t)\mathbb E\psi(t).\] Boundary terms must vanish and the derivative convention must be the Hermitian one. This is an identity of expectations, not an identity for each matrix.}
\GSfield{Fate}{The loop-equation method is already present in \citet{brezinitzyksonparisizuber1978planar}. Its exact finite-size identity precedes any use of large-N factorization. The notebook\textquoteright s later correction is preserved rather than reported as a new editorial discovery.}
\GSfield{Present status}{As of 12 November 2026, the precise statement above is established; the qualifications distinguish it from the abbreviated notebook formulation.}
\GSfield{Assessment}{literal statement: correct with missing hypotheses; publication match: exact; verification: read the result; development: annotated; confidence: high}
\end{GSitem}
\begin{GSitem}{1986-706.3}{E}{Planes through triples of cube vertices}
\GSfield{Kind; source}{problem; \GSpage{P001364}; I. M. Gelfand, attributed to T. Alekseevskaya}
\GSfield{The notes}{Gelfand attributes an elegant problem to T. Alekseevskaya: draw a plane through every three vertices of a three-dimensional cube and count the pieces. The notebook already gives the answer 96.}
\GSfield{Statement}{Take the eight vertices of a nondegenerate cube in \(\mathbb R^3\). Form all distinct affine planes through triples of distinct vertices, and count the connected components of the open cube after deleting these planes. There are 20 distinct planes; six are supporting face planes and fourteen cut the interior. The number of components is 96.}
\GSfield{Fate}{\oursComputed{} The chapter independently verifies the already recorded answer by exact rational cell enumeration and by the intersection-poset count \(1+14+47+34=96\). It supplies a rational witness for each region. This is a reproducible annotation of an answered notebook problem; no claim of a new open problem or new numerical discovery is made.}
\GSfield{Present status}{As of 12 November 2026, the notebook answer is confirmed by two independent exact methods in this companion. Historical publication priority for this elementary puzzle has not been determined.}
\GSfield{Assessment}{literal statement: correct; publication match: none found; verification: read the result; development: annotated; confidence: high}
\end{GSitem}
\begin{GSitem}{1986-706.4}{E}{The planar sector for classical matrix symmetry classes}
\GSfield{Kind; source}{claim; \GSpage{P001364}}
\GSfield{The notes}{The notes write \(O(\infty)=SU(\infty)=Sp(\infty)\) and say the large-size limit is the same for all classical groups.}
\GSfield{Statement}{With matching matrix and coupling normalizations, the leading oriented planar diagrams of real, complex and quaternionic invariant matrix models are related. The full theories differ in symmetry factors and nonorientable or subleading terms. There is no equality or isomorphism of the displayed infinite classical groups, and equality of arbitrary matrix-model observables is not implied.}
\GSfield{Fate}{The graphical expansion in \citet{bessisitzyksonzuber1980graphical} explains the intended leading-diagram comparison. The literal group equality is shorthand that becomes false if read as a group-theoretic theorem.}
\GSfield{Present status}{As of 12 November 2026, the precise statement above is established; the qualifications distinguish it from the abbreviated notebook formulation.}
\GSfield{Assessment}{literal statement: false as written; publication match: partial; verification: read the result; development: annotated; confidence: high}
\end{GSitem}
\begin{GSitem}{1986-706.5}{E}{Large-N factorization and closure of the planar equation}
\GSfield{Kind; source}{theorem; \GSpage{P001365}}
\GSfield{The notes}{Normalized trace correlations are replaced by the product of their expectations up to \(O(N^{-2})\).}
\GSfield{Statement}{In the unitary invariant formal topological expansion, connected correlators of two normalized traces begin at order \(N^{-2}\). In convergent analytic models this requires hypotheses that ensure the relevant planar limit and fluctuations; it is not a statement about arbitrary random matrices. Under those assumptions the finite Schwinger--Dyson equation closes for the limiting one-trace function.}
\GSfield{Fate}{\citet{brezinitzyksonparisizuber1978planar} develops the planar limit; \citet{bessisitzyksonzuber1980graphical} gives the Euler-characteristic bookkeeping. The same factorization principle appears in earlier occurrence 1982-565.3 and in the reduction discussion of record 698.}
\GSfield{Present status}{As of 12 November 2026, the precise statement above is established; the qualifications distinguish it from the abbreviated notebook formulation.}
\GSfield{Assessment}{literal statement: correct with missing hypotheses; publication match: partial; verification: read the result; development: annotated; confidence: high}
\end{GSitem}
\begin{GSitem}{1986-706.6}{E}{The planar resolvent and its quadratic equation}
\GSfield{Kind; source}{construction; \GSpage{P001365-P001366}}
\GSfield{The notes}{The Laplace transform of the trace exponential becomes \(f(z)=N^{-1}\mathbb E\operatorname{Tr}(z-A)^{-1}\), with \(R=f^2-\lambda^{-1}V^{\prime}f\); the next page solves the quadratic using \(-R\) under the root.}
\GSfield{Statement}{For a limiting compactly supported eigenvalue measure \(\rho\), \(W(z)=\int(z-x)^{-1}d\rho(x)\) is analytic off its support and satisfies \(W(z)\sim z^{-1}\). The planar equation gives \(R(z)=W(z)^2-\lambda^{-1}V^{\prime}(z)W(z)\), a polynomial fixed by moments and support conditions. Therefore \[W(z)=\frac{V^{\prime}(z)}{2\lambda}-\sqrt{\left(\frac{V^{\prime}(z)}{2\lambda}\right)^2+R(z)},\] with the branch fixed at infinity. A Laplace representation requires its own convergence half-plane.}
\GSfield{Fate}{\citet[planar eigenvalue solution]{brezinitzyksonparisizuber1978planar} supplies the established algebraic-resolvent method. Relative to the notebook\textquoteright s own definition of \(R\), its radicand sign is incorrect.}
\GSfield{Present status}{As of 12 November 2026, the precise statement above is established; the qualifications distinguish it from the abbreviated notebook formulation.}
\GSfield{Assessment}{literal statement: false as written; publication match: exact; verification: read the result; development: annotated; confidence: high}
\end{GSitem}
\begin{GSitem}{1986-706.7}{E}{Boundary values, density, and the one-cut ansatz}
\GSfield{Kind; source}{method; \GSpage{P001365-P001366}}
\GSfield{The notes}{The discontinuity of \(f\) is called its imaginary part without a factor \(i\); a single square root \(\sqrt{(z-a)(z-b)}\) is then used.}
\GSfield{Statement}{For a real density on its support, \(W_+-W_-=-2\pi i\rho\), \(\operatorname{Im}W_+=(W_+-W_-)/(2i)\), and \(W_+^2-W_-^2=4i\operatorname{Re}W_+\operatorname{Im}W_+\). The form \(W=V^{\prime}/(2\lambda)+M(z)\sqrt{(z-a)(z-b)}\) assumes a one-cut solution. Polynomial \(M\), endpoints, normalization and positivity are constrained by the potential; multi-cut phases need more branch points.}
\GSfield{Fate}{\citet[Sections~2--3]{brezinitzyksonparisizuber1978planar} develops the one-cut solution and its selection conditions. Zakharevich\textquoteright s boundary-value intervention on P001366 fits this scalar Riemann--Hilbert reformulation; it does not remove the support assumptions.}
\GSfield{Present status}{As of 12 November 2026, the precise statement above is established; the qualifications distinguish it from the abbreviated notebook formulation.}
\GSfield{Assessment}{literal statement: false as written; publication match: exact; verification: read the result; development: annotated; confidence: high}
\end{GSitem}
\begin{GSitem}{1986-706.8}{E}{Wilson loops as gauge-invariant parallel transport traces}
\GSfield{Kind; source}{definition; \GSpage{P001366}}
\GSfield{The notes}{The matrix exponential is replaced by \(W(c)=N^{-1}\operatorname{Tr}P\exp\oint_c A_\mu dx_\mu\).}
\GSfield{Statement}{For a smooth connection in a unitary representation and a piecewise smooth closed path, path ordering defines its parallel transport. Its normalized trace is invariant under gauge transformations because the holonomy transforms by conjugation at the basepoint. For quantum fields this formal observable requires a regulator and, in continuum limits, renormalization depending on path singularities.}
\GSfield{Fate}{\citet[Sections~1 and~3]{makeenko2008cusped} reviews this construction and the additional ultraviolet issues at cusps and intersections. It is established gauge-theory notation, not a claim that an unregularized four-dimensional path integral has been constructed.}
\GSfield{Present status}{As of 12 November 2026, the precise statement above is established; the qualifications distinguish it from the abbreviated notebook formulation.}
\GSfield{Assessment}{literal statement: correct with missing hypotheses; publication match: exact; verification: read the result; development: annotated; confidence: high}
\end{GSitem}
\begin{GSitem}{1986-706.9}{E}{The Makeenko--Migdal loop equation and its scope}
\GSfield{Kind; source}{method; \GSpage{P001366-P001367}}
\GSfield{The notes}{The lecture proposes repeating the matrix integration-by-parts argument for the gauge-field functional integral and ends with an abbreviated functional derivative.}
\GSfield{Statement}{The formal gauge-field Schwinger--Dyson identity relates an area derivative of a Wilson loop to a contact term splitting the contour at coincident points. In the planar limit, factorization replaces the two-loop expectation by a product. A precise continuum equation must specify regularization, renormalization and treatment of cusps and intersections; deriving this identity is not a solution of four-dimensional Yang--Mills theory.}
\GSfield{Fate}{The original equation was published in \citet{makeenkomigdal1979exact}; \citet[Section~3.1, equation~(3.2), and Section~3.3]{makeenko2008cusped} gives the regularized interpretation and its limitations. The lecture is an exposition of the established loop-equation method.}
\GSfield{Present status}{As of 12 November 2026, the precise statement above is established; the qualifications distinguish it from the abbreviated notebook formulation.}
\GSfield{Assessment}{literal statement: correct with missing hypotheses; publication match: partial; verification: read the result; development: annotated; confidence: high}
\end{GSitem}
\begin{GSitem}{1986-706.10}{F}{The proposed noncommutative analogue of Karasev symbols}
\GSfield{Kind; source}{remark; \GSpage{P001363; P001366}}
\GSfield{The notes}{The talk is introduced as a noncommutative analogue of Karasev\textquoteright s symbol. Later it says that the components of a gauge field cannot be diagonalized simultaneously.}
\GSfield{Statement}{The obstruction to simultaneous diagonalization of noncommuting matrices is definite. The intended analogue of a Karasev symbol, however, has no stated domain, product, map or asymptotic parameter beyond large matrix size. It cannot be uniquely identified with a particular deformation-quantization construction from this phrase.}
\GSfield{Fate}{The matrix and loop methods are identified in the preceding items. They supply thematic context, but no exact symbol calculus is reconstructed or credited on the strength of this introductory phrase.}
\GSfield{Present status}{As of 12 November 2026, the particular proposed symbol calculus is not identifiable from the record; the matrix loop equations themselves are established.}
\GSfield{Assessment}{literal statement: not determinable; publication match: none found; verification: unresolved lead; development: none; confidence: low}
\end{GSitem}

\GStalk{708}{1987-04-13}{Ya. G. Sinai}{P001367-P001372}{Spin glasses and ultrametricity}
Sinai\textquoteright s discussion on P001368--P001369 distinguishes finite-range spin glasses from the mean-field model, raises the thermodynamic ground-state problem, and describes the Parisi picture of hierarchical overlaps. Later theorems settle carefully specified mean-field assertions, but do not justify the literal finite-size statement of many exactly degenerate ground states.

\begin{GSitem}{1986-708.1}{E}{Finite-range and mean-field spin-glass Hamiltonians}
\GSfield{Kind; source}{definition; \GSpage{P001368}}
\GSfield{The notes}{The lecture starts with nearest-neighbour random couplings and then passes to all-pairs Gaussian couplings scaled by \(N^{-1/2}\).}
\GSfield{Statement}{The finite-range Edwards--Anderson model uses a specified lattice and neighbouring pairs. The zero-field Sherrington--Kirkpatrick model has \(\sigma\in\{-1,1\}^N\), independent standard Gaussian \(J_{ij}\), and \(H_N(\sigma)=N^{-1/2}\sum_{i<j}J_{ij}\sigma_i\sigma_j\). Independence, centering and the summation convention determine the normalization. The thermodynamic ground energy is normalized by \(N\).}
\GSfield{Fate}{The precise Gaussian mean-field setting is the one treated in \citet[Introduction]{auffingerchen2017ground} and \citet{talagrand2006parisi}. Their theorems do not automatically solve the short-range model introduced in the first line.}
\GSfield{Present status}{As of 12 November 2026, the precise statement above is established; the qualifications distinguish it from the abbreviated notebook formulation.}
\GSfield{Assessment}{literal statement: correct with missing hypotheses; publication match: exact; verification: read the result; development: annotated; confidence: high}
\end{GSitem}
\begin{GSitem}{1986-708.2}{B}{The thermodynamic ground-state energy of the SK model}
\GSfield{Kind; source}{problem; \GSpage{P001368}}
\GSfield{The notes}{The notes ask to study \(N\to\infty\) and \(\lim\min H\), without the extensive normalization.}
\GSfield{Statement}{For the Gaussian SK Hamiltonian just specified, \(N^{-1}\max_\sigma H_N(\sigma)\) has a deterministic limit, and the minimum has the negative limit by symmetry. For covariance function \(\xi(q)=q^2/2\), the positive ground-energy constant equals \[\inf_{\gamma}\left\{\Psi_\gamma(0,0)-\frac12\int_0^1 t\xi^{\prime\prime}(t)\gamma(t)\,dt\right\},\] where \(\gamma\ge0\) is nondecreasing, right-continuous and integrable, and \(\Psi_t=-\tfrac12\xi^{\prime\prime}(\Psi_{xx}+\gamma\Psi_x^2)\), \(\Psi(1,x)=|x|\), in the appropriate weak solution sense.}
\GSfield{Fate}{\citet[Theorem~1]{auffingerchen2017ground} proves the zero-temperature Parisi formula, first circulated in June 2016 and published in November 2017. It gives a variational characterization rather than an elementary closed numerical constant. The finite-temperature Parisi formula was proved earlier by \citet{talagrand2006parisi}, published in 2006.}
\GSfield{Present status}{As of 12 November 2026, the normalized mean-field ground-energy problem has a rigorous Parisi variational solution. The omitted factor 1/N is essential; this conclusion does not cover arbitrary finite-range spin glasses.}
\GSfield{Assessment}{literal statement: correct with missing hypotheses; publication match: exact; verification: read the result; development: annotated; confidence: high; problem \hyperref[prob:GP-1986-65]{GP-1986-65}}
\end{GSitem}
\begin{GSitem}{1986-708.3}{E}{Exact finite-volume ground-state degeneracy and near-ground states}
\GSfield{Kind; source}{claim; \GSpage{P001368}}
\GSfield{The notes}{The notes describe many ground states with energies differing by about \(1/N\), and assert positive limiting entropy of their set.}
\GSfield{Statement}{For finite \(N\ge2\), independent continuous Gaussian couplings in the zero-field complete-graph SK model give almost surely exactly one minimizing pair \(\{\sigma,-\sigma\}\). Consequently the entropy per spin of exact minimizers is \(N^{-1}\log2\to0\). States with slightly different energies are near-ground states; their count depends on a specified energy window and normalization, which the notebook does not provide.}
\GSfield{Fate}{\ours{} For \(\tau\ne\pm\sigma\), some product \(\tau_i\tau_j\) differs from \(\sigma_i\sigma_j\), so the energy difference is a nondegenerate Gaussian and vanishes with probability zero. There are finitely many pairs; spin-flip symmetry supplies the only forced equality. Independently, equality of every pair product forces \(\tau_i/\sigma_i\) to be constant in \(i\), checking the algebraic step. This elementary annotation corrects the literal sentence without replacing the much richer asymptotic-state problem treated in \citet{panchenko2013ultrametricity}.}
\GSfield{Present status}{As of 12 November 2026, the finite-volume claim is false under the stated Gaussian SK assumptions. The notebook\textquoteright s unspecified near-ground energy window is insufficient to define a separate quantitative entropy problem.}
\GSfield{Assessment}{literal statement: false as written; publication match: partial; verification: read the result; development: annotated; confidence: high}
\end{GSitem}
\begin{GSitem}{1986-708.4}{B}{Ultrametricity of the limiting overlap array}
\GSfield{Kind; source}{conjecture; \GSpage{P001368}}
\GSfield{The notes}{For three states, two overlaps are equal and the third is larger; the statement is attributed to Parisi and Virasoro and phrased in terms of ground states.}
\GSfield{Statement}{Let a random probability measure on the unit ball of a separable Hilbert space satisfy the Ghirlanda--Guerra identities for all bounded overlap tests. For independent replicas with overlaps \(R_{ab}\), \(R_{12}\ge\min(R_{13},R_{23})\) almost surely. Permuting the replicas gives the ultrametric triple property. The hypothesis is a statement about the thermodynamic Gibbs or asymptotic overlap measure, not about every finite set of near-minimizers.}
\GSfield{Fate}{\citet[Theorem~1]{panchenko2013ultrametricity}, first circulated in December 2011 and published in 2013, proves this implication without a finite-support assumption on the overlap distribution. Its introduction explains the averaging or small-perturbation mechanisms producing the identities in spin-glass models. These precise hypotheses prevent an unjustified universal finite-ground-state claim.}
\GSfield{Present status}{As of 12 November 2026, the ultrametricity theorem under the Ghirlanda--Guerra identities is established. The notebook\textquoteright s informal state terminology is interpreted through this rigorous limiting framework, not as an exact statement for all finite samples.}
\GSfield{Assessment}{literal statement: correct with missing hypotheses; publication match: partial; verification: read the result; development: annotated; confidence: high; problem \hyperref[prob:GP-1986-66]{GP-1986-66}}
\end{GSitem}
\begin{GSitem}{1986-708.5}{B}{The spin-glass transition and the scope of the Parisi formula}
\GSfield{Kind; source}{remark; \GSpage{P001369}}
\GSfield{The notes}{A final line announces a transition to a spin-glass phase without specifying temperature, field, dimension or order parameter.}
\GSfield{Statement}{In mean-field SK theory the free-energy variational problem distinguishes replica-symmetric and symmetry-broken regimes. A precise transition assertion requires the covariance normalization, external field and criterion for the order parameter. The notebook line alone does not fix a numerical critical temperature or state a finite-dimensional transition theorem.}
\GSfield{Fate}{\citet{talagrand2006parisi} proves the mean-field Parisi formula and makes its thermodynamic content rigorous. This identifies a substantial later development of the topic, but does not justify treating an unspecified phase-transition sentence as a solved universal claim.}
\GSfield{Present status}{As of 12 November 2026, rigorous mean-field free-energy theory is established. No exact finite-range phase-transition question is assigned to this unqualified line.}
\GSfield{Assessment}{literal statement: correct with missing hypotheses; publication match: partial; verification: publisher metadata checked; development: annotated; confidence: medium}
\end{GSitem}
\begin{GSitem}{1986-708.6}{E}{Stochastic quantization attribution}
\GSfield{Kind; source}{remark; \GSpage{P001369}}
\GSfield{The notes}{The notes credit Parisi with inventing stochastic quantization.}
\GSfield{Statement}{The named construction introduces an auxiliary stochastic time and a Langevin process whose equilibrium formally reproduces the Euclidean field-theory measure. The original proposal is jointly associated with Parisi and Wu; defining a process and proving equilibrium for a singular field theory requires additional analysis.}
\GSfield{Fate}{This is the established proposal of \citet{parisiwu1981perturbation}, and recurs from the matrix discussion in occurrence 1982-565.2. It is a historical remark, not an assertion that every interacting field theory has already been constructed by this method.}
\GSfield{Present status}{As of 12 November 2026, the precise statement above is established; the qualifications distinguish it from the abbreviated notebook formulation.}
\GSfield{Assessment}{literal statement: correct with missing hypotheses; publication match: partial; verification: indirect; development: annotated; confidence: high}
\end{GSitem}

\GStalk{707}{1987-04-13}{A. S. Schwarz}{P001367-P001369}{Super Riemann surfaces and string measures}
The Schwarz heading occurs on P001369 after Sinai\textquoteright s discussion and continues into P001371. The talk treats superconformal distributions, moduli, periods, super Selberg formulas and determinants in string theory. Statements about uniformization, free modules and determinant powers need qualifications; the catalogue record number, date and interval are retained, with the manuscript order identified explicitly.

\begin{GSitem}{1986-707.1}{E}{Super Riemann surfaces as nonintegrable odd distributions}
\GSfield{Kind; source}{definition; \GSpage{P001369}}
\GSfield{The notes}{The local odd derivative is \(D=\partial_\theta+\theta\partial_z\); superconformal coordinate changes preserve it up to multiplication by a function.}
\GSfield{Statement}{A complex super Riemann surface is a complex supermanifold of dimension \(1|1\) with a rank \(0|1\) subbundle \(\mathcal D\subset T\) such that the bracket induces an isomorphism \(\mathcal D\otimes\mathcal D\to T/\mathcal D\). Locally \(D^2=\partial_z\), and the rescaling factor for a change of generator is invertible and even.}
\GSfield{Fate}{This definition was established in the pre-1987 superstring literature. \citet[Section~2.1]{witten2012superriemann} gives the invariant definition and its local normal form; \citet[Section~2]{cranerabin1988uniformization} states the coordinate version used by the lecture.}
\GSfield{Present status}{As of 12 November 2026, the precise statement above is established; the qualifications distinguish it from the abbreviated notebook formulation.}
\GSfield{Assessment}{literal statement: correct; publication match: exact; verification: read the result; development: annotated; confidence: high}
\end{GSitem}
\begin{GSitem}{1986-707.2}{E}{The coordinate condition for superconformal maps}
\GSfield{Kind; source}{construction; \GSpage{P001369}}
\GSfield{The notes}{One even and one odd analytic function parametrize the coordinate changes. The ordinary analogy writes \(\widetilde\partial=F\partial\), \(F=\partial\widetilde z/\partial z\).}
\GSfield{Statement}{For a supercoordinate change \((z,\theta)\mapsto(z_1,\theta_1)\), the superconformal condition is \(Dz_1=\theta_1D\theta_1\), with invertible even \(D\theta_1\); then \(D=(D\theta_1)D_1\). In the ordinary analogy \(\partial_{z_1}=(\partial z_1/\partial z)^{-1}\partial_z\). The notebook has inverted the ordinary Jacobian factor.}
\GSfield{Fate}{\citet[Section~2, equations~(2.5)--(2.8)]{cranerabin1988uniformization} gives the exact transformation law and the even/odd parametrization. The missing reciprocal in the ordinary comparison is a faithful notebook error.}
\GSfield{Present status}{As of 12 November 2026, the precise statement above is established; the qualifications distinguish it from the abbreviated notebook formulation.}
\GSfield{Assessment}{literal statement: false as written; publication match: exact; verification: read the result; development: annotated; confidence: high}
\end{GSitem}
\begin{GSitem}{1986-707.3}{E}{The local dimension of supermoduli}
\GSfield{Kind; source}{theorem; \GSpage{P001369}}
\GSfield{The notes}{The corrected dimension in the notebook is \((3g-3|2g-2)\).}
\GSfield{Statement}{For a smooth unpunctured super Riemann surface of genus \(g\ge2\), the supermoduli stack has local dimension \(3g-3|2g-2\). Its reduced data include a spin structure, and finite automorphisms require a stack or orbifold interpretation. Genus 0 and 1 require automorphism terms; punctures alter the dimension.}
\GSfield{Fate}{\citet[Section~2.2, equations~(2.24)--(2.27)]{witten2012superriemann} derives the count from deformation cohomology. \citet[Introduction and Section~3]{cranerabin1988uniformization} distinguishes the already known local index count from its global super-Teichm\"uller construction, published in 1988. The bare dimension assertion is therefore E.}
\GSfield{Present status}{As of 12 November 2026, the precise statement above is established; the qualifications distinguish it from the abbreviated notebook formulation.}
\GSfield{Assessment}{literal statement: correct with missing hypotheses; publication match: exact; verification: read the result; development: annotated; confidence: high}
\end{GSitem}
\begin{GSitem}{1986-707.4}{B}{Uniformization of super Riemann surfaces}
\GSfield{Kind; source}{question; \GSpage{P001369}}
\GSfield{The notes}{Asked for a Riemann mapping analogue, Schwarz answers that a simply connected underlying surface reduces to a standard model.}
\GSfield{Statement}{For the DeWitt-type super Riemann surfaces of the cited theorem, with trivial topology in the nilpotent directions, uniformization uses the supersphere, superplane or upper half-superplane, together with superconformal covering transformations. The topology assumption and supermanifold category matter: simple connectedness of the body alone is not a complete theorem for every broader Rogers supermanifold.}
\GSfield{Reconstruction}{The DeWitt topology and the distinction between individual surfaces and families are supplied from the theorem; the notebook does not specify them.}
\GSfield{Fate}{\citet[Sections~2--5]{cranerabin1988uniformization} proves uniformization and the associated Teichm\"uller description in the stated category, publishing in 1988; its Section~6 explicitly separates a more general Rogers-topology conjecture. This precisely qualifies the short answer rather than asserting a universal category-independent result.}
\GSfield{Present status}{As of 12 November 2026, the precise statement above is established; the qualifications distinguish it from the abbreviated notebook formulation.}
\GSfield{Assessment}{literal statement: correct with missing hypotheses; publication match: partial; verification: read the result; development: annotated; confidence: high}
\end{GSitem}
\begin{GSitem}{1986-707.5}{A}{Riemann--Roch and nonfree cohomology modules}
\GSfield{Kind; source}{theorem; \GSpage{P001369-P001370}}
\GSfield{The notes}{Gelfand asks for the function ring on a compact supercurve. Schwarz warns that kernels of equations over a Grassmann algebra need not be free, and only then discusses Riemann--Roch dimensions.}
\GSfield{Statement}{For a family over an odd base, cohomology is a module over the base algebra and need not be free. Ranks can be used only on loci where the relevant direct images are locally free and commute with base change. The correct general invariant is the cohomological index or determinant/Berezinian of cohomology, with parity retained; a numerical dimension formula does not automatically describe the full function algebra.}
\GSfield{Fate}{The warning in the notebook is substantive and should be preserved. The general supercurve theory is identified by \citet{roslyschwarzvoronov1988geometry}, published in March 1988; its original full text was not recovered, so an exact theorem-number match is not claimed. \citet[Section~3.4]{witten2012superriemann} verifies the cohomological formulation and its qualifications.}
\GSfield{Present status}{As of 12 November 2026, the precise statement above is established; the qualifications distinguish it from the abbreviated notebook formulation.}
\GSfield{Assessment}{literal statement: correct with missing hypotheses; publication match: partial; verification: publisher metadata checked; development: annotated; confidence: medium}
\end{GSitem}
\begin{GSitem}{1986-707.6}{A}{Super Abelian differentials, primitives, and periods}
\GSfield{Kind; source}{construction; \GSpage{P001370}}
\GSfield{The notes}{Sections transforming like \(D\Phi\) are called Abelian differentials. Primitives on the universal cover give periods, followed by Jacobian and Picard constructions.}
\GSfield{Statement}{The superconformal distribution determines the appropriate line bundle of differentials and its Berezinian interpretation. Periods are defined from compatible closed differential representatives and cycles of the reduced surface. A normalized super period matrix has additional spin and nonsingularity conditions and can develop poles when the relevant Dirac operator acquires zero modes; it is not an everywhere regular coordinate on all supermoduli.}
\GSfield{Fate}{\citet{roslyschwarzvoronov1988geometry} is the near-contemporary publication match for this theory. \citet[Section~8, especially the end of~8.3]{witten2012superriemann} supplies a readable exact account of periods and the pole qualification. The 1988 paper is verified at metadata/abstract level; the later technical account was opened and read.}
\GSfield{Present status}{As of 12 November 2026, the precise statement above is established; the qualifications distinguish it from the abbreviated notebook formulation.}
\GSfield{Assessment}{literal statement: correct with missing hypotheses; publication match: partial; verification: read the result; development: annotated; confidence: high}
\end{GSitem}
\begin{GSitem}{1986-707.7}{A}{The super Selberg trace formula}
\GSfield{Kind; source}{theorem; \GSpage{P001370}}
\GSfield{The notes}{The string partition function is related to a superanalogue of Selberg zeta, described as joint work with Manin.}
\GSfield{Statement}{The identifiable result is a Selberg-type trace formula for discrete groups acting on the complex upper half-superplane, relating spectral data of super differential operators to geometric conjugacy-class data. Applications concern determinant factors in fermionic-string measures. The notebook specifies neither the test-function class nor the precise super zeta product, so those formulas are not reconstructed from this sentence.}
\GSfield{Fate}{\citet{baranovmaninfrolovschwarz1987selberg}, published in September 1987, is an exact subject and collaboration match. The publisher abstract explicitly states the trace formula and its measure applications. Verification here is bibliographic and abstract-level; no uninspected detailed trace identity is asserted.}
\GSfield{Present status}{As of 12 November 2026, the announced trace-formula construction is published. The original full proof was not inspected in this run, and the notebook contains too little to compare normalizations.}
\GSfield{Assessment}{literal statement: correct with missing hypotheses; publication match: partial; verification: publisher metadata checked; development: annotated; confidence: high}
\end{GSitem}
\begin{GSitem}{1986-707.8}{E}{Odd frames and the superconformal metric gauge}
\GSfield{Kind; source}{construction; \GSpage{P001370}}
\GSfield{The notes}{An odd field \(\widehat e\) satisfying \(\{\widehat e,\overline{\widehat e}\}=\tau\widehat e+\overline\tau\overline{\widehat e}\) is said to admit local coordinates with \(\widehat e=\lambda D\).}
\GSfield{Statement}{The bracket condition is a local integrability condition for the appropriate complex superconformal distribution. A nondegeneracy condition on its self-bracket and an invertible scale \(\lambda\) are also needed. The resulting frame description is local and carries gauge freedom; a single arbitrary odd vector field on a \(2|2\) supermanifold is not automatically a supermetric.}
\GSfield{Fate}{\citet[Sections~3.5--3.6]{witten2012superriemann} explains the relation between superconformal geometry, supergravity and frame gauges. The local formulation is part of the pre-existing supergravity framework; the notebook suppresses the nondegeneracy and reality conventions.}
\GSfield{Present status}{As of 12 November 2026, the precise statement above is established; the qualifications distinguish it from the abbreviated notebook formulation.}
\GSfield{Assessment}{literal statement: correct with missing hypotheses; publication match: partial; verification: read the result; development: annotated; confidence: high}
\end{GSitem}
\begin{GSitem}{1986-707.9}{A}{Determinant expressions for bosonic and fermionic string measures}
\GSfield{Kind; source}{construction; \GSpage{P001370-P001371}}
\GSfield{The notes}{The notes contrast a scalar determinant power \(-13\) for the bosonic string with \(-5\) for the fermionic string, followed by ghost determinants and a reference to Baranov, Schwarz and Voronov.}
\GSfield{Statement}{The exponents reflect Gaussian matter-field counts in the respective critical dimensions, but a genuine string measure also specifies ghost complexes, fermionic Pfaffians or superdeterminants, zero-mode removal, spin structures, determinant lines and integration over the appropriate moduli. A bare product of unprimed scalar determinants is only schematic. The local quadratic action also needs its pairing and conjugation convention.}
\GSfield{Fate}{\citet{roslyschwarzvoronov1988geometry} is a direct near-contemporary match for the super-Laplacian and measure program. The established bosonic determinant-line background includes \citet{beilinsonmanin1986mumford}. The notebook does not state enough to identify every determinant factor or to infer convergence of the integrated string amplitude.}
\GSfield{Present status}{As of 12 November 2026, the determinant-line and supergeometric constructions are established. The notebook formulas remain schematic and do not constitute a complete measure or a finiteness theorem.}
\GSfield{Assessment}{literal statement: correct with missing hypotheses; publication match: partial; verification: publisher metadata checked; development: annotated; confidence: medium}
\end{GSitem}
\begin{GSitem}{1986-707.10}{F}{A proposed noncommutative Ray--Singer torsion}
\GSfield{Kind; source}{remark; \GSpage{P001371}}
\GSfield{The notes}{The final proposal is a noncommutative analogue of Ray--Singer torsion, explicitly accompanied by ``no promise yet.}
\GSfield{Statement}{No coefficient algebra, trace, elliptic complex, determinant convention or invariance statement is supplied. Several inequivalent noncommutative torsion theories could fit the phrase; none is singled out by the record.}
\GSfield{Fate}{The sentence is preserved as a proposal without assigning it to an arbitrary later torsion construction. The nearby string determinant theory is thematic context, not an exact solution.}
\GSfield{Present status}{As of 12 November 2026, this particular proposal is not precisely identifiable from the notebook. It is not classified D merely because the exact referent was not found.}
\GSfield{Assessment}{literal statement: not determinable; publication match: none found; verification: unresolved lead; development: none; confidence: low}
\end{GSitem}

\GStalk{709}{1987-04-13}{D. V. Alekseevskii}{P001370-P001375}{Quaternionic structures and integrability}
The Alekseevskii heading on P001371 introduces the failure of a naive quaternionic differentiability condition, quaternionic projective changes of coordinates, and the invariant rank-three bundle of almost complex structures. The discussion distinguishes torsion-free quaternionic geometry from a chosen hypercomplex triple and includes a twistor-style integrability proposal.

\begin{GSitem}{1986-709.1}{E}{The rigidity of naive quaternionic differentiability}
\GSfield{Kind; source}{claim; \GSpage{P001371}}
\GSfield{The notes}{Transition maps have differentials that are morphisms of left quaternionic modules. The notes conclude that all automorphisms are linear and that simple connectedness gives a linear space.}
\GSfield{Statement}{For a connected open set, sufficiently smooth maps with everywhere quaternionic-linear differential are affine in fixed quaternionic coordinates. Translations must be allowed. An affine atlas on a simply connected manifold supplies a developing map into \(\mathbb H^n\), not automatically a global identification with all of \(\mathbb H^n\); completeness is an additional global requirement.}
\GSfield{Fate}{\citet[Section~3, Theorem~1]{sudbery1979quaternionic} gives the precise one-variable affine-rigidity result. \citet[Introduction and Section~1]{salamon1986quaternionic} explains why such affine structures are too restrictive for quaternionic geometry. The notebook\textquoteright s local point is right, but ``linear omits translations and the global simply-connected conclusion omits completeness.}
\GSfield{Present status}{As of 12 November 2026, the precise statement above is established; the qualifications distinguish it from the abbreviated notebook formulation.}
\GSfield{Assessment}{literal statement: false as written; publication match: partial; verification: read the result; development: annotated; confidence: high}
\end{GSitem}
\begin{GSitem}{1986-709.2}{E}{Quaternionic projective charts and inversion}
\GSfield{Kind; source}{construction; \GSpage{P001371-P001372}}
\GSfield{The notes}{The space of quaternionic lines is labelled \(\mathbb CP^n\); the two charts on its projective line are related by \(q\mapsto q^{-1}\), with derivative \(-q^{-1}(dq)q^{-1}\).}
\GSfield{Statement}{The space is \(\mathbb HP^n\). On \(\mathbb HP^1\), inversion is a transition map and its differential has multiplication on both sides. It generally fails to be a morphism of one fixed left \(\mathbb H\)-module structure. Its differential does preserve the appropriate quaternionic structure up to rotation, which motivates the enlarged structure group.}
\GSfield{Fate}{\citet[Section~1, projective-coordinate construction]{salamon1986quaternionic} gives exactly this motivation and the group \(GL(n,\mathbb H)Sp(1)\). The manuscript\textquoteright s use of \(\mathbb C\) is a faithful notation error, not a translation issue.}
\GSfield{Present status}{As of 12 November 2026, the precise statement above is established; the qualifications distinguish it from the abbreviated notebook formulation.}
\GSfield{Assessment}{literal statement: false as written; publication match: exact; verification: read the result; development: annotated; confidence: high}
\end{GSitem}
\begin{GSitem}{1986-709.3}{F}{The ambiguous sphere obstruction}
\GSfield{Kind; source}{remark; \GSpage{P001372}}
\GSfield{The notes}{The topological explanation says that \(S^3\) would have three complex structures.}
\GSfield{Statement}{An almost complex structure requires even real tangent dimension, so the stated \(S^3\) cannot carry even one. The surrounding projective-line discussion suggests \(S^4=\mathbb HP^1\), but replacing 3 by 4 would be an editorial conjecture, and the intended obstruction to a global quaternionic frame would still need to be stated.}
\GSfield{Fate}{The projective geometry in \citet{salamon1986quaternionic} supplies context but does not license silently changing the handwritten superscript. The recorded sentence is mathematically impossible as it stands; its intended topological proof is not identifiable.}
\GSfield{Present status}{As of 12 November 2026, the literal dimensional error is clear. The intended sphere and obstruction remain an unresolved reading/reconstruction issue, not an open mathematical problem.}
\GSfield{Assessment}{literal statement: false as written; publication match: thematic only; verification: unresolved lead; development: none; confidence: medium}
\end{GSitem}
\begin{GSitem}{1986-709.4}{E}{Almost quaternionic structures without a distinguished triple}
\GSfield{Kind; source}{definition; \GSpage{P001372}}
\GSfield{The notes}{A quaternionic subalgebra of \(\operatorname{End}(V)\) is specified, but its isomorphism with \(\mathbb H\) and the individual \(I,J,K\) are not fixed.}
\GSfield{Statement}{On a real \(4n\)-manifold, an almost quaternionic structure is equivalently a rank-three subbundle \(Q\subset\operatorname{End}(TM)\) locally spanned by \(I,J,K\) satisfying \(I^2=J^2=K^2=-1\), \(IJ=K=-JI\), together with the associated rank-four algebra \(\mathbb R\operatorname{Id}\oplus Q\). Changing an oriented orthonormal local triple rotates it by \(SO(3)\). A global chosen triple is extra hypercomplex-type data.}
\GSfield{Fate}{\citet[Section~1]{salamon1986quaternionic} formulates the same invariant reduction and distinguishes it from an integrable \(GL(n,\mathbb H)\)-structure. The notebook\textquoteright s explicit warning not to choose \(I,J,K\) globally is essential.}
\GSfield{Present status}{As of 12 November 2026, the precise statement above is established; the qualifications distinguish it from the abbreviated notebook formulation.}
\GSfield{Assessment}{literal statement: correct; publication match: exact; verification: read the result; development: annotated; confidence: high}
\end{GSitem}
\begin{GSitem}{1986-709.5}{E}{The general and metric quaternionic structure groups}
\GSfield{Kind; source}{theorem; \GSpage{P001372}}
\GSfield{The notes}{The notes display both \(GL(1,\mathbb H)GL(n,\mathbb H)\) and \(Sp(1)Sp(n)\), while attaching the latter to the earlier module-linearity definition.}
\GSfield{Statement}{The stabilizer of the quaternionic subalgebra is \(GL(n,\mathbb H)Sp(1)\), with the common central \(\{\pm1\}\) identified; positive scalar factors can be absorbed into \(GL(n,\mathbb H)\). After a compatible positive definite metric is chosen, the group reduces to \(Sp(n)Sp(1)\). Preserving a fixed left module structure alone gives \(GL(n,\mathbb H)\), not the compact metric group.}
\GSfield{Fate}{\citet[Section~1]{salamon1986quaternionic} explicitly identifies these groups and their relation. The compact-group line in the notebook has a missing metric condition and an incorrect reference to definition (1).}
\GSfield{Present status}{As of 12 November 2026, the precise statement above is established; the qualifications distinguish it from the abbreviated notebook formulation.}
\GSfield{Assessment}{literal statement: false as written; publication match: exact; verification: read the result; development: annotated; confidence: high}
\end{GSitem}
\begin{GSitem}{1986-709.6}{E}{The special role of four dimensions}
\GSfield{Kind; source}{construction; \GSpage{P001372}}
\GSfield{The notes}{Every four-dimensional Riemannian manifold is said to have two quaternionic structures because \(SO(4)\) is a product.}
\GSfield{Statement}{For an oriented Riemannian four-manifold, the self-dual and anti-self-dual rank-three bundles of two-forms give the two pointwise quaternionic families. The group relation is \(SO(4)\cong(Sp(1)\times Sp(1))/\{(1,1),(-1,-1)\}\). Reversing orientation exchanges the choices. A quaternionic structure in dimension four encodes an oriented conformal class, not a unique metric.}
\GSfield{Fate}{\citet[Section~1 and the four-dimensional discussion]{salamon1986quaternionic} gives the exceptional relation \(Sp(1)Sp(1)=SO(4)\). The quotient and orientation data must accompany the notebook\textquoteright s shorthand ``product.}
\GSfield{Present status}{As of 12 November 2026, the precise statement above is established; the qualifications distinguish it from the abbreviated notebook formulation.}
\GSfield{Assessment}{literal statement: correct with missing hypotheses; publication match: partial; verification: read the result; development: annotated; confidence: high}
\end{GSitem}
\begin{GSitem}{1986-709.7}{E}{Torsion-free quaternionic connections and twistor integrability}
\GSfield{Kind; source}{definition; \GSpage{P001373}}
\GSfield{The notes}{The condition for 1-integrability is a torsion-free connection preserving the quaternionic subalgebra; parallel transport preserves it, and the connection need not be metric.}
\GSfield{Statement}{For \(n>1\), a quaternionic manifold is an almost quaternionic \(4n\)-manifold admitting a torsion-free connection preserving \(Q\). This is weaker than a flat projective quaternionic atlas and does not choose a metric. Its associated sphere bundle of complex structures has an integrable twistor complex structure. In dimension four the analogous twistor integrability additionally requires the appropriate self-duality condition on the conformal curvature.}
\GSfield{Fate}{\citet[Sections~2--3 and Corollary~7.4]{salamon1986quaternionic} establishes the torsion/curvature framework and the twistor construction, published before the lecture in 1986. The notebook\textquoteright s explicit ``not necessarily metric is correct and prevents a conflation with quaternionic K"ahler geometry.}
\GSfield{Present status}{As of 12 November 2026, the precise statement above is established; the qualifications distinguish it from the abbreviated notebook formulation.}
\GSfield{Assessment}{literal statement: correct with missing hypotheses; publication match: partial; verification: read the result; development: annotated; confidence: high}
\end{GSitem}
\begin{GSitem}{1986-709.8}{F}{Goncharov on complexified quaternionic cones}
\GSfield{Kind; source}{remark; \GSpage{P001373}; A. B. Goncharov}
\GSfield{The notes}{For real-analytic \(M\), Goncharov forms a cone from vectors \(v+iJv\) in a complexification and says that there must be integrability.}
\GSfield{Statement}{The family of complex eigenspaces as \(J\) ranges over the quaternionic sphere is a natural twistor-type object. A precise theorem would specify whether the proposed condition is involutivity of a distribution on an incidence space, existence of complex submanifolds, or vanishing of intrinsic torsion. The notebook gives no such condition or equivalence.}
\GSfield{Fate}{\citet[twistor construction, Corollary~7.4]{salamon1986quaternionic} is a close mathematical antecedent. It does not identify the exact extra assertion Goncharov intended; the item remains a partially articulated proposal rather than a newly formulated general theorem.}
\GSfield{Present status}{As of 12 November 2026, the cone construction is recognizable, but the intended integrability criterion is not sufficiently specified for an exact fate assignment.}
\GSfield{Assessment}{literal statement: not determinable; publication match: thematic only; verification: unresolved lead; development: none; confidence: medium}
\end{GSitem}

\GStalk{710}{1987-04-20}{A. A. Kirillov}{P001373-P001377}{Hecke algebras, braids, and link invariants}
After Gelfand\textquoteright s incompletely specified question about commuting partially defined isometries, Kirillov explains Hecke algebras, braid groups and Markov traces. The notes conflate algebraic traces, positive traces and factor representations, and the observation about roots of unity should be read against Jones\textquoteright s pre-existing index theorem.

\begin{GSitem}{1986-710.1}{F}{Commuting partially defined isometries and unitary completion}
\GSfield{Kind; source}{question; \GSpage{P001374}; I. M. Gelfand}
\GSfield{The notes}{Gelfand asks how two commuting isometries in an n-dimensional space can be turned into commuting unitaries, adding explicitly that they are not defined everywhere; Kasparov is said to have done part of it.}
\GSfield{Statement}{A precise extension problem must state each initial domain and range, the meaning of commutation on partially defined compositions, whether the ambient space may grow, and whether the goal is an extension or a compression dilation. In finite dimension an everywhere-defined isometry is already unitary. The handwritten domain qualification is therefore decisive.}
\GSfield{Fate}{The notebook is not an exact formulation of the theorem on two commuting everywhere-defined contractions and cannot be replaced by it. Searches combining Kasparov, commuting partial isometries and unitary completion did not identify the referenced partial result with its domains. The question is preserved without inventing a missing formulation.}
\GSfield{Present status}{As of 12 November 2026, the exact problem and cited Kasparov contribution remain unidentified. This is a limitation of the source data, not a claim that all such extension problems are open.}
\GSfield{Assessment}{literal statement: not determinable; publication match: none found; verification: unresolved lead; development: none; confidence: medium}
\end{GSitem}
\begin{GSitem}{1986-710.2}{E}{Type-A Hecke algebras and braid-group representations}
\GSfield{Kind; source}{definition; \GSpage{P001374-P001375}}
\GSfield{The notes}{The generators satisfy \((T_i+1)(T_i-q)=0\), the adjacent braid relation and distant commutativity; inverse braid generators are sent to \((T_i+1-q)/q\).}
\GSfield{Statement}{For \(q\ne0\), the type-A Hecke algebra \(H_n(q)\) is generated by \(T_1,\ldots,T_{n-1}\) with these relations. They make \(T_i\) invertible, with \(T_i^{-1}=q^{-1}(T_i+1-q)\), and define a representation of the Artin braid group. The algebras form a compatible tower under the usual inclusions.}
\GSfield{Fate}{\citet[Sections~3--4]{jones1987hecke} uses exactly this presentation. The theory and its link-invariant application were already circulating by the June 1985 Bourbaki lecture \citet{connes1986indice}; the later 1987 article is the detailed source, not the first appearance.}
\GSfield{Present status}{As of 12 November 2026, the precise statement above is established; the qualifications distinguish it from the abbreviated notebook formulation.}
\GSfield{Assessment}{literal statement: correct with missing hypotheses; publication match: exact; verification: read the result; development: annotated; confidence: high}
\end{GSitem}
\begin{GSitem}{1986-710.3}{F}{The unnamed Leningrad type-B extension}
\GSfield{Kind; source}{remark; \GSpage{P001374}}
\GSfield{The notes}{The notes say that a type-B Hecke algebra had likewise been applied in Leningrad.}
\GSfield{Statement}{Type-B Hecke and braid constructions have additional generators, relations and parameters. The sentence gives neither authors nor the target links, boundary conditions or trace, so it cannot identify one exact extension of the type-A invariant.}
\GSfield{Fate}{\citet[Section~5, concluding discussion]{jones1987hecke} explicitly distinguishes other Coxeter types as further directions. This thematic connection does not establish which Leningrad result the notebook refers to.}
\GSfield{Present status}{As of 12 November 2026, the intended local attribution remains unidentified. No new open type-B problem is assigned to this passing report.}
\GSfield{Assessment}{literal statement: not determinable; publication match: thematic only; verification: unresolved lead; development: none; confidence: low}
\end{GSitem}
\begin{GSitem}{1986-710.4}{E}{Links as braid closures}
\GSfield{Kind; source}{construction; \GSpage{P001375}}
\GSfield{The notes}{A link is called an immersed compact oriented one-manifold in three-space, and closing a braid gives a link.}
\GSfield{Statement}{A tame oriented link is an embedded finite disjoint union of circles in oriented three-space, considered up to ambient isotopy. Closing an oriented braid gives such a link, and Alexander\textquoteright s theorem represents every oriented link by a braid closure. General immersions with double points are singular links and require different equivalence data.}
\GSfield{Fate}{\citet[Section~2]{jones1987hecke} states the braid-closure setting. The manuscript\textquoteright s ``immersed must be strengthened to ``embedded for the knot invariant discussed here.}
\GSfield{Present status}{As of 12 November 2026, the precise statement above is established; the qualifications distinguish it from the abbreviated notebook formulation.}
\GSfield{Assessment}{literal statement: false as written; publication match: exact; verification: read the result; development: annotated; confidence: high}
\end{GSitem}
\begin{GSitem}{1986-710.5}{E}{Markov equivalence of braid representatives}
\GSfield{Kind; source}{theorem; \GSpage{P001375}}
\GSfield{The notes}{Conjugation inside \(B_n\) and stabilization \(b\mapsto b\sigma_n^{\pm1}\) are said to generate equivalence of closures.}
\GSfield{Statement}{Two braids have ambient-isotopic oriented closures exactly when they are related by a finite sequence of conjugations, positive or negative stabilizations, and inverse destabilizations, allowing the braid index to change. Therefore a compatible trace alone handles conjugation, while stabilization needs an additional normalization.}
\GSfield{Fate}{\citet[Section~2]{jones1987hecke} invokes Markov\textquoteright s theorem as the topological foundation of the trace construction. This was established before the lecture; the notebook correctly identifies the difficult converse, not just the easy invariance of individual moves.}
\GSfield{Present status}{As of 12 November 2026, the precise statement above is established; the qualifications distinguish it from the abbreviated notebook formulation.}
\GSfield{Assessment}{literal statement: correct; publication match: exact; verification: read the result; development: annotated; confidence: high}
\end{GSitem}
\begin{GSitem}{1986-710.6}{E}{The algebraic Ocneanu trace has a second parameter}
\GSfield{Kind; source}{theorem; \GSpage{P001375}}
\GSfield{The notes}{For every \(q\), the notebook asserts a unique factor representation with trace, whose values are Laurent polynomials in \(q\).}
\GSfield{Statement}{For fixed Hecke parameter \(q\), the compatible algebraic trace is characterized only after a second parameter \(z\) is supplied: \(\operatorname{tr}(1)=1\), \(\operatorname{tr}(ab)=\operatorname{tr}(ba)\), and \(\operatorname{tr}_{n+1}(aT_n)=z\operatorname{tr}_n(a)\). This is an algebraic Markov trace, not automatically a positive state or a unique von Neumann factor representation. Values on braid words also involve the inverse parameter \(q^{-1}\).}
\GSfield{Fate}{\citet[Theorem~5.1]{jones1987hecke} proves precisely the trace theorem and explicitly makes \(z\) independent of \(q\). The earlier exposition \citet{connes1986indice} explains the relation to subfactors without asserting the notebook\textquoteright s universal factor uniqueness.}
\GSfield{Present status}{As of 12 November 2026, the precise statement above is established; the qualifications distinguish it from the abbreviated notebook formulation.}
\GSfield{Assessment}{literal statement: false as written; publication match: exact; verification: read the result; development: annotated; confidence: high}
\end{GSitem}
\begin{GSitem}{1986-710.7}{E}{From Markov traces to two-variable link polynomials}
\GSfield{Kind; source}{construction; \GSpage{P001375-P001376}}
\GSfield{The notes}{The desired function of braids is invariant under both Markov moves and becomes a Laurent-polynomial knot invariant.}
\GSfield{Statement}{The normalized trace expression includes factors depending on braid index and exponent sum so that both stabilization signs act trivially. The general Hecke trace gives the two-variable HOMFLY-type invariant; the Jones polynomial is a particular specialization/quotient. Equality of these invariants is a necessary condition for isotopy, not a complete classifier of links.}
\GSfield{Fate}{\citet[Section~6]{jones1987hecke} gives the normalized construction and skein relation, tracing the two-variable discovery to the 1985 work of Freyd, Yetter, Hoste, Lickorish, Millett, Ocneanu and Przytycki--Traczyk. \citet{connes1986indice} shows that this was established before Kirillov\textquoteright s lecture.}
\GSfield{Present status}{As of 12 November 2026, the precise statement above is established; the qualifications distinguish it from the abbreviated notebook formulation.}
\GSfield{Assessment}{literal statement: correct with missing hypotheses; publication match: exact; verification: read the result; development: annotated; confidence: high}
\end{GSitem}
\begin{GSitem}{1986-710.8}{E}{Specialization at plus and minus one}
\GSfield{Kind; source}{question; \GSpage{P001376}}
\GSfield{The notes}{The exercise asks for the invariant obtained by setting \(q=\pm1\), mentioning the ratio of Jones-polynomial values.}
\GSfield{Statement}{After choosing the usual Jones variable \(t\) and normalization \(V_{\text{unknot}}(t)=1\), a knot satisfies \(V_K(1)=1\) and \(|V_K(-1)|=|\Delta_K(-1)|\), its determinant. The Hecke parameter \(q\), trace parameter and link-polynomial variable must first be related; without that choice the raw exercise does not determine one numerical specialization.}
\GSfield{Reconstruction}{The Jones normalization and the distinction between its variable and the Hecke parameter are supplied editorially.}
\GSfield{Fate}{\citet[Section~12, special values]{jones1987hecke} gives the specialization results. They are established identities rather than an open question. The statement here is an explicit modern normalization of the exercise, not an assertion that the notebook\textquoteright s variables were already fixed this way.}
\GSfield{Present status}{As of 12 November 2026, the precise statement above is established; the qualifications distinguish it from the abbreviated notebook formulation.}
\GSfield{Assessment}{literal statement: correct with missing hypotheses; publication match: partial; verification: read the result; development: annotated; confidence: high}
\end{GSitem}
\begin{GSitem}{1986-710.9}{E}{Roots of unity and the existence of subfactor representations}
\GSfield{Kind; source}{claim; \GSpage{P001376}}
\GSfield{The notes}{The notes say Jones excludes roots of unity; Gelfand suggests studying precisely them and wonders whether factors disappear there. Vershik is mentioned without a theorem.}
\GSfield{Statement}{Generic semisimplicity and generic representation formulas may exclude roots of unity, but the Hecke algebra and Markov trace do not thereby disappear. Roots of unity are central to the discrete subfactor-index regime \(4\cos^2(\pi/m)\), \(m\ge3\), and yield important quotient representations. Positivity and factor realizability require specified involutions, parameters and quotients.}
\GSfield{Fate}{\citet[Section~4 after Theorem~4.5 and later root-of-unity sections]{jones1987hecke} explicitly says these parameters are important and repeatedly studies them. \citet{connes1986indice} already relates the discrete indices to the braid construction. The broad exclusion attributed to Jones is false; the unnamed Vershik contribution is not uniquely identified.}
\GSfield{Present status}{As of 12 November 2026, roots-of-unity representation and knot theories are established. The notebook\textquoteright s suggestion is historically interesting, but it is not an undeveloped general direction.}
\GSfield{Assessment}{literal statement: false as written; publication match: partial; verification: read the result; development: annotated; confidence: high}
\end{GSitem}

\GStalk{711}{1987-04-20}{A. A. Beilinson}{P001375-P001379}{Holomorphic Chern classes and their integral comparison}
The explicit Beilinson heading on P001376 starts with Chern classes and their Hodge and integral structures, then formulates the comparison problem using transition functions and logarithms. This is the introduction to the following week\textquoteright s Deligne-cohomology lecture, not a second copy of that lecture.

\begin{GSitem}{1986-711.1}{E}{Integral Chern classes and the holomorphic Hodge filtration}
\GSfield{Kind; source}{theorem; \GSpage{P001376}}
\GSfield{The notes}{For a holomorphic bundle on a compact complex analytic manifold, Chern classes are integral and lie in the indicated Hodge-filtration term; cohomology is computed by the holomorphic de Rham complex.}
\GSfield{Statement}{For a complex manifold \(X\) and holomorphic bundle \(E\), \(c_i(E)\in H^{2i}(X,\mathbb Z)\) has complex image in \(F^iH^{2i}(X,\mathbb C)\). Here \(F^i\) means the image of hypercohomology of the truncated holomorphic de Rham complex, not automatically an injectively embedded hypercohomology group. A Hodge direct-sum decomposition requires additional hypotheses, such as compact K\"ahler.}
\GSfield{Fate}{\citet[Sections~1--2 and~8]{esnaultviehweg1987deligne} gives the Deligne-cohomology setting and the compatible Chern classes; the MPI 87-2 preprint is explicitly cited in the following lecture. The underlying integral and filtration properties are established antecedents, not new conclusions of the 1987 announced work.}
\GSfield{Present status}{As of 12 November 2026, the precise statement above is established; the qualifications distinguish it from the abbreviated notebook formulation.}
\GSfield{Assessment}{literal statement: correct with missing hypotheses; publication match: exact; verification: read the result; development: annotated; confidence: high}
\end{GSitem}
\begin{GSitem}{1986-711.2}{E}{Atiyah classes, Chern characters, and invariant polynomials}
\GSfield{Kind; source}{construction; \GSpage{P001376-P001377}}
\GSfield{The notes}{The transition form \(A=g_{ij}^{-1}dg_{ij}\) is treated as an ordinary matrix-valued Cech cocycle, and \(c_i(E)=\operatorname{Tr}(A^i)\) is written.}
\GSfield{Statement}{The transition forms represent the Atiyah class in \(H^1(X,\Omega_X^1\otimes\operatorname{End}E)\), with the correct conjugation identifications between trivializations. After the chosen \(2\pi i\) and sign normalization, \(\operatorname{tr}(A^i)/i!\) represents the degree-i Chern character, not the Chern class. Chern classes are the elementary invariant polynomials, for example \(c_2=\tfrac12((\operatorname{tr}a)^2-\operatorname{tr}(a^2))\) for normalized \(a\).}
\GSfield{Fate}{The distinction between universal Chern classes and trace-polynomial representatives is part of the characteristic-class construction in \citet[Section~8]{esnaultviehweg1987deligne}. The notebook\textquoteright s trace-power equality is false for general \(i>1\); fixing it also requires the twisted endomorphism coefficients, not only a factorial.}
\GSfield{Present status}{As of 12 November 2026, the precise statement above is established; the qualifications distinguish it from the abbreviated notebook formulation.}
\GSfield{Assessment}{literal statement: false as written; publication match: partial; verification: read the result; development: annotated; confidence: high}
\end{GSitem}
\begin{GSitem}{1986-711.3}{A}{An explicit cochain comparison with integral characteristic classes}
\GSfield{Kind; source}{problem; \GSpage{P001377}}
\GSfield{The notes}{The problem is to construct explicitly a homotopy between the differential-form characteristic cocycle and an integral characteristic cocycle. The lecture announces Beilinson, MacPherson and Schechtman\textquoteright s Notes on motivic cohomology.}
\GSfield{Statement}{The comparison asks for cochains witnessing equality of the complex image of an integral Chern cocycle and the chosen filtered de Rham representative. A Deligne Chern class packages precisely the integral cocycle, filtered representative and comparison homotopy. An explicit low-weight polylogarithmic formula is a stronger construction than the abstract existence of this lift.}
\GSfield{Fate}{The announced work was published as \citet{beilinsonmacphersonschechtman1987motivic} in 1987, but its original proof text was not recovered here, so a general formula-by-formula match is not asserted. \citet[Proposition~8.2]{esnaultviehweg1987deligne} gives existence and functorial uniqueness, while \citet[Section~1.7 and Section~4]{goncharov2016exponential} provides an inspected explicit second-Chern construction. The final lecture records its weight-two ingredients.}
\GSfield{Present status}{As of 12 November 2026, Deligne Chern classes and explicit low-weight comparison cocycles are established. The source-specific match to every higher-weight formula of the announced 1987 paper remains limited by access; no general unsolved problem is inferred from that limit.}
\GSfield{Assessment}{literal statement: correct with missing hypotheses; publication match: partial; verification: indirect; development: annotated; confidence: medium; problem \hyperref[prob:GP-1986-67]{GP-1986-67}}
\end{GSitem}
\begin{GSitem}{1986-711.4}{E}{The first Chern class from logarithms of transition functions}
\GSfield{Kind; source}{construction; \GSpage{P001377}}
\GSfield{The notes}{For \(c_1\), the notes use \(d\log\det g_{ij}\) and propose dropping the exterior derivative to recover an integral discrepancy.}
\GSfield{Statement}{On a sufficiently fine cover, choose holomorphic logarithms \(\ell_{ij}\) of the determinant transitions. On triple overlaps, \(\ell_{ij}+\ell_{jk}+\ell_{ki}\in2\pi i\mathbb Z\), with signs adjusted to the chosen Cech convention. Dividing that discrepancy by \(2\pi i\) gives the integral first Chern cocycle, and \(d\ell_{ij}\) its de Rham comparison. Changing logarithm branches changes the cocycle by a coboundary.}
\GSfield{Fate}{\citet[weight-one complex and Chern classes]{esnaultviehweg1987deligne} identifies this construction with the exponential sequence. It is the classical complete model for the higher comparison question, and recurs in record 712.}
\GSfield{Present status}{As of 12 November 2026, the precise statement above is established; the qualifications distinguish it from the abbreviated notebook formulation.}
\GSfield{Assessment}{literal statement: correct with missing hypotheses; publication match: exact; verification: read the result; development: annotated; confidence: high}
\end{GSitem}

\GStalk{712}{1987-04-27}{A. A. Beilinson}{P001377-P001384}{Deligne cohomology, motivic complexes, and the dilogarithm}
The lecture defines the Deligne complex, its products and characteristic classes, explains the intermediate-Jacobian exact sequence, and sketches motivic refinements and the weight-two dilogarithmic complex. It cites the MPI 87-2 preprint explicitly. The last cross-ratio construction is best matched to the low-weight motivic characteristic-class program; its abbreviation is not a proof of the general motivic conjectures.

\begin{GSitem}{1986-712.1}{E}{The lecture bibliography and its chronology}
\GSfield{Kind; source}{remark; \GSpage{P001378}}
\GSfield{The notes}{The notes cite Soul\'e\textquoteright s February 1985 Bourbaki talk, Esnault--Viehweg\textquoteright s preprint, Beilinson\textquoteright s Current Problems volume 24 paper, and a tentative Contemporary Mathematics volume 25.}
\GSfield{Statement}{The directly identified sources are Soul\'e\textquoteright s Bourbaki expos\'e 644, delivered in February 1985 and printed in Ast\'erisque 133--134 (1986), and Esnault--Viehweg, MPI 87-2. Beilinson\textquoteright s Notes on absolute Hodge cohomology is in Contemporary Mathematics 55, Part~I (1986), not 25. These references show that much of the Deligne-complex exposition is pre-existing material.}
\GSfield{Fate}{\citet{soule1986regulateurs} and \citet{esnaultviehweg1987deligne} were opened as original texts. The latter\textquoteright s bibliography confirms the absolute-Hodge reference. The specific 1987 announced motivic construction should be separated from these background definitions and theorems.}
\GSfield{Present status}{As of 12 November 2026, the precise statement above is established; the qualifications distinguish it from the abbreviated notebook formulation.}
\GSfield{Assessment}{literal statement: correct with missing hypotheses; publication match: exact; verification: read the result; development: annotated; confidence: high}
\end{GSitem}
\begin{GSitem}{1986-712.2}{E}{The analytic Deligne complex and its degree convention}
\GSfield{Kind; source}{definition; \GSpage{P001378}}
\GSfield{The notes}{The complex is \(\mathbb Z\to\mathcal O\to\Omega^1\to\cdots\to\Omega^{p-1}\), with the first arrow multiplying by \((2\pi i)^p\).}
\GSfield{Statement}{For a complex manifold \(X\) and integer \(p\ge1\), let \(\mathbb Z(p)=(2\pi i)^p\mathbb Z\subset\mathbb C\), and put \(\mathbb Z(p)_{\mathcal D}=[\mathbb Z(p)\to\mathcal O_X\to\Omega_X^1\to\cdots\to\Omega_X^{p-1}]\) in cohomological degrees \(0,\ldots,p\). Its hypercohomology is \(H^n_{\mathcal D}(X,\mathbb Z(p))\). Weight zero is the constant sheaf \(\mathbb Z\). The algebraic open-variety version requires the corresponding logarithmic construction.}
\GSfield{Fate}{\citet[Section~1.1 and Section~2]{esnaultviehweg1987deligne} fixes these conventions and distinguishes analytic from algebraic Deligne--Beilinson cohomology. The notebook\textquoteright s ordinary \(H\) notation must be read as hypercohomology when the argument is a complex.}
\GSfield{Present status}{As of 12 November 2026, the precise statement above is established; the qualifications distinguish it from the abbreviated notebook formulation.}
\GSfield{Assessment}{literal statement: correct with missing hypotheses; publication match: exact; verification: read the result; development: annotated; confidence: high}
\end{GSitem}
\begin{GSitem}{1986-712.3}{E}{Weight-one and weight-two cohomology sheaves}
\GSfield{Kind; source}{theorem; \GSpage{P001378}}
\GSfield{The notes}{The exponential map identifies the weight-one complex with \(\mathcal O^*[-1]\); for weight two the notes give \(\mathbb C/\mathbb Z(2)\) and \(\Omega^1/d\mathcal O\).}
\GSfield{Statement}{The exponential sequence gives \(\mathbb Z(1)_{\mathcal D}\simeq\mathcal O_X^*[-1]\). In weight two, the cohomology sheaves are \(\mathcal H^1=\mathbb C/\mathbb Z(2)\) and \(\mathcal H^2=\Omega_X^1/d\mathcal O_X\). They are sheaves, not the global Deligne groups of \(X\); the latter require hypercohomology. Locally the holomorphic Poincar\'e lemma identifies the second quotient with the image of \(d:\Omega^1\to\Omega^2\).}
\GSfield{Fate}{\citet[Section~1.1 and the weight-one comparison in Section~2]{esnaultviehweg1987deligne} verifies the exponential quasi-isomorphism. The weight-two assertions are the cohomology sheaves of the displayed three-term complex, with the notebook\textquoteright s Tate factor retained.}
\GSfield{Present status}{As of 12 November 2026, the precise statement above is established; the qualifications distinguish it from the abbreviated notebook formulation.}
\GSfield{Assessment}{literal statement: correct with missing hypotheses; publication match: exact; verification: read the result; development: annotated; confidence: high}
\end{GSitem}
\begin{GSitem}{1986-712.4}{E}{Topological and curvature maps from the Deligne complex}
\GSfield{Kind; source}{construction; \GSpage{P001379}}
\GSfield{The notes}{One map projects to the integer component; a second takes the exterior derivative of the final differential form.}
\GSfield{Statement}{There is a forgetful morphism \(\mathbb Z(p)_{\mathcal D}\to\mathbb Z(p)\) and a curvature morphism \(\mathbb Z(p)_{\mathcal D}\to\Omega_{X,\mathrm{cl}}^p[-p]\), zero except on the last component, where it is \(d\). Calling the first target \(\mathbb Z\) is harmless only after identifying the Tate lattice with the abstract integers; comparison in complex cohomology must retain \((2\pi i)^p\).}
\GSfield{Fate}{\citet[Sections~1--3]{esnaultviehweg1987deligne} constructs the corresponding projections and compatible products. These are pre-existing structural maps, and the last-component derivative automatically takes values in closed forms.}
\GSfield{Present status}{As of 12 November 2026, the precise statement above is established; the qualifications distinguish it from the abbreviated notebook formulation.}
\GSfield{Assessment}{literal statement: correct with missing hypotheses; publication match: exact; verification: read the result; development: annotated; confidence: high}
\end{GSitem}
\begin{GSitem}{1986-712.5}{E}{A chain-level Deligne product and commutativity up to homotopy}
\GSfield{Kind; source}{construction; \GSpage{P001379}}
\GSfield{The notes}{The notes give ordinary multiplication when either factor is an integer, then \(x\wedge dy\) when the right form has top degree, and claim commutativity up to homotopy.}
\GSfield{Statement}{For positive weights \(p,q\), one standard asymmetric product of the truncated complexes is \(x\cup y=xy\) when \(\deg x=0\); it is \(x\wedge dy\) when \(\deg x>0\) and \(\deg y=q\); it is zero otherwise. In particular a positive-degree left factor multiplied by a degree-zero right factor is zero in this model. This product is a chain map; its opposite is chain-homotopic to it, inducing graded commutativity on hypercohomology.}
\GSfield{Fate}{\citet[Section~1.1 and Section~3, especially Proposition~3.5 and Theorem~3.9]{esnaultviehweg1987deligne} gives product conventions and their homotopies. The notebook combines the asymmetric top-form rule with two-sided integer multiplication and thereby violates the chain-map convention; one cannot keep all three literal rules simultaneously.}
\GSfield{Present status}{As of 12 November 2026, the precise statement above is established; the qualifications distinguish it from the abbreviated notebook formulation.}
\GSfield{Assessment}{literal statement: false as written; publication match: exact; verification: read the result; development: annotated; confidence: high}
\end{GSitem}
\begin{GSitem}{1986-712.6}{E}{Functorial Deligne Chern classes and the meaning of uniqueness}
\GSfield{Kind; source}{theorem; \GSpage{P001379}}
\GSfield{The notes}{Holomorphic bundles have classes \(c_p(E)\in H^{2p}_{\mathcal D}(X,\mathbb Z(p))\), determined by functoriality and their ordinary Chern-class images.}
\GSfield{Statement}{There is a universal natural system of Deligne Chern classes compatible with pullback, the topological Chern classes and the usual characteristic-class identities. The uniqueness is a universal functorial statement. On a fixed \(X\), specifying a topological image alone does not select a unique Deligne class because the kernel can be a nonzero intermediate Jacobian.}
\GSfield{Fate}{\citet[Proposition~8.2 and its proof using the universal bundle]{esnaultviehweg1987deligne} proves the stated existence and uniqueness. The preprint was available as MPI 87-2 and is explicitly named at the lecture\textquoteright s start, so this is E rather than a new 1987 conjecture.}
\GSfield{Present status}{As of 12 November 2026, the precise statement above is established; the qualifications distinguish it from the abbreviated notebook formulation.}
\GSfield{Assessment}{literal statement: correct with missing hypotheses; publication match: exact; verification: read the result; development: annotated; confidence: high}
\end{GSitem}
\begin{GSitem}{1986-712.7}{E}{Cech hypercohomology and the signed total differential}
\GSfield{Kind; source}{construction; \GSpage{P001379-P001380}}
\GSfield{The notes}{The double complex is totalized using \(d+\delta\), with component relations linking an integral cocycle and forms on multiple overlaps.}
\GSfield{Statement}{For Cech degree \(r\), a standard total differential is \(D=\delta+(-1)^r d\), unless one has already signed one of the two component differentials. On an acyclic cover or a suitable resolution, total cohomology computes the hypercohomology of the Deligne complex. The component equations must include the same signs and the initial Tate-lattice inclusion.}
\GSfield{Fate}{\citet[Sections~1--2]{esnaultviehweg1987deligne} uses the complex-of-sheaves framework. The notebook\textquoteright s abbreviated unsigned formula is acceptable only with an unstated anticommuting convention; with the ordinary commuting Cech and de Rham differentials it is incorrect.}
\GSfield{Present status}{As of 12 November 2026, the precise statement above is established; the qualifications distinguish it from the abbreviated notebook formulation.}
\GSfield{Assessment}{literal statement: correct with missing hypotheses; publication match: exact; verification: read the result; development: annotated; confidence: high}
\end{GSitem}
\begin{GSitem}{1986-712.8}{E}{The Deligne exact sequence and the intermediate Jacobian}
\GSfield{Kind; source}{theorem; \GSpage{P001380-P001381}}
\GSfield{The notes}{For compact \(X\), Hodge decomposition is asserted and Deligne cohomology is an extension of integral Hodge classes by a quotient called the Griffiths Jacobian.}
\GSfield{Statement}{For compact K\"ahler \(X\) and \(p\ge1\), put \[J^p(X)=\frac{H^{2p-1}(X,\mathbb C)}{F^pH^{2p-1}(X,\mathbb C)+\operatorname{im}H^{2p-1}(X,\mathbb Z(p))}.\] There is an exact sequence \[0\to J^p(X)\to H^{2p}_{\mathcal D}(X,\mathbb Z(p))\to\{c\in H^{2p}(X,\mathbb Z(p)):c_{\mathbb C}\in F^p\}\to0.\] The left group is the intermediate Jacobian. The right description retains torsion; a literal intersection inside complex cohomology would lose it.}
\GSfield{Fate}{\citet[Section~7 and the basic exact sequences]{esnaultviehweg1987deligne} gives the Jacobian interpretation. Compactness alone does not guarantee the Hodge decomposition used in the notebook, and all complex comparisons need the Tate lattice.}
\GSfield{Present status}{As of 12 November 2026, the precise statement above is established; the qualifications distinguish it from the abbreviated notebook formulation.}
\GSfield{Assessment}{literal statement: false as written; publication match: exact; verification: read the result; development: annotated; confidence: high}
\end{GSitem}
\begin{GSitem}{1986-712.9}{E}{A Deligne cocycle as an explicit comparison homotopy}
\GSfield{Kind; source}{method; \GSpage{P001380-P001381}}
\GSfield{The notes}{The notes describe an integral cocycle, a filtered de Rham cocycle and a homotopy between them, then write a differential of a sum beginning with \(\alpha_0\), although the preceding homotopy sum begins with \(\alpha_1\).}
\GSfield{Statement}{The intermediate components form a total cochain of degree one less than the compared cocycles. With consistent total-complex signs, its differential is their difference. The degree-n integral cocycle \(\alpha_0\) is one endpoint and cannot be included as an ordinary degree-\(n-1\) homotopy term. This is the mapping-cone interpretation of Deligne cohomology, not an equality between unsorted forms of different total degrees.}
\GSfield{Fate}{\citet[cone construction in Section~2]{esnaultviehweg1987deligne} provides the precise framework. The notebook\textquoteright s reference to Gabrielov--Gelfand--Losik points to earlier combinatorial characteristic-class comparisons; it is an antecedent, not an identification of every Deligne homotopy with one of their formulas.}
\GSfield{Present status}{As of 12 November 2026, the precise statement above is established; the qualifications distinguish it from the abbreviated notebook formulation.}
\GSfield{Assessment}{literal statement: false as written; publication match: partial; verification: read the result; development: annotated; confidence: high}
\end{GSitem}
\begin{GSitem}{1986-712.10}{E}{Picard group and the Deligne first Chern class}
\GSfield{Kind; source}{theorem; \GSpage{P001381}}
\GSfield{The notes}{The determinant bundle controls \(c_1\); logarithms of transition functions give the cocycle, but the final formula reads \(\log(\deg g_{ij})\).}
\GSfield{Statement}{For a complex manifold, \(H^2_{\mathcal D}(X,\mathbb Z(1))\cong H^1(X,\mathcal O_X^*)=\operatorname{Pic}(X)\). Under this identification a line bundle gives its Deligne first Chern class; for a higher-rank bundle it is the class of \(\det E\). The logarithm is \(\log\det g_{ij}\), and its Cech discrepancy lies in \(2\pi i\mathbb Z\).}
\GSfield{Fate}{\citet[weight-one exponential comparison and Section~8]{esnaultviehweg1987deligne} gives the result. The manuscript\textquoteright s \(\deg\) is a faithful slip; it is not a meaningful degree of a variable transition matrix in this formula.}
\GSfield{Present status}{As of 12 November 2026, the precise statement above is established; the qualifications distinguish it from the abbreviated notebook formulation.}
\GSfield{Assessment}{literal statement: false as written; publication match: exact; verification: read the result; development: annotated; confidence: high}
\end{GSitem}
\begin{GSitem}{1986-712.11}{B}{Motivic complexes and the regulator refinement}
\GSfield{Kind; source}{construction; \GSpage{P001382}}
\GSfield{The notes}{Chern classes should first take values in \(\mathbb Z(p)_M\) and then map through a polylogarithmic regulator to Deligne cohomology. Weight zero is \(\mathbb Z\), weight one is \(\mathcal O^*[-1]\), and negative effective weights are set to zero.}
\GSfield{Statement}{For smooth algebraic varieties over a field, motivic cohomology and its complexes now have rigorous constructions, including cycle and transfer models. In weight one, \(\mathbb Z(1)\simeq\mathbb G_m[-1]\); motivic Chern classes admit realization/regulator maps to the corresponding complex-analytic cohomology. The convention of zero negative weights belongs to the effective complex formulation, not to the assertion that negative Tate objects vanish in a stable motivic category.}
\GSfield{Reconstruction}{The switch from arbitrary analytic manifolds to smooth algebraic varieties, the topology of the motivic sheaves and the effective/stable distinction are necessary modern specifications.}
\GSfield{Fate}{\citet[Lectures~3--4,~15 and~19]{mazzavoevodskyweibel2006motivic} develops the established theory from Voevodsky\textquoteright s 1999--2000 lectures, published in 2006. The explicit low-weight regulator is inspected in \citet[Proposition~1.7]{goncharov2016exponential}. Existence of this theory does not prove every vanishing property proposed in the next notebook line.}
\GSfield{Present status}{As of 12 November 2026, the precise statement above is established; the qualifications distinguish it from the abbreviated notebook formulation.}
\GSfield{Assessment}{literal statement: correct with missing hypotheses; publication match: partial; verification: read the result; development: annotated; confidence: high}
\end{GSitem}
\begin{GSitem}{1986-712.12}{C}{Beilinson--Soule vanishing in nonpositive degrees}
\GSfield{Kind; source}{conjecture; \GSpage{P001382}}
\GSfield{The notes}{For positive weight \(p\), the proposed motivic complex has cohomology only in degrees \([1,p]\).}
\GSfield{Statement}{The unresolved lower-degree assertion includes the Beilinson--Soul\'e vanishing conjecture: for every field \(F\), \(p>0\), and \(j\le0\), \(H^j(F,\mathbb Q(p))=0\). The notebook\textquoteright s integral sheaf-amplitude formulation is stronger in its coefficient and sheaf details. Upper-degree bounds and the existence of motivic complexes must be separated from this nonpositive-degree vanishing.}
\GSfield{Fate}{\citet[Section~4 and Theorem~4.1]{garkusha2023semilocal} states the amplitude conjecture and proves equivalent criteria, not the general conjecture (2020 preprint, 2022 online publication, 2023 journal issue). \citet{soule2026interview} explicitly describes the conjecture as still open in an interview dated 4 November 2026. Known arithmetic cases do not establish it for all fields.}
\GSfield{Present status}{Open in general as of 12 November 2026, with direct near-cutoff confirmation from Soul\'e. The construction of motivic cohomology itself is established and must not be reported as solving this vanishing conjecture.}
\GSfield{Assessment}{literal statement: correct with missing hypotheses; publication match: exact; verification: read the result; development: annotated; confidence: high; problem \hyperref[prob:GP-1986-68]{GP-1986-68}}
\GSfield{Worth developing}{Worth preserving with its exact coefficient and degree scope: it is a definite surviving part of the notebook program, but no bounded proof of the general vanishing assertion is attempted.}
\end{GSitem}
\begin{GSitem}{1986-712.13}{B}{Milnor K-theory and the top motivic cohomology sheaf}
\GSfield{Kind; source}{definition; \GSpage{P001382-P001383}}
\GSfield{The notes}{The notes generate the Milnor algebra by units with the Steinberg relation \(\{a,1-a\}=0\), mention additivity in a product, and identify the top weight-p sheaf with Milnor K-theory.}
\GSfield{Statement}{For a field \(F\), \(K_*^M(F)\) is the tensor algebra over \(\mathbb Z\) on the abelian group \(F^*\), modulo the two-sided ideal generated by \(a\otimes(1-a)\), \(a\ne0,1\). Symbols are multiplicatively additive in each entry and the resulting algebra is graded-commutative. For smooth varieties the appropriate unramified sheaf is used; the field comparison is \(H^p(F,\mathbb Z(p))\cong K_p^M(F)\). Graded commutativity does not mean that every odd square vanishes integrally.}
\GSfield{Fate}{\citet[Lecture~5, Theorem~5.1]{mazzavoevodskyweibel2006motivic} proves the field comparison and fixes the Milnor symbols. The tensor presentation already includes additivity in units; a free algebra on the underlying set of units without those relations would be the wrong object.}
\GSfield{Present status}{As of 12 November 2026, the precise statement above is established; the qualifications distinguish it from the abbreviated notebook formulation.}
\GSfield{Assessment}{literal statement: correct with missing hypotheses; publication match: partial; verification: read the result; development: annotated; confidence: high}
\end{GSitem}
\begin{GSitem}{1986-712.14}{E}{The weight-two Bloch complex and the five-term relation}
\GSfield{Kind; source}{construction; \GSpage{P001383}}
\GSfield{The notes}{The degree-one group is generated by functions avoiding \(0,1,\infty\), modulo the alternating cross-ratio relation from five points; the differential is \(t\mapsto t\wedge(1-t)\).}
\GSfield{Statement}{Over a field, let the pre-Bloch group be generated by \([z]\), \(z\ne0,1\), modulo the five-term cross-ratio relations. The map \([z]\mapsto z\wedge(1-z)\) kills those relations and defines a two-term weight-two complex. The opposite sign convention \((1-z)\wedge z\) is equally valid if used consistently. Integral comparisons require careful antisymmetric-tensor and 2-torsion conventions; the exterior-square model is safest after inverting 2 or rationalizing.}
\GSfield{Fate}{\citet[Section~1.5]{goncharov2016exponential} constructs the sheafified Bloch complex and its regulator. \citet[Introduction and Theorem~5.2]{suslin1991bloch} explains the exact integral K-theory relation and its torsion terms. The notebook itself already warns about 2-torsion, which should not be dropped.}
\GSfield{Present status}{As of 12 November 2026, the precise statement above is established; the qualifications distinguish it from the abbreviated notebook formulation.}
\GSfield{Assessment}{literal statement: correct with missing hypotheses; publication match: partial; verification: read the result; development: annotated; confidence: high}
\end{GSitem}
\begin{GSitem}{1986-712.15}{E}{Suslin and the indecomposable K3 group}
\GSfield{Kind; source}{theorem; \GSpage{P001383}}
\GSfield{The notes}{Suslin is credited with connecting the Bloch complex to K-theory, and the construction is called Bloch--Suslin.}
\GSfield{Statement}{For an infinite field \(F\), the Bloch group obtained as the kernel of the weight-two symbol boundary is closely related to \(K_3^{\mathrm{ind}}(F)=\operatorname{coker}(K_3^M(F)\to K_3(F))\). After tensoring with \(\mathbb Q\), Suslin\textquoteright s comparison identifies these groups. Integrally there is a roots-of-unity torsion term and a convention-sensitive modification; one must not identify the whole pre-Bloch group or all of \(K_3(F)\) with the kernel.}
\GSfield{Reconstruction}{The rational comparison is stated without the convention-sensitive integral extension. The classification records prior circulation, rather than inferring a post-1987 discovery from the printed date.}
\GSfield{Fate}{\citet[Introduction, p.~217, and Theorem~5.2, p.~237]{suslin1991bloch} is the precise later publication: Russian original 1990, English issue 1991 (translation copyright 1992). Suslin explicitly says that most results of the paper were obtained in 1982. The April 1987 notebook independently witnesses the prepublication circulation of the Bloch--K-theory connection. It is therefore treated as exposition of an already circulating result, while the introduction does not date every refinement of the final theorem separately. The theorem was read directly and specifies an infinite field, the indecomposable quotient and the modified roots-of-unity torsion term.}
\GSfield{Present status}{As of 12 November 2026, the precise field comparison is established. The notebook omits the distinction among pre-Bloch, Bloch and indecomposable K-theory groups, which the statement restores.}
\GSfield{Assessment}{literal statement: correct with missing hypotheses; publication match: partial; verification: read the result; development: annotated; confidence: high}
\end{GSitem}
\begin{GSitem}{1986-712.16}{B}{The dilogarithmic regulator map}
\GSfield{Kind; source}{construction; \GSpage{P001383-P001384}}
\GSfield{The notes}{The final diagram maps the weight-two motivic complex to the Deligne complex using logarithms and the dilogarithm.}
\GSfield{Statement}{A chain-level regulator must include logarithmic correction terms and branch data so that the differential of the dilogarithmic term agrees with the symbol boundary and the five-term relation is respected in the target complex. A single uncorrected branch of \(\operatorname{Li}_2(z)\) is not a globally single-valued map to the Deligne complex. The real Bloch--Wigner version gives a single-valued real regulator after the appropriate real realization.}
\GSfield{Fate}{\citet[Proposition~1.7 and its preceding formulas]{goncharov2016exponential} supplies an explicit weight-two morphism of complexes. The classical five-term and single-valued real formulas are treated in \citet{zagier2007dilogarithm}. Earlier seminar occurrences 1974-342 and 1979-477 concern dilogarithms/regulators but are not the same full motivic-complex conjecture.}
\GSfield{Present status}{As of 12 November 2026, the precise statement above is established; the qualifications distinguish it from the abbreviated notebook formulation.}
\GSfield{Assessment}{literal statement: correct with missing hypotheses; publication match: partial; verification: read the result; development: annotated; confidence: high}
\end{GSitem}
\begin{GSitem}{1986-712.17}{A}{A cross-ratio cocycle for the second Chern class of a rank-two bundle}
\GSfield{Kind; source}{construction; \GSpage{P001384}}
\GSfield{The notes}{For a rank-two bundle and local sections \(s_i\), the notes place one component on triple overlaps in units wedged with units and another on quadruple overlaps in the Bloch group, obtained from four projective points.}
\GSfield{Statement}{Choose local nonzero sections in sufficiently general position so that all required pairwise determinants and cross-ratios are defined and invertible on the refined intersections. A triple of sections supplies an exterior product of determinant functions; a quadruple supplies a Bloch symbol of their cross-ratio. The Pl\"ucker identity identifies the symbol boundary with the Cech boundary of the triple component, and the five-term relation supplies the next cocycle identity. Degenerate configurations require refinement or the appropriate resolution, not evaluation at forbidden cross-ratios.}
\GSfield{Reconstruction}{The general-position, nonvanishing and refinement conditions are supplied from the explicit cocycle construction.}
\GSfield{Fate}{\citet[Section~1.7, equations~(1.34)--(1.44), and Section~4]{goncharov2016exponential} provides an inspected explicit construction matching the notebook\textquoteright s last sketch, including the regulator to Deligne cohomology. The lecture announces the earlier \citet{beilinsonmacphersonschechtman1987motivic}; its exact original formula was not inspected, so the later source is not presented as first discovery.}
\GSfield{Present status}{As of 12 November 2026, the rank-two weight-two construction is explicit and established. This is a completion of the recorded sketch, not a claim that a polylogarithmic formula in every weight follows from it.}
\GSfield{Assessment}{literal statement: correct with missing hypotheses; publication match: partial; verification: read the result; development: annotated; confidence: medium; problem \hyperref[prob:GP-1986-67]{GP-1986-67}}
\end{GSitem}

\section{Developments}\label{dev:1986}
\noindent\textbf{Scope and credit.}
The first two sections develop bounded parts of the class C problems
\hyperref[prob:GP-1986-33]{GP-1986-33} and
\hyperref[prob:GP-1986-60]{GP-1986-60}. The last two are verification
annotations to questions already answered in the literature or notebook.
The markers \ours\ and \oursComputed\ identify proofs and computations
supplied here; they do not assert historical novelty. No finite calculation
below is promoted to an all-level spectral theorem or a large-\(N\) result.

\subsection{Quantum KdV: a level-two calculation and a Jordan obstruction}
\label{dev:kdv}

\paragraph{The question and its scope.}
At \GSpage{P001318} the notebook asks to quantize the KdV integrals
and then compute the spectrum in irreducible Virasoro representations.
Reading the latter as the spectrum of the quantized integrals, rather
than of \(L_0\) alone, is an editorial reconstruction based on the
preceding clause. The notes specify no central charge, highest weight,
positivity, or completeness criterion. A precise modern target fixes
\(c,h\in\mathbb C\), the irreducible highest-weight module \(L(c,h)\),
charge normalizations, and a grade \(n\); one seeks the joint generalized
eigenspaces and nilpotent actions, including multiplicities. An ordinary
joint eigenbasis requires an additional semisimplicity assertion.

The construction of commuting charges and their spectral classification
are different problems.
\citet[eq.~(11)]{bazhanovlukyanovzamolodchikov1996integrable}
gives the normalization used below (December 1994 preprint; 1996
publication). The opened May 2026 version of
\citet[Definitions~6.5--6.6 and Theorem~6.8]{masoeroruzza2026odeim}
proves a completeness statement at \(c=-2\),
\(h=(\beta^2-1)/8\), \(\beta\notin\mathbb Z\), for the specified
admissible QQ solutions modulo their entire-function gauge.
This is not a theorem for every \(c,h\). The calculation below
does not depend on that preprint.

\paragraph{Conventions.}
Let \(M(c,h)\) be the Verma module generated by \(v\), where
\[
 L_mv=0\ (m>0),\quad L_0v=hv,\quad
 [L_m,L_n]=(m-n)L_{m+n}
       +\frac c{12}(m^3-m)\delta_{m+n,0}.
\]
Its grade-two PBW basis is \(v_1=L_{-2}v,\ v_2=L_{-1}^2v\). Set
\[
 I_3=2\sum_{n\geq1}L_{-n}L_n+
       L_0^2-\frac{c+2}{12}L_0+\frac{c(5c+22)}{2880}.
\]
The sum is finite on every vector of the algebraic module.

\begin{proposition}[The level-two charge]\ours\ 
For
\[
 q_2=(h+2)^2-\frac{(c+2)(h+2)}{12}+\frac{c(5c+22)}{2880},
\]
the matrix of \(I_3\) on \(M(c,h)[2]\) is
\[
 q_2\,1+\begin{pmatrix}8h+c&12h\\6&8h+4\end{pmatrix}.
\]
Its characteristic roots, with algebraic multiplicity, are
\[
 q_2+8h+\frac{c+4}{2}
       \ \pm\ \frac12\sqrt{(c-4)^2+288h}.
\]
It is diagonalizable on this two-dimensional subspace exactly when
\((c-4)^2+288h\ne0\).
\end{proposition}
\begin{proof}
The commutators give
\[
\begin{aligned}
 L_1v_1&=3L_{-1}v,&L_1v_2&=(4h+2)L_{-1}v,\\
 L_2v_1&=(4h+c/2)v,&L_2v_2&=6hv.
\end{aligned}
\]
For example,
\(L_1L_{-1}^2v=2L_0L_{-1}v+2L_{-1}L_0v
=(4h+2)L_{-1}v\).
All \(L_n\), \(n>2\), annihilate this grade. Multiplication by
\(2L_{-1}\) and \(2L_{-2}\) gives the matrix.
Its nonscalar part \(B\) has trace \(16h+c+4\) and discriminant
\((\operatorname{tr}B)^2-4\det B=(c-4)^2+288h\).
If this vanishes, the repeated-root matrix is not scalar, since its
lower-left entry is \(6\); hence it is not diagonalizable.
\end{proof}

\begin{proposition}[A Jordan block in an irreducible quotient]\ours\ 
For \(c=1,\ h=-1/32\), grade two survives unchanged in \(L(c,h)\).
On that grade,
\[
 I_3=\frac{28893}{5120}\,1+
       \begin{pmatrix}-3/2&-3/8\\6&3/2\end{pmatrix}.
\]
The second matrix \(N\) is nonzero and satisfies \(N^2=0\).
Thus there is no ordinary simultaneous eigenbasis for the local charges
on this irreducible module.
\end{proposition}
\begin{proof}
Substitution gives the displayed matrix and \(N^2=0\).
Every submodule is graded: it is invariant under the diagonal operator
\(L_0\), and each vector has only finitely many weight components.
A proper submodule meeting grade two would therefore have a nonzero
singular vector in grade one or two. In grade one the condition is
\(2h=0\), which fails. In grade two, annihilation by \(L_1,L_2\)
requires
\[
 \begin{pmatrix}3&4h+2\\4h+c/2&6h\end{pmatrix}\binom ab=0.
\]
At the stated parameters its determinant is \(-81/64\ne0\).
There is consequently no such singular vector, and no proper
submodule meets grade two. The Jordan block persists in the
irreducible quotient. A simultaneous eigenbasis would diagonalize
\(I_3\), a contradiction.
\end{proof}

\paragraph{Interpretation and checks.}
This example is nonunitary: the highest-weight contravariant form
is negative on \(L_{-1}v\). It does not challenge finite-grade
diagonalizability with a positive-definite unitary inner product.
It also does not refute a notebook conjecture: the notebook asks for
a spectrum, without asserting semisimplicity.

The independent program \texttt{code/root/check\_\allowbreak{}quantum\_\allowbreak{}kdv.py}
constructs the action by recursive PBW rewriting, without entering
the proposed matrix. A separate Gram calculation gives
\[
 G=\begin{pmatrix}4h+c/2&6h\\6h&4h(2h+1)\end{pmatrix},
 \quad
 \det G=2h\{16h^2+(2c-10)h+c\}.
\]
It checks \(GI_3=I_3^{\mathsf T}G\); at the exceptional parameters
\(\det G=-81/1024\). The program also checks Gram symmetry in
grades one, three, and four at two other rational parameter choices.
All exact checks passed; the output is supplied.

\paragraph{Remaining obstruction.}
This calculation does not organize arbitrary grades, settle every
singular-vector quotient, or prove ODE/IM completeness. The reviewed
sources establish particular constructions and regimes, without
documenting a reason that the general problem was abandoned.
A further bounded run could construct the charge algebra in grades
three through six and separate its Jordan loci from degeneration
of the Verma form. It would give a finite classification, not an
all-grade theorem.

\subsection{The multi-source unitary integral: evaluation, error control, and a failed extension}
\label{dev:unitary}

\paragraph{Supported and reconstructed formulations.}
On \GSpage{P001330}, Kazakov asks for the integral with two diagonal
matrix pairs, writing identical pairs first. We use probability
Haar measure, \(\int_{U(N)}dU=1\), and write
\[
 Z_N(P,Q)=\int_{U(N)}
 \exp\!\left(\sum_{\mu=1}^r
   \operatorname{Tr}(U^\ast P^\mu UQ^\mu)\right)dU,\qquad
 P^\mu=\operatorname{diag}(p_i^\mu),\quad
 Q^\mu=\operatorname{diag}(q_j^\mu).
\]
The notebook supports \(r=2\); arbitrary finite \(r\), complex entries,
and the explicit Haar normalization are editorial extensions.
Put \(a_{ij}=\sum_\mu p_i^\mu q_j^\mu\). The exponent becomes
\(X_A(U)=\sum_{ij}a_{ij}|U_{ij}|^2\).
The diagonal entries are fixed parameters, so no contour for them is
needed.

\paragraph{Literature and the remaining target.}
The \(N=2\) reduction and the failure of an immediate localization
extension to \(N>2\) are treated in
\citet[Sections~2.2--2.3 and Appendix~A]{holguin2024multimatrix}
(July 2023 preprint; 2024 publication). Their displayed angular measure
and our probability Haar measure must be compared after fixing the
overall normalization. All-rank polynomial Haar moments are given by
\citet[Corollary~2.4]{collinssniady2006integration}
(February 2004 preprint; 2006 publication).
Thus an assertion that the integral has never been evaluated would
be wrong if evaluation includes a convergent explicit series.
A meaningful remaining target is a tractable formula or asymptotic
description for general multi-source data in a stated regime.

\begin{proposition}[The normalized \(U(2)\) integral]\ours\ 
For \(A\in M_2(\mathbb C)\), put
\(\alpha=a_{11}+a_{22}\), \(\beta=a_{12}+a_{21}\).
Then
\[
 Z_2(A)=
 \begin{cases}
 (e^\alpha-e^\beta)/(\alpha-\beta),&\alpha\ne\beta,\\
 e^\alpha,&\alpha=\beta.
 \end{cases}
\]
These are the values of one entire function. Equivalently,
\(Z_2=e^s\operatorname{sinhc}(z)\), where
\(s=(\alpha+\beta)/2\), \(z=(\alpha-\beta)/2\),
\(\operatorname{sinhc}(z)=\sinh(z)/z\) for \(z\ne0\), and
\(\operatorname{sinhc}(0)=1\).
\end{proposition}
\begin{proof}
For Haar \(U\in U(2)\), \(t=|U_{11}|^2\) is uniform on \([0,1]\).
Indeed a first column is obtained by normalizing two independent
complex Gaussians. Their squared moduli are independent exponentials.
The change \((x,y)=(rt,r(1-t))\), with Jacobian \(r\), gives
density \(\int_0^\infty re^{-r}\,dr=1\) for \(t\).
Unitarity then yields
\[
 (|U_{ij}|^2)=\begin{pmatrix}t&1-t\\1-t&t\end{pmatrix},
 \qquad Z_2=\int_0^1e^{\beta+(\alpha-\beta)t}\,dt.
\]
Integration gives the formula; uniform convergence of its exponential
series on compact parameter sets gives entireness and the diagonal value.
\end{proof}

\begin{proposition}[All-rank convergent evaluation with an explicit tail]\ours\ 
For \(A\in M_N(\mathbb C)\), set
\(M=\sum_i\max_j|a_{ij}|\). Then
\[
 Z_N(A)=\sum_{m=0}^\infty\frac{\mathbb E X_A^m}{m!},\qquad
 \left|Z_N(A)-\sum_{m=0}^K\frac{\mathbb E X_A^m}{m!}\right|
 \leq\frac{e^M M^{K+1}}{(K+1)!}.
\]
Each moment is an explicit finite sum, even for \(m>N\).
Index a matrix by \(S_m\), and put
\[
 G_{\sigma,\tau}=N^{\#\operatorname{cycles}(\sigma^{-1}\tau)},\qquad
 d_{\boldsymbol i}(\sigma)=
     \prod_{k=1}^m\delta_{i_k,i_{\sigma(k)}}.
\]
Writing \(G^+\) for the Moore--Penrose inverse, for \(m\geq1\) one has
\[
 \mathbb E X_A^m=
 \sum_{\boldsymbol i,\boldsymbol j\in\{1,\ldots,N\}^m}
 \left(\prod_{k=1}^m a_{i_kj_k}\right)
 \sum_{\sigma,\tau\in S_m}
 d_{\boldsymbol i}(\sigma)(G^+)_{\sigma,\tau}
 d_{\boldsymbol j}(\tau).
\]
For rational \(A\), all moments are rational.
\end{proposition}
\begin{proof}
Each row of \((|U_{ij}|^2)\) sums to one, giving \(|X_A|\leq M\).
The exponential series converges absolutely and uniformly on the
compact integration group. Termwise integration and the scalar bound
\[
 \left|e^z-\sum_{m=0}^Kz^m/m!\right|
 \leq e^{|z|}|z|^{K+1}/(K+1)!
\]
prove the series and error bound.

On \(\mathcal H=(\mathbb C^N)^{\otimes m}\), averaging conjugation
by \(U^{\otimes m}\) is the Hilbert--Schmidt orthogonal projection
onto the commutant. Schur--Weyl duality says that this commutant is
spanned by the permutation operators \(P_\sigma\).
Their Gram matrix is
\(\operatorname{Tr}(P_\sigma^\ast P_\tau)=G_{\sigma,\tau}\).
The orthogonal projection onto a possibly dependent family uses its
Gram pseudoinverse. Its diagonal coefficients give
\[
 \mathbb E\prod_{k=1}^m|U_{i_kj_k}|^2
 =\sum_{\sigma,\tau}
 d_{\boldsymbol i}(\sigma)(G^+)_{\sigma,\tau}
 d_{\boldsymbol j}(\tau).
\]
Expansion of \(X_A^m\) proves the stated formula.
The pseudoinverse of the integer matrix \(G\) is rational, by
rational rank factorization and inversion of nonsingular rational
Gram matrices. This proves the rationality assertion. A usual
inverse would be invalid when the permutation operators are
dependent.
\end{proof}

\paragraph{Independent low-degree identities.}\ours\ 
Let \(S=\sum_{ij}a_{ij}\), \(F=\sum_{ij}a_{ij}^2\),
\(R=\sum_i(\sum_j a_{ij})^2\), and
\(C=\sum_j(\sum_i a_{ij})^2\). For \(N\geq2\),
\[
 \mathbb EX_A=S/N,\qquad
 \mathbb EX_A^2=
 \frac{S^2+F}{N^2-1}-\frac{R+C}{N(N^2-1)}.
\]
The squares are algebraic, not absolute squares.
For \(N=2\), independent integration gives
\[
 \mathbb EX_A^m=\frac1{m+1}\sum_{j=0}^m\alpha^j\beta^{m-j},
\]
including coincident endpoints.

\begin{example}[Failure of a specific permutation-only extension]\oursComputed\ 
The following stronger formula is constructed for this audit.
The notebook does not assert that its saddle approximation is exact.
Take the two-component diagonal data to be the rows of
\[
 P=\begin{pmatrix}0&0\\1&0\\0&1\end{pmatrix},\quad
 Q=\begin{pmatrix}0&0\\2&1\\1&3\end{pmatrix},\quad
 A=PQ^{\mathsf T}=
   \begin{pmatrix}0&0&0\\0&2&1\\0&1&3\end{pmatrix}.
\]
For \(\sigma\in S_3\), define
\[
 S_\sigma=\sum_i a_{i,\sigma(i)},\quad
 D_\sigma=\prod_{i<j}
  (a_{i,\sigma(i)}+a_{j,\sigma(j)}
                  -a_{i,\sigma(j)}-a_{j,\sigma(i)}).
\]
The naive extension of the ordinary rank-one fixed-point expression is
\[
 F(t)=\frac2{t^3}\sum_{\sigma\in S_3}
                   \frac{e^{tS_\sigma}}{D_\sigma}.
\]
Its exact data are
\[
\begin{array}{c|rrrrrr}
 \sigma&123&132&213&231&312&321\\ \hline
 S_\sigma&5&2&3&1&1&2\\
 D_\sigma&18&-3&-8&-2&-2&-3
\end{array}
\]
so that
\[
 \sum_\sigma D_\sigma^{-1}=-125/72,\qquad
 F(t)=-\frac{125}{36}t^{-3}+O(t^{-2}).
\]
The true probability-normalized Haar integral is entire and has
\[
 Z_3(tA)=1+\frac73t+\frac{71}{24}t^2+\frac{65}{24}t^3+O(t^4).
\]
Consequently this permutation-only expression cannot be the general
two-source formula.
\end{example}
\begin{proof}
The six substitutions give the table. All denominator factors are
nonzero. Their interpretation is checked independently: rotating
rows \(i,j\) of the permutation matrix by a real Givens rotation
changes the phase to
\[
 S_\sigma-
 (a_{i,\sigma(i)}+a_{j,\sigma(j)}
      -a_{i,\sigma(j)}-a_{j,\sigma(i)})\sin^2\theta.
\]
The code differentiates all eighteen curves to verify the quadratic
factors. Summing the reciprocals gives \(-125/72\), proving the pole.
The genuine integral is entire by the preceding proposition, and
its first moments are computed independently from the Gram formula.
\end{proof}

\paragraph{Verification and practical limits.}
The program \texttt{check\_\allowbreak{}unitary\_\allowbreak{}integral.py} compares the exact
Gram calculation with the \(U(2)\) integral for real, complex, and
coincident endpoints and with the ordinary rank-one HCIZ determinant
at \(N=3\). It tests \(m>N\): the Gram ranks are \(5\) for
\((N,m)=(2,3)\) and \(14\) for \((2,4)\).
All comparisons pass.
For \(A=\left(\begin{smallmatrix}1&2\\3&5\end{smallmatrix}\right)\),
centering at \(11/2\) leaves radius \(1/2\).
At order 20 the computed error is about \(2.26\cdot10^{-27}\),
below the analytic bound \(3.77\cdot10^{-24}\).
These decimals illustrate the proved bound; they are not an
interval-arithmetic certificate.
The separate \texttt{check\_\allowbreak{}stationary\_\allowbreak{}formula.py} checks the
counterexample by both the quadratic variations and exact moments.

The procedure is convergent but factorial in degree. It does not
supply an efficient large-\(N\) algorithm or a compact generic
formula. A further bounded run could exploit permutation symmetries,
compute a prescribed set of \(N=3\) moments, and test a precisely
stated approximation with rigorous tails. The later literature
explicitly revisits the problem, so there is no basis to call it
abandoned.

\subsection{An annotation on simultaneous complex structures}
\label{dev:complex}

\paragraph{Source and later result.}
Gelfand's question at \GSpage{P001276} asks for a complex structure
making all the real codimension-two subspaces complex. The notebook
already warns that even-dimensional intersections are insufficient.
Here a complex structure means a single \emph{linear}
\(J\in\operatorname{End}_{\mathbb R}(\mathbb R^4)\) with \(J^2=-1\).
Nonlinear equivalence of complements is a different question.
\citet[Theorems~1 and~3]{batzies2007planes} gives the four-plane
classification and the transverse normal form (received 2001,
revised 2005, published 2007). The following is an elementary
rederivation, not a new classification.

\begin{proposition}[A finite-family criterion in the transverse case]\ours\ 
Let \(P_1,\ldots,P_k\), \(k\geq3\), be real two-dimensional linear
subspaces of \(\mathbb R^4\), with \(P_i\cap P_j=\{0\}\) for
\(i\ne j\). There are real coordinates with
\[
 P_1=\mathbb R^2\oplus0,\quad P_2=0\oplus\mathbb R^2,\quad
 P_3=\operatorname{graph}(1),\quad
 P_j=\operatorname{graph}(A_j)\ (j\geq4).
\]
A common linear complex structure exists exactly when either every
\(A_j\) is scalar, or a nonscalar \(A\) among them satisfies
\[
 (\operatorname{tr}A)^2-4\det A<0,\qquad
 [A,A_j]=0\quad\hbox{for every }j.
\]
For \(k=4\), the condition is that \(A_4\) is scalar or has negative
discriminant.
\end{proposition}
\begin{proof}
In \(P_1\oplus P_2\), transversality makes \(P_3\) the graph of an
isomorphism \(T:P_1\to P_2\). A coordinate change in \(P_2\)
makes \(T=1\); every other plane is then a graph.
The hypotheses imply
\(\det A_j\ne0\), \(\det(A_j-1)\ne0\), and
\(\det(A_i-A_j)\ne0\) for distinct graph planes.

A linear \(J\) preserving \(P_1,P_2\) is block diagonal.
Preserving \(\operatorname{graph}(1)\) makes its blocks equal:
\(J=\operatorname{diag}(B,B)\). The remaining conditions are
\(B^2=-1\) and \(BA_j=A_jB\).
In a basis \(e,Be\), the centralizer of \(B\) consists exactly of
\(a\,1+bB\). A nonscalar such matrix has discriminant \(-4b^2<0\),
and all such matrices commute. This proves necessity.
Conversely, if every \(A_j\) is scalar, any complex structure \(B\)
works. Otherwise, for a chosen \(A\) with discriminant \(\Delta<0\),
Cayley--Hamilton gives
\[
 B=\frac{2A-(\operatorname{tr}A)1}{\sqrt{-\Delta}},
 \qquad B^2=-1.
\]
Every matrix commuting with \(A\) commutes with \(B\).
Thus \(\operatorname{diag}(B,B)\) has all the required invariances.
\end{proof}

\begin{corollary}[The intersection lattice is insufficient]\ours\ 
Arrangements of four pairwise transverse two-planes in \(\mathbb R^4\)
can have the same intersection lattice, and be joined by a transverse
continuous deformation, while only one admits a common linear
complex structure.
\end{corollary}
\begin{proof}
Keep the first three planes in the normal form and use
\(A_t=\operatorname{diag}(2,2+t)\), \(0\leq t\leq1\).
The determinants \(2(2+t)\) and \(1+t\) of \(A_t\) and \(A_t-1\)
stay positive. All pairwise intersections stay zero and the
intersection lattice stays the rank-two uniform lattice.
At \(t=0\) the matrix is scalar; at \(t>0\) its discriminant is
\(t^2>0\). Apply the proposition.
\end{proof}

\paragraph{Independent check and boundary of the result.}
The code \texttt{check\_\allowbreak{}complex\_\allowbreak{}structures.py} independently solves
the full four-by-four invariance equations. It checks scalar,
nonreal, distinct-real, and Jordan cases and commuting and
noncommuting families. In the last case, the invariant algebra is
only the real scalars, excluding \(J^2=-1\) directly. All checks
agree. This settles the stated transverse four-dimensional linear
subcase; it is not a combinatorial characterization of arbitrary
higher-dimensional or nonlinear configurations. No unspecified
extension is labelled an open conjecture here.

\subsection{An exact audit of the cube's 96 regions}
\label{dev:cube}

\paragraph{The recorded answer.}
The Gelfand--T.~Alekseevskaya question at \GSpage{P001364},
item \hyperref[item:1986-706.3]{1986-706.3}, already has the written
answer \(96\). We count the connected components of the
\emph{interior} of a cube after removing every distinct affine
plane determined by three vertices. This is verification of the
recorded answer, not a new problem.

\begin{proposition}[The cube region count]\oursComputed\ 
For \([-1,1]^3\), the 56 vertex triples determine 20 distinct planes.
Six are faces and do not cut the interior. The other 14 planes cut
the interior into exactly 96 regions.
\end{proposition}
\begin{proof}
The cross-product calculation in
\texttt{check\_\allowbreak{}cube\_\allowbreak{}regions.py} enumerates and deduplicates the planes.
The interior-cutting ones are
\[
 x=\pm y,\quad x=\pm z,\quad y=\pm z,\qquad
 \varepsilon_1x+\varepsilon_2y+\varepsilon_3z=1
     \quad(\varepsilon_i\in\{-1,1\}).
\]
The last eight equations describe distinct planes.
The code computes every triple intersection of the full 20-plane
family in the closed cube and includes its vertices, giving
35 candidate points. Every vertex of every intermediate cell
belongs to this set: a vertex of a full-dimensional bounded
three-dimensional polyhedron lies on at least three independent
bounding planes.

Process the 14 cutting planes. A convex cell splits exactly when
its candidate points have both a strict positive and a strict
negative value of the new plane equation. Linear functions attain
their extrema at vertices, so this is necessary and sufficient.
Points on the plane are kept in both new closed cells.
Exact rational arithmetic gives the successive counts
\[
 1,2,4,6,8,12,15,18,26,36,46,56,72,84,96.
\]
The supplied \texttt{cube\_\allowbreak{}witnesses.json} gives a rational interior
point of each final cell. All 96 sign vectors are distinct, with
no zero coordinate. The cells are convex and hence connected.
\end{proof}

\paragraph{Independent count.}
The second calculation uses the intersection poset rather than
cell splitting. For an arrangement restricted to an open convex
set, the region count is the sum of \(|\mu|\) over flats meeting
that set. The usual formula is stated in
\citet[Theorem~A, p.~100]{greenezaslavsky1983whitney}.
The restriction follows by inclusion--exclusion with compactly
supported Euler characteristic: each nonempty flat--cube
intersection is open convex in its affine span, contributing
\((-1)^{\dim}\); each region contributes \((-1)^3\).
Grouping identical intersections gives the Möbius sum, with its
alternating codimension signs.

Exact interval tests find 43 line flats meeting the cube interior,
with total \(|\mu|=47\), and 21 point flats with total \(|\mu|=34\).
Thus the second count is
\[
                    1+14+47+34=96.
\]
Flats only meeting the boundary are excluded. Both independent
counts agree with the notebook; no correction to its answer is
needed.

\section{Problem index}\label{problems:1986}
Earlier-year occurrence identifiers are retained as text; current-year identifiers link to their entries.

\subsection{GP-1982-40: Multiplicativity of regularized determinants}\label{prob:GP-1982-40}
For specified admissible elliptic operators and a fixed regularization, determine the defect \(\log\det(AB)-\log\det A-\log\det B\), and conditions under which it vanishes. The operator order, spectral cuts, commutativity, and treatment of kernels must be fixed.
\par\noindent\textit{Occurrences:} 1982-587.7, \hyperref[item:1986-668.8]{1986-668.8}.
\par\noindent\textit{Status (2026-11):} As of 12 November 2026, the multiplicative anomaly and its local-residue description are established in the settings treated by \citet{kontsevichvishik1994geometry}. Unqualified multiplicativity is false; this occurrence reuses the earlier precise problem.

\subsection{GP-1986-01: Odd-moduli integration and the higher-genus RNS measure}\label{prob:GP-1986-01}
For super Riemann surfaces of a fixed genus and spin parity, determine the odd-moduli pushforward of the RNS measure on any specified projection domain, its modular transformation, and its degeneration behavior. Separately determine whether a global holomorphic projection to reduced spin moduli exists.
\par\noindent\textit{Occurrences:} \hyperref[item:1986-670.6]{1986-670.6}, \hyperref[item:1986-670.7]{1986-670.7}.
\par\noindent\textit{Status (2026-11):} As of 12 November 2026, the global-projection premise is obstructed for \(g\geq5\) by Donagi--Witten and for both parity components in every \(g\geq3\) by the theorem stated in the September 2026 Correa--Donagi--Noja preprint. Local odd integration remains meaningful. The 2025 genus-three analysis proves a three-term form but leaves its exact coefficient values conjectural and its numerical extraction without a rigorous error bound. Thus that coefficient determination remains a precise partial remainder; no general unsolved global projection is asserted.

\subsection{GP-1986-02: Which lattice theta sums have the required modular transformation?}\label{prob:GP-1986-02}
For a positive-definite lattice, determine the modular transformation of its theta series, specifying the modular group and scalar, multiplier, or vector-valued representation. The seminar asks whether the desired full-group scalar transformation forces even unimodularity.
\par\noindent\textit{Occurrences:} \hyperref[item:1986-671.7]{1986-671.7}.
\par\noindent\textit{Status (2026-11):} As of 12 November 2026, Poisson summation gives the precise dual-lattice transformation, and even unimodular lattices supply the full modular-group scalar form in the stated convention. Theta series of other integral lattices remain modular for suitable subgroups or vector-valued representations, so the unrestricted only-if assertion is false. This is a settled clarification, not a new open problem.

\subsection{GP-1986-03: A common linear complex structure for real codimension-two arrangements}\label{prob:GP-1986-03}
Given real codimension-two linear subspaces \(L_i\subset\mathbb R^{2m}\), decide whether there is a real linear operator \(J\) with \(J^2=-I\) and \(J(L_i)=L_i\) for every \(i\). A 2-arrangement also imposes even-codimension intersection conditions; these are additional hypotheses, not an automatic consequence of the individual codimensions.
\par\noindent\textit{Occurrences:} \hyperref[item:1986-672.8]{1986-672.8}.
\par\noindent\textit{Status (2026-11):} As of 12 November 2026, the proposed automatic complexification is false. For pairwise transverse planes in \(\mathbb R^4\), \citet[Theorems 1 and 3]{batzies2007planes} gives simultaneous-conjugacy and four-plane criteria. The chapter rederives a bounded linear-algebra formulation as an annotation to known work, with no new research claim and no inference that a higher-dimensional classification is an open problem. Topological equivalence is not settled by the linear criterion.

\subsection{GP-1986-04: Polynomial quotients of the affine type-A Hecke algebra}\label{prob:GP-1986-04}
For \(q\ne0\), identify the universal quotient of the affine type-A Hecke algebra by \(p(X_n)=0\), where \(p\) is monic of degree \(r\) with nonzero constant term, and determine its dimension and a basis; the notebook singles out \(p(z)=z^2-1\).
\par\noindent\textit{Occurrences:} \hyperref[item:1986-681.6]{1986-681.6}.
\par\noindent\textit{Status (2026-11):} As of 12 November 2026, this is the cyclotomic Hecke algebra of type \(G(r,1,n)\), of dimension \(r^n n!\); \(r=2\) yields type B. The Ariki--Koike basis theorem holds without semisimplicity, as read in \citet[\S2, (2.2)]{mathas2004matrix}. The universal quotient target answers the precise notebook reading adopted here; it is not a claim of an evaluation into the original finite type-A algebra for every polynomial.

\subsection{GP-1986-30: Faithfulness of universal Burau representations}\label{prob:GP-1986-30}
Determine the ranks \(n\) for which the Burau representation over \(\mathbb Z[t,t^{-1}]\) is faithful.
\par\noindent\textit{Occurrences:} \hyperref[item:1986-686.4]{1986-686.4}.
\par\noindent\textit{Status (2026-11):} As of 12 November 2026, faithful for \(n\le3\), unfaithful for \(n\ge5\). A July 2026 preprint, revised September 2026, claims faithfulness for \(n=4\); its main theorem was read but not independently checked here.

\subsection{GP-1986-31: Three- and four-dimensional reduced planar gauge models}\label{prob:GP-1986-31}
For a specified one-site unitary plaquette model, determine its nonperturbative planar free energy and Wilson loops, with the center-symmetry conditions needed for volume reduction made explicit.
\par\noindent\textit{Occurrences:} \hyperref[item:1986-688.7]{1986-688.7}, \hyperref[item:1986-695.14]{1986-695.14}.
\par\noindent\textit{Status (2026-11):} As of 12 November 2026 the general all-coupling planar solution is unavailable in the sources checked; exact two-dimensional reductions and higher-dimensional numerical methods are established. Unconditional equivalence of the original untwisted model with four-dimensional weak-coupling gauge theory is false.

\subsection{GP-1986-32: Commuting quantizations of local KdV Hamiltonians}\label{prob:GP-1986-32}
Construct commuting normally ordered deformations of the complete family of classical local KdV integrals.
\par\noindent\textit{Occurrences:} \hyperref[item:1986-693.6]{1986-693.6}.
\par\noindent\textit{Status (2026-11):} As of 12 November 2026 solved as an algebraic/formal and generic-parameter construction; spectral classification and exceptional representation specializations are distinct questions.

\subsection{GP-1986-33: Quantum KdV spectrum in irreducible Virasoro representations}\label{prob:GP-1986-33}
For fixed central charge and highest weight, determine the joint spectral and Jordan data of the commuting quantum KdV charges on the irreducible highest-weight module.
\par\noindent\textit{Occurrences:} \hyperref[item:1986-693.15]{1986-693.15}.
\par\noindent\textit{Status (2026-11):} As of 12 November 2026 partly solved, with explicit constructions and special-regime results. The opened 2026 Masoero--Ruzza v1 proves a particular c=-2 admissible-QQ completeness statement for nonintegral momentum; it is not an arbitrary-parameter theorem. The chapter provides a checked level-two Jordan obstruction in a nonunitary irreducible quotient.

\subsection{GP-1986-34: Quantization of second Adler--Gelfand--Dickey structures}\label{prob:GP-1986-34}
Construct a quantum algebra deforming the second Gelfand--Dickey bracket on monic scalar operators of arbitrary order, in the principal Drinfeld--Sokolov setting.
\par\noindent\textit{Occurrences:} \hyperref[item:1986-693.16]{1986-693.16}.
\par\noindent\textit{Status (2026-11):} As of 12 November 2026 the quantum W-algebra/BRST construction resolves the algebraic problem, with precise generic-parameter identifications in the cited sources.

\subsection{GP-1986-35: Finite equational basis for subspace-lattice identities}\label{prob:GP-1986-35}
Do identities valid in subspace lattices of all finite-dimensional rational vector spaces follow from finitely many lattice identities? This is the explicit equational reading of the notebook question about linear equivalence.
\par\noindent\textit{Occurrences:} \hyperref[item:1986-696.6]{1986-696.6}.
\par\noindent\textit{Status (2026-11):} As of 12 November 2026 the stated rational formulation and the analogous all-field variety have a negative answer by Herrmann (1984), Theorem 1.4. The field and type of axioms were not stated in the notebook; the decidability problem GP-1985-20 is different.

\subsection{GP-1986-36: Preprojective reduction of the free modular lattice}\label{prob:GP-1986-36}
Describe the quotient of \(D^r\) by equality in all indecomposable preprojective rational representations and characterize its represented subspace families.
\par\noindent\textit{Occurrences:} \hyperref[item:1986-696.11]{1986-696.11}.
\par\noindent\textit{Status (2026-11):} As of 12 November 2026 the 1987 Gelfand--Lidskii--Ponomarev theorem gives a precise large-rank answer through compatible eventually generated families and graded right ideals. The present chapter uses the safely verified subrange r>=7; it does not turn the historical r=5 exception into a present-day open claim.

\subsection{GP-1986-60: The two-source unitary orbital integral}\label{prob:GP-1986-60}
Evaluate \(\int_{U(N)}\exp\sum_{\mu=1}^2\operatorname{Tr}(U^*P_\mu UQ_\mu)\,dU\) for diagonal source matrices, specifying whether the target is a convergent series, a compact finite formula or large-\(N\) information.
\par\noindent\textit{Occurrences:} \hyperref[item:1986-688.5]{1986-688.5}, \hyperref[item:1986-695.11]{1986-695.11}, \hyperref[item:1986-698.1]{1986-698.1}.
\par\noindent\textit{Status (2026-11):} As of 12 November 2026, the general integral has exact convergent Haar expansions; an explicit rank-two formula and rank-three expansions appear in \citet{holguin2024multimatrix}. The chapter develops a checked finite-rank formula and controlled series. No general compact HCIZ-type determinant or general large-\(N\) solution is claimed.

\subsection{GP-1986-61: Riemann hypothesis in the toroidal and horocycle formulations}\label{prob:GP-1986-61}
Do all nontrivial zeros of \(\zeta(s)\) satisfy \(\operatorname{Re}s=1/2\)? In the recorded approach, can toroidal spectral parameters be put in a suitable positive unitary realization, or equivalently can the requisite RH-strength horocycle remainders be proved under the stated test-function hypotheses?
\par\noindent\textit{Occurrences:} \hyperref[item:1986-700.15]{1986-700.15}, \hyperref[item:1986-700.17]{1986-700.17}, \hyperref[item:1986-703.8]{1986-703.8}.
\par\noindent\textit{Status (2026-11):} Open as of 12 November 2026. The toroidal structure theorems and qualitative horocycle equidistribution are established, but the positivity and zero-location conclusions are not.

\subsection{GP-1986-62: Orientable gluings of a labelled polygon by genus}\label{prob:GP-1986-62}
For a labelled oriented \(2n\)-gon, count orientation-reversing pairings of its sides whose quotient has genus \(g\), with cyclic rotations kept distinct.
\par\noindent\textit{Occurrences:} \hyperref[item:1986-700.2]{1986-700.2}, \hyperref[item:1986-702.1]{1986-702.1}, \hyperref[item:1986-702.7]{1986-702.7}.
\par\noindent\textit{Status (2026-11):} Solved before the recorded discussions by Harer and Zagier (1986), with an exact coefficient formula, a recurrence, and a Gaussian-matrix derivation. As of 12 November 2026 this is an established enumeration theorem.

\subsection{GP-1986-63: Coloured planar maps with controlled face counts and interfaces}\label{prob:GP-1986-63}
Enumerate rooted orientable planar maps with specified numbers of faces of the prescribed degrees and colours, and a specified number of edges between unlike colours; specify primal/dual and rooting conventions before extracting coefficients.
\par\noindent\textit{Occurrences:} \hyperref[item:1986-701.1]{1986-701.1}, \hyperref[item:1986-702.8]{1986-702.8}.
\par\noindent\textit{Status (2026-11):} As of 12 November 2026, formal multi-matrix generating functions and the planar Ising solution provide the recorded framework. The notebook does not select a unique finer counting convention, so no unsolved numerical problem is inferred.

\subsection{GP-1986-64: Artin holomorphy for nontrivial irreducible complex representations}\label{prob:GP-1986-64}
For every nontrivial irreducible continuous finite-dimensional complex representation \(\rho\) of the absolute Galois group of a number field, is its Artin L-function entire?
\par\noindent\textit{Occurrences:} \hyperref[item:1986-703.12]{1986-703.12}.
\par\noindent\textit{Status (2026-11):} Open in general as of 12 November 2026; Brauer induction proves meromorphic continuation, and abelian, monomial and numerous other cases are known. Recent partial holomorphy criteria do not establish the universal conjecture.

\subsection{GP-1986-65: The normalized mean-field spin-glass ground energy}\label{prob:GP-1986-65}
Determine the deterministic limit of \(N^{-1}\min_{\sigma\in\{-1,1\}^N}N^{-1/2}\sum_{i<j}J_{ij}\sigma_i\sigma_j\) for independent standard Gaussian couplings.
\par\noindent\textit{Occurrences:} \hyperref[item:1986-708.2]{1986-708.2}.
\par\noindent\textit{Status (2026-11):} Solved in the variational sense as of 12 November 2026 by the zero-temperature Parisi formula of Auffinger and Chen (2016 preprint; 2017 publication). This is a mean-field theorem, with the normalization and Gaussian assumptions explicit.

\subsection{GP-1986-66: Parisi ultrametricity of asymptotic spin-glass overlaps}\label{prob:GP-1986-66}
Identify a rigorous thermodynamic interpretation of the hierarchical overlap assertion and prove ultrametricity of the resulting replica array.
\par\noindent\textit{Occurrences:} \hyperref[item:1986-708.4]{1986-708.4}.
\par\noindent\textit{Status (2026-11):} As of 12 November 2026, Panchenko\textquoteright s theorem proves ultrametricity for random overlap measures satisfying the Ghirlanda--Guerra identities (2011 preprint; 2013 publication). It does not assert ultrametricity for arbitrary finite near-ground-state sets.

\subsection{GP-1986-67: Explicit integral and de Rham comparison for Chern classes}\label{prob:GP-1986-67}
Construct explicit Cech/de Rham comparison data, preferably through low-weight motivic complexes, for the integral Chern classes of a holomorphic vector bundle.
\par\noindent\textit{Occurrences:} \hyperref[item:1986-711.3]{1986-711.3}, \hyperref[item:1986-712.17]{1986-712.17}.
\par\noindent\textit{Status (2026-11):} As of 12 November 2026, functorial Deligne Chern classes exist, and inspected explicit weight-two formulas construct the comparison through the Bloch complex and dilogarithm. The announced 1987 paper is bibliographically identified, but its complete higher-weight formulas were not inspected. This access limitation is not classified as mathematical openness.

\subsection{GP-1986-68: Beilinson--Soule motivic vanishing}\label{prob:GP-1986-68}
For every field \(F\), positive weight \(p\), and \(j\le0\), prove \(H^j(F,\mathbb Q(p))=0\); compare with the stronger integral sheaf-amplitude requirement \([1,p]\) stated in the notebook.
\par\noindent\textit{Occurrences:} \hyperref[item:1986-712.12]{1986-712.12}.
\par\noindent\textit{Status (2026-11):} Open in general as of 12 November 2026. Modern motivic complexes exist, the upper-degree comparison with Milnor K-theory is established, and important arithmetic cases are known. Garkusha\textquoteright s criteria and Soul\'e\textquoteright s 4 November 2026 interview explicitly retain the general lower-degree conjecture.

\section{Corrections to the transcription and translation}\label{corr:1986}
The evidence distinguishes a transcription correction from a mathematical emendation of wording faithfully preserved in the notebook.
\subsection*{P001257 (both; provenance): Boundary fragment before the Kontsevich heading}
\noindent\textit{Transcription reads.} Page included in the catalogue range of record 667\par
\noindent\textit{Critical reading.} Preserve as a Beilinson continuation, indexed at 1986-667.4 only for catalogue containment; do not attribute it to Kontsevich.\par
\noindent\textit{Evidence.} The next heading begins on P001258; the determinant-line discussion continues 1985-666.21 and is separately attributed in the item.\par
\subsection*{P001258 (both; provenance): Kontsevich session heading}
\noindent\textit{Transcription reads.} 24/II-85; catalogue record 667: 1987-02-24\par
\noindent\textit{Critical reading.} Retain both readings and flag the conflict; the neighboring session headings read 1986. Do not silently select a priority date.\par
\noindent\textit{Evidence.} Original scan directly inspected; the year 85 is visible. The catalogue date is preserved in talks.json.\par
\subsection*{P001259 (both; provenance): Date of records 668}
\noindent\textit{Transcription reads.} Catalogue: 1987-02-24; manuscript: 24.2.86\par
\noindent\textit{Critical reading.} Preserve the catalogue date but display the manuscript date separately, with a discrepancy flag.\par
\noindent\textit{Evidence.} Russian and English page headings agree on the manuscript year; 1986 dates are Mondays. Olshanskii 1987 independently records presentation by Gelfand on 5 May 1986.\par
\subsection*{P001262 (both; provenance): Date of records 669 and 670}
\noindent\textit{Transcription reads.} Catalogue: 1987-03-03; manuscript: 3/III-86\par
\noindent\textit{Critical reading.} Preserve the catalogue date but display the manuscript date separately, with a discrepancy flag.\par
\noindent\textit{Evidence.} Russian and English page headings agree on the manuscript year; 1986 dates are Mondays. Olshanskii 1987 independently records presentation by Gelfand on 5 May 1986.\par
\subsection*{P001268 (both; provenance): Date of records 671}
\noindent\textit{Transcription reads.} Catalogue: 1987-03-10; manuscript: 10/III-86\par
\noindent\textit{Critical reading.} Preserve the catalogue date but display the manuscript date separately, with a discrepancy flag.\par
\noindent\textit{Evidence.} Russian and English page headings agree on the manuscript year; 1986 dates are Mondays. Olshanskii 1987 independently records presentation by Gelfand on 5 May 1986.\par
\subsection*{P001274 (both; provenance): Date of records 672 and 673}
\noindent\textit{Transcription reads.} Catalogue: 1987-03-24; manuscript: 24/III-86\par
\noindent\textit{Critical reading.} Preserve the catalogue date but display the manuscript date separately, with a discrepancy flag.\par
\noindent\textit{Evidence.} Russian and English page headings agree on the manuscript year; 1986 dates are Mondays. Olshanskii 1987 independently records presentation by Gelfand on 5 May 1986.\par
\subsection*{P001281 (both; provenance): Date of records 674 and 675}
\noindent\textit{Transcription reads.} Catalogue: 1987-04-07; manuscript: 7/IV-86\par
\noindent\textit{Critical reading.} Preserve the catalogue date but display the manuscript date separately, with a discrepancy flag.\par
\noindent\textit{Evidence.} Russian and English page headings agree on the manuscript year; 1986 dates are Mondays. Olshanskii 1987 independently records presentation by Gelfand on 5 May 1986.\par
\subsection*{P001287 (both; provenance): Date of records 676}
\noindent\textit{Transcription reads.} Catalogue: 1987-04-28; manuscript: 28.4.86\par
\noindent\textit{Critical reading.} Preserve the catalogue date but display the manuscript date separately, with a discrepancy flag.\par
\noindent\textit{Evidence.} Russian and English page headings agree on the manuscript year; 1986 dates are Mondays. Olshanskii 1987 independently records presentation by Gelfand on 5 May 1986.\par
\subsection*{P001290 (both; provenance): Date of records 677 and 678}
\noindent\textit{Transcription reads.} Catalogue: 1987-05-05; manuscript: 5/V-86\par
\noindent\textit{Critical reading.} Preserve the catalogue date but display the manuscript date separately, with a discrepancy flag.\par
\noindent\textit{Evidence.} Russian and English page headings agree on the manuscript year; 1986 dates are Mondays. Olshanskii 1987 independently records presentation by Gelfand on 5 May 1986.\par
\subsection*{P001257-P001298 (both; provenance): Overlapping catalogue page ranges}
\noindent\textit{Transcription reads.} Adjacent catalogue ranges include parts of neighboring talks\par
\noindent\textit{Critical reading.} Use the named manuscript headings to assign items; retain the catalogue ranges as provenance only.\par
\noindent\textit{Evidence.} Per-record summaries state the actual boundaries; in particular P001282 changes from Cherednik to Kogan, P001291 changes from Olshanskii to Cherednik, and P001298 changes to Ovsienko and Kazakov.\par
\subsection*{P001258 (both; faithful error): Tangent-space quotient}
\noindent\textit{Transcription reads.} \(\mathfrak g/\mathfrak k\)\par
\noindent\textit{Critical reading.} \(\mathfrak g/\gamma\) for the decorated moduli problem, with the curve-dependent stabilizer retained.\par
\noindent\textit{Evidence.} Scan and following Kodaira--Spencer formula; Kontsevich 1987, section 1; item 1986-667.2.\par
\subsection*{P001260 (both; faithful error): Equation (16), determinant variation}
\noindent\textit{Transcription reads.} \(\delta_y\delta_{\bar y}[\det\Delta_j/(\det N^{(j)}\det N^{(1-j)})]=c_j\)\par
\noindent\textit{Critical reading.} Interpret as curvature: mixed derivatives of the logarithm of the determinant-line norm, with the universal two-form and normalization supplied.\par
\noindent\textit{Evidence.} Scan directly inspected; Belavin--Knizhnik 1986, section 3; item 1986-668.7. The missing logarithm is absent in the manuscript too.\par
\subsection*{P001263 (both; faithful error): Example (9), pure weight filtration}
\noindent\textit{Transcription reads.} \(W_a=H^a(C,\mathbb Z),\ W_{a+1}=0\)\par
\noindent\textit{Critical reading.} \(W_{a-1}=0\), with the filtration considered over \(\mathbb Q\); preserve the written plus sign in the diplomatic text.\par
\noindent\textit{Evidence.} Scan directly inspected; the very next definition (10) explicitly has the correct subscript a-1; item 1986-669.3.\par
\subsection*{P001263 (both; faithful error): Example (9), purity hypothesis}
\noindent\textit{Transcription reads.} A smooth complex manifold has pure Hodge structure in every degree\par
\noindent\textit{Critical reading.} For this standard assertion require a smooth proper complex algebraic variety, or an appropriate compact Kahler hypothesis; a general smooth open variety has mixed cohomology.\par
\noindent\textit{Evidence.} Scan and Russian text; Deligne 1971, sections 2.1 and 3.2; item 1986-669.3.\par
\subsection*{P001265 (both; faithful error): Particle-size comparison}
\noindent\textit{Transcription reads.} Proton size approximately \(10^{-35}\,\mathrm{cm}\)\par
\noindent\textit{Critical reading.} Flag the numerical entry as inconsistent with the notebook comparison on P001269; do not silently manufacture a corrected measurement.\par
\noindent\textit{Evidence.} The English and Russian entries agree; item 1986-670.1 records the discrepancy as a notebook error.\par
\subsection*{P001266 (both; ambiguity): Equation (10), Green kernel}
\noindent\textit{Transcription reads.} One Abelian-integral factor and a theta function are written schematically\par
\noindent\textit{Critical reading.} Complete the harmonic correction as a quadratic form in \(\operatorname{Im}\int\omega\), and use the prime form or specify the half-differential factors and normalization.\par
\noindent\textit{Evidence.} Scan directly inspected; the second vector factor is suppressed. Belavin--Knizhnik 1986, section 2; item 1986-670.3.\par
\subsection*{P001270 (both; faithful error): Equation for the harmonic winding representative}
\noindent\textit{Transcription reads.} Both holomorphic and antiholomorphic terms appear with plus signs\par
\noindent\textit{Critical reading.} \(x_{\mathrm{cl}}(z)=((\Lambda_1-\bar\tau\Lambda_2)z-(\Lambda_1-\tau\Lambda_2)\bar z)/(2i\operatorname{Im}\tau)\) for the two written shifts.\par
\noindent\textit{Evidence.} Original scan checked; direct substitution in \texttt{code/check\_editorial.py} verifies shifts by tau and 1 and reality. Item 1986-671.5.\par
\subsection*{P001270 (both; faithful error): Equation (12), sign of the functional-integral action}
\noindent\textit{Transcription reads.} A positive kinetic-action exponent is displayed in the Euclidean functional integral\par
\noindent\textit{Critical reading.} Use \(e^{-S_E}\) for a positive Euclidean action. The later winding sum (21) already has the correct negative exponent; do not change that line.\par
\noindent\textit{Evidence.} Russian and English formulas (12) and (21) have different sign conventions; compact-boson construction in Gross et al. 1985; item 1986-671.6.\par
\subsection*{P001271 (both; ambiguity): Question about modularity and self-duality}
\noindent\textit{Transcription reads.} Modularity without naming the modular group or transformation representation\par
\noindent\textit{Critical reading.} Distinguish a scalar full-group transformation from modularity for a congruence subgroup or a vector-valued representation.\par
\noindent\textit{Evidence.} Poisson summation and the heterotic lattice construction; item 1986-671.7 and GP-1986-02. This is a hypothesis clarification, not a new open problem.\par
\subsection*{P001272-P001273 (both; faithful error): Root-operator commutator}
\noindent\textit{Transcription reads.} Zero when the sum of the two roots is not a root\par
\noindent\textit{Critical reading.} Exclude opposite roots from the zero case; their bracket has Cartan and central terms, with the required cocycle conventions.\par
\noindent\textit{Evidence.} Russian and English text; identified Frenkel--Kac construction, with direct-original access limitation retained in item 1986-671.9.\par
\subsection*{P001276 (both; ambiguity): Question about complexifying real codimension-two arrangements}
\noindent\textit{Transcription reads.} Real codimension two without further conditions on intersections\par
\noindent\textit{Critical reading.} Keep individual real codimension two distinct from the extra even-codimension intersection hypothesis of a 2-arrangement, and distinguish a linear complex structure from topological equivalence.\par
\noindent\textit{Evidence.} Russian text and original scan checked by the integrating researcher; Ziegler 1993 and Batzies 2007; item 1986-672.8 and GP-1986-03.\par
\subsection*{P001278 (both; faithful error): Plane-curve Milnor-fibre formula}
\noindent\textit{Transcription reads.} \(\mu=2g-(\text{holes})+1\)\par
\noindent\textit{Critical reading.} \(\mu=2g+b-1\) when holes denotes the number \(b\) of boundary components of the Milnor fibre.\par
\noindent\textit{Evidence.} Scan directly inspected and Russian text agrees. Milnor 1968; item 1986-673.4. This is an emendation of a faithful reading.\par
\subsection*{P001279 (both; faithful error): Sentence following equation (7)}
\noindent\textit{Transcription reads.} The logarithmic exponent k is the size of a Jordan block\par
\noindent\textit{Critical reading.} A block of size \(s\) permits logarithmic degree at most \(s-1\); a given period need not attain that maximum.\par
\noindent\textit{Evidence.} Original scan shows the insertion for block size; Varchenko 1982, sections 1--2; item 1986-673.7.\par
\subsection*{P001284 (both; faithful error): Equation (10), Virasoro commutator}
\noindent\textit{Transcription reads.} \(L_{n-m}\)\par
\noindent\textit{Critical reading.} \(L_{n+m}\)\par
\noindent\textit{Evidence.} Scan directly inspected; the mode grading and Furuuchi--Ohta 2006 oscillator conventions determine the sum; item 1986-675.5.\par
\subsection*{P001285 (both; faithful error): Equation (21), level-two constraint}
\noindent\textit{Transcription reads.} \(\partial_\mu h_{\mu\nu}=A_\nu\)\par
\noindent\textit{Critical reading.} The right side involves the level-two vector \(B_\nu\) of equation (16), up to its stated normalization, rather than the separate level-one vector \(A_\nu\).\par
\noindent\textit{Evidence.} Original scan directly inspected. Equation (16) uses B for the coefficient of the mode -2; item 1986-675.7 renames it C for clarity.\par
\subsection*{P001287 (both; faithful error): RTT relation, left side only}
\noindent\textit{Transcription reads.} \(L^2(-\lambda_2)\)\par
\noindent\textit{Critical reading.} \(L^2(\lambda_2)\) in the convention of the same equation on its right side.\par
\noindent\textit{Evidence.} Original scan directly inspected; the unmatched minus sign is present in both transcriptions and incompatible with the coefficient presentation; item 1986-676.1.\par
\subsection*{P001287 (both; faithful error): Quantum-determinant antisymmetrizer}
\noindent\textit{Transcription reads.} Tensor degree n used where the matrix size is N\par
\noindent\textit{Critical reading.} Use the top antisymmetrizer on \((\mathbb C^N)^{\otimes N}\); the finite level \(n\) is a different parameter.\par
\noindent\textit{Evidence.} Original scan checked; quantum determinant definition in Molev--Nazarov--Olshanski 1996; item 1986-676.2.\par
\subsection*{P001287 (both; faithful error): Possible quantum-determinant poles}
\noindent\textit{Transcription reads.} \(i=0,\ldots,n-1\)\par
\noindent\textit{Critical reading.} \(i=0,\ldots,N-1\) for the N shifted matrix factors, with orders controlled by the finite level.\par
\noindent\textit{Evidence.} Original scan checked; item 1986-676.3 separates matrix size N from level n.\par
\subsection*{P001289 (both; faithful error): Last line, symmetric-group order}
\noindent\textit{Transcription reads.} \(S_n\), ord. \(n\)\par
\noindent\textit{Critical reading.} \(S_n\) has order \(n!\); retain the written ord. n as notebook notation.\par
\noindent\textit{Evidence.} Original scan directly inspected; the Russian abbreviation is por. n. Item 1986-676.6.\par
\subsection*{P001291 (both; ambiguity): Highest-weight-vector uniqueness}
\noindent\textit{Transcription reads.} A unique highest-weight vector\par
\noindent\textit{Critical reading.} A unique highest-weight line in an irreducible highest-weight module; a vector requires a normalization.\par
\noindent\textit{Evidence.} Russian and English highest-weight setup; item 1986-678.1.\par
\subsection*{P001292 (both; faithful error): Two-dimensional antisymmetrization coefficient}
\noindent\textit{Transcription reads.} Coefficient \(1/4\) in the scalar mixed determinant\par
\noindent\textit{Critical reading.} Coefficient \(1/2\) for the normalized projector \((I-P)/2\).\par
\noindent\textit{Evidence.} Original scan directly inspected. \texttt{code/check\_editorial.py} computes the projector matrix exactly and also checks A=B=I; item 1986-678.3.\par
\subsection*{P001293 (both; faithful error): First closure axiom}
\noindent\textit{Transcription reads.} \(\overline{A\cap B}=\bar A\cap\bar B\)\par
\noindent\textit{Critical reading.} Monotonicity: \(A\subseteq B\Rightarrow\operatorname{cl}(A)\subseteq\operatorname{cl}(B)\). Preserve the displayed stronger false identity as notebook text.\par
\noindent\textit{Evidence.} Original scan directly inspected; the following span example and Oxley closure axioms determine the intended definition; item 1986-680.1.\par
\subsection*{P001293 (both; ambiguity): Closed singleton assertion}
\noindent\textit{Transcription reads.} \(\bar e_i=e_i\) for a vector configuration\par
\noindent\textit{Critical reading.} Add simplicity: exclude zero vectors and parallel pairs. Looplessness alone is insufficient.\par
\noindent\textit{Evidence.} Scan and vector-span example; Oxley and Gelfand--Serganova; item 1986-680.2.\par
\subsection*{P001294 (both; faithful error): Two-generator Coxeter example}
\noindent\textit{Transcription reads.} A finite integer product-order relation may still give an infinite rank-two group\par
\noindent\textit{Critical reading.} The rank-two group has order \(2m\) for finite \(m\). Infinity requires the omitted product-order relation, written \(m=\infty\).\par
\noindent\textit{Evidence.} Scan directly inspected; Coxeter 1934, introduction; item 1986-680.4.\par
\subsection*{P001294 (both; ambiguity): Icosahedral group discussion}
\noindent\textit{Transcription reads.} The successive groups of order 120 are not distinguished\par
\noindent\textit{Critical reading.} Distinguish the full icosahedral reflection group H3 from the binary icosahedral subgroup of SU(2); the 600-cell has tetrahedral three-dimensional facets.\par
\noindent\textit{Evidence.} Original scan inspected; Coxeter 1934 order table on p.618; item 1986-680.5.\par
\subsection*{P001294 (both; ambiguity): Bruhat-order definition}
\noindent\textit{Transcription reads.} A reduced word occurs inside another reduced word\par
\noindent\textit{Critical reading.} Read as a subsequence of a reduced expression, not necessarily a consecutive substring; translated orders require a least element, not merely some minimal element.\par
\noindent\textit{Evidence.} Scan and Gelfand--Serganova 1987, section 8; item 1986-680.6.\par
\subsection*{P001295 (both; ambiguity): Seven-point nonrepresentability example}
\noindent\textit{Transcription reads.} Not every matroid is representable, illustrated by the seven listed triples\par
\noindent\textit{Critical reading.} The listed Fano matroid is nonrepresentable over the real or complex field but representable exactly in characteristic two; specify the field.\par
\noindent\textit{Evidence.} Original scan inspected; the written triples identify the Fano lines. Oxley, Theorem 4.1; item 1986-680.8.\par
\subsection*{P001295 (both; faithful error): Polynomial/exterior duality example near page foot}
\noindent\textit{Transcription reads.} \(x_ix_n-x_kx_i\)\par
\noindent\textit{Critical reading.} \(x_ix_k-x_kx_i\), with the same two indices in both products.\par
\noindent\textit{Evidence.} Original scan and Russian/English formula; the immediately preceding commutator definition and paired dual monomial fix the intended index; item 1986-680.10.\par
\subsection*{P001296 (both; faithful error): Coefficient sum in the quadratic-dual calculation}
\noindent\textit{Transcription reads.} \(\sum a_{ij}x_ix_k\)\par
\noindent\textit{Critical reading.} \(\sum a_{ij}x_ix_j\)\par
\noindent\textit{Evidence.} Original scan inspected; repeated coefficient indices and the following symmetric-matrix condition determine the correction; item 1986-680.11.\par
\subsection*{P001297 (both; ambiguity): Definition of the subgroup in the symmetric space}
\noindent\textit{Transcription reads.} Compact subgroup\par
\noindent\textit{Critical reading.} In the cited algebraic compactification theorem use the symmetric algebraic subgroup from an involution, with the specified normalizer. The abbreviation does not license an arbitrary compact real subgroup.\par
\noindent\textit{Evidence.} Scan inspected; De Concini--Procesi setup explicitly recalled in Chirivi--Maffei 2004, section 1; item 1986-681.1.\par
\subsection*{P001297 (both; faithful error): Normal-crossings parenthesis}
\noindent\textit{Transcription reads.} Boundary codimension k and coordinate vanishing equations are conflated\par
\noindent\textit{Critical reading.} The union is \(z_1\cdots z_k=0\), of codimension one; the intersection is \(z_1=\cdots=z_k=0\), of codimension k.\par
\noindent\textit{Evidence.} Original scan and Chirivi--Maffei 2004, Proposition 1.1; item 1986-681.2.\par
\subsection*{P001297 (both; faithful error): Picard restriction}
\noindent\textit{Transcription reads.} An isomorphism onto the Picard group of the closed orbit\par
\noindent\textit{Critical reading.} Restriction is injective; its image is the specified subgroup generated by the relevant weights, not generally the entire Picard group.\par
\noindent\textit{Evidence.} Original scan inspected; Chirivi--Maffei 2004, Propositions 1.2--1.4, citing De Concini--Procesi 1983, Proposition 8.1; item 1986-681.3.\par
\subsection*{P001297 (both; faithful error): Commuting finite Hecke generators}
\noindent\textit{Transcription reads.} \(j\ne i,i+1\)\par
\noindent\textit{Critical reading.} \(|i-j|>1\)\par
\noindent\textit{Evidence.} Original scan inspected; the written asymmetric exclusion incorrectly includes the adjacent generator i-1; Mathas 2004, section 2; item 1986-681.4.\par
\subsection*{P001298 (both; faithful error): Proposed affine evaluation}
\noindent\textit{Transcription reads.} \(X_{n-1}\mapsto qT_{n-1}^{2}\)\par
\noindent\textit{Critical reading.} \(X_{n-1}\mapsto qT_{n-1}^{-2}\), and use inverse products recursively with the written Bernstein relation.\par
\noindent\textit{Evidence.} Original scan has the positive square. Exact noncommutative reduction and independent regular-representation check are in \texttt{code/check\_editorial.py}; item 1986-681.5.\par
\subsection*{P001298 (both; provenance): Transition after Cherednik polynomial-relation question}
\noindent\textit{Transcription reads.} The full page is inside more than one catalogue range\par
\noindent\textit{Critical reading.} Assign the upper Hecke question to record 681 with Cherednik as intervention speaker; assign the lower named Ovsienko and Kazakov contributions to their own records.\par
\noindent\textit{Evidence.} Original scan inspected; the speaker headings and change of subject are explicit; no page-level attribution is imposed.\par
\subsection*{P001299-P001302 (both; faithful error): Matrix-model normalization}
\noindent\textit{Transcription reads.} Quadratic term U\textasciicircum{}2, followed by the standard 72g\textasciicircum{}2 series\par
\noindent\textit{Critical reading.} Use U\textasciicircum{}2/2 for the displayed standard series, or rescale g consistently.\par
\noindent\textit{Evidence.} Comparison with BIPZ (1978), Eqs. (42), (51), (52); these are mathematical normalization repairs, not asserted paleographic corrections.\par
\subsection*{P001302 (both; faithful error): Critical coupling}
\noindent\textit{Transcription reads.} g\_c=1/(108 sqrt(3))\par
\noindent\textit{Critical reading.} g\_c\textasciicircum{}2=1/(108 sqrt(3)), in the BIPZ quadratic normalization.\par
\noindent\textit{Evidence.} BIPZ (1978), Eq. (52) explicitly gives g\_c squared.\par
\subsection*{P001303 (both; provenance): Talk heading and catalogue attribution}
\noindent\textit{Transcription reads.} Catalogue: V. L. Ginzburg, 29 September 1986\par
\noindent\textit{Critical reading.} Manuscript: D. Gurevich, 13 October 1986. Preserve catalogue metadata and explicitly attribute the surviving items.\par
\noindent\textit{Evidence.} Original scan viewed: handwritten heading and date are clear; generalized Lie algebra content matches Gurevich (1986).\par
\subsection*{P001304-P001306 (both; provenance): Date of Kazakov continuation}
\noindent\textit{Transcription reads.} Catalogue: 29 September 1986\par
\noindent\textit{Critical reading.} The notes follow the heading dated 13 October 1986; keep the catalogue date and flag the manuscript sequence.\par
\noindent\textit{Evidence.} Viewed P001303 heading and checked source continuity into Kazakov section.\par
\subsection*{P001308 (both; ambiguity): Exponent versus Hecke parameter}
\noindent\textit{Transcription reads.} f\textasciicircum{}alpha; (T+1)(T-alpha)=0; alpha=-1 is hyperelliptic\par
\noindent\textit{Critical reading.} Write t=exp(2 pi i alpha); the Hecke relation uses t, and the double cover has t=-1.\par
\noindent\textit{Evidence.} Original scan viewed: the same handwritten alpha is indeed used in all three positions; this is a faithful mathematical ambiguity, not a mistranscription.\par
\subsection*{P001312 (both; faithful error): 1000-dimensional equatorial layer}
\noindent\textit{Transcription reads.} Thickness R/100 contains almost all volume\par
\noindent\textit{Critical reading.} About 0.126 for full thickness R/100, or about 0.248 for half-width R/100.\par
\noindent\textit{Evidence.} Russian transcription explicitly flags R/100 as manuscript reading; direct cross-sectional integral and beta-function check are saved in code/check\_slab.py.\par
\subsection*{P001313 (both; faithful error): Duistermaat--Heckman display}
\noindent\textit{Transcription reads.} (2 pi)\textasciicircum{}(-n) times a phase-free Hessian sum with omega\textasciicircum{}n\par
\noindent\textit{Critical reading.} For Liouville measure omega\textasciicircum{}n/n!, use (2 pi/t)\textasciicircum{}n and signature phases, or the equivalent oriented-weight formula.\par
\noindent\textit{Evidence.} Original scan directly viewed: inverse prefactor and omitted phase are present; comparison with DH addendum Theorem (1.5).\par
\subsection*{P001313 (both; provenance): Attribution of nilpotent-action conjecture}
\noindent\textit{Transcription reads.} Hypothesis (mine)\par
\noindent\textit{Critical reading.} Attribute to the notebook author, not automatically to I. M. Gelfand or V. L. Ginzburg.\par
\noindent\textit{Evidence.} Original scan and Russian transcription read «Гипотеза (моя)».\par
\subsection*{P001316 (both; faithful error): Quadratic oscillator modes}
\noindent\textit{Transcription reads.} Sum only over p>q\par
\noindent\textit{Critical reading.} Include the diagonal half-square when p=q, and specify normal ordering and zero modes.\par
\noindent\textit{Evidence.} Russian transcription explicitly confirms strict p>q; compare standard free-boson Virasoro modes and BLZ Eq. (31).\par
\subsection*{P001318 (both; faithful error): Riccati sweep equation}
\noindent\textit{Transcription reads.} alpha prime - alpha squared = 0\par
\noindent\textit{Critical reading.} alpha prime - alpha squared = -a(x), for y prime + alpha y=0 and y double prime-a y=0.\par
\noindent\textit{Evidence.} Original scan viewed; Russian source explicitly confirms the zero and exclamation mark. Direct differentiation provides the stated editorial repair.\par
\subsection*{P001319-P001321 (both; provenance): Missing catalogue page range}
\noindent\textit{Transcription reads.} Record 694 has no catalogue pages\par
\noindent\textit{Critical reading.} Retain empty catalogue pages but index surviving P001319-P001321 under record 694 as an explicit manuscript attribution.\par
\noindent\textit{Evidence.} Original P001319 scan has clear Beilinson heading and date 10.11.86; the coherent continuation ends before the dated Kazakov report.\par
\subsection*{P001319 (both; ambiguity): Dimension of moduli}
\noindent\textit{Transcription reads.} Dimension 6g-6\par
\noindent\textit{Critical reading.} Real dimension 6g-6, or complex dimension 3g-3, for g>=2.\par
\noindent\textit{Evidence.} The note is in a complex-algebraic context but does not label real dimension.\par
\subsection*{P001320 (both; faithful error): First tangent quotient}
\noindent\textit{Transcription reads.} T(K backslash X)\par
\noindent\textit{Critical reading.} T(K backslash x), vector fields regular away from the marked point.\par
\noindent\textit{Evidence.} English and Russian evidence blocks identify capital X in the first formula, while its defining diagram and second formula use x.\par
\subsection*{P001322 (both; faithful error): Vertex allocation}
\noindent\textit{Transcription reads.} sigma\_i=n\_i/3 is called a defect\par
\noindent\textit{Critical reading.} n\_i/3 is a unit-triangle area allocation; angular defect is 2 pi-n\_i pi/3.\par
\noindent\textit{Evidence.} The following page explains area division into thirds; this is a mathematical terminology correction.\par
\subsection*{P001323 (both; faithful error): Gaussian surface integral}
\noindent\textit{Transcription reads.} Integrate all vertex positions without fixing translation\par
\noindent\textit{Critical reading.} Pin one vertex or divide by translation volume before taking determinant powers.\par
\noindent\textit{Evidence.} Every edge energy is translation invariant; the constant graph-Laplacian mode has zero eigenvalue.\par
\subsection*{P001324 (both; faithful error): Planar limit Eq. (11)}
\noindent\textit{Transcription reads.} Z=lim log I\_N\par
\noindent\textit{Critical reading.} Z=lim N\textasciicircum{}(-2) log[I\_N(g)/I\_N(0)], as a formal normalized series.\par
\noindent\textit{Evidence.} Compare BIPZ (1978) and the N squared power of connected planar vacuum graphs.\par
\subsection*{P001325 (both; faithful error): Three-matrix angular Eq. (14)}
\noindent\textit{Transcription reads.} Three independent angular matrices; final spectral factor mu\par
\noindent\textit{Critical reading.} Impose U12 U23 U31=I, integrate only two independent relative matrices, and use lambda in the third pair.\par
\noindent\textit{Evidence.} Both source transcriptions identify final mu; Daul (1995), Sec. 1 explicitly explains the nonindependence of cyclic relative angles.\par
\subsection*{P001327 (both; faithful error): Preprojective definition}
\noindent\textit{Transcription reads.} Only finitely many morphisms rho to rho\_0\par
\noindent\textit{Critical reading.} Only finitely many isomorphism classes of indecomposable rho have Hom(rho,rho\_0) nonzero.\par
\noindent\textit{Evidence.} Original scan explicitly says morphisms; compare the representation-type finiteness in Gelfand--Lidskii--Ponomarev (1987), Sec. 1.\par
\subsection*{P001327 (both; faithful error): Perfect element definition}
\noindent\textit{Transcription reads.} For every representation rho, rho(a)=V or 0\par
\noindent\textit{Critical reading.} For every indecomposable representation rho, rho(a)=V or 0.\par
\noindent\textit{Evidence.} Original scan omits indecomposable here; the 1974 and 1987 published definitions require it.\par
\subsection*{P001328 (both; ambiguity): Complementary-series modulus}
\noindent\textit{Transcription reads.} Ordinary-looking absolute determinant raised to alpha, 0<alpha<1/2\par
\noindent\textit{Critical reading.} Use normalized local modulus |det|\_C=ordinary |det| squared for the interval 0<alpha<1/2.\par
\noindent\textit{Evidence.} Badulescu--Renard (2010), Sec. 1.2 explicitly states the square-modulus convention.\par
\subsection*{P001329 (both; faithful error): Discrete-series definition}
\noindent\textit{Transcription reads.} Square-integrable matrix coefficients on GL(k,F)\par
\noindent\textit{Critical reading.} Square integrability modulo the center, with unitary central character.\par
\noindent\textit{Evidence.} The center is noncompact; source definitions in Tadic (1986) and Badulescu--Renard (2010) include the quotient.\par
\subsection*{P001329 (both; faithful error): Speh segment exponents}
\noindent\textit{Transcription reads.} Second exponent n/2-1; no delta on final factor\par
\noindent\textit{Critical reading.} Second exponent (n-3)/2; final factor nu\textasciicircum{}(-(n-1)/2) delta.\par
\noindent\textit{Evidence.} Original scan viewed; the overwritten second exponent reads n/2-1, in agreement with both transcriptions. Compare Badulescu--Renard Eq. (7.1).\par
\subsection*{P001329 (both; faithful error): Field scope}
\noindent\textit{Transcription reads.} For every field\par
\noindent\textit{Critical reading.} For nondiscrete locally compact fields in the relevant topological representation category.\par
\noindent\textit{Evidence.} Theorems of Tadic and Vogan assume local fields and normalized smooth/Harish-Chandra induction.\par
\subsection*{P001331 (both; faithful error): ordinary angular integral}
\noindent\textit{Transcription reads.} determinant divided by two Vandermondes\par
\noindent\textit{Critical reading.} Multiply by the normalized-Haar factor product(j!,j=1,...,N-1)/(2c)\textasciicircum{}(N(N-1)/2).\par
\noindent\textit{Evidence.} The original scan contains no such factor; the missing normalization is a faithful mathematical abbreviation, not a translation error.\par
\subsection*{P001332 (both; faithful error): monic moment determinant}
\noindent\textit{Transcription reads.} moment determinant without denominator\par
\noindent\textit{Critical reading.} Divide by the leading moment minor to make the polynomial monic.\par
\noindent\textit{Evidence.} Russian and original scan agree; expansion in the leading power gives that minor, not 1.\par
\subsection*{P001332 (both; faithful error): product of biorthogonal norms}
\noindent\textit{Transcription reads.} product h\_i, i=1,...,N\par
\noindent\textit{Critical reading.} For labelled eigenvalues use N! product h\_i, i=0,...,N-1, together with the declared angular normalization.\par
\noindent\textit{Evidence.} The scan does not declare ordered eigenvalues; the necessary convention is absent.\par
\subsection*{P001329-P001334 (both; provenance): catalogue boundary of record 698}
\noindent\textit{Transcription reads.} whole range attributed to Kazakov\par
\noindent\textit{Critical reading.} Kazakov heading on P001330; P001329 is the preceding discussion and 9 February 1987 starts on P001333.\par
\noindent\textit{Evidence.} Explicit speaker and date headings in the Russian manuscript.\par
\subsection*{P001333 (both; faithful error): modular transformation law}
\noindent\textit{Transcription reads.} determinant plus or minus 1\par
\noindent\textit{Critical reading.} Use determinant +1 for the stated holomorphic action on the upper half-plane.\par
\noindent\textit{Evidence.} Matrices of negative determinant exchange the half-planes; compare the subsequent slash action and Zagier 1977.\par
\subsection*{P001334 (both; faithful error): Fourier expansion}
\noindent\textit{Transcription reads.} exp(-2 pi i n z)\par
\noindent\textit{Critical reading.} Use exp(+2 pi i n z) for holomorphy at infinity and the later q definition.\par
\noindent\textit{Evidence.} The later q convention and all displayed modular-form expansions fix the sign.\par
\subsection*{P001334 (both; faithful error): Eisenstein lattice sum}
\noindent\textit{Transcription reads.} denominator squared\par
\noindent\textit{Critical reading.} The denominator exponent must equal the modular weight k.\par
\noindent\textit{Evidence.} The manuscript examples E4 and E6 and their transformation laws require exponent k.\par
\subsection*{P001335 (both; faithful error): Hecke coset formula}
\noindent\textit{Transcription reads.} extra factor 1/2\par
\noindent\textit{Critical reading.} Remove it if cosets already identify opposite matrices; alternatively specify doubled representatives.\par
\noindent\textit{Evidence.} T1=identity and the displayed T2 formula fix the normalization.\par
\subsection*{P001337 (both; faithful error): spectral trace sum}
\noindent\textit{Transcription reads.} sum of eigenforms through n\par
\noindent\textit{Critical reading.} Sum through dim S\_k.\par
\noindent\textit{Evidence.} The Hecke index n and the dimension of the cusp-form space are distinct.\par
\subsection*{P001339 (both; faithful error): class-number table}
\noindent\textit{Transcription reads.} H(-32)=2\par
\noindent\textit{Critical reading.} H(-32)=3 under the Hurwitz convention used in the trace formula.\par
\noindent\textit{Evidence.} The weight includes the primitive discriminant -32 contribution 2 and the imprimitive discriminant -8 contribution 1.\par
\subsection*{P001341 (both; faithful error): Eisenstein completion}
\noindent\textit{Transcription reads.} (2 pi)\textasciicircum{}s Gamma(s) E\par
\noindent\textit{Critical reading.} pi\textasciicircum{}(-s) Gamma(s) zeta(2s) E for the primitive sum.\par
\noindent\textit{Evidence.} Zagier 1981 Section 1 gives the symmetric completion.\par
\subsection*{P001342 (both; faithful error): horocycle limit}
\noindent\textit{Transcription reads.} integral over fundamental domain without factor\par
\noindent\textit{Critical reading.} Multiply the unnormalized hyperbolic integral by 3/pi.\par
\noindent\textit{Evidence.} The hyperbolic area of the modular fundamental domain is pi/3; the preceding formula has the correct factor.\par
\subsection*{P001343 (both; faithful error): negative zeta-values}
\noindent\textit{Transcription reads.} zeta(-3)=1/252\par
\noindent\textit{Critical reading.} zeta(-3)=1/120; zeta(-5)=-1/252.\par
\noindent\textit{Evidence.} Bernoulli formula zeta(1-2n)=-B\_(2n)/(2n).\par
\subsection*{P001344 (both; faithful error): E4 Fourier series}
\noindent\textit{Transcription reads.} coefficient 28 at q\textasciicircum{}2\par
\noindent\textit{Critical reading.} The q\textasciicircum{}2 coefficient inside the divisor-sum series is 9; 28 belongs to q\textasciicircum{}3.\par
\noindent\textit{Evidence.} sigma\_3(2)=9 and sigma\_3(3)=28.\par
\subsection*{P001345 (both; faithful error): Eisenstein constant}
\noindent\textit{Transcription reads.} -B\_k/(2k) labelled zeta(1-k)\par
\noindent\textit{Critical reading.} It is zeta(1-k)/2 for even k at least 2.\par
\noindent\textit{Evidence.} The Bernoulli formula has denominator k, not 2k.\par
\subsection*{P001345 (both; faithful error): range of Eisenstein product weights}
\noindent\textit{Transcription reads.} k/2 < n <= k-2\par
\noindent\textit{Critical reading.} For holomorphic weights at least 4 use 2 <= n <= k/2-2.\par
\noindent\textit{Evidence.} The stated product has weights 2n and k-2n.\par
\subsection*{P001346 (both; faithful error): Gamma integral factor}
\noindent\textit{Transcription reads.} (k-1)!\par
\noindent\textit{Critical reading.} (k-2)! for the displayed y\textasciicircum{}(k-2) integral.\par
\noindent\textit{Evidence.} The following page uses the corrected factorial; Gamma(k-1)=(k-2)!.\par
\subsection*{P001347 (both; faithful error): critical value conclusion}
\noindent\textit{Transcription reads.} L(f,2n)/zeta(2n) independent of n\par
\noindent\textit{Critical reading.} The relevant ratios become algebraic after division by one parity period; they are not all equal.\par
\noindent\textit{Evidence.} Kohnen--Zagier exact period theorem and its weight-12 table.\par
\subsection*{P001348 (both; provenance): bibliographic title}
\noindent\textit{Transcription reads.} Modular forms with rational coefficients\par
\noindent\textit{Critical reading.} Modular forms with rational periods.\par
\noindent\textit{Evidence.} Author bibliography and opened 1984 original.\par
\subsection*{P001353 (both; faithful error): moment integrand}
\noindent\textit{Transcription reads.} lambda\textasciicircum{}(2k) for the n-th even moment\par
\noindent\textit{Critical reading.} Use lambda\textasciicircum{}(2n), with the Gaussian scaling fixed consistently.\par
\noindent\textit{Evidence.} The moment being calculated is indexed by n; k is the orthogonal-polynomial summation index.\par
\subsection*{P001353 (both; faithful error): Harer--Zagier generating series}
\noindent\textit{Transcription reads.} normalization using (2n)!!\par
\noindent\textit{Critical reading.} For C\_n(N), use 1+2 sum C\_n(N)x\textasciicircum{}(n+1)/(2n-1)!!=((1+x)/(1-x))\textasciicircum{}N.\par
\noindent\textit{Evidence.} The zeroth moment is N and the first is N\textasciicircum{}2; the opened original gives the exact normalization.\par
\subsection*{P001354 (both; faithful error): hyperbolic-tangent coefficient formula}
\noindent\textit{Transcription reads.} square-root scaling inconsistent with u/2\par
\noindent\textit{Critical reading.} Use coefficient [u\textasciicircum{}(2g)] of ((u/2)/tanh(u/2))\textasciicircum{}(n+1).\par
\noindent\textit{Evidence.} Harer--Zagier Theorem 2; the n=2,g=1 coefficient must equal 1.\par
\subsection*{P001355-P001358 (both; provenance): records 703 and 704}
\noindent\textit{Transcription reads.} catalogue spans overlap; record 703 stops at P001356\par
\noindent\textit{Critical reading.} Attribute the explicit Beilinson heading and continuation through P001358 to 703, and the preceding Gelfand proposals to 704.\par
\noindent\textit{Evidence.} Original scan P001355 clearly separates the two named speakers. Catalogue metadata are retained in talks.json.\par
\subsection*{P001356 (both; faithful error): p-adic Gaussian analogue}
\noindent\textit{Transcription reads.} Z\_p minus zero called an open subgroup\par
\noindent\textit{Critical reading.} It is an open multiplicative submonoid of Q\_p\textasciicircum{}*; the open compact subgroup is Z\_p\textasciicircum{}*. The zeta test function is 1\_(Z\_p).\par
\noindent\textit{Evidence.} The manuscript writes Z\_p without a unit star; negative valuations show failure of closure under inversion.\par
\subsection*{P001356 (both; ambiguity): regularized power sum}
\noindent\textit{Transcription reads.} evaluate the first regular term at t=1\par
\noindent\textit{Critical reading.} Specify the parameter t=e\textasciicircum{}(-u) and take the constant term in u, or declare another regularization.\par
\noindent\textit{Evidence.} A Laurent finite part changes under nonlinear change of parameter; ordinary substitution diverges.\par
\subsection*{P001356 (both; faithful error): p-adic zeta measure}
\noindent\textit{Transcription reads.} one Z\_p-valued measure without regularizing factor\par
\noindent\textit{Critical reading.} Use a pseudomeasure for the zeta branch or a bounded Mazur measure with its auxiliary factor and character conventions.\par
\noindent\textit{Evidence.} Gelbart et al. equations (2.1), (2.12), (2.13) explicitly distinguish bounded measures from Bernoulli distributions.\par
\subsection*{P001356 (both; faithful error): unit decomposition}
\noindent\textit{Transcription reads.} Z\_p\textasciicircum{}*=(1+pZ\_p) times roots of unity of order p-1\par
\noindent\textit{Critical reading.} State p odd; at p=2 use {+1,-1} times (1+4Z\_2).\par
\noindent\textit{Evidence.} The two-adic unit group has a separate order-two factor.\par
\subsection*{P001355 (both; ambiguity): factor type}
\noindent\textit{Transcription reads.} III\_infinity\par
\noindent\textit{Critical reading.} The recognizable uniqueness result concerns type III\_1 with separable predual; retain III\_infinity as the recorded label.\par
\noindent\textit{Evidence.} Original scan confirms infinity. Haagerup 1987 Section 2 proves uniqueness for III\_1, conventionally denoted R\_infinity.\par
\subsection*{P001359 (both; faithful error): KdV and the Lax pair}
\noindent\textit{Transcription reads.} u\_t+6u u\_x+u\_xxx=0 with the displayed positive A\par
\noindent\textit{Critical reading.} For the displayed L and A, u\_t=6u u\_x+u\_xxx; reverse t consistently to obtain the first convention.\par
\noindent\textit{Evidence.} Original scan shows the positive third-order operator; compatibility supplies the sign check.\par
\subsection*{P001359 (both; faithful error): vacuum eigenfunction}
\noindent\textit{Transcription reads.} exp(ikx+4ik\textasciicircum{}3 t), eigenvalue k\textasciicircum{}2\par
\noindent\textit{Critical reading.} For the displayed pair use eigenvalue -k\textasciicircum{}2 and exp(ikx-4ik\textasciicircum{}3 t).\par
\noindent\textit{Evidence.} Direct substitution into both vacuum differential equations; original scan inspected.\par
\subsection*{P001360 (both; faithful error): KP Lax pair}
\noindent\textit{Transcription reads.} coefficient 2 and no nonlocal time term\par
\noindent\textit{Critical reading.} Use a consistent alpha coefficient and include the auxiliary term 3 alpha w with w\_x=u\_y.\par
\noindent\textit{Evidence.} Original scan confirms 2, distinct from alpha; the displayed pair otherwise forces extra restrictions.\par
\subsection*{P001362 (both; faithful error): zero curvature equivalence}
\noindent\textit{Transcription reads.} [L\_i,L\_j]=0 iff [M\_i,M\_j]=0\par
\noindent\textit{Critical reading.} Include partial\_i M\_j-partial\_j M\_i, unless these cross derivatives cancel.\par
\noindent\textit{Evidence.} Composition of parameter-dependent differential operators.\par
\subsection*{P001362 (both; faithful error): example matrix A}
\noindent\textit{Transcription reads.} A=[[0,1],[1,0]] while A\textasciicircum{}2=-1\par
\noindent\textit{Critical reading.} A=[[0,1],[-1,0]] is a consistent correction, retaining I=diag(1,-1).\par
\noindent\textit{Evidence.} Direct scan confirms two positive ones. Two independent exact calculations are preserved in check\_spectral\_operators.py and its output.\par
\subsection*{P001363 (both; ambiguity): matrix reality convention}
\noindent\textit{Transcription reads.} A\textasciicircum{}*=-A followed by real-eigenvalue resolvent formulas\par
\noindent\textit{Critical reading.} Declare an anti-Hermitian contour and transformed potential, or explicitly reconstruct the Hermitian model.\par
\noindent\textit{Evidence.} The manuscript anti-Hermitian sign is faithful; later real-cut formulas require a convention change.\par
\subsection*{P001364 (both; transcription): cube-problem parenthesis}
\noindent\textit{Transcription reads.} vypisal?\par
\noindent\textit{Critical reading.} The scan reads iziaschnaia, meaning elegant, in the parenthesis after the attribution.\par
\noindent\textit{Evidence.} Direct inspection of P001364; the word is a description of the problem, not a copying instruction.\par
\subsection*{P001365 (both; faithful error): imaginary part and jump}
\noindent\textit{Transcription reads.} Im f=(f\_+-f\_-)/2; f\_+\textasciicircum{}2-f\_-\textasciicircum{}2=2 Im f Re f\par
\noindent\textit{Critical reading.} Use Im f=(f\_+-f\_-)/(2i) and f\_+\textasciicircum{}2-f\_-\textasciicircum{}2=4i Im f Re f under usual conventions.\par
\noindent\textit{Evidence.} The two notebook equations are mutually inconsistent even if Im is redefined as a half-jump.\par
\subsection*{P001366 (both; faithful error): resolvent square root}
\noindent\textit{Transcription reads.} sqrt((Vprime/(2lambda))\textasciicircum{}2-R)\par
\noindent\textit{Critical reading.} With \(R=f^2-(V'/\lambda)f\), the radicand is \[\left(V'/(2\lambda)\right)^2+R.\]\par
\noindent\textit{Evidence.} The preceding definition fixes the sign algebraically; the original evidence block confirms minus R is in the manuscript.\par
\subsection*{P001369 (both; faithful error): ordinary conformal derivative}
\noindent\textit{Transcription reads.} partial\_tilde=F partial with F=d z\_tilde/dz\par
\noindent\textit{Critical reading.} Use partial\_tilde=F\textasciicircum{}(-1) partial, or reverse the definition of F.\par
\noindent\textit{Evidence.} The evidence block and manuscript preserve both formulas; the ordinary chain rule supplies the reciprocal.\par
\subsection*{P001369 (both; ambiguity): supermoduli dimension}
\noindent\textit{Transcription reads.} (3g-3|2g-2) without genus qualification\par
\noindent\textit{Critical reading.} State unpunctured genus g>=2; include automorphism and puncture corrections in other cases.\par
\noindent\textit{Evidence.} Witten equations (2.26)-(2.27) distinguish the dimension from the automorphism-corrected formula.\par
\subsection*{P001371 (both; provenance): Schwarz versus Sinai catalogue sequence}
\noindent\textit{Transcription reads.} record 707 before record 708\par
\noindent\textit{Critical reading.} Retain catalogue record numbers and dates, while placing the Sinai material on P001368-P001369 before the explicit Schwarz heading on P001369.\par
\noindent\textit{Evidence.} Manuscript speaker headings determine the actual order.\par
\subsection*{P001368 (both; faithful error): ground energy limit}
\noindent\textit{Transcription reads.} lim min H\par
\noindent\textit{Critical reading.} For the displayed SK scaling study lim N\textasciicircum{}(-1) min H.\par
\noindent\textit{Evidence.} Auffinger-Chen Theorem 1 concerns energy per spin; total extremal energy is extensive.\par
\subsection*{P001368 (both; faithful error): ground-state entropy}
\noindent\textit{Transcription reads.} many ground states differing in energy and positive exact-ground-state entropy\par
\noindent\textit{Critical reading.} Distinguish exact minimizers from near-minimizers in a specified energy window; Gaussian SK has a unique minimizing spin-flip pair almost surely.\par
\noindent\textit{Evidence.} Finite Gaussian nondegeneracy argument recorded with an independent algebraic check in item 708.3.\par
\subsection*{P001371 (both; faithful error): quaternionic projective space}
\noindent\textit{Transcription reads.} CP\textasciicircum{}n and CP\textasciicircum{}1 for quaternionic lines\par
\noindent\textit{Critical reading.} HP\textasciicircum{}n and HP\textasciicircum{}1.\par
\noindent\textit{Evidence.} The evidence block confirms C is written, while the definition explicitly uses quaternionic subspaces.\par
\subsection*{P001371 (both; faithful error): affine rigidity}
\noindent\textit{Transcription reads.} all transformations linear; simply connected gives the linear space\par
\noindent\textit{Critical reading.} Locally affine, including translations; a global identification with H\textasciicircum{}n also requires suitable completeness.\par
\noindent\textit{Evidence.} Sudbery Section 3 Theorem 1 includes the constant term; Salamon distinguishes locally affine structures.\par
\subsection*{P001372 (both; ambiguity): sphere dimension}
\noindent\textit{Transcription reads.} S\textasciicircum{}3 with three complex structures\par
\noindent\textit{Critical reading.} Retain S\textasciicircum{}3 as written and flag impossibility; S\textasciicircum{}4 is a plausible contextual conjecture, not a secure transcription correction.\par
\noindent\textit{Evidence.} Original evidence block identifies superscript 3; odd-dimensional real tangent spaces have no complex structure.\par
\subsection*{P001372 (both; faithful error): compact structure group}
\noindent\textit{Transcription reads.} Sp(1)Sp(n) for definition (1)\par
\noindent\textit{Critical reading.} Sp(1)Sp(n) after imposing a compatible metric; definition (1) alone has GL(n,H).\par
\noindent\textit{Evidence.} Salamon Section 1 identifies the general stabilizer and its maximal compact subgroup.\par
\subsection*{P001373 (both; faithful error): quaternion algebra in tangent automorphisms}
\noindent\textit{Transcription reads.} subalgebra in Aut T\_xM\par
\noindent\textit{Critical reading.} subalgebra in End T\_xM, whose nonzero quaternionic elements act invertibly.\par
\noindent\textit{Evidence.} Aut is a group and does not contain the zero endomorphism; the earlier definition on P001372 correctly uses End.\par
\subsection*{P001375 (both; faithful error): link definition}
\noindent\textit{Transcription reads.} immersed compact oriented one-manifold in R\textasciicircum{}3\par
\noindent\textit{Critical reading.} embedded tame compact oriented one-manifold, unless a theory of singular links is intended.\par
\noindent\textit{Evidence.} The manuscript evidence block confirms immersed; Alexander-Markov theory in this lecture uses embeddings.\par
\subsection*{P001375 (both; faithful error): factor trace claim}
\noindent\textit{Transcription reads.} unique factor representation for every q\par
\noindent\textit{Critical reading.} Unique algebraic Markov trace after fixing its additional parameter z; positivity and factor realizations are separate questions.\par
\noindent\textit{Evidence.} Jones 1987 Theorem 5.1 states the two-parameter trace theorem.\par
\subsection*{P001376 (both; faithful error): roots of unity}
\noindent\textit{Transcription reads.} Jones excludes roots of unity from consideration\par
\noindent\textit{Critical reading.} Only specified generic semisimplicity formulas exclude them; Jones explicitly emphasizes and studies them.\par
\noindent\textit{Evidence.} Jones 1987 discussion following Theorem 4.5 and the special-value/root-of-unity sections directly contradict the blanket attribution.\par
\subsection*{P001376 (both; faithful error): Hodge filtration}
\noindent\textit{Transcription reads.} F\textasciicircum{}j H\textasciicircum{}k equals hypercohomology of the truncated complex\par
\noindent\textit{Critical reading.} Use the image in de Rham hypercohomology; injectivity and a Hodge decomposition require additional conditions.\par
\noindent\textit{Evidence.} The definition of a filtration is by this image; compactness alone does not imply the Kahler Hodge decomposition.\par
\subsection*{P001377 (both; faithful error): trace-power characteristic class}
\noindent\textit{Transcription reads.} c\_i(E)=Tr(A\textasciicircum{}i)\par
\noindent\textit{Critical reading.} Use the normalized Chern-character component (1/i!)Tr(a\textasciicircum{}i), or the elementary invariant polynomial for c\_i.\par
\noindent\textit{Evidence.} The manuscript evidence block confirms c\_i; already for i=2 the elementary symmetric polynomial differs from the power sum.\par
\subsection*{P001378 (both; ambiguity): Contemporary Mathematics reference}
\noindent\textit{Transcription reads.} approximately volume 25\par
\noindent\textit{Critical reading.} Beilinson, Notes on absolute Hodge cohomology, Contemporary Mathematics 55, Part I (1986), 35-68.\par
\noindent\textit{Evidence.} Bibliography of the explicitly cited Esnault-Viehweg preprint identifies the work.\par
\subsection*{P001379 (both; faithful error): Deligne product with an integer on the right}
\noindent\textit{Transcription reads.} omega\_1 star z\_2=z\_2 omega\_1, together with the asymmetric top-form rule\par
\noindent\textit{Critical reading.} For the displayed asymmetric chain product, positive-degree x cup degree-zero y is zero.\par
\noindent\textit{Evidence.} Esnault-Viehweg Section 1.1 and the homotopic product models in Section 3; retaining both integer rules violates the chain-map identity.\par
\subsection*{P001380 (both; faithful error): Hodge decomposition hypothesis}
\noindent\textit{Transcription reads.} For compact X, H\textasciicircum{}n is the direct sum of H\textasciicircum{}(a,b)\par
\noindent\textit{Critical reading.} Add compact Kahler, or another hypothesis giving the asserted Hodge decomposition.\par
\noindent\textit{Evidence.} Compact complex manifolds need not have a degenerating Hodge-to-de Rham spectral sequence.\par
\subsection*{P001380-P001381 (both; faithful error): Deligne exact sequence lattice}
\noindent\textit{Transcription reads.} ordinary Z in every weight and an intersection inside complex cohomology\par
\noindent\textit{Critical reading.} Use the Tate lattice Z(p), and use the preimage of F\textasciicircum{}p in integral cohomology to retain torsion.\par
\noindent\textit{Evidence.} Esnault-Viehweg basic exact sequences; torsion maps to zero in complex cohomology.\par
\subsection*{P001381 (both; faithful error): comparison homotopy, equation 19}
\noindent\textit{Transcription reads.} d(alpha\_0+alpha\_1+...+alpha\_i)=alpha\_0-d alpha\_i\par
\noindent\textit{Critical reading.} Use the signed total differential of the degree-(n-1) cochain built from alpha\_1,...,alpha\_i; alpha\_0 is an endpoint.\par
\noindent\textit{Evidence.} The manuscript itself starts the homotopy sum at alpha\_1 on the previous line; alpha\_0 has a different total degree.\par
\subsection*{P001381 (both; faithful error): first Chern logarithm}
\noindent\textit{Transcription reads.} log(deg g\_ij)\par
\noindent\textit{Critical reading.} log(det g\_ij), with the 2pi i lattice normalization.\par
\noindent\textit{Evidence.} The evidence block confirms deg in equation 24 and det in equation 21; only the determinant is a transition function of det E.\par
\subsection*{P001382 (both; faithful error): motivic Chern class membership}
\noindent\textit{Transcription reads.} c\_i,alg(E)=H\textasciicircum{}(2i)(X,Z(i)\_M)\par
\noindent\textit{Critical reading.} c\_i,alg(E) belongs to H\textasciicircum{}(2i)(X,Z(i)\_M).\par
\noindent\textit{Evidence.} The source evidence block confirms an equals sign; the left side is a class and the right side a group.\par
\subsection*{P001383 (en; translation): historical name}
\noindent\textit{Transcription reads.} First devised by Spencer, Bloch\par
\noindent\textit{Critical reading.} First devised by Spencer Bloch, one mathematician; the given name is not a separate attribution.\par
\noindent\textit{Evidence.} The manuscript name Spencer is confirmed; Bloch is Spencer Bloch, as in the cited original K-theory literature.\par
\subsection*{P001342-P001343 (both; faithful error): toroidal spectral criterion}
\noindent\textit{Transcription reads.} unitarity of the representation would imply RH\par
\noindent\textit{Critical reading.} The relevant spectral conclusion is temperedness; unitarity alone includes complementary series.\par
\noindent\textit{Evidence.} Cornelissen-Lorscheid Introduction states the RH strategy with tempered constituents; the positive pairing alone is not enough.\par
